\documentclass[11pt,a4paper]{article}
\usepackage[margin=27mm]{geometry}
\usepackage{fontspec}
\setmainfont{Latin Modern Roman}
\usepackage{amsmath,amssymb,amsthm,mathtools}
\usepackage{microtype,booktabs,longtable,array}
\usepackage{xcolor}
\usepackage[unicode,colorlinks=true,linkcolor=blue!45!black,citecolor=blue!45!black,urlcolor=blue!45!black]{hyperref}
\usepackage{xurl}
\usepackage[numbers,sort&compress]{natbib}
\setlength{\emergencystretch}{2em}
\renewcommand{\thesection}{\Alph{section}}
\newcommand{\Fp}{\mathbf F_p}
\newcommand{\Z}{\mathbf Z}
\newcommand{\Q}{\mathbf Q}
\newcommand{\F}{\mathbf F}
\newcommand{\G}{\mathcal G}
\newcommand{\cl}{\mathrm{cl}}
\newcommand{\pr}{\mathrm{pr}}
\newcommand{\et}{\text{ét}}
\newcommand{\Gal}{\operatorname{Gal}}
\newcommand{\Hom}{\operatorname{Hom}}
\newcommand{\Aut}{\operatorname{Aut}}
\newcommand{\Out}{\operatorname{Out}}
\newcommand{\im}{\operatorname{im}}
\newcommand{\Spec}{\operatorname{Spec}}
\newcommand{\Pic}{\operatorname{Pic}}
\newcommand{\Inn}{\operatorname{Inn}}
\newcommand{\crr}{\mathrm{c}r}
\providecommand{\mathscr}{\mathcal}
\newenvironment{claim}[3]{\par\medskip\noindent\hypertarget{#1}{\textbf{#1. #2.}}\ \textbf{#3.}\par\nopagebreak\smallskip\begingroup\itshape}{\endgroup\par\smallskip}
\newcommand{\statementref}[1]{\hyperlink{#1}{#1}}
\title{ANABELIAN\_A1\\Prime-to-characteristic point separation:\\a proof audit and its scope}
\author{AI-produced by GPT-6 Astra\\OpenAI, ChatGPT Pro mode}
\date{17 November 2026}
\begin{document}
\maketitle
\begin{abstract}
This manuscript gives an independently audited proof of universal injectivity
of point sections for proper hyperbolic curves over finite fields, using
the whole arithmetic Galois quotient. Every assertion is assigned an
explicit status. The manuscript also records the separate existence,
completion, reconstruction, and characteristic-zero questions requested
with the proof audit. No priority is claimed.
\end{abstract}
\noindent\textbf{Overall task status: INCOMPLETE.}
Unsettled statements are not asserted as theorems.
\tableofcontents

\section{The point-separation argument}
\subsection{Conventions, conclusion, and the logical chain}
We retain the whole arithmetic quotient. If $k=\mathbf F_q$ has
characteristic $p$, $C/k$ is smooth, proper and geometrically connected,
and $\Sigma$ is a set of primes different from $p$, our notation is
\[
 \Delta_C^\Sigma=\pi_1(C_{\bar k})^\Sigma,\qquad
 1\longrightarrow\Delta_C^\Sigma\longrightarrow
 \Pi_C^\Sigma\longrightarrow G_k\longrightarrow1.
\]
For $\Sigma=\{\ell:\ell\ne p\}$ we omit the superscript. All
homomorphisms of profinite groups are continuous. Point sections are
classes under conjugation by the geometric kernel, not by a finite
approximation to that kernel. We reserve $E_C$ for the image on closed
points of the exceptional set of Sa\"idi--Tamagawa.

\begin{claim}{A.1}{Proved}{The main theorem, including every constant extension}
Let $C/k$ be a smooth proper geometrically connected curve of genus at
least two over a finite field of characteristic $p$. The map
\[
 C(k)\longrightarrow
 \{s:G_k\to\Pi_C\mid \operatorname{pr}\circ s=\operatorname{id}\}/\Delta_C
\]
is injective. The same assertion holds for $C_L/L$ for every finite
extension $L/k$. Thus AS-A.3 (AR.3), with precisely these hypotheses,
is correct. This answers cote 149, scan 36, for proper hyperbolic curves.
\end{claim}
\begin{proof}
We give two complete chains below, and specify their dependencies here.
Fix any $r\ne p$. The relative configuration exactness and boundary
transport in \statementref{A.3}--\statementref{A.5}
turn a geometric conjugacy of the sections of $P,Q$ into an isomorphism
of the arithmetic cuspidally pro-$r$ groups of $C-P$ and $C-Q$ over the
identity of $\Pi_C$, preserving the cusp conjugacy class. The criterion
\statementref{A.19} proves $P=Q$. Its inputs are proved separately:
the same-prime open-subgroup lemma, integral Frobenius and point-count
comparison, compactness over the entire system of covers, the uniform
finite exceptional image, self-normalization, and the orientation
extension and Kummer difference cocycle. No full mixed-prime
configuration-fibre injection is used.

There is also a proof avoiding the finite exceptional-image theorem:
the proper-field support and Kummer-cover arguments proved below
replace the finite-exception argument by the Weil bound along a
constant-field tower. The same relative configuration transport then
applies. Every hypothesis and comparison map in these two chains is
supplied in the statements that follow. Applying either chain over the
finite field $L$ proves the asserted uniformity in constant extensions;
it is not an assertion that a finite-level conjugacy lifts to the full group.
\end{proof}

The attached register identifies AS-A.87--A.90 with AG.39--AG.42,
and AS-A.97--A.98 with AR.8--AR.9. We use the working labels only
when discussing that earlier manuscript. Our statements are numbered
independently by Part. The proof audit and the computational checks
are AI-produced; we claim no priority.



\subsection{The relative configuration sequence, with all maps specified}

Throughout this subsection, a finite $\Sigma$-group means a finite group
whose order has all prime divisors in $\Sigma$. Completion means the
inverse limit of all such quotients. A prime $r\in\Sigma$ is fixed.
For a complex smooth proper curve $C$ of genus $g\geq2$, put
\[
 \Lambda=\pi_1(C-P),\quad \Gamma=\pi_1((C\times C)-\operatorname{diag}),
 \quad B=\pi_1(C),\quad \Delta=B^{\Sigma}.
\]
The second projection is used, so that $\Lambda$ is its discrete normal
fibre subgroup. This exact discrete sequence is the
Fadell--Neuwirth sequence; the precise surface hypotheses and the
vanishing of $\pi_2(C)$ are given in
\citet[Section 2.1]{Bellingeri2004}. Superscripts $\Sigma$ below denote
pro-$\Sigma$ completion, not a quotient of the arithmetic Galois group.

\begin{claim}{A.2}{Proved}{The free-fibre monodromy lemma}
Let $1\to L\to H\to B_0\to1$ be an exact sequence of discrete groups,
where $L$ is free of finite rank at least two. If conjugation by $H$
acts through an $r$-group on $L/[L,L]L^r$, the natural continuous map
$L^{(r)}\to H^{(r)}$ is injective. Every prescribed finite $r$-quotient
of $L$ is consequently dominated by the restriction of a finite
$r$-quotient of $H$.
\end{claim}
\begin{proof}
Write $L_r=L^{(r)}$. Its iterated Frattini subgroups are open and
characteristic, and they are cofinal: their images eventually vanish in
each finite $r$-quotient. For a finite $r$-group $Q$, the kernel of
$\Aut(Q)\to\Aut(Q/\Phi(Q))$ is an $r$-group. Indeed, it acts freely
on the $|\Phi(Q)|^d$ ordered lifts of a fixed basis of $Q/\Phi(Q)$,
and every such lift generates $Q$. Passing to characteristic finite
quotients proves that the kernel of
\[
 \Aut(L_r)\longrightarrow\Aut(L_r/\Phi(L_r))
\]
is pro-$r$. Thus the closure of the conjugation image of $H$ in the
profinite group $\Aut(L_r)$ is pro-$r$. The action extends continuously
to $H^{(r)}$.

For completeness, $L_r$ has trivial centre. The Magnus map sends its
free generators to $1+X_i$ in the units of the noncommutative formal
power-series algebra $\F_r\langle\!\langle X_1,\ldots,X_d\rangle\!\rangle$.
It is injective: for each finite $r$-quotient $Q$ of $L_r$, substitution
$X_i\mapsto q_i-1$ is defined because the augmentation ideal of
$\F_r[Q]$ is nilpotent, and it recovers the quotient map to $Q$.
Nilpotence follows by induction on $|Q|$, using a central subgroup of
order $r$ and $(z-1)^r=0$. A power series commuting with both $X_1$
and $X_2$ is scalar. To see this, compare homogeneous terms: a
homogeneous polynomial commuting with $X_1$ has only the monomial
$X_1^n$, by comparing initial and terminal letters successively; then
commutation with $X_2$ eliminates every positive-degree term. A central
Magnus unit of augmentation one is therefore one.

The composite $L_r\to H^{(r)}\to\Aut(L_r)$ is the inner action,
so it is injective. A continuous injective map from a compact group
into a Hausdorff profinite group is a closed embedding. Given an open
normal subgroup $W$ of $L_r$, there is an open normal subgroup $N$ of
$H^{(r)}$ with $N\cap L_r\subset W$: use the induced topology and
then take the finite normal core of an open neighbourhood. The quotient
$H^{(r)}/N$ gives the stated domination.
\end{proof}

\begin{claim}{A.3}{Proved}{Exactness for a relative single-prime cusp kernel}
Set $F=\Lambda^\Sigma$, $E=\Gamma^\Sigma$, and
\[
 \begin{split}
 K_F&=\ker(F\to\Delta),\qquad
 K_E=\ker(E\to\Delta\times\Delta),\\
 F_{\crr}&=F/\ker(K_F\to K_F^{(r)}),\qquad
 E_{\crr}=E/\ker(K_E\to K_E^{(r)}).
 \end{split}
\]
The induced continuous sequence
\[
 1\longrightarrow F_{\crr}\longrightarrow E_{\crr}
       \xrightarrow{\pr_2}\Delta\longrightarrow1
\]
is exact. Every finite quotient of $F_{\crr}$ is dominated on the
fibre by a finite quotient of $E_{\crr}$ which has a normal $r$-subgroup
and a quotient $A\times A$, for a finite $\Sigma$-quotient $A$ of
$\Delta$.
\end{claim}
\begin{proof}
Let $q:F_{\crr}\twoheadrightarrow Q$ be continuous with $Q$ finite.
The image $R$ of $K_F^{(r)}$ is a normal $r$-subgroup, and
$A=Q/R$ is a quotient of $\Delta$. Let $a:C'\to C$ be the connected
Galois proper cover corresponding to this quotient. Put
\[
 Y=(C'\times C')-(a\times a)^{-1}(\operatorname{diag} C).
\]
Its removed divisor is the disjoint union of the graphs of the deck
transformations of $a$, so $Y\to C'$ has individually labelled
punctures. Choose a point of $C'$ above $P$. The fibre is
$V'=C'-a^{-1}(P)$. The surface fibration gives a discrete exact sequence
\[
 1\to L=\pi_1(V')\to H=\pi_1(Y)\to\pi_1(C')\to1.
\]
Here $L$ is free of rank $2g(C')+|A|-1\geq2$. In
$H_1(V',\F_r)$ the boundary classes span a subspace $W$, and there is
an exact sequence
\[
 0\to W\to H_1(V',\F_r)\to H_1(C',\F_r)\to0.
\]
Monodromy acts trivially on $W$, because the boundary sections are
labelled and oriented. It acts trivially on the quotient, because the
compactified family is the product $C'\times C'\to C'$. Choosing any
vector-space splitting shows that its image is contained in the additive
$r$-group $\Hom_{\F_r}(H_1(C',\F_r),W)$. This does not require
$r\nmid|A|$. By \statementref{A.2}, $L^{(r)}$ embeds in $H^{(r)}$.

The subgroup $H$ is normal in $\Gamma$, with quotient $A\times A$,
and $\Lambda\cap H=L$. The restriction of $q$ to $L$ takes values in
$R$. Choose a finite $r$-quotient of $H$, with kernel $J$, whose
restriction dominates this quotient of $L$. Thus
$J\cap L\subset\ker(q)$. Take
\[
 J^*=\bigcap_{\gamma\in\Gamma/H}\gamma J\gamma^{-1}.
\]
This is a finite intersection since $J$ is normal in $H$. The group
$H/J^*$ embeds in a finite product of $r$-groups. Consequently
$\Gamma/J^*$ is a finite $\Sigma$-group, an extension of $A\times A$
by an $r$-group, and $J^*\cap\Lambda\subset\ker(q)$. This last
containment remains true on completed fibre groups, as is seen in the
joint finite image of the two quotient maps.

The quotient of $\Gamma$ factors through $E_{\crr}$. Its image of
$K_E$ lies in $H/J^*$ because the $A\times A$ quotient comes from
the proper product; it therefore kills $\ker(K_E\to K_E^{(r)})$.
Conversely, the fibre map $F\to E_{\crr}$ has pro-$r$ image on
$K_F$, so it factors through $F_{\crr}$. Detection of every finite
quotient proves that this induced map is injective.

Finally, right exactness of pro-$\Sigma$ completion follows directly
from its universal property: quotienting $E$ by the closed image of
$F$ represents precisely the finite $\Sigma$-quotients of $B$, and
hence gives $\Delta$. The subgroup killed to form $E_{\crr}$ is
inside $K_E$, which is inside this fibre kernel. Thus the image of
$F_{\crr}$ is the whole kernel of $\pr_2$, and $\pr_2$ is onto.
\end{proof}

\begin{claim}{A.4}{Proved}{Specialization and the arithmetic fibre}
Let $C$ be smooth proper geometrically connected of genus at least two
over an algebraically closed field of characteristic $p$, and let
$\Sigma$ be any set of primes not containing $p$, with $r\in\Sigma$.
The preceding relative exact sequence holds for geometric etale groups.
If $C/k$ is defined over a finite field and $P\in C(k)$, its arithmetic
cuspidally pro-$r$ group is the inverse image of a chosen point section
$s_P(G_k)$ under the second projection of the arithmetic configuration
quotient. The comparison is over the identity of the whole $G_k$.
\end{claim}
\begin{proof}
We record the exact specialization input. SGA~1, Expose XIII,
Corollary~2.9, assumes a smooth proper morphism of finite presentation
$X\to S$, the complement $U$ of a relative normal-crossings divisor,
and the set of primes different from every residual characteristic of
$S$. For a locally constant constructible sheaf of finite groups with
these primes and tame along the boundary, it gives locally constant
$R^1$ and cohomological properness through degree one. Section~2.10
constructs the specialization homomorphism by strict localization and
a path; its ambiguity is inner conjugation. Applied to constant finite
groups, Corollary~2.9 identifies their torsors on the two geometric
fibres, so the resulting maximal prime-to-$p$ specialization map is an
isomorphism. It is functorial in the pair. Lemma~2.11 identifies the
images of inertia groups under this specialization. These statements
are in the opened text at pages 288--290, corresponding to original
margin pages 389--392 of \cite{LA_SGA1}.

Here a suitable lift exists. SGA~1, Expose III, Theorem~7.3 assumes
$H^2(C,T_C)=H^2(C,\mathcal O_C)=0$ for a smooth projective special
fibre over the residue field of a complete local ring and gives a
smooth projective lift; Corollary~7.4 states the curve case. Both
vanishings hold for a curve. A chosen point lifts successively by
formal smoothness, and the compatible points give a section over the
complete ring. Equivalently, the marked deformation obstruction lies
in $H^2(C,T_C(-P))=0$. Take the Witt ring of a perfect algebraically
closed field of definition. The pairs
\[
 (\mathcal C,\mathcal P),\qquad
 (\mathcal C\times\mathcal C,\operatorname{diag}\mathcal C)
\]
are smooth proper pairs with smooth relative boundary. Their fibre
inclusion is transverse with multiplicity one. Functorial
specialization therefore identifies the groups and maps used in
\statementref{A.3}; the characteristic relative kernels and their
pro-$r$ quotients are identified as well.

For general algebraically closed fields, first descend the pointed
curve to a finitely generated field over $\F_p$, and take its
algebraic closure there. This is countable. Its Witt fraction field has
cardinality at most that of $\mathbb C$, so an algebraic closure
embeds into $\mathbb C$. Invariance of tame prime-to-$p$ groups of
smooth proper pairs under algebraically closed field extension, also
obtained from the same proper-pair comparison of finite torsors,
permits this descent and the return to the original field. Taking a
further maximal pro-$\Sigma$ quotient transfers the result for every
$\Sigma$ under consideration.

For the arithmetic assertion, write $\Pi_C$ for the geometrically
pro-$\Sigma$ group over the unchanged $G_k$. The configuration group
maps onto $\Pi_C\times_{G_k}\Pi_C$. Its kernel relative to this
product is geometric; replacing it by its pro-$r$ quotient gives a
group $\mathcal E_{\crr}$ with geometric subgroup $E_{\crr}$.
The actual fibre $C-P$ induces a map of extensions
\[
 \Pi_{C-P,\crr}\longrightarrow
 \mathcal E_{\crr}\times_{\Pi_C,s_P}G_k.
\]
On geometric kernels it is the isomorphism of
\statementref{A.3}, and on arithmetic quotients it is the identity
of $G_k$. Injectivity and surjectivity follow by the elementary
comparison of short exact sequences. Since $s_P$ is injective, the
right-hand group identifies with $\pr_2^{-1}(s_P(G_k))$.
Paths give exactly geometric inner conjugacy, not a change of the
arithmetic quotient.
\end{proof}

\begin{claim}{A.5}{Proved}{The common diagonal boundary and cusp transport}
In the arithmetic situation of \statementref{A.4}, the image
$\mathcal D$ of a fixed diagonal boundary component fits into a
continuous exact sequence
\[
 1\to\Z_r(1)\to\mathcal D\to\Pi_C\to1,
\]
and maps diagonally to $\Pi_C\times_{G_k}\Pi_C$. With paths chosen
through this boundary, the actual cusp subgroup of $C-P$ is
$\mathcal D\cap\pr_2^{-1}(s_P(G_k))$. Consequently, geometrically
conjugate proper point sections give an isomorphism of their actual
cuspidally pro-$r$ punctured groups over the identity of $\Pi_C$,
carrying the cusp subgroup to a geometric conjugate of the other cusp.
\end{claim}
\begin{proof}
Put $X=C\times C$, $D=\operatorname{diag} C$, and
$N=\mathcal O_X(D)|_D$. For an integer $n$ prime to $p$, the boundary
of $\sqrt[n]{(X,D)}$ is the gerbe of $n$-th roots of $N$ on $D$.
Here is the tame comparison being used, rather than a configuration
fibre exactness assertion. SGA~1 XIII, Proposition~5.2, assumes a
regular local scheme, a normal-crossings divisor with local equations
$f_i$, and a finite cover tame along that divisor. If $n_i$ is the
least common multiple of the relevant inertia orders, these integers
are invertible and the cover extends uniquely to a finite etale cover
after adjoining $T_i$ with $T_i^{n_i}=f_i$. The conclusion is over the
entire Kummer chart; the proof invokes purity. This is the opened
statement at pages 315--316, original margin page 429, of \cite{LA_SGA1}.
Thus every finite tame cover of $X-D$ extends to a sufficiently high
root stack. Restriction to its boundary gerbe defines compatible maps
from the boundary fundamental groups to the configuration group.

A finite etale cover of the boundary gerbe with trivial stabilizer
action descends to $D$, and covers of $D$ pull back fully faithfully.
Etale locally the gerbe is $B\mu_n$; consequently the kernel of its
map to $\pi_1(D)$ is the image of $\mu_n$. An injection of this
finite stabilizer has not been assumed. Taking the inverse limit over
root orders, and then the geometrically pro-$\Sigma$ quotient, gives
a right exact sequence
\[
 \widehat\Z^{\Sigma}(1)\longrightarrow\mathcal T
                 \longrightarrow\Pi_C\longrightarrow1.
\]
Exactness follows also directly from the same trivial-stabilizer
descent criterion. The transition maps of the root gerbes are
surjective on fundamental groups: their maps on the base are the
identity and the maps of stabilizer images are onto. Compact inverse
limits preserve these right exact sequences. Right exactness of
pro-$\Sigma$ completion shows that discarding the geometric base
kernel introduces no further kernel over $\Pi_C$.

Let $\mathcal D$ be the image of $\mathcal T$ in
$\mathcal E_{\crr}$. Its map to the proper product is the injective
diagonal on $\Pi_C$, so its kernel over $\Pi_C$ is precisely the
image of inertia, a quotient of $\Z_r(1)$. It is the full
$\Z_r(1)$. Indeed, the geometric group of the fibre $C-P$ maps
onto the free pro-$r$ group on $2g$ generators. Its boundary loop is
the product of their commutators, up to orientation. Sending one pair
to $1+e_{12},1+e_{23}$ in the upper-unitriangular group over $\Z_r$
and the others to one sends that boundary loop to $1+e_{13}$, of
infinite $r$-adic order. The fibre injection of
\statementref{A.3} retains this inertia in the total quotient.
Transversality identifies it with diagonal inertia with multiplicity
one. This argument includes $r=2$.

The local transverse Kummer chart sends the actual cusp decomposition
group into $\mathcal D$ over $s_P(G_k)$, identically on inertia and
on the arithmetic residue group. Comparing the two short exact
sequences therefore identifies it with the whole inverse image of
$s_P(G_k)$ in $\mathcal D$. The paths for $P$ and $Q$ may both be
chosen in the same connected boundary tower, so this uses the same
$\mathcal D$ without assuming any prior separation of points.

Suppose $\gamma\in\Delta$ conjugates $s_P(G_k)$ to $s_Q(G_k)$.
Choose a geometric $d\in\mathcal D$ over $\gamma$. Conjugation by
$d$ identifies the two fibre subgroups and their cusp subgroups; on
the first proper coordinate it induces $\Inn(\gamma)$. The first
projection from the geometric part of the $Q$-fibre onto $\Delta$
is surjective. Choose there an $e$ mapping to $\gamma^{-1}$. Then
$\Inn(e)\Inn(d)$ is over the identity of $\Pi_C$ and $G_k$, and
carries the first cusp to a geometric conjugate of the second. No
lifting of a point section to a section of the configuration fibration
has been assumed.
\end{proof}

\subsection{Finite positive and negative tests in genus two}

\begin{claim}{A.6}{Proved}{A mixed-prime finite quotient detecting the relative fibre}
For genus two, $p=2$, $r=3$, and
$A=(\Z/3)^4\times\Z/5$, there is a quotient of the fibre of
order $1215$, with central cuspidal subgroup of order three and proper
quotient $A$, which embeds in a quotient of the configuration group of
order $492075$ with central kernel of order three over $A\times A$.
The defining surface braid relations and the fibre injection have been
checked by the retained program.
\end{claim}
\begin{proof}
Let $V=\F_3^4$ have standard alternating form
$J(v,w)=v_1w_2-v_2w_1+v_3w_4-v_4w_3$. On
\[
 G=(V\oplus V)\times\F_3\times\F_5^2
\]
define multiplication by adding the noncentral coordinates and adding
$2J(v_1+v_2,w_1+w_2)$ to the central coordinate. This is associative
because the added function is bilinear; $2$ is $1/2$ in $\F_3$.
Its order is $3^9\cdot25=492075$. Factor interchange preserves the
law. Form $G\rtimes\langle\tau\rangle$, where $\tau$ interchanges
the two $V$ coordinates and the two $\F_5$ coordinates.

Use \citet[Theorem 1.2]{Bellingeri2004} with $n=g=2$. Write its first
strand generators as $a,b,c,d$, put
$t=[a,b^{-1}][c,d^{-1}]$, and denote the interchange generator by
$\sigma$. Send $a,b,c,d$ to the four standard basis vectors in the
first $V$, giving $a$ also the coordinate $(1,0)$ in $\F_5^2$;
all central coordinates are zero. Send $\sigma$ to the product of
central coordinate $2$ with $\tau$. The eleven relations of the
presentation become identities: the four same-generator relations
use zero self-pairing, the four different-handle relations use zero
cross-handle pairing, the two same-handle mixed relations use the fact that
$[a,b]$ and $[c,d]$ both have central coordinate $1$ in $\F_3$,
and the surface relation
uses $-2=1$ in $\F_3$, giving $t=\sigma^2$.
The code checks each relation directly in this finite group.

The pure braid image lies in $G$ and is all of $G$. The first and
second strand images give eight independent $\F_3$ coordinates;
$a^3$ and its $\sigma$-conjugate give the two $\F_5$ coordinates,
and $[a,b]$ gives the central coordinate. The first fibre image is
\[
 Q=V\times\F_3\times\F_5\hookrightarrow G,
 \qquad (v,z,f)\longmapsto((v,0),z,(f,0)).
\]
This is injective by its coordinates and has order $3^5\cdot5=1215$.
The normal closure of $t$ is the central $\F_3$, in both groups.
Thus these finite maps factor through the respective cuspidally
pro-three groups over the full pro-odd proper groups, and their
proper quotients are exactly $A$ and $A\times A$. The program also
enumerates the $1215$ fibre elements and checks closure under its
specified generators. By \statementref{A.4}, the same geometric
finite test applies to the genus-two curve $y^2+y=x^5$ over
$\overline{\F}_2$.
\end{proof}

\begin{claim}{A.7}{Proved}{An explicit finite obstruction to full pro-odd exactness}
Let $F_4=\langle a,b,c,d\rangle$ be the first fibre in the genus-two
pure surface braid group $P_2$. Map $F_4$ onto the Heisenberg group
$H=\operatorname{UT}_3(\F_3)$ by sending $a,b$ to its standard
generators and $c,d$ to one. Let $K$ be the kernel and
$N=[K,K]K^5$. Then the finite odd quotient $F_4/N$, of order
$27\cdot5^{82}$, is not dominated on the fibre by any finite odd
quotient of $P_2$. In particular $F_4^{(2')}\to P_2^{(2')}$ is not
injective. The conclusion applies to the geometric groups of
$y^2+y=x^5$ in characteristic two by specialization.
\end{claim}
\begin{proof}
This is a finite certificate, supplemented by the proof that it
obstructs every possible finite odd total quotient. In Bellingeri's
presentation used above, set
$C=\sigma c\sigma^{-1}$ and $D=\sigma d\sigma^{-1}$.
Conjugation by them fixes $a,b$. On the other generators it is
\[
 \begin{array}{ll}
 \varphi_C(c)=tct^{-1},&
 \varphi_C(d)=tdt^{-1}(tc^{-1}tct^{-1}),\\
 \varphi_D(c)=ct^{-1},&
 \varphi_D(d)=tdt^{-1}.
 \end{array}                                                    \tag{CF.1}
\]
These are the actual conjugation automorphisms, not guessed
endomorphisms. For example the same-generator relation and its
$\sigma$-conjugate give
$CcC^{-1}=tct^{-1}$ and
$\varphi_C(t)=tc^{-1}tct^{-1}$. The mixed relation gives
$\varphi_C(dt^{-1})=tdt^{-1}$. Its conjugate gives
$\varphi_D(c)=ct^{-1}$; the analogous same-generator relation gives
the formula for $d$. The different-handle relations give the asserted
fixing of $a,b$. Normality of the discrete fibre follows from the
Fadell--Neuwirth sequence. The code independently verifies these
rewritten relations and the images of $t$ by free reduction.

The Heisenberg group is represented by triples with product
\[
 (x,y,z)(u,v,w)=(x+u,y+v,z+w+xv)\quad\text{over }\F_3.
\]
Both $\varphi_C^3$ and $\varphi_D^3$ preserve the quotient map to
$H$, as direct evaluation of all four generator images verifies.
Thus they normalize $K$, and they normalize $N$. A spanning tree in
the Cayley graph of this quotient gives a free Schreier basis of $K$
of rank $1+27(4-1)=82$. Its homology coordinates are computed by
recording oriented non-tree edges; the program verifies that its
basis-coordinate matrix is the identity.

Let $z=[a,b]\in H$. In the integral homology of $K$, write $c_s$
and $d_s$ for the lifted $c,d$ loops at vertex $s$. They are non-tree
edges, since $c,d$ map to one. On the independent vectors
\[
                 x=c_1-c_z,\qquad y=d_1-d_z
\]
the complete integral Schreier calculation gives
\[
 M_C=\begin{pmatrix}1&3\\0&1\end{pmatrix},\qquad
 M_D=\begin{pmatrix}1&0\\-3&1\end{pmatrix}.                    \tag{CF.2}
\]
Here $M_C$ is the action of $C^3$ and $M_D$ that of $D^3$.
The program exports the entire two $82$ by $82$ integral matrices
and checks all four equalities defining (CF.2), so this is not a
numerical approximation or a guessed invariant subspace. Hence
\[
 M_CM_D=\begin{pmatrix}-8&3\\-3&1\end{pmatrix},\qquad
 (M_CM_D)(x-y)=-(x-y)\pmod5.
\]
The class $x-y$ in $K/N$ is nonzero, having four nonzero Schreier
coordinates.

Suppose a finite odd quotient $P_2\to T$ with kernel $J$ dominated
$F_4/N$. Then $J\cap F_4\subset N$. Let
$H_0=\langle F_4,C^3,D^3\rangle$. It normalizes $K$ and $N$.
The quotient of the image of $K$ by the image of $N$ in $T$ is
exactly $K/N$, since $J\cap K\subset N$. The conjugation action of
$H_0$ on it therefore factors through the finite odd group
$\im(H_0\to T)$. In particular $C^3D^3$ must act with odd order.
Its action on the nonzero line generated by $x-y$ has order two,
a contradiction. No invariance of $K$ under the entire $P_2$ was
needed. Finally, an injective map between the two compact pro-odd
groups would induce the full topology on the fibre and would supply
a dominating finite total quotient for this specified finite fibre
quotient. Thus it cannot be injective.

The retained files are
\path{code/configuration_finite_tests.py}, its exact output, and
\path{code/configuration_finite_certificates.json}; the latter
contains the presentation inputs, the free words, the Schreier basis,
the integral matrices, and the sparse vector. The factorized arithmetic size is $27\cdot5^{82}$; its exact decimal
value is also retained in the program output.
\end{proof}

\subsection{Which sets of primes preserve the full configuration fibre?}

\begin{claim}{A.8}{Proved}{Exactness would pass to proper covers}
Suppose $C/\mathbb C$ has genus at least two and the full pro-$\Sigma$
fibre map for $(C\times C)-\operatorname{diag}$ is injective.
For every connected finite Galois etale $\Sigma$-cover $C'\to C$,
the analogous full fibre map for $C'$ is injective as well.
\end{claim}
\begin{proof}
First, if a subgroup $H$ of a discrete group $G$ contains a normal
subgroup of finite $\Sigma$-index, it has its full pro-$\Sigma$
topology induced from $G$. Given a normal $\Sigma$-index subgroup
of $H$, intersect with the fixed normal subgroup and take the
$G$-core. The quotient of the latter normal subgroup embeds in a
finite product of $\Sigma$-groups, and the resulting quotient of
$G$ is an extension of $\Sigma$-groups. It is therefore a
$\Sigma$-group. This proves the topology assertion and the
identification of the corresponding completed open subgroup.

Let $a:C'\to C$ have group $A$, and set
$Y=(C'\times C')-\bigcup_{\alpha\in A}\operatorname{graph}(\alpha)$.
It is a normal $A\times A$ cover of the original configuration
space. Choose $P'\in C'$ above $P$. Its fibre is
$V'=C'-a^{-1}(P)$, the corresponding proper-cover subgroup of
$C-P$. The preceding topology assertion and the assumed injection
identify $(\pi_1V')^\Sigma$ with a closed normal subgroup of
$(\pi_1Y)^\Sigma$; the quotient is $(\pi_1C')^\Sigma$.

Let $N$ be the closed normal subgroup of this fibre generated by
inertias at the points of $a^{-1}(P)$ other than $P'$. Monodromy
preserves these individual puncture labels. Thus the discrete total
group normalizes $N$, and so does its completion, since the
normalizer of a closed subgroup is closed. Filling the nonidentity
graphs in $Y$ kills exactly the closed total normal subgroup generated
by these same meridians. This elementary assertion follows by putting
a transverse disk through each smooth divisor and applying van
Kampen; every meridian along that divisor is a total conjugate of
its meridian in the chosen fibre. Since $N$ is already total-normal,
that total filling kernel equals $N$.

It follows that the quotient of the completed total group is
$(\pi_1((C'\times C')-\operatorname{diag}))^\Sigma$, and that its
injected fibre is $(\pi_1(C'-P'))^\Sigma$. Here we used right
exactness of pro-$\Sigma$ completion for both fillings; it follows
from the finite-quotient universal property exactly as in
\statementref{A.3}. This proves the assertion without assuming
that arbitrary completion commutes with kernels.
\end{proof}

\begin{claim}{A.9}{Proved}{A finite symplectic monodromy obstruction}
Let $C/\mathbb C$ have genus $h\geq2$, and let $r,\ell$ be distinct
primes. There is a connected degree-$r^3$ cover of $C-P$, with
$r$-group Galois group and cusp inertia of order $r$, whose compact
homology modulo $\ell$ has a subspace of dimension $2h-2$ on which
a subgroup of the configuration group preserving the cover acts
onto $\operatorname{Sp}_{2h-2}(\F_\ell)$. Consequently, if $r,\ell\in\Sigma$
and the full pro-$\Sigma$ configuration fibre were injective, every
prime divisor of $|\operatorname{Sp}_{2h-2}(\F_\ell)|$ would belong
to $\Sigma$.
\end{claim}
\begin{proof}
Choose a one-holed torus containing $P$ and a complementary connected
subsurface $M$ of genus $h-1$ with one boundary component. Map the
free group of $C-P$ onto $\operatorname{UT}_3(\F_r)$ by sending
the first two handle generators to its standard generators and all
generators supported on $M$ to one. Its boundary commutator is a
nontrivial central element $z$ of order $r$ (up to inversion).
Let $D^\circ\to C-P$ be the resulting regular cover and $D$ its
compactification. The restriction over $M$ is the disjoint union of
$r^3$ copies $M_s$, labelled by deck elements $s$. For a fixed $s$,
put
\[
 W=\{u_s-u_{zs}:u\in H_1(M,\F_\ell)\}\subset H_1(D,\F_\ell).
\]
The map defining $W$ is injective. Indeed, the usual symplectic handle
cycles in the disjoint subsurfaces $M_s$ have block-diagonal
intersection pairings with their respective dual cycles, so all their
classes are independent in $D$. This argument also works when
$\ell=2$, where the restricted pairing on $W$ itself can be
 degenerate. We use its identified vector space, not a false claim
that this restricted pairing is nondegenerate.

We give the topological calculation of the action explicitly. Let
$\alpha$ be an oriented simple nonseparating curve in $M$, joined to
$P$ by an embedded arc avoiding a fixed first-strand basepoint.
Pushing $P$ once round this based simple loop is the product of two
opposite Dehn twists on its left and right push-offs, denoted
$\alpha_0,\alpha_1$. To verify the identity, identify an annular
neighbourhood of the based loop with angular coordinate $\theta$ and
transverse coordinate $t$, and move the marked point by
$\theta\mapsto\theta+2\pi u$ during the isotopy parameter $u$.
At its two boundary collars the resulting map is one positive and
one negative twist. It is the identity outside this neighbourhood.
The two push-offs and the puncture bound a pair of pants. Their
cover monodromies are respectively $1,z$ (possibly after replacing
$z$ by $z^{-1}$); centrality makes the joining path immaterial.

The $r$-th power of each twist lifts and fixes the chosen sheet
outside the twisting annuli. On an annular cover of degree $n$,
where $n\mid r$, a twist of exponent $r$ lifts to exponent $r/n$:
in coordinates $\theta_{\rm base}=n\theta_{\rm cover}$ the lift is
$\theta_{\rm cover}\mapsto\theta_{\rm cover}+2\pi(r/n)f(t)$.
It is the identity on the boundary circles, so it glues to the
identity outside all annuli. This proves existence of the lift,
its sheet convention, and preservation of the cover without an
unproved assertion about outer actions.

The lifts of $\alpha_0$ have degree one and are the curves
$\alpha_s$ in the separate copies of $M$. Their twists have
exponent $r$. The lifts of $\alpha_1$ have degree $r$, one for
each central $\langle z\rangle$-orbit, and their twists have
exponent one. Each such long lift has homology class
\[
                        \sum_{j=0}^{r-1}[\alpha_{z^j s}]
                                                               \tag{CF.3}
\]
up to its common orientation. For this equality, lift the pair of
pants: each connected component is a degree-$r$ cyclic cover with
$r$ short boundaries, one long boundary, and one boundary above the
puncture. After compactification the last boundary bounds a disk.
The oriented boundary relation gives exactly (CF.3), with no extra
homology component.

The intersection of (CF.3) with $u_s-u_{zs}$ is zero: the two
coefficients are equal and cancel. Thus the long-lift twists act
trivially on $W$. The short-lift twists send
\[
 u_s-u_{zs}\longmapsto
 u_s-u_{zs}\ \mathbin{\pm}\ r\langle u,\alpha\rangle
                          (\alpha_s-\alpha_{zs}).              \tag{CF.4}
\]
The whole lifted point-push therefore preserves $W$ and induces the
usual symplectic transvection of exponent $\pm r$ on its identified
space $H_1(M,\F_\ell)$. All these lifts have the same sheet labels,
because they fix every point above the chosen basepoint outside the
annuli.

Since $r$ is invertible modulo $\ell$, powers give the transvection
of any coefficient. Every primitive integral homology class of $M$
is represented by a simple nonseparating curve: Euclidean handle
slides reduce a primitive vector to a handle basis vector, and
reversing the slides supplies the curve. Every nonzero vector modulo
$\ell$ has such a primitive integral lift (adjust one coordinate
by multiples of $\ell$ if necessary). We therefore obtain all
symplectic transvections on $H_1(M,\F_\ell)$.

These generate its full symplectic group. Here is a linear-algebra
proof, including characteristic two. A transvection of suitable
coefficient about $v-u$ sends $u$ to $v$ whenever
$\langle u,v\rangle\ne0$. If the pairing is zero, an intermediate
vector with both pairings nonzero exists, since two proper
hyperplanes do not exhaust the vector space. Thus the transvection
group is transitive on nonzero vectors. After fixing the first vector
$e$ of a symplectic basis, write its partner as $f+w+ae$, with $w$
in the symplectic complement. The product
$T_{e+w}T_w^{-1}T_e^{-1}$ fixes $e$ and removes $w$, with signs
chosen according to $\langle e,f\rangle=1$; a power of $T_e$
then removes $ae$. Induction in the symplectic complement fixes the
whole basis. This proves generation of $\operatorname{Sp}_{2h-2}$.

To obtain the last assertion, suppose the completed fibre is a
closed subgroup of the pro-$\Sigma$ configuration group. The
$r^3$-index subgroup defining $D^\circ$ has its full pro-$\Sigma$
topology there by the normal-core argument. Its quotient
$H_1(D,\F_\ell)$ is finite. The chosen point-push powers normalize
both the covering subgroup and the closed kernel defining this
homology quotient. Their closures still do so, because normalizers
of closed subgroups are closed. Conjugation on a finite quotient is
continuous, as one checks on finitely many coset representatives.
It therefore has finite $\Sigma$-image. Restricting to $W$ gives
a quotient of that image, but the preceding calculation makes this
quotient the full $\operatorname{Sp}_{2h-2}(\F_\ell)$. Its prime
divisors must lie in $\Sigma$.
\end{proof}

\begin{claim}{A.10}{Proved}{Classification of prime sets for the full fibre sequence}
For a smooth proper complex curve of genus at least two, the full
pro-$\Sigma$ configuration fibre sequence is exact if and only if
$\Sigma$ is empty, consists of one prime, or is the set of all
primes. If $C$ instead has characteristic $p$ and
$\Sigma\subseteq\{\ell:\ell\ne p\}$, it is exact if and only if
$|\Sigma|\leq1$. In particular the full prime-to-$p$ fibre map is
noninjective for every $p$ and every such curve, as claimed in R4.
The relative cuspidally pro-$r$ sequence of
\statementref{A.3} remains exact for every $r\in\Sigma$.
\end{claim}
\begin{proof}
The empty case is trivial. For $\Sigma=\{r\}$, monodromy on
$H_1(C-P,\F_r)=H_1(C,\F_r)$ is trivial, because the compactified
family is the product and there is only one boundary class, which
vanishes in abelian homology. Apply \statementref{A.2} and right
exactness. For the set of all primes, we use
\citet[Proposition 2.2]{LA_MT2008}, with a characteristic-zero
hyperbolic curve and the projection from its second to its first
configuration space: the proposition gives the full profinite fibre
exact sequence. Its prime-set hypothesis allows the total set (or a
singleton); it does not assert the sequence for a proper multi-prime
set.

Suppose now that $\Sigma$ is proper and contains distinct primes
$r,\ell$, and choose $p\notin\Sigma$. If the fibre sequence were
exact for $C$, it would be exact for every connected cyclic proper
cover of degree $r^m$, by \statementref{A.8}. Such covers exist by
the surjection $H_1(C,\Z)\twoheadrightarrow\Z/r^m$ and have genus
$h=1+r^m(g-1)$, hence arbitrarily large genus. Let $d$ be the
multiplicative order of $\ell$ modulo $p$ (for $p=2$, take $d=1$).
Choose $m$ such that $h-1\geq d$. Then $p$ divides
$\ell^{2d}-1$, and therefore divides
\[
 |\operatorname{Sp}_{2h-2}(\F_\ell)|
       =\ell^{(h-1)^2}\prod_{i=1}^{h-1}(\ell^{2i}-1).
\]
For reference, this order formula follows inductively by choosing a
nonzero first vector, its partner of pairing one, and a symplectic
basis of the complement; the successive factors are
$(\ell^{2i}-1)\ell^{2i-1}$. The divisibility contradicts
\statementref{A.9}, since $p\notin\Sigma$. Thus left injectivity
fails. Middle and right exactness still hold by the universal
finite-quotient argument, so this identifies the precise lost part
of the sequence.

For characteristic $p$, the functorial proper-pair specialization in
\statementref{A.4} identifies the pro-$\Sigma$ fibre map with the
complex one. As $p\notin\Sigma$, a nonempty such set is always
proper, giving exactly the stated classification. The mechanism also
explains the relative repair: a single retained cuspidal prime makes
all finite covering monodromy kernels used in
\statementref{A.3} $r$-groups; keeping two independent cusp primes
permits the symplectic finite actions above and their omitted-prime
divisors.
\end{proof}

\begin{claim}{A.11}{Proved}{The labelled several-section extension}
Let $C_0,B/\mathbb C$ be smooth proper curves, with $B$ of genus at
least one and $C_0$ of genus at least two, and let
$f_1,\ldots,f_m:B\to C_0$ have pairwise disjoint graphs. For
$Y=(C_0\times B)-\bigcup_i\operatorname{graph}(f_i)$, the relative
cuspidally pro-$r$ fibre sequence over pro-$\Sigma$ proper groups
is exact for every $r\in\Sigma$. The same assertion holds for simultaneously liftable
smooth proper pairs in characteristic prime to every member of
$\Sigma$ by the preceding specialization. In particular it applies
to the diagonal families after finite proper covers used in
cuspidally abelian reconstruction.
\end{claim}
\begin{proof}
For a finite fibre quotient, take its proper quotient $A$ and the
associated Galois cover $C_1\to C_0$. Pull this cover back by every
$f_i$, and take a common finite Galois $\Sigma$-cover $B_1\to B$
which trivializes all these $A$-torsors. Such a cover is given by
the intersection of their kernels; its group is a subgroup of a
finite product of copies of $A$. Now the inverse boundary in
$C_1\times B_1$ consists of disjoint labelled sections. Its compact
family is a product, its base is aspherical, and the homology action
modulo $r$ is the same shear group used in
\statementref{A.3}. Apply \statementref{A.2} and take the
normal core in $\pi_1(Y)$. The resulting total finite quotient is
an $r$-group by $A\times A_B$, where $A_B$ is the chosen base-cover
group; this is a $\Sigma$-group. Every finite fibre quotient is
detected, and the middle and right exactness proof is unchanged.
For the specific characteristic-$p$ diagonal-cover families, lift
the proper etale covers together with the smooth curve (finite
prime-to-$p$ etale covers lift uniquely on the proper-pair comparison)
and specialize. For a general characteristic-$p$ family of graphs,
the same conclusion requires a simultaneous lift of the indicated
smooth proper pair; the statement is understood with that hypothesis,
not as an assertion that arbitrary morphisms of curves lift.
\end{proof}

\begin{claim}{A.12}{Proved}{A finite two-prime test beyond the cofinite case}
For a genus-three curve, the full pro-$\{2,3\}$ configuration fibre
map is noninjective. A particular obstructing fibre quotient has order
$8\cdot3^{41}$; an exact computation of the preserving braid action
on a four-dimensional homology subspace gives
$\operatorname{Sp}_4(\F_3)$, of order $51840$, and an explicit
acting element of order five.
\end{claim}
\begin{proof}
Let $F_6$ have handle generators $a,b,c,d,e,f$, and let $K$ be the
kernel of its quotient onto $\operatorname{UT}_3(\F_2)$ obtained
from $a,b$ alone. Then $K$ is free of rank $1+8(6-1)=41$.
Use $N=[K,K]K^3$. The actual point-push lifts
$C,D,E,F$ are the $\sigma$-conjugates of $c,d,e,f$ in Bellingeri's
genus-three presentation. The program
\path{code/configuration_symplectic_crosscheck.py} derives their
conjugation words by the same relations as (CF.1), squares the words,
and checks preservation of the quotient to $\operatorname{UT}_3(\F_2)$.
It computes their action on the full integral Schreier homology and
on the submodule spanned by the four differences of the $c,d,e,f$
loops at the identity and at the nontrivial central vertex.
The five actions of $C^2,D^2,E^2,F^2,(CE)^2$ are exactly the
symplectic transvections of coefficient two, including the
transvection in direction $c+e$. Exhaustive multiplication of the
resulting four-by-four matrices modulo three gives $51840$ distinct
matrices. It also checks that
\[
                    (CE)^2 F^2 D^2 D^2 C^2
\]
acts on this subspace by a matrix of exact order five.
The matrices, group order, and word are retained in the JSON
certificate and exact output alongside the program. Any finite
$\{2,3\}$ total quotient dominating $F_6/N$ would make this action
factor through a finite $\{2,3\}$ group, contradicting that order.
The normalizer argument is exactly the one in
\statementref{A.7}. This finite computation verifies an instance
of the higher-dimensional mechanism in \statementref{A.9}; the
general proof of that statement does not rely on the computation.
\end{proof}

\subsection{Direct characteristic-positive proof, including proper wild covers}

\begin{claim}{A.13}{Proved}{Relative cuspidal exactness for arbitrary prime sets containing a non-residue prime}
Let \(C\) be a smooth proper geometrically connected curve of genus at
least two over an algebraically closed field of characteristic \(p>0\).
Let \(\Sigma\) be any set of primes and choose
\(r\in\Sigma\setminus\{p\}\). In the definitions of A.3, now formed
with maximal pro-\(\Sigma\) groups even when \(p\in\Sigma\), the
relative cuspidally pro-\(r\) configuration sequence is exact. The
boundary comparison and arithmetic transport of A.5 remain valid.
More generally, the finite-detection proof works for a constant proper
compactification with disjoint labelled sections in characteristic
\(p\), without a simultaneous lifting hypothesis.
\end{claim}
\begin{proof}
We replace the specialization step by a direct theorem whose arithmetic
base group is retained in full. In the opened text of
\citep[Proposition~2.7]{LA_Stix2005}, the data are a connected locally noetherian scheme \(S\),
a smooth proper \(n\)-pointed genus-\(g\) curve over \(S\),
\(2-2g-n<0\), and a set \(L\) of primes invertible on \(S\).
Writing \(C^0\) for the complement of its sections, the asserted
sequence is
\[
 1\longrightarrow \pi_1(C^0_{\bar s})^{(L)}
  \longrightarrow \pi_1^{(L)}(C^0)
  \longrightarrow \pi_1(S)\longrightarrow1.             \tag{A.CF.12}
\]
The middle term is the quotient which pro-\(L\)-completes the fibre
kernel and retains the full group \(\pi_1(S)\). The proposition also
gives its continuous conjugation map to the automorphism group of the
fibre. Its proof constructs that map using the diagonal section after
base change by \(C^0\); the centrelessness of the hyperbolic pro-\(L\)
fibre proves injectivity. Thus the quoted statement is exactly a
relative completion assertion, and does not assert exactness after
pro-\(\Sigma\) completion of the total group.

Take any finite proper connected Galois \(\Sigma\)-cover
\(C'\to C\) with group \(A\), allowing \(p\mid\#A\).
The finite étale map is still a cover of smooth proper curves. Put
\[
 V=(C'\times C')\setminus\bigcup_{a\in A}\Gamma_a,
 \qquad S=C',\qquad L=\{r\}.
\]
Projection to the second factor has constant proper compactification
and disjoint labelled sections. Its geometric fibre has finitely
generated nonabelian free pro-\(r\) group \(P\).
The exact residue sequence for its mod-\(r\) homology has boundary
submodule generated by the labelled punctures, with their one sum
relation, and quotient \(H_1(C',\F_r)\). Both modules have trivial
base monodromy: the labels are fixed, the ground field is algebraically
closed, and the compactification is constant. Consequently the image
on \(P/\Phi(P)\) is contained in the additive group of shears from
the quotient module to the boundary module. It is a finite
\(r\)-group.

Compose \(\pi_1(V)\) with the conjugation map supplied by (A.CF.12).
The kernel of
\(\operatorname{Aut}(P)\to\operatorname{Aut}(P/\Phi(P))\)
is pro-\(r\), as proved in A.2. The image of this composite is
therefore pro-\(r\), so the homomorphism factors through
\(\pi_1(V)^{(r)}\). On the fibre it is the inner-automorphism
embedding of \(P\), which is faithful because \(P\) is centreless.
This proves the required finite-detection injection directly in
characteristic \(p\).

The remainder of A.3 only takes finite quotients and normal cores.
The resulting groups are an \(r\)-group by \(A\times A\), hence
are \(\Sigma\)-groups also when \(A\) has a \(p\)-part. All the
relative quotient maps and kernel identifications are unchanged.
For the boundary, proper covers remain étale along the diagonal,
whereas the new inertia has \(r\)-power order and is tame since
\(r\ne p\). The same diagonal root-gerbe construction and local
intersection argument from A.5 therefore apply. The arithmetic
assertion follows by retaining the entire absolute Galois quotient,
exactly as there.

For the asserted constant-compactification variant, first take a
common finite proper base cover which trivializes the pullbacks of
the chosen finite proper fibre cover along the finitely many section
maps. After this cover the lifted punctures are disjoint and labelled.
The same trivial-module residue sequence, Stix proposition, and
normal-core argument apply. No lifting of the family to characteristic
zero enters this proof.
\end{proof}


\subsection{The cuspidal criterion and its arithmetic inputs}

We retain the whole arithmetic Galois quotient throughout this subsection.

The manuscript question addressed is cote~149, scan~36: can the covering condition force the equality of the two points?

\begin{claim}{A.14}{Proved}{The same-prime quotient lemma}

Let $F$ be profinite, $K\triangleleft F$ closed, $r$ prime, and
$N=\ker(K\to K^{(r)})$. For every open $I\le F$ containing $N$, the
continuous quotient map induces a natural isomorphism
$I^{(r)}\simeq(I/N)^{(r)}$. No restriction on $[F:I]$ is imposed.

\end{claim}

\begin{proof}

Let F be a profinite group, K a closed normal subgroup, r a prime, and N the kernel of the maximal pro-r quotient \(K{\to}K^{(r)}\). Because N is characteristic in K, it is normal in F. Let \(I{\le}F\) be open and contain N. We prove that the continuous epimorphism \(I{\to}I/N\) induces an isomorphism \(I^{(r)}\to(I/N)^{(r)}\).

Take an arbitrary continuous quotient I/J which is a finite r-group, where J is open and normal in I. Put \(I_0=\operatorname{core}_F(I)\); it is open normal in F, contained in I, and contains N. Put \(J_0=J{\cap}I_0\). Then \(J_0\) is open normal in \(I_0\) and \(I_0/J_0\hookrightarrow I/J\), so \(I_0/J_0\) is an r-group. Take \(M=\operatorname{core}_F(J_0)\). There are only finitely many F-conjugates of \(J_0\), because \(J_0\) is open. Since \(I_0\) is normal in F, \(I_0/M\) embeds in the product of the quotients \(I_0/{}^fJ_0\), each isomorphic to \(I_0/J_0\). Thus \(I_0/M\) is an r-group.

The image of K in \(F/I_0\) is an r-group: \(N{\subset}I_0\), so this image is a finite quotient of K/N, a pro-r group. The image of K in F/M is an extension of that image by a subgroup of \(I_0/M\). It too is an r-group. By the universal defining property of N, N is killed by \(K{\to}F/M\). Therefore \(N\subset M\subset J_0\subset J\).

Every finite r-quotient of I consequently kills N. This is exactly the asserted isomorphism of maximal pro-r quotients, with its natural map, not merely an abstract isomorphism. No finite-generation assumption, normality of I, or r-power-index assumption on I was used.


\end{proof}

\begin{claim}{A.15}{Proved}{Compactification and Frobenius are retained at the cusp prime}

Let $C/k$ be smooth proper geometrically connected of genus at least two,
$k$ finite of characteristic $p$, $P\in C(k)$, and $r\ne p$.
Every open geometric subgroup of the maximal cuspidally pro-$r$
quotient of $\pi_1(C-P)^{(p')}$, after its cusp inertia is killed and
its maximal pro-$r$ quotient is taken, gives the full geometric
pro-$r$ fundamental group of its actual smooth compactification.
The comparison is compatible with every arithmetic action preserving
the cover. Thus a cusp-preserving isomorphism over $G_k$ preserves
the point counts of every corresponding open cover over every
finite constant extension.

\end{claim}

\begin{proof}

Take \(F=\pi_1((C-P)_{\bar k})^{(p')}\), \(K=\ker(F\to\Delta_C)\), \(r{\ne}p\). An open I/N in F/N corresponds to an actual connected finite étale cover \(D^\circ\to(C-P)_{\bar k}\) whose Galois closure has prime-to-p order. The full pro-prime-to-p topology on its geometric group is induced from F: for an open subgroup with a prime-to-p Galois core, the normal-core construction again gives a prime-to-p group, since extensions and subgroups of products of prime-to-p groups have prime-to-p order. Thus I is the full prime-to-p geometric group of \(D^{\circ}\).

Let D be its smooth compactification. Let \(T{\trianglelefteq}I\) be the closed normal subgroup generated by the inertia groups of the points in \(D-D^{\circ}\). Then \(I/T=\pi_1(D)^{(p')}\). The image of T in I/N is precisely the closed normal subgroup generated by the images of those inertia groups. Taking a maximal pro-r quotient commutes with quotient by a closed normal subgroup, by the universal property of continuous maps to finite r-groups. Therefore

\[
 ((I/N)/\overline{TN}/N)^{(r)}
 \simeq (I/T)^{(r)}
 =\pi_1(D)^{(r)}.
\]

This isomorphism is induced by the actual punctured-cover and compactification maps. Every arithmetic conjugation preserving the cover preserves its puncture inertia system and all kernels in this construction. It is therefore equivariant for the outer action of \(G_L\) when the cover descends to a finite field L. Consequently it identifies \(H^1(D_{\bar L},\mathbf Z_r)\) with its actual Frobenius action, not only its rank or characteristic polynomial modulo r.

If two such arithmetic relative groups are isomorphic over \(G_L\) and their cusp decomposition groups correspond, corresponding covers have isomorphic compactified r-adic \(H^1\) representations. Hence they have equal integral Weil polynomials and equal compactified point counts over each finite extension of L. The missing points are indexed by the cusp double cosets, and their degrees are the indices of the projected decomposition groups in \(G_L\). Thus these degrees also agree. Subtracting the missing points gives equal point counts for the open covers over every constant extension.

The point-count comparison uses the trace formula and the integral
Frobenius polynomial for a smooth proper curve over a finite field:
these are \cite[formula (1.5.1) and Theorem (1.6)]{WeilCurves}.
Their hypotheses apply to both smooth compactifications and to every
finite extension of their common constant field. Equality of their
\(r\)-adic characteristic polynomials is equality of integer polynomials,
since \(\mathbf Z[T]\hookrightarrow\mathbf Z_r[T]\).


\end{proof}

\begin{claim}{A.16}{Known, with derivation}{The exact uniform exceptional-set input}

For a proper hyperbolic curve $C/k$ over a finite field of
characteristic $p$, and $\Sigma=\mathrm{Primes}-\{p\}$, there is a
finite set $E_C\subset C^{\cl}$ such that outside its inverse image
a decomposition subgroup has a unique geometric point label, even
up to commensurability, over every finite constant extension at once.
The numbering of Definitions~1.3 and 1.5, Remark~1.2, and
Proposition~1.8(vi) is the same in the attached arXiv text and the
journal version of \cite{LA_ST2009}.

\end{claim}

\begin{proof}

Let k be finite of characteristic p, X/k a smooth proper geometrically connected curve, \(U=X-S\) hyperbolic, and \({\Sigma}\) a set of primes containing at least one prime different from p.~Let \({\Pi}_U\) be the arithmetic group retaining the whole \(G_k\) and the maximal geometric pro-\({\Sigma}\) quotient of the tame group. Its universal pro-cover compactification is denoted by \(\widetilde X_U\).

The attached ST2009 Definition 1.3(iii) calls a point a of a set map f exceptional when \(\#f^{-1}(f(a))>1\). Definition 1.5 applies this to the map \(\overline D\) from closed points upstairs to commensurability classes of decomposition subgroups. Definition 1.10 puts \(E_U\) equal to its image in \(X^{\rm cl}\).

Proposition 1.8(iii) states injectivity of \(\overline D\) over a fixed geometric point; (iv) gives finite fibres over one closed point and injectivity over one rational point; (v) says the exceptional set is \({\Pi}_U\)-stable; under the additional hypothesis that \({\Sigma}\) is cofinite, (vi) states that \(E_{\widetilde U}/\Pi_U\) is finite. This quotient is \(E_U\). The same numbering and hypotheses occur in the opened journal PDF.

Remark 1.2 says that the commensurability maps are unchanged upon passing to any open subgroup. More explicitly, if H is open in \({\Pi}_U\) and \(D_1,D_2\) are closed subgroups, then \(D_1\) and \(D_2\) are commensurate if and only if \(D_1{\cap}H\) and \(D_2{\cap}H\) are commensurate. One direction is immediate. For the other, \(D_i{\cap}H\) has finite index in \(D_i\); the common intersection has finite index in both. Universal pro-cover points are the same on replacing U by its H-cover. In particular, when H is the preimage of \(G_L\) for a finite constant extension L/k, the exceptional points upstairs are exactly the same points and the exceptional image on \(X_L\) is the finite inverse image of \(E_U\).

For \({\Sigma}=\mathrm{Primes}-\{p\}\), all hypotheses are met when C is proper of genus at least two: U=C is hyperbolic and \({\Sigma}\) is nonempty and cofinite. Consequently there is one finite set \(E_C\subset C^{\rm cl}\) whose geometric inverse image excludes all exceptional labels over all finite extensions. If R is a closed point outside \(E_C\), a decomposition subgroup at R cannot be commensurate, even after conjugation, with a subgroup at another geometric point. Conjugation moves the lift and preserves the underlying closed-point label. This is exactly the use in AS-AG.40 Step 3.

\textbf{Check of their proof and its hypotheses.}

Their proof first replaces U by an open-subgroup cover whose compactification has genus \textgreater1 and is nonhyperelliptic. Commensurability invariance just proved makes finiteness upstairs descend under this finite cover. For a nonhyperelliptic compactification X, the difference morphism \(\delta:X\times X\to J_X\), \((P,Q)\mapsto[P-Q]\), is injective on geometric points away from the diagonal. Indeed, equality of two differences gives \(P+Q'\sim P'+Q\); distinct linearly equivalent effective degree-two divisors would supply a degree-two map to the projective line, contrary to nonhyperellipticity. Equality of the divisors then gives the equal ordered pair, unless both differences vanish.

The image W of \({\delta}\) contains no translate of a positive-dimensional abelian subvariety. A two-dimensional such translate would be W itself, hence an abelian variety containing zero and generating \(J_X\). Then \(J_X\) would have dimension two, contrary to nonhyperellipticity. A one-dimensional translate would, by the quasi-finite bijection off zero, admit a finite radicial cover mapping nontrivially to X. Its normalization has genus one (purely inseparable extensions of the elliptic function field preserve genus), so it cannot map nontrivially to a genus\textgreater1 curve; after factoring off inseparability, the separable part contradicts Riemann--Hurwitz.

An exceptional pair has zero abelian Kummer difference after a finite constant extension. Therefore its nonzero difference class lies in \(J_X(\bar k)\{\Sigma'\}\), the torsion supported on the complementary primes. Here \({\Sigma}'\) is finite. The opened proof applies Boxall's finite-prime torsion intersection theorem to W and concludes that its intersection with that subgroup is finite. Injectivity of \({\delta}\) away from the diagonal then gives finitely many exceptional geometric pairs on X. For the prime-to-p choice the complementary prime set is exactly \{p\}, with no missing restriction to odd p, p-rank, ordinarity, or nonhyperellipticity of the original curve. The cover reduction is what handles genus two and hyperelliptic original curves.

We have checked the derivation and the hypotheses of the stated ST theorem. We have not independently reproved Boxall's theorem or the nonhyperelliptic-cover existence result cited within ST's proof.


\end{proof}

\begin{claim}{A.17}{Proved}{A rational point section is self-normalizing}

Let $C/L$ be smooth proper geometrically connected of genus at least
two over a finite field, and let $\Sigma$ contain a prime
$r\ne\operatorname{char}L$. If $D\subset\Pi_C^\Sigma$
is a continuous section subgroup mapping isomorphically onto $G_L$,
then $N_{\Pi_C^\Sigma}(D)=D$.

\end{claim}

\begin{proof}

The exact ST2009 Lemma 1.7 hypotheses are: \(Z{\subset}{\Pi}_U\) is closed, pr(Z) is open in \(G_k\), and \(pr|_Z\) is injective. Its conclusion is that pr embeds the commensurator in \(G_k\); commensurator, normalizer, and centralizer coincide and contain Z with finite index. If Z is the section of an L-rational point and we work in \({\Pi}_{U_L}\), its projection is the whole \(G_L\). Therefore the conclusion specializes to

\[N_{\Pi_{U_L}}(Z)=Z_{\Pi_{U_L}}(Z)=Z.\]

This does not assert self-normalization of every closed-point subgroup in the original \({\Pi}_U\) without enlarging constants. It is exactly sufficient for AS-AG.40.

For clarity, the key proof of the lemma can be re-derived. A geometric commensurator \({\sigma}\) centralizes the open intersection of Z and \({\sigma}Z{\sigma}^{-1}\), because both project injectively to the same abelian arithmetic group. If \({\sigma}{\ne}1\), choose an open characteristic \(N{\trianglelefteq}{\Delta}_U\) not containing \({\sigma}\). The open subgroup \(\bar H=\langle N,\sigma\rangle\) has a nonzero image of \({\sigma}\) in its abelianization, detected by the cyclic quotient \(\bar H/N\). The open part of Z centralizing \({\sigma}\) stabilizes \(\bar H\), so this abelian image is fixed by the Frobenius action on the corresponding genuine cover. The abelian geometric group is a product of prime-adic Tate modules of the generalized Jacobian (proper Jacobian for proper covers); none has a nonzero Frobenius-fixed element because its eigenvalues have nonzero Weil weights, and its Tate module is torsion-free. This contradiction proves that the geometric part of the commensurator is trivial. Its injection into the procyclic arithmetic group gives the rest.

The opened source used in the first argument is
\cite[Lemma 1.7]{LA_ST2009}; it has no cofiniteness assumption on \(\Sigma\).


\end{proof}

\begin{claim}{A.18}{Proved}{Point recognition in the whole cuspidal quotient}

Under the hypotheses of AS-AG.40, a cusp-preserving isomorphism
$\beta:G_P\simeq G_Q$ over $\mathrm{id}_{\Pi_C}$ carries every
noncuspidal point section, after its finite residue-field extension,
to an actual noncuspidal point section. Outside $E_C\cup\{P,Q\}$ the
target point is the same point of $C$. This conclusion holds in the
whole quotient, and not merely at a finite level.

\end{claim}

\begin{proof}

Suppose \(\beta:G_P\xrightarrow\sim G_Q\) is the isomorphism in AS-AG.40, over \(\operatorname{id}_{{\Pi}_C}\), preserving the cusp decomposition conjugacy class. For a finite extension L/k, restrict both groups to the preimage of \(G_L\); \({\beta}\) remains an isomorphism over \(\operatorname{id}_{{\Pi}_{C_L}}\) and carries the base-changed cusps. Choose \(R{\in}(C-P)(L)\) and a decomposition subgroup \(D_R^P\). Let \(S={\beta}(D_R^P)\), a closed subgroup mapping isomorphically onto \(G_L\).

Let \(\mathcal H\) be the directed set of all open subgroups \(H{\le}(G_Q)_L\) containing S, ordered by reverse inclusion. All these subgroups have full arithmetic projection, so their connected covers \(X_H\) of \(C_L-Q\) are geometrically connected over L. The source cover for \({\beta}^{-1}(H)\) has a rational point above R because it contains \(D_R^P\). By Statement \statementref{A.15}'s point-count equality, \(X_H(L)\) is a nonempty finite set.

For \(H'{\subset}H\) the covering map gives a map \(X_{H'}(L){\to}X_H(L)\). The inverse limit of these sets is nonempty even though no individual transition map is asserted surjective. To prove this, take their product, a compact space. Compatibility along each transition defines a closed condition. Any finite set of compatibility conditions is satisfied by choosing a point on the cover for the intersection of the finitely many subgroups involved and projecting it to the relevant covers. Thus the closed conditions have the finite-intersection property, and compactness supplies a compatible family \((x_H)\).

The component at \(H=(G_Q)_L\) is a point \(T{\in}(C-Q)(L)\). Choose a geometric lift \(\widetilde T\) over T in the universal pro-cover with every projection equal to \(x_H\). This is possible by the same compactness argument in the profinite fibre above T. Its decomposition subgroup \(D_{\widetilde T}\) maps isomorphically to \(G_L\): all these covers are étale at T and the residue field upstairs is \(\bar L\). For every H, rationality of \(x_H\) says that \(D_{\widetilde T}{\cap}H\) surjects onto \(G_L\). Since \(D_{\widetilde T}\) itself projects injectively, this implies \(D_{\widetilde T}{\subset}H\). Therefore

\[
D_{\widetilde T}\subset\bigcap_{H\supset S}H=S.
\]

Both sides project isomorphically onto \(G_L\), so they are equal. Thus \({\beta}\) carries a genuine noncuspidal point decomposition subgroup to the decomposition subgroup of a genuine noncuspidal point of the target, in the whole profinite quotient.

If R lies outside \(E_C\) and P,Q, its proper image \(D_R\) in \({\Pi}_{C_L}\) equals the proper image \(D_T\). The finite-exception uniqueness from Statement \statementref{A.16} forces T=R. This works for every closed R of C outside the one fixed finite set, upon taking L=k(R). It does not assume an a priori bijection of rational point sets, and does not identify one finite quotient with the inverse limit.


\end{proof}

\begin{claim}{A.19}{Proved}{The complete conditional cuspidal criterion}

Let $k$ be finite of characteristic $p$, $C/k$ smooth proper
geometrically connected of genus at least two, $P,Q\in C(k)$, and
$r\ne p$. Let $G_P,G_Q$ be the maximal cuspidally pro-$r$ quotients
of the arithmetic geometrically prime-to-$p$ groups of $C-P,C-Q$,
relative to their maps to the full proper group $\Pi_C$. If a
continuous isomorphism $\beta:G_P\to G_Q$ over
$\mathrm{id}_{\Pi_C}$ carries the cusp decomposition subgroup to
a conjugate of the cusp decomposition subgroup, then $P=Q$.

\end{claim}

\begin{proof}

For clarity about the orientation-class input, the precise finite-coefficient
facts are as follows. For a smooth proper connected curve over an
algebraically closed field and \(r\) invertible in that field,
\(H^2_{\et}(C,\mu_{r^a})\simeq\mathbf Z/r^a\) via degree, with a point
class sent to one; see \cite[Proposition 14.2 and Theorem 24.1(a)]{MilneEtale}.
The pro-\(\Sigma\) group comparison when \(r\in\Sigma\) can be verified
directly from \cite[Theorem 14.9 and Remark 14.10]{MilneEtale}.
In the tower of all connected geometric finite \(\Sigma\)-covers,
every \(H^1(-,\mathbf Z/r^a)\) class is killed by its corresponding
cyclic torsor, and every \(H^2(-,\mu_{r^a})\) class is killed by an
étale cover of degree divisible by \(r^a\). Such covers exist since
the Jacobian has nonzero \(r\)-adic Tate module. The Galois closures of
the composite covers remain \(\Sigma\)-covers, by the normal-core
argument. Cohomology above degree two vanishes. The spectral sequence
therefore identifies group cohomology with curve cohomology for these
coefficients. The same comparison over the finite ground field follows
by comparing the two Hochschild--Serre spectral sequences.
Inverse limits give the required \(\mathbf Z_r(1)\) statement.
The tame surface presentation, including its canonical inertia words,
is \cite[Expos\'{e} XIII, Corollary 2.12]{LA_SGA1}; its hypotheses are a
smooth proper connected curve over a separably closed field and a
finite set of distinct punctures, with the full prime-to-characteristic
quotient taken. Passing to any \(\Sigma\) containing \(r\) preserves the
pro-\(r\) test used here.

Let k be finite of characteristic p, C/k smooth proper geometrically connected with \(g{\ge}2\), \(P,Q{\in}C(k)\), \(r{\ne}p\). Let \(G_P,G_Q\) be the maximal cuspidally pro-r quotients relative to the full proper geometrically prime-to-p arithmetic group \({\Pi}_C\). Suppose \({\beta}:G_P{\to}G_Q\) is a continuous isomorphism over \(\operatorname{id}_{{\Pi}_C}\) carrying the cusp decomposition group to a conjugate of the cusp decomposition group. We prove P=Q.

\textbf{Principal p-primary difference.}

Projection of the cuspidal decomposition groups gives conjugate point decomposition groups at P,Q in \({\Pi}_C\). A conjugator can be made geometric: multiply it by an element of the target decomposition subgroup with the same arithmetic projection. After passing to the geometric abelianization, the Kummer difference \([P-Q]\) has zero class in \(H^1(G_k,T_\ell\) \(J_C)\) for every \(\ell{\ne}p\). Since \(J_C(k)\) is finite, its ell-adic completion is its ell-primary subgroup and the Kummer map is injective on that subgroup. Thus \([P-Q]\) has p-power order. There is \(e{\ge}0\) with \(p^e(P-Q)\) principal. A rational generator over \(\bar k\) descends to \(k(C)^{\times}\) by Hilbert 90: its Galois ratios are constants, and killing that scalar cocycle gives f with

\[\operatorname{div}(f)=p^e(P-Q).\]

\textbf{The central extensions and the multiplier.}

Let \(M=\mathbf{Z}_r(1)\). Divide the geometric cuspidal kernel in \(G_P\) by its commutators with the full geometric subgroup. The resulting arithmetic extension \(E_P{\to}{\Pi}_C\) is the maximal geometrically cuspidally central quotient. Its kernel is procyclic, generated by the puncture inertia, because the original cuspidal kernel is normally generated by the conjugate puncture inertia groups. It is exactly M, not a finite quotient: in the geometric punctured surface presentation the puncture word is the product of the g handle commutators. Map the first pair of handles to \(1+e_{12},1+e_{23}\) in the Heisenberg pro-r group \(\operatorname{UT}_3(\mathbf{Z}_r)\) and all other handles to 1. The puncture word maps to \(1+e_{13}\), of infinite r-adic order, and its normal closure is central in that target. This proves that the central inertia generator has infinite r-adic order. The arithmetic action is the cyclotomic action, as on tame inertia. The same construction gives \(E_Q\).

The class of the geometric central extension is the degree-one orientation class in \(H^2(\Delta_C,M)\simeq Z_r\), with the sign fixed by the Kummer residue convention. One equivalent construction is the Kummer first Chern class of \(O_C(P)\); its canonical trivialization on \(C-P\) gives the homomorphism from the punctured group to the extension and maps local inertia with residue +1. Its geometric restriction is the class of a point, hence a generator. This is the familiar punctured-surface orientation extension; no configuration-group injection enters.

Because \({\beta}\) preserves the proper quotient and the geometric subgroup, it descends to an isomorphism \(\bar\beta:E_P\to E_Q\) over \(\operatorname{id}_{{\Pi}_C}\). On M it acts by a unit \(u{\in}\mathbf{Z}_r^{\times}\). Restricting the extension isomorphism to the geometric group gives u times the degree-one class equal to the same degree-one class. Thus u=1.

\textbf{The difference cocycle and its residues.}

On the continuous fibre product \(E_P{\times}_{{\Pi}_C}E_Q\) define

\[\chi(u,v)=\bar\beta(u)v^{-1}\in M.\]

The common arithmetic projection acts on M cyclotomically. Multiplication verifies \(\chi(ee')=\chi(e)+e\cdot\chi(e')\); therefore \({\chi}\) is a continuous crossed homomorphism. Pull back to the arithmetic geometrically prime-to-p group \({\Gamma}\) of \(C-\{P,Q\}\), using its natural maps to both central extensions. On inertia at P the second map is trivial and the first has residue +1; on inertia at Q the first map is trivial and the inverse of the second gives residue \(-1\). Thus \({\chi}\) has residues \((+1,-1)\).

The Kummer class \({\kappa}_r(f)\) has residues \((p^e,-p^e)\). Hence \({\kappa}_r(f)-p^e{\chi}\) has zero restriction to both inertia groups. Since M is trivial under the geometric group, a cocycle vanishing on these inertia groups vanishes on their closed normal closure. The kernel of \({\Gamma}{\to}{\Pi}_C\) is precisely that closure. Thus this difference descends to \(H^1({\Pi}_C,M)\).

Inflation-restriction gives an isomorphism

\[H^1(G_k,M)\xrightarrow\sim H^1(\Pi_C,M),\]

because \(H^1(\Delta_C,M)^{G_k}=0\). Explicitly \(H^1(\Delta_C,M)\) is \(\operatorname{Hom}(T_rJ_C,\mathbf{Z}_r(1))\); the arithmetic Frobenius eigenvalues on its \(\mathbf{Q}_r\)-extension have complex absolute value \(\sqrt{\#k}\), so none is 1, and the \(\mathbf{Z}_r\)-module is torsion-free. Finite-field Kummer theory identifies \(H^1(G_k,M)\) with the r-primary completion of \(k^{\times}\). Therefore there is \(a{\in}k^{\times}\) representing \({\kappa}_r(f)-p^e{\chi}\).

\textbf{Evaluation at all good closed points.}

Take any closed point R outside \(E_C{\cup}\{P,Q\}\), and put L=k(R). Statement \statementref{A.18} gives a conjugacy in \((G_Q)_L\) between \({\beta}(D_R^P)\) and \(D_R^Q\). The projection of a conjugator to \({\Pi}_{C_L}\) normalizes \(D_R\). Statement \statementref{A.17} says this normalizer is exactly \(D_R\). Multiply the conjugator by the inverse of the element of \(D_R^Q\) with that projection. This leaves its conjugation of the subgroup unchanged and makes it belong to the cuspidal kernel.

In \(E_Q\) its image therefore lies in M. Conjugating a section by an element of M changes its M-valued cocycle by a coboundary. It follows that \({\chi}\) restricts to zero in \(H^1(G_L,M)\) along the genuine section of R. Evaluation of the preceding Kummer identity gives

\[\kappa_{r,L}(f(R)/a)=0.\]

The point R is allowed to have arbitrary degree. The field L is chosen separately for each R, and \({\beta}\) and the extension comparison are restricted along \(G_L{\to}G_k\); no incompatible Frobenius convention is introduced.

\textbf{Infinitely many usable points and contradiction.}

Assume \(P{\ne}Q\). Then f is nonconstant. Over \(\bar k\) it gives a surjective finite map to the projective line. The excluded closed points have finitely many geometric inverse images, and their images under f form a finite set. Because \(r{\ne}p\), \(\mu_{r^\infty}(\bar k)\) is infinite. Choose a nontrivial \({\zeta}\) in that group with \(a{\zeta}\) outside the finite set of excluded values. Every geometric point over \(a{\zeta}\) is defined over a finite field. Choose one and regard it as a closed point R of C. Then R avoids all exclusions, and \(f(R)/a={\zeta}\) belongs to \(k(R)^{\times}\).

The finite-field Kummer map to \(H^1(G_{k(R)},\mathbf{Z}_r(1))\) is injective on its r-primary subgroup. Consequently the nontrivial r-power root \({\zeta}\) has nonzero Kummer class, contrary to the evaluation identity. This proves P=Q.

This last argument does not require choosing all roots of unity in the original k, and a root becoming an r-th power in some finite extension does not make its entire r-adic Kummer class vanish. The completion of a finite cyclic group preserves its r-primary subgroup.


\end{proof}

\begin{claim}{A.20}{Proved}{Proper-field support suffices for the criterion}

Let $\Sigma$ be any set of primes containing $r\ne p$.
Use the full proper geometric pro-$\Sigma$ group in the preceding
criterion. It suffices that the noncuspidal point transport induced
by $\beta$ have unchanged labels outside $C(F)$, for one proper
algebraic subfield $k\subset F\subsetneq\bar k$. In particular,
an exceptional image supported over such an $F$ suffices, even if it
is infinite. The annihilating integer for $[P-Q]$ need not be
coprime to $r$.

\end{claim}

\begin{proof}

This strengthening also applies to the \(J_C\)-large prime sets of \cite{LA_ST2017}.

Fix any prime set \({\Sigma}\) containing a prime \(r\ne p\), and use the full proper geometric pro-\({\Sigma}\) group in the definitions of \(G_P,G_Q\). The same-prime lemma, compactification comparison, central extension construction, and self-normalizer argument remain valid. Assume \({\beta}\) is a cusp-preserving isomorphism over \(\operatorname{id}_{{\Pi}_C^{\Sigma}}\). Suppose that there is a proper algebraic subfield \(F{\subsetneq}\bar k\) containing k such that, for every closed R whose geometric points are not F-rational, the point transported by \({\beta}\) has the same label R. Then P=Q.

Only the final Kummer argument changes. Let n be the exact order of \([P-Q]\) in \(J_C(k)\), and choose \(f{\in}k(C)^{\times}\) with \(\operatorname{div}(f)=n(P-Q)\). This order exists because \(J_C(k)\) is finite. The argument for the central extensions uses neither that n is p-primary nor that n is coprime to r: the difference \({\kappa}_r(f)-n{\chi}\) has zero residues, hence is a constant class \({\kappa}_r(a)\), and \({\chi}\) evaluates to zero at every point with unchanged label. Thus \({\kappa}_r(f(R)/a)=0\) for every such R.

The function f/a is not an r-th power in \(\bar k(C)^{\times}\). If \(r\nmid n\) its valuations at P,Q already rule this out. If \(r\mid n\) and \(f/a=b^r\), then \(\operatorname{div}(b)=(n/r)(P-Q)\), contradicting the exact order of \([P-Q]\). The geometric divisor class has the same order as its k-rational class, since \(J_C(k){\to}J_C(\bar k)\) is injective. Equivalently, a rational principal divisor descends by Hilbert 90.

Let D be the smooth compactification of the Kummer cover \(z^r=f/a\). It is geometrically connected, since r is prime and f/a is not a geometric r-th power; \(r{\ne}p\) makes it generically separable. Take a finite L/k containing \({\mu}_r\). Because F is a proper subfield of \(\bar k\), some prime ell has finite exponent in the supernatural degree {[}F:k{]}. To see this directly, the nonzero closed subgroup \(G_F{\subset}\widehat{\mathbf Z}\) has a nonzero ell-adic projection for some ell, and a nonzero closed subgroup of \(\mathbf{Z}_\ell\) is open. Put \(L_m\) equal to the degree-\(\ell^m\) extension of L. Then

\[F_0=F\cap\bigcup_m L_m\]

is a finite extension of k: the degrees outside ell are bounded by {[}L:k{]}, and the ell-part of {[}F:k{]} is bounded. Thus all excluded F-rational points lying in \(C(L_m)\) belong to the fixed finite set \(C(F_0)\). Write \(B=\#C(F_0)+2\), including P,Q.

For every remaining \(R{\in}C(L_m)\), the zero r-adic Kummer class implies that f(R)/a is an r-th power in \(L_m\). Since \({\mu}_r{\subset}L_m\) and f has no zero or pole away from P,Q, the fibre of D above R consists of exactly r \(L_m\)-rational points. Hence

\[\#D(L_m)\ge r\bigl(\#C(L_m)-B\bigr).\]

Writing \(Q_m=\#L_m\) and applying the Weil bounds to the geometrically connected proper curves C and D gives

\[(r-1)(Q_m+1)\le (2g_D+2r g_C)\sqrt{Q_m}+rB.\]

This is impossible as \(m{\to}{\infty}\). Thus P=Q. The auxiliary ell is allowed to equal p or r; the extension is an extension of finite constants and no ramification issue occurs. Only the Kummer degree r must be different from p.

In particular, if \(E_C^{\Sigma}\) is supported over a proper algebraic subfield, the needed same-label assertion follows from compactness recognition and the definition of E. Thus, once the generalized configuration transport is established, almost \({\Sigma}\)-separated in the sense of ST2017 implies \({\Sigma}\)-separated, i.e.~\(E_C^{\Sigma}=\varnothing\). This conclusion uses no finiteness assertion about \(E_C\) itself.

The finite-exception case is a special case: finitely many closed-point residue fields generate one finite F, which is proper in \(\bar k\). In fact the original roots-of-unity proof also works directly with arbitrary n, so no coprimality condition was hidden there.


\end{proof}

\begin{claim}{A.21}{Proved}{An independent elementary support-field proof}

For a proper hyperbolic curve $C/k$ over a finite field of
characteristic $p$, put
\[
B_C=\{R\in C(\bar k):\exists S\ne R,
       [R-S]\in J_C(\bar k)[p^\infty]\}.
\]
There is a proper algebraic subfield $F\subsetneq\bar k$ containing
$k$ for which $B_C\subset C(F)$. Hence the conditional criterion
\statementref{A.19} admits a proof using
\statementref{A.20} and this support field in place of the
exceptional-set finiteness theorem.

\end{claim}

\begin{proof}

We can remove the Boxall finiteness input entirely from an alternative proof of the main point-separation theorem. Define

\[B_C=\{R\in C(\bar k):\exists S\ne R,\ [R-S]\in J_C(\bar k)[p^\infty]\}.\]

We prove that \(B_C\) is supported over one proper algebraic subfield of \(\bar k\). This is a statement about divisor classes, not a replacement of group collisions by a claimed converse.

\textbf{The fibres of the difference map.}

For every nonzero \(t{\in}J_C(\bar k)\), the set

\[\{(R,S)\in C(\bar k)^2:R\ne S,[R-S]=t\}\]

has at most two elements. Indeed, for two such pairs, R+S' is linearly equivalent to R'+S. If these effective degree-two divisors are equal, the nonzero-difference condition forces the ordered pairs to be equal. Otherwise they define a degree-two pencil on C. There is at most one such pencil when \(g{\ge}2\): two distinct degree-two maps would give a birational map of C to a curve of bidegree (2,2) in \(P^1{\times}P^1\), whose arithmetic genus is one, impossible. The degree-two pencil is separable even in characteristic two, because a purely inseparable degree-two map to \(P^1\) would identify C with a Frobenius twist of a rational curve, again impossible in positive genus. Denote its involution by iota. Equality of the two fibre divisors then forces (R',S')=(iota(S),iota(R)). This is the sole possible second pair.

\textbf{A nontrivial arithmetic subgroup fixing all p-primary torsion.}

Let f be the p-rank of \(J_C\). Then \(T_pJ_C\cong Z_p^f\) and its finite quotients identify with the geometric p-power torsion groups. Choose an odd prime \(\ell{\ne}p\) not dividing \(|\operatorname{GL}_f(\mathbf{F}_p)|\), with \(|\operatorname{GL}_0(\mathbf{F}_p)|=1\). The kernel of \(\operatorname{GL}_f(\mathbf{Z}_p){\to}\operatorname{GL}_f(\mathbf{F}_p)\) is pro-p, as its congruence filtration has elementary abelian p-group successive quotients. The ell-adic factor \(H=\mathbf{Z}_\ell{\subset}G_k\) has trivial image in \(\operatorname{GL}_f(\mathbf{F}_p)\), since the latter order is prime to ell; its image in the pro-p congruence kernel is also trivial. Thus H fixes all of \(J_C(\bar k)[p^{\infty}]\).

If \(t=[R-S]\) is a nonzero p-primary difference, H acts on the at-most-two-element fibre over t. A pro-ell group with ell odd cannot act nontrivially on a set of cardinality at most two. Hence H fixes R and S individually. Therefore

\[B_C\subset C(F),\qquad F=\bar k^H\subsetneq\bar k.\]

If f=0, the p-primary torsion group itself is zero, so \(B_C\) is empty because the Abel map of a positive-genus curve is injective.

\textbf{Application to the conditional criterion.}

From the compactness proof, \({\beta}\) carries the point R to an actual target point S over the same finite field. Because \({\beta}\) is over \(\operatorname{id}_{{\Pi}_C}\), their full proper point sections have equal geometric conjugacy class. Passing to the Jacobian Kummer map proves \([R-S]\) is p-primary. This is only the necessary implication from group collision to divisor-class torsion, which has already been proved in the Kummer step of \statementref{A.19}. Therefore R outside \(B_C\) must have S=R.

Statement \statementref{A.20} now applies with the explicit proper field F above and proves P=Q. This independently establishes the conditional cuspidal criterion without ST2009's exceptional-set finiteness, Boxall's theorem, or an assumption that a p-primary divisor difference produces a group collision. It retains only the exact cusp isomorphism as input, the full same-prime compactification lemma, the compactness recognition, the elementary degree-two-pencil bound, standard curve cohomology, finite-field Kummer theory, and the Weil bound.

Consequently, after the independent verification of AG.41 and AG.42, this is an alternative complete route to the main theorem. The original K1 audit remains useful for checking the earlier proof exactly as written, but it is not an indispensable external premise of this replacement proof.


\end{proof}

\begin{claim}{A.22}{Proved}{Exceptional emptiness and all-extension point separation}

For a proper hyperbolic $C/k$ over a finite field of characteristic
$p$ and every prime set $\Sigma$ containing some $r\ne p$,
$E_C^\Sigma=\varnothing$ is equivalent to injectivity of the
geometric pro-$\Sigma$ point-section map after every finite
constant extension of $k$.

\end{claim}

\begin{proof}

For every prime set \({\Sigma}\) containing some prime \(r\ne p\), let \({\Pi}_C^{\Sigma}\) retain the full arithmetic quotient. Assume C is proper hyperbolic over finite k. Then \(E_C^{\Sigma}\) is empty if and only if the point-section map is injective after every finite constant extension.

If distinct L-rational points have geometrically conjugate sections, choose their lifts so that the subgroups are equal. The universal pro-cover for \(C_L\) is the same universal pro-cover with the open arithmetic group \({\Pi}_{C_L}^{\Sigma}\). Equal decomposition groups are commensurate. The lifts remain distinct because their geometric images are distinct. Thus the original exceptional set is nonempty.

Conversely, take two distinct upstairs points with commensurate decomposition groups \(D_1,D_2\). Proposition 1.8(iii), whose proof only needs Lemma 1.7 and no cofiniteness, shows that their geometric images are distinct. The intersection \(D_1{\cap}D_2\) is open in both. Its arithmetic projection is open in \(G_k\) because the \(D_i\) project injectively onto open subgroups. Choose a finite L/k containing both point residue fields and with \(G_L\) contained in that projection. Since each \(D_i\) projects injectively, their restrictions over \(G_L\) are precisely the same inverse image of \(G_L\) in \(D_1{\cap}D_2\). They give equal sections of distinct L-points. This is a collision in the whole \({\Pi}_{C_L}^{\Sigma}\).

ST2017 Proposition 2.11(i), read in full, states \(E_{C'}=f^{-1}(E_C)\) for any étale \({\Sigma}\)-covering f, including constant extensions. This is compatible with the equivalence. Its ramified statement is an inclusion with the ramification locus, so must not be promoted to equality. Proposition 2.12, under \(J_C\)-largeness, only confines exceptional points to a proper algebraic subfield, finite of degree at most two over the \({\Sigma}'\)-torsion field (degree one for nonhyperelliptic C); it does not exhibit a nonempty exceptional set or contradict emptiness for the full prime-to-p set.


\end{proof}

\begin{claim}{A.23}{Proved}{Large prime sets and the finite-or-infinite dichotomy}

Assume the cuspidal transport construction from equal point sections
has been established for a prime set $\Sigma$ containing some
$r\ne p$, allowing $p\in\Sigma$. Then
$E_C^\Sigma=\varnothing$ whenever $\Sigma$ is $J_C$-large in the
sense of \cite{LA_ST2017}; this includes every cofinite $\Sigma$.
More generally, almost $\Sigma$-separation implies
$E_C^\Sigma=\varnothing$, and a finite $E_C^\Sigma$ is necessarily
empty. Thus the exceptional image is either empty or infinite.
This dichotomy alone does not classify every prime set.

\end{claim}

\begin{proof}

For any prime set \({\Sigma}\) containing some \(r\ne p\), put

\[T_{\Sigma'}J_C=\prod_{s\in\Sigma'}T_sJ_C,\qquad H=\ker(G_k\to\operatorname{Aut}(T_{\Sigma'}J_C)).\]

Here the p-adic Tate module means the étale p-adic module, of rank equal to the p-rank. Under the ST2017 definition, \({\Sigma}\) is \(J_C\)-large exactly when H is nontrivial. Every \({\Sigma}'\)-primary geometric torsion point is fixed by H. The fibre argument of the difference-map argument in \statementref{A.21} shows that 2H fixes each ordered pair realizing a nonzero \({\Sigma}'\)-primary difference. Since multiplication by two is injective on the closed subgroup \(H{\subset}\widehat{\mathbf Z}\), 2H is still nonzero. Thus

\[B_C^\Sigma\subset C(\bar k^{2H}),\qquad \bar k^{2H}\ne\bar k.\]

This recovers and makes explicit the support mechanism of ST2017 Proposition 2.12: for a nonhyperelliptic curve the fibre is a singleton, so H itself suffices; for a hyperelliptic curve one may need the index-at-most-two subgroup 2H.

The Kummer assertion used here includes the characteristic prime when
$p\in\Sigma$. Over a finite field $L$, multiplication by $p^n$ is
surjective on $J_C(\bar L)$, with kernel $J_C[p^n](\bar L)$.
It gives an injective map
\[
 J_C(L)/p^nJ_C(L)\longrightarrow
 H^1(G_L,J_C[p^n](\bar L)).
\]
These are the same classes as those read by the abelian geometric
pro-$p$ fundamental group. Indeed, factor $[p^n]$ through the
quotient by the connected component $J_C[p^n]^0$: the first
isogeny is radicial and bijective on geometric points, and the
remaining isogeny is finite \emph{\'etale}, with geometric kernel
$J_C[p^n](\bar L)$. Its lifting cocycle therefore equals the
multiplication-$p^n$ Kummer cocycle. The Abel--Jacobi map identifies
the abelian curve group with the Jacobian group at this prime as
well, by \cite[Proposition 9.1]{MilneJac}. Passing to the inverse
limit is legitimate for these finite groups and detects the entire
$p$-primary part of the finite group $J_C(L)$. Thus equality of
point sections kills the components at every prime in $\Sigma$,
including $p$ when it is retained; the remaining difference has
precisely complementary-prime torsion. No ordinary-Jacobian
hypothesis is required.

In the generalized criterion, compactness carries each R to S with \([R-S]\) \({\Sigma}'\)-primary by Kummer theory at all primes in \({\Sigma}\). Hence outside \(B_C^{\Sigma}\) the labels agree. Statement \statementref{A.20} then proves the generalized conditional criterion for every \(J_C\)-large \({\Sigma}\). If the configuration transport is verified for every such \({\Sigma}\), it follows that \(E_C^{\Sigma}=\varnothing\) for every \(J_C\)-large \({\Sigma}\) and every proper hyperbolic C over a finite field. Every cofinite \({\Sigma}\) is \(J_C\)-large: choose an auxiliary odd ell outside the finite set \({\Sigma}'\) and outside the prime divisors of all finite groups \(\operatorname{GL}_{\operatorname{rank}T_s}(\mathbf F_s)\), \(s{\in}{\Sigma}'\); then \(\mathbf{Z}_\ell\) acts trivially on the complementary Tate module.

Even without \(J_C\)-largeness, finiteness of \(E_C^{\Sigma}\) supplies finite-field support and Statement \statementref{A.20} applies. Therefore, again under the generalized transport theorem, a finite exceptional set must be empty. Every \(E_C^{\Sigma}\) is consequently either empty or infinite. This is a dichotomy, not a classification of which alternative occurs for every arbitrary \({\Sigma}\). The singleton theorem in Statement \statementref{A.26} supplies explicit infinite cases; the \(J_C\)-large theorem supplies empty cases. We have not proved that these exhaust all possible prime sets.


\end{proof}

\begin{claim}{A.24}{Proved}{Genus-one negative control}

The elliptic curve $E/\mathbf F_2:y^2+xy=x^3+1$ has distinct
rational points $O$ and $T=(0,1)$ whose sections collide in the
entire geometrically prime-to-two arithmetic group. Its exceptional
image is all of $E^{\cl}$. The finite-exception hypothesis used for
proper hyperbolic curves fails here.

\end{claim}

\begin{proof}

Take \(E/\mathbf{F}_2\) given by \(y^2+xy=x^3+1\), with origin O at infinity. Its affine \(\mathbf{F}_2\)-points are (0,1),(1,0),(1,1), so it has four rational points. It is nonsingular: a potential affine singularity must have x=0 from \(\mathbf{F}_y=x\), but at (0,1) the derivative \(\mathbf{F}_x=y+x^2\) is 1; the standard projective infinity point is also nonsingular. The inverse of (x,y) is (x,y+x). Hence T=(0,1) is a nonzero point of exact order two.

The entire geometric fundamental group of an elliptic curve is abelian. Its prime-to-two quotient is the product of the odd Tate modules. Using O as splitting, the class of the point section of T is exactly its Kummer class in that product. The two-primary point T is divisible by every odd integer in the finite group \(E(\mathbf{F}_2)\); consequently every odd-adic Kummer component vanishes. Thus O,T collide in the entire geometrically prime-to-two group, not just in an abelianized approximation of a larger nonabelian group.

For every \(x\in E(\bar F_2)\), the same argument over a finite field containing x gives a collision between x and x+T. These are distinct. Thus the exceptional image is all of \(E^{cl}\); it is infinite. The exact failed AS-AG.40 hypothesis is the finite-exception statement in Step 3. ST2009 Proposition 1.8 assumes hyperbolicity of U; the proper genus-one curve is not hyperbolic. The rest of the Kummer argument does not repair that missing input.


\end{proof}

\begin{claim}{A.25}{Proved}{The abelianized negative control in genus two}

The smooth proper genus-two curve
$C/\mathbf F_2:y^2+x(x+1)y=x^5+1$ has distinct rational points
$P=(0,1)$ and $Q=(1,0)$ with $[P-Q]$ of exact order two. They
collide in the entire arithmetic extension by the geometric
abelian prime-to-two quotient. Both the cusp-residue and the
compactified-cover cohomology arguments cease to apply to that
abelianized replacement.

\end{claim}

\begin{proof}

Take \(C/\mathbf{F}_2\) given by

\[y^2+x(x+1)y=x^5+1.\]

At finite ramification points x=0,1 there are respectively P=(0,1),Q=(1,0), and \(\mathbf{F}_x=y+x^4\) is 1 at both. The infinity chart \(u=1/x,v=y/x^3\) is

\[v^2+(u+u^2)v=u+u^6,\]

whose derivative in u is 1 at (0,0). This is a smooth genus-two hyperelliptic curve, with three simple Artin--Schreier branch points \((0,1,{\infty})\). For an elementary genus check, put z=y/{[}x(x+1){]} in the function field. The right-hand Artin--Schreier function is \((x^5+1)/(x^2(x+1)^2)\). At 0, subtract the Artin--Schreier coboundary \((1/x)^2+1/x\); the remaining pole has order one. At 1 the original pole has order one, and at infinity it has order one. The degree-two separable map therefore has different contribution two at each of three points, giving \(2g-2=-4+6=2\).

The rational function x/(x+1) has divisor \(2P-2Q\). The class \([P-Q]\) is nonzero, because a function with divisor \(P-Q\) would give a degree-one map to the projective line, impossible in positive genus. Hence its exact order is two. Both point sections become equal up to geometric conjugation in the entire extension by \(\Delta_C^{\rm ab,(2')}\), by finite-field Kummer theory. The program gives \(N_1=N_2=3\) and \(\#J_C(\mathbf{F}_2)=4\), consistently.

The full AS-AG.40 input does not survive replacing all groups by their geometric abelianizations. At a once-punctured curve the sole puncture inertia is zero in the geometric abelianization, since the puncture word is a product of commutators. Thus there is no central cusp kernel \(\mathbf{Z}_r(1)\) in that abelianized punctured group and no cocycle with residue +1 at that cusp. Moreover, even the cohomology-retention statement needed for point recognition fails for general open subgroups of the proper abelian quotient: an index-r geometric subgroup has r-adic abelian rank 2g, while its actual degree-r étale curve cover has genus \(1+r(g-1)\), hence r-adic rank \(2+2r(g-1)>2g\). Therefore the relevant failure occurs before using finite exceptions or Kummer evaluations; AG.39.15 must be applied to the full punctured pro-prime-to-p group and its actual cuspidal kernel, not to an abelianized replacement.


\end{proof}

\begin{claim}{A.26}{Proved}{Empty and singleton prime sets have infinitely many exceptions}

For every proper hyperbolic curve $C/k$ over a finite field and every
prime $r$, including $r=\operatorname{char}k$, the exceptional
image $E_C^{\{r\}}$ is infinite. Here the full arithmetic quotient
$G_k$ is retained, and for $r=\operatorname{char}k$ we use the
same definition by commensurable point decomposition groups.
For $\Sigma=\varnothing$, the group is $\Pi_C^\varnothing=G_k$;
all rational point sections are equal and every closed point is
exceptional. Thus retaining at most one geometric prime does not
separate points over every finite constant extension.

\end{claim}

\begin{proof}

\textbf{Elementary automorphism lemma.} Let D be a finitely generated pro-r group, with Frattini quotient \(D/\Phi(D)\cong F_r^d\). Then the kernel of \(\operatorname{Aut}(D){\to}\operatorname{GL}_d(\mathbf{F}_r)\) is pro-r.

Here is a proof avoiding an unopened automorphism theorem. For a finite r-group H, every maximal subgroup has index r, so a tuple whose images generate \(H/{\Phi}(H)\) generates H. Fix one basis of \(H/{\Phi}(H)\), and consider the set of all ordered lifts of that basis to H. It has cardinality \(|{\Phi}(H)|^d\), an r-power. The kernel of \(\operatorname{Aut}(H){\to}\operatorname{Aut}(H/{\Phi}(H))\) acts freely on this set: fixing a generating tuple fixes H pointwise. Therefore the kernel's order divides this r-power, so it is an r-group. A finitely generated pro-r group has a cofinal system of open characteristic subgroups (intersect all subgroups of bounded index). Its finite quotients are r-groups; the image of \({\Phi}(D)\) in each is the finite quotient's Frattini subgroup. Passing to the inverse limit of the preceding finite kernels proves the assertion.

Now assume C/k proper hyperbolic has a rational point O, write \({\Delta}_r\) for its geometric pro-r group, and use \(s_O\) to express \({\Pi}_C^{\{r\}}={\Delta}_r{\rtimes}G_k\). Its Frattini rank is \(d=2g\) when \(r\ne p\), and \(d=f_C\), the geometric p-rank of the Jacobian, when \(r=p\). The latter equality follows from the Artin--Schreier comparison and the all-prime Abel--Jacobi comparison in \cite[Proposition 9.1 and Lemma 9.4]{MilneJac}. In particular it is finite. For a pro-r group, finitely many lifts of a basis of its Frattini quotient generate the group: otherwise their closed generated subgroup lies in a maximal open subgroup of index r, contradicting that they generate the quotient. Thus the automorphism lemma applies at the characteristic prime as well. If \(f_C=0\), the pro-p group itself is trivial. Choose a prime q different from r and not dividing \(|\operatorname{GL}_d(\mathbf{F}_r)|\), with \(\operatorname{GL}_0\) interpreted as the trivial group. The q-adic direct factor \(Z_q\subset\widehat Z=G_k\) has trivial image under \(G_k{\to}\operatorname{Aut}({\Delta}_r)\): its image in \(\operatorname{GL}_d(\mathbf{F}_r)\) is a q-group, hence trivial, and its remaining image in the pro-r kernel is both pro-q and pro-r, hence trivial. Consequently \(s_O(\mathbf{Z}_q)\) is central in the entire arithmetic group.

For every closed point x of degree divisible by q, choose \(h{\in}\mathbf{Z}_q\) with q-adic coordinate 1. This h acts nontrivially on a geometric point over x, because its image modulo deg(x) is nonzero. Let \(z=s_O(h)\), and choose any lift \(\tilde{x}\) in the universal pro-cover. Then \(z\tilde{x}\) differs from \(\tilde{x}\) already on the geometric curve, whereas

\[D_{z\tilde{x}}=zD_{\tilde{x}}z^{-1}=D_{\tilde{x}}.\]

Thus x is exceptional. There are infinitely many such closed points: for sufficiently large j the Weil bound implies \(\#C(k_{q^j})>\#C(k_{q^{j-1}})\), so there is a point of exact degree \(q^j\). The equality of the displayed decomposition groups, restricted to a finite field defining both geometric points, gives equal whole section subgroups for two distinct rational points. No finite-level approximation is involved.

If C(k) is empty, first take a finite constant extension with a point. Equality of two decomposition groups in this open arithmetic subgroup makes their decomposition groups in the original arithmetic group commensurable: each original group contains that restricted group with finite index. The geometric points are still distinct. Infinitely many closed points over the extension map to infinitely many closed points over k because the extension is finite. This proves descent of infiniteness directly, without invoking ST2009's additional prime-set hypotheses. Therefore \(E_C^{\{r\}}\) is infinite for every proper hyperbolic C over every finite field and every prime r.

When \(\Sigma=\varnothing\), the geometric quotient is trivial. A geometric point over a closed point of degree n has decomposition group \(G_{k_n}\subset G_k\). Any two such open subgroups of the procyclic group are commensurable. Since the curve has infinitely many geometric points, every closed point is exceptional, and all rational points over any fixed finite extension give the same identity section of that extension's whole arithmetic group.

This proof gives actual all-level collisions after finite extensions. It also shows exactly why ST2009's cofiniteness hypothesis cannot be omitted from AS-AG.40's finite-exception argument. We make no inference that arbitrary finite sets with more than one prime have virtually pro-\({\Sigma}\) automorphism groups; that would require a separate theorem and is not supplied here.


\end{proof}

\begin{claim}{A.27}{Proved}{A completely specified pro-three collision}

Let $L=\mathbf F_2[t]/(t^7+t+1)$ and $C/L:y^2+y=x^5$.
The distinct $L$-points
$P=(t,t^6)$ and $Q=(t^2,t^6+t^5)$ have conjugate sections in
the entire arithmetic extension by the geometric pro-three group.
The assertion is not merely a finite-quotient collision.

\end{claim}

\begin{proof}

Let \(L=F_2[t]/(t^7+t+1)\), and take \(C/L:y^2+y=x^5\) with \(O={\infty}\). It is smooth of genus two. The program verifies the field identity for every nonzero element of the 128-element ring, so irreducibility of the chosen modulus is checked rather than assumed. Put

\[P=(t,t^6),\qquad Q=(t^2,t^6+t^5).\]

They satisfy the equation and Q is the \(\mathbf{F}_2\)-Frobenius image of P. The orbit has seven distinct points; the code gives their complete coordinates. Work in the geometrically pro-three quotient over \(\mathbf{F}_2\). We have \(|GL_4(F_3)|=24\,261\,120\), which is not divisible by seven. The preceding argument gives a central image of \(\mathbf{Z}_7{\subset}G_{\mathbf{F}_2}\). Its element with 7-adic coordinate 1 has image 1 in \(\operatorname{Gal}(L/\mathbf{F}_2)\), so sends a lift of P to a lift of Q while leaving the decomposition group unchanged. Intersect that common decomposition subgroup with the preimage of \(G_L\). This gives equal section subgroups for the two L-rational points in the entire geometrically pro-three arithmetic group over L.

The program finds \(\#C(L)=129\), of which 126 points have exact degree seven over \(\mathbf{F}_2\), making 18 exceptional degree-seven closed points. Infinitely many other degrees are handled by the preceding theorem; no bounded computation is being used to prove infiniteness.


\end{proof}

\begin{claim}{A.28}{Open}{The remaining arbitrary-prime-set classification}

The assertions above prove emptiness for $J_C$-large sets (once
the stated transport theorem is inserted), infiniteness for
singletons, and a finite-or-infinite dichotomy. They do not decide
the alternative for every arbitrary nonempty prime set that is
neither a singleton nor $J_C$-large. That additional classification
asked in A12 remains unproved in this audit.

\end{claim}

\paragraph{State of the unresolved assertion.}

The remaining obstruction is exact: the Kummer comparison confines a
possible exceptional pair to a fibre over complementary-prime torsion,
but when the complementary Tate representation is faithful its kernel
gives no proper field supporting those pairs. The single-prime central
arithmetic subgroup proof cannot be applied to a multi-prime completion:
its automorphism group need not have a pro-\(\Sigma\) kernel on a finite
Frattini quotient. Neither assertion is silently assumed. The next
approach would be to classify the kernel of the full outer Frobenius
action for these prime sets, or to establish sparse support of the
possible point-pair differences along a constant-field tower.



\subsection{A concrete obstruction to the finite-Frattini extension}

\begin{claim}{A.29}{Proved}{Mixed-prime surface groups have infinite Frattini quotient}
Let $\Sigma$ contain two distinct primes $r,s$, and let $\Delta$
be the pro-$\Sigma$ completion of a surface group of genus at least
two. Its Frattini quotient $\Delta/\Phi(\Delta)$ is infinite,
where $\Phi$ is the intersection of the maximal open subgroups.
In particular this applies to the geometric pro-$\Sigma$ group
of a proper curve in characteristic outside $\Sigma$.

For $\Sigma=\{2,3\}$ there is, moreover, an explicit two-generated
group $G$ of order $576$, with $\Phi(G)=1$, and an automorphism of
order seven acting trivially on $G^{\mathrm{ab}}$.
Neither assertion is a counterexample to mixed-prime point separation.
They show why the finite-Frattini and abelianization arguments used
for a single prime do not extend by the same reasoning.
\end{claim}
\begin{proof}
For each positive integer $n$, put $d_n=\operatorname{ord}_{s^n}(r)$.
Choose an element $\zeta_n$ of exact order $s^n$ in
$\mathbf F_{r^{d_n}}^\times$, and form
\[
 G_n=(\mathbf F_{r^{d_n}},+)\rtimes\langle\zeta_n\rangle,
 \qquad (a,i)(b,j)=(a+\zeta_n^i b,i+j).
\]
The action on the additive normal subgroup is irreducible over
$\mathbf F_r$, because $\mathbf F_r[\zeta_n]=\mathbf F_{r^{d_n}}$.
It is faithful. Consequently the complement $H_n=\langle\zeta_n\rangle$
is a maximal subgroup: a subgroup containing it intersects the
additive group in an invariant subspace. Its core is trivial,
since a nonidentity scalar fails to fix some additive element.
Thus $\Phi(G_n)\subset\bigcap_{g\in G_n}gH_ng^{-1}=1$.
The elements $(1,0)$ and $(0,1)$ generate $G_n$, since the scalar
orbit of $1$ spans the whole additive group.

The surface group surjects onto the free group on two generators
by sending two of its $a$-generators to free generators and every
$b$-generator, and all remaining $a$-generators, to one. Its
pro-$\Sigma$ completion therefore surjects onto each $G_n$.
The image of a Frattini subgroup under a surjection is contained
in the target Frattini subgroup: pull back every maximal subgroup.
Hence $\Delta/\Phi(\Delta)$ surjects onto every $G_n$.
Their orders are unbounded, since they are divisible by $s^n$;
the quotient is infinite. The characteristic-$p$ interpretation
uses the geometric specialization isomorphism already checked in
A.4, with all primes of $\Sigma$ invertible.

For the finite example use
\[
 \mathbf F_{64}=\mathbf F_2[t]/(t^6+t+1),\quad
 \zeta=t^7,\quad\lambda=t^9,
 \qquad G=(\mathbf F_{64},+)\rtimes\langle\zeta\rangle.
\]
Here $\zeta$ has order nine and generates the degree-six field;
$\lambda$ has order seven. Multiplication by $\lambda$ on the
additive factor commutes with multiplication by $\zeta$, so
\[
                    (a,i)\longmapsto(\lambda a,i)
\]
is an automorphism of order seven. Since $\zeta-1$ is invertible,
the commutators with $\zeta$ fill the additive factor. Thus
$G^{\mathrm{ab}}=\mathbf Z/9$, on which this automorphism is
the identity. The program \path{code/frattini_two_prime.py}
checks the field, the root orders, all $576^2$ automorphism
compatibilities, the two-generator orbit, the commutator image,
and the trivial complement core. This supplies the exact finite
counterexample to replacing the Frattini quotient by abelianization
in the attempted extension of the single-prime argument.
\end{proof}


\subsection{The covering condition and finitely generated fields}

\begin{claim}{A.30}{Proved}{Prime-to-characteristic Kummer kernels}
Let $K$ be a field finitely generated over $\mathbf F_p$, and let $G/K$
be a semiabelian variety, an extension of an abelian variety by a torus.
For the prime-to-$p$ Kummer map on rational points, the kernel is
\[
 \ker\left(G(K)\longrightarrow
       \prod_{\ell\ne p}H^1(G_K,T_\ell G)\right)=G(K)[p^\infty].
\]
For a field finitely generated over $\mathbf Q$, the analogous map
using every prime has zero kernel. No finite-generation assertion
about $K^\times$ is needed or made.
\end{claim}
\begin{proof}
For each integer $n$ invertible in $K$, the finite Kummer sequence
injects $G(K)/nG(K)$ into $H^1(G_K,G[n])$. Compatibility in $n$
identifies the displayed kernel with
$\bigcap_{(n,p)=1}nG(K)$; in characteristic zero use all positive $n$.
Passing to a finite separable extension splitting the torus preserves
any equality $x=ny$. We may therefore establish the required
implication there; the assertion that a rational point is killed by
a power of $p$ descends through the injective inclusion on points.

For a split torus, the divisor sequence on a normal proper integral
model shows that $K^\times$ modulo its constant-field units embeds in
the free abelian group on prime divisors. Its image is free abelian.
The constants are finite in positive characteristic. In characteristic
zero they form a number field; the unit theorem and the principal
fractional ideal group show again that its multiplicative group is
finite torsion plus a free abelian group. Thus $T(K)$ is free abelian
plus finite torsion. The image of $G(K)$ in the abelian quotient $A(K)$
is finitely generated. Here we use the Lang--N\'eron theorem
\cite[Theorem 2.1]{Conrad2006}: for a finitely generated regular
extension of a constant field, rational points on an abelian variety
modulo the points of its trace form a finitely generated group. The
constant-field trace has finite rational points over a finite field,
and finitely generated points over a number field.

An extension of a free abelian group plus finite torsion by a finitely
generated abelian group is itself a free abelian group plus a finitely
generated group. Indeed choose finitely many lifts of generators and
relations of the quotient. The finitely many relation values involve
only finitely many elements of a basis of the free kernel; all remaining
basis elements split off. In such a group the intersection of all
prime-to-$p$ multiples is exactly its $p$-primary torsion; the
intersection of all multiples is zero. This proves the assertion.
The converse for a $p$-primary point is immediate because every integer
prime to $p$ acts invertibly on its finite cyclic subgroup.
\end{proof}

\begin{claim}{A.31}{Proved}{The covering condition is equivalent to collision}
Let $C/k$ be as in \statementref{A.1}, and let $P,Q\in C(k)$.
Their sections are geometrically conjugate if and only if the
following condition holds. For every finite extension $L/k$ and every
connected finite \emph{geometrically} prime-to-$p$ Galois cover
$D\to C_L$ with a specified $P'\in D(L)$ over $P$, there is
$Q'\in D(L)$ above $Q$ such that
$[Q'-P']\in J_D(L)[p^\infty]$. When this condition holds, the lift
$Q'$ is unique. This is the condition of cote 149, scan 36, with
\emph{$p$-torsion} understood as $p$-primary torsion.
\end{claim}
\begin{proof}
Equality of section classes identifies the Galois actions on the
prime-to-$p$ fibre functor. A fixed point $P'$ in the fibre therefore
gives a fixed point $Q'$ with equal lifted section classes. For a
geometric Galois cover of order prime to $p$, the geometric open
subgroup has its full prime-to-$p$ topology: a core is obtained by
intersecting finitely many conjugate subgroups, and the resulting
quotient embeds in a product of groups of prime-to-$p$ order.
Thus the lifted conjugacy is in the correct full arithmetic quotient
of $D$. Abelianization and \statementref{A.30} give the torsion.

We justify uniqueness independently of the theorem being proved.
Over an algebraic closure suppose two such lifts differ by a
nonidentity deck transformation $a$ of order $m$ prime to $p$.
Then $n(aR-R)$ is principal for some power $n$ of $p$. Choose $f$
with that divisor. Its norm under $\langle a\rangle$ is a constant;
scale $f$ to make the norm one. Hilbert 90 writes
$f=a(u)/u$, with a convention for the action on divisors.
Consequently $\operatorname{div}(u)-nR$ is $a$-invariant (replacing
$a$ by its inverse changes no conclusion). Every orbit of this free
action has length $m$, so the degree of an invariant divisor is
divisible by $m$. Its degree is $-n$, a contradiction.

Conversely use a cofinal system of open characteristic normal
subgroups $N$ of $\Delta_C$. They are stable under the action defined
by $s_P$, and the subgroup $N\rtimes s_P(G_k)$ defines a connected
pointed cover $D_N/k$. Its geometric deck group need not be constant.
After a finite constant extension making it constant, the assumed
condition supplies a lift $Q_N$. Uniqueness makes this lift invariant
under the original $G_k$, hence rational over $k$. Norm maps on
Jacobians send the distinguished difference class to one at every
coarser level; uniqueness gives compatibility. Compatible rational
lifts determine a $G_k$-fixed element of the profinite path torsor
from $P$ to $Q$, equivalently a geometric conjugacy of their sections.
This proves the converse in the whole group. By \statementref{A.1}
the condition in fact forces $P=Q$.
\end{proof}

\begin{claim}{A.32}{Proved}{Specialization from finite type to a finite field}
If proper hyperbolic curves over all finite fields of characteristic
$p$ have injective geometrically prime-to-$p$ point-section maps,
then so do proper hyperbolic curves over every field $K$ finitely
generated over $\mathbf F_p$.
\end{claim}
\begin{proof}
Suppose $P\ne Q$ collide on $C/K$. Spread the smooth proper curve,
the two sections, and its constant genus over a smooth connected
finite-type scheme $S/\mathbf F_q$ with function field $K$, shrinking
$S$ so that the sections are disjoint. Choose a closed point $s$,
let $R=\mathcal O^h_{S,s}$ and $L=\operatorname{Frac}(R)$. Pullback
along $G_L\to G_K$ preserves the proposed geometric conjugacy.

Smooth proper specialization identifies the geometric prime-to-$p$
fundamental groups of the generic and special fibres over this
henselian base. We use \cite[Expos\'e X, Corollary 3.9, and
Expos\'e XIII, Corollaries 2.9--2.10]{LA_SGA1}, with a smooth proper
curve and no boundary or wild inertia. The comparison is equivariant
for the arithmetic map
$G_L\to\pi_1(\operatorname{Spec}R)=G_{k(s)}$ and gives a map of
arithmetic extensions. This arithmetic map is surjective: a connected
finite \emph{\'etale} cover of the normal henselian domain $R$ has
connected generic fibre, so every finite quotient of $G_{k(s)}$
is already seen by $G_L$.

The specialization map carries the sections supplied by $P,Q$ to
those supplied by $P_s,Q_s$ after this surjection. A single geometric
conjugator therefore conjugates their values for every element of
$G_{k(s)}$. Since $P_s\ne Q_s$, this contradicts the finite-field
hypothesis. This explains the direction of specialization and the
arithmetic surjectivity needed in the reduction.
\end{proof}

\begin{claim}{A.33}{Proved}{Every hyperbolic curve over a finitely generated field}
Let $K$ be finitely generated over $\mathbf F_p$, and let $X/K$ be
a smooth geometrically connected hyperbolic curve. Its map from
$X(K)$ to geometric conjugacy classes of sections of its arithmetic
geometrically prime-to-$p$ group is injective. For $K$ finitely
generated over $\mathbf Q$, the same assertion holds for the full
geometric fundamental group. The positive-characteristic assertion
is equivalent to \statementref{A.1}.
\end{claim}
\begin{proof}
First make a finite algebraic extension over which the smooth
geometric compactification and its distinct boundary points descend.
This is possible because their equations and morphisms use finitely
many algebraic coefficients. The purely inseparable part induces
universal homeomorphisms and equivalences on finite \emph{\'{e}tale}
covers, and identifies the absolute Galois groups. The separable part
restricts both sections along the same Galois homomorphism. Either
part therefore preserves a collision, and distinct points remain
distinct. The enlarged field is still finitely generated. This step
handles imperfect fields and boundary residue fields that cannot be
split by a separable extension alone.

For compactified genus at least two, the quotient induced by
$X\hookrightarrow\overline X$ carries a collision to one on the
proper curve. Statements \statementref{A.1} and
\statementref{A.32} exclude it.

If the compactification has genus zero, make a finite separable
extension splitting three boundary points, and choose their
coordinates to be $0,1,\infty$. The coordinate is a morphism
$X\to\mathbf G_m$ injective on geometric points. A collision
has equal Kummer coordinates. By \statementref{A.30} their ratio
is $p$-primary torsion in $K^\times$, hence one in characteristic $p$.

For compactified genus one, a putative rational point $P$ makes
the compactification an elliptic curve $E$, with origin $P$ and
nonempty boundary $B$. A second colliding point $Q$ has class of
exact order $p^a$ in $E(K)$, by \statementref{A.30}. Choose $m$
prime to $p$ with $m^2>p^a$. Pull back the punctured curve by
$[m]:E\to E$, and make a finite separable extension splitting the
boundary and deck group. The lifted collision supplies $P'=0,Q'$
in $D=E-[m]^{-1}B$. Its class in the generalized Jacobian $G_D$ is
$p$-primary. Its projection to $E$ has exact order $p^a$, because
$[m]$ is invertible on geometric $p$-primary torsion. The torus
kernel of $G_D\to E$ has no nonzero geometric $p$-primary points,
so the generalized class itself has exact order $p^a$.
The definition of the generalized Jacobian with reduced modulus
then supplies a function $f$ with
\[
 \operatorname{div}(f)=p^a(Q'-P'),\qquad
 f(T)=1\quad(T\in[m]^{-1}B).
\]
But $f-1$ has at least $m^2$ distinct zeros and pole divisor of
degree $p^a$, impossible. The boundary points are distinct because
$[m]$ is \emph{\'etale}, and none is $P'$ or $Q'$.

In characteristic zero use the Abel map into the generalized
Jacobian and the all-prime case of \statementref{A.30}. Its
injectivity on points follows directly: equality would give a
function with divisor $Q-P$ and value one on the reduced boundary.
For positive compactified genus a degree-one function is impossible;
for genus zero hyperbolicity supplies at least three boundary points,
while a degree-one function takes a specified value only once.
Finally finite fields are finitely generated and proper genus at
least two is hyperbolic, proving the equivalence. This answers
cote 149, scans 35--36, including the affine cases.
\end{proof}



\subsection{The configuration statements in the cited literature}

\begin{claim}{A.34}{Proved: the literal scope of the disputed assertion}{A10, A14}
Mochizuki's original author Proposition~1.8(ii), equivalently
Proposition~1.6(ii) in the separate LaTeX version, asserts the full
prime-to-$p$ punctured-fibre identification when the ground field has
characteristic $p$ and $\Sigma$ consists of all primes different from $p$.
Consequently the noninjectivity established in the configuration audit
contradicts this specialization of the assertion itself. The February
2014 Comments do not amend this assertion.
\end{claim}
\begin{proof}
The conventions preceding original Proposition~1.6 define the geometric
configuration group by taking the maximal pro-$\Sigma$ quotient of the
tame fundamental group and then the maximal cuspidally pro-$\Sigma^\dagger$
quotient over the proper product. For the present finite-field choice,
$\Sigma^\dagger=\Sigma$. A closed subgroup of a pro-$\Sigma$ group is
already pro-$\Sigma$, so the latter quotient makes no change. The same
observation applies to the punctured fibre. Original Proposition~1.8(ii)
then identifies the inverse image of a rational-point decomposition
group under a configuration projection with that punctured group.
It is therefore not a statement about a singleton prime or about only
the completion of the fibre kernel.

The proof cites Stix's monodromy paper. Its relative completion retains
the entire base fundamental group and completes the fibre kernel.
The arXiv version's \S2.3 and Proposition~2.7 give that construction.
They do not justify additionally completing the base. The four corrections
in the 2014 Comments concern descended forms in decomposition recognition,
targets of diagonal covers, a Kummer injection, and infinitude of a
cyclotomic image, respectively. None replaces the fibre identification.
These readings establish the precise conflict; they do not establish a
verdict about every theorem of the paper.
\citep[original pp.~11,14--15]{LA_Mochizuki2007}
\citep{LA_Mochizuki2014,LA_Stix2005}
\end{proof}

\paragraph{Version control.}
The following numbering distinctions matter. In Sa\"idi--Tamagawa 2009,
the supplied arXiv Remark~3.13 is journal Remark~3.14, and the supplied
arXiv Theorem~2.7 is journal Theorem~2.6. The former is the explicitly
conditional configuration-space route to elimination of exceptional
points; journal Remark~3.13 concerns a different extension of the Isom
theorem. Mochizuki's original Theorem~3.10 is Theorem~3.1 in the LaTeX
version. Statements below specify the original numbering when necessary.
\citep{LA_ST2009,LA_Mochizuki2007}

\begin{claim}{A.35}{Proved: a limited repair of the fibre input}{A10}
Let $k$ be a perfect field, let $C/k$ be a proper smooth
geometrically connected curve of genus at least two, and let
$\Sigma^\dagger=\{r\in\Sigma:r\text{ is invertible in }k\}$
be nonempty. Thus $\Sigma^\dagger=\Sigma$ in characteristic
zero and $\Sigma^\dagger=\Sigma\setminus\{p\}$ in
characteristic $p>0$.
Let $F\to E\to\Delta_2\to1$ be the geometric right-exact configuration
sequence, let $F\twoheadrightarrow\Delta_1$ be the proper fibre quotient,
and let $E\twoheadrightarrow\Delta_1\times\Delta_2$ be the proper product
quotient. Write $K_F,K_E$ for their cuspidal kernels. Suppose the relative
groups are formed from the \emph{tame} boundary groups, with Mochizuki's
maximal cuspidally pro-$\Sigma^\dagger$ convention. The proper groups
may retain $p$ in $\Sigma$. The relative cuspidally pro-$r$ fibre
theorem of Part~A, including \statementref{A.13}, holds for every
$r\in\Sigma^\dagger$.
Then
\[
 K_F^{(r)}\xrightarrow{\sim}K_E^{(r)}\quad(r\in\Sigma^\dagger),
 \qquad K_F^{\mathrm{ab}}\xrightarrow{\sim}K_E^{\mathrm{ab}}.
\]
Consequently the configuration fibre sequence is left exact after
cuspidal abelianization, as well as after the specified fixed-prime
cuspidal completion. The one-puncture geometrically cuspidally central
fibre identification also holds. For a $k$-rational point, the
corresponding arithmetic inverse-image identifications hold over
the unchanged quotient $G_k$.
\end{claim}
\begin{proof}
Right exactness gives a surjection $K_F\twoheadrightarrow K_E$:
an element of $K_E$ lifts to $F$, and its trivial first proper projection
forces the lift into $K_F$. The relative pro-$r$ theorem identifies the
two cuspidal pro-$r$ kernels by taking kernels of the identified fibre
groups over $\Delta_1$. Every abelian pro-$\Sigma$ group is the product
of its pro-$r$ factors. In positive characteristic there is no $p$
factor in the present cuspidal
kernels: each is topologically normally generated by conjugates of
tame inertia, and every map to a finite $p$-group kills each such
generator. In particular $K_F^{(p)}=K_E^{(p)}=1$ when $p\in\Sigma$.
This assertion would be false for the full \emph{wild} punctured
\'etale group, and that group is not the convention used here.
Moreover the maximal abelian pro-$r$ quotient of
a group is the abelianization of its maximal pro-$r$ quotient. Taking
the product of the resulting isomorphisms proves the assertion about
abelianizations. Applying the elementary extension isomorphism criterion
over $\Delta_1$ proves the cuspidally abelian fibre statement.

For one puncture, taking $\Delta_1$-coinvariants of the abelian cuspidal
kernel leaves the orientation module
$\widehat{\Z}^{\Sigma^\dagger}(1)$, generated by the boundary loop. The other
geometric proper factor acts trivially on this orientation: the family
has one labelled boundary and its local parameter gives the same tame
inertia orientation. Hence taking its coinvariants adds no relation.
This proves the geometrically central assertion. The arithmetic action
on the surviving module is the cyclotomic action. Equivalently this
last extension is constructed directly from the Kummer cycle class of
the diagonal, whose restriction to the fibre is the degree-one point
class. Neither argument requires injection of the full fibre.

The relative pro-$r$ proof is geometric and uses no finite-field
count or Frobenius argument. In positive characteristic it is
exactly the direct Stix argument of \statementref{A.13}; in
characteristic zero the same argument applies to every $r$, or
one uses the topological construction after descending the finite
data. For a rational point, the natural map from the punctured
arithmetic quotient to the inverse image of its decomposition
group is an isomorphism on the geometric kernels just proved and
on the arithmetic quotients, both $G_k$. The elementary extension
criterion makes it an isomorphism. Finite sets of closed points
are handled after a finite separable splitting extension; all
maps are natural geometric maps over the unchanged arithmetic
quotient, so their Galois descent data commute and the
identification descends. This justifies the stated field scope.
\end{proof}

\paragraph{Exact repair boundaries.}
This proves a replacement for the false intermediate assertion over
every perfect ground field with these tame cuspidal conventions,
including proper $\Sigma$-groups which retain the positive
characteristic. It covers all fields of the original paper:
its \S0 defines a nonarchimedean local field to mean a finite
extension of $\mathbb Q_p$, so its local convention
$\Sigma^\dagger=\Sigma$ is exactly the characteristic-zero
case here. It does not turn the original
full sequence into an exact sequence. The following dependency loci were
checked in the original author text. Proposition~1.8(iii),(iv) uses
1.8(ii) for central quotients. The paragraph preceding Proposition~1.12
(original p.~20) cites 1.8(ii),(iii) for surjectivity of the symmetrized
fundamental section. Proposition~1.14 builds cuspidally abelian quotients
as inverse limits. Before Proposition~3.8 (original p.~63), a punctured
fibre is inserted in the configuration group by 1.8(ii); Proposition~3.9(v)
identifies the fixed-$l$ cuspidal quotient, which enters Theorem~3.10.

At these loci one may use the central or abelian fibre identification
when the subsequent construction is central or abelian, and the relative
pro-$l$ identification when it is cuspidally pro-$l$. This specifies the
replacement actually proved here. Constructions on finite proper covers
retain their labelled boundary components. Here is the needed extension
of the relative proof. For a product compactification $C_0\times B$
with finitely many disjoint section graphs, choose the finite Galois
$A$-cover $C_1\to C_0$ required by a finite fibre quotient. Take a
finite Galois $\Sigma$-cover of $B$ dominating its pullbacks along
every section; their finite intersection and core give such a cover.
All lifted punctures are then labelled, the compactification remains
a product, and mod-$r$ monodromy is a shear between proper homology
and the boundary module. The same relative pro-$r$ and normal-core
argument applies, with finite proper deck group $A\times A_B$.
When this group's order is divisible by $p$, the direct algebraic
argument of \statementref{A.13} uses Stix's relative pro-$r$
completion over the full base; it requires no equivariant lift of the
proper cover to characteristic zero.
This proves the required compatibility for these finite proper covers.
For naturality under a finite proper pullback not itself visible
as a $\Sigma$-cover, apply the construction to that curve with
its own proper pro-$\Sigma$ group: naturality of the restriction
maps and the diagonal cycle class gives the comparison square.
One does not pretend that this cover is an open subgroup of the
original pro-$\Sigma$ group. Any conclusion retaining the entire
mixed-prime punctured kernel still loses its asserted fibre
identification. The dependency census below records that loss
separately from the repaired quotient constructions.
\citep[Propositions~1.8,1.12,1.14,3.9 and Theorem~3.10]{LA_Mochizuki2007}

\begin{claim}{A.36}{Proved: the primary and total formation distinction}{A10}
Mochizuki--Tamagawa's configuration exactness theorem and Wakabayashi's
configuration theorem do not assert the mixed prime-to-$p$ exactness
refuted in Part~A.
\end{claim}
\begin{proof}
Mochizuki--Tamagawa Definition~1.1 restricts the formation to a
PT-formation, namely primary or total. Proposition~2.2 additionally
requires invertibility on the connected base and the stated formation
conditions on the cyclotomic and compactified-fibre mod-$l$ monodromy
images. For classes of all finite $\Sigma$-groups, primary means one
prime; total means all finite groups. The proof uses that the kernel of
the automorphism action of a finitely generated pro-$l$ group on its
mod-$l$ abelianization is pro-$l$. A set of several primes omitting the
characteristic is not one of these two formations. The 2015 Comments
do not enlarge the formation scope. Wakabayashi fixes one invertible
prime throughout. Thus neither theorem supplies the disputed
mixed-prime completion step.
\citep[Definition~1.1 and Proposition~2.2]{LA_MT2008}
\citep{LA_MT2015,LA_Wakabayashi2016}
\end{proof}

\begin{claim}{A.37}{Proved: Wakabayashi's actual input and output}{A11}
Fix a prime $l$ invertible in the two finite fields. Wakabayashi's
Theorem~A, proved in the more general logarithmic form as Theorem~3.9,
starts with a Frobenius-preserving isomorphism of the geometrically
pro-$l$ groups of once-punctured curves, carrying the specified new
cusp decomposition groups to one another. It extends this isomorphism
to the two-point configuration groups, compatibly with the proper
quotient and switching, up to the stated inner ambiguity. Corollary~3.10
transports punctures at other points with corresponding proper
decomposition groups. Theorem~C, namely Theorem~4.4, reconstructs the
higher configuration groups from the geometrically pro-$l$ curve group.
These are reconstruction results in the singleton-prime category.
\end{claim}
\begin{proof}
These are the hypotheses and conclusions of the supplied source. Its
proof uses the inverse image of a point decomposition group under the
configuration projection, so the fibre method itself is already present.
Remark~3.10.1 explicitly permits nonuniqueness of the corresponding
point. The theorem retains a pro-$l$ proper group, rather than the full
prime-to-$p$ proper group with only a pro-$r$ cuspidal kernel. No operation
in the theorem reconstructs mixed finite quotients from separate
pro-$l$ quotients. Nor may the cofinite-prime exceptional-set theorem
be applied to a singleton set without verifying its hypothesis. Thus
these results do not conflict with AG.40--AG.42 and, with the stated
Sa\"idi--Tamagawa finiteness input alone, do not supply their universal
point-separation conclusion.
\citep[Theorems~3.9,4.4 and Corollary~3.10]{LA_Wakabayashi2016}
\end{proof}

\paragraph{Sa\"idi--Tamagawa dependencies.}
The 2009 configuration route in journal Remark~3.14 is expressly
conditional. Its failure does not by itself invalidate that paper's
Isom theorem. Theorem~2.2 reconstructs the configuration cuspidally
abelian group; Corollary~2.3 restricts the labelled puncture comparison
to the complements of exceptional sets; journal Theorem~2.6 gives
fixed-$l$ cuspidal reconstruction. The 2017 text repeats these as
Theorem~3.2, Corollary~3.3 and Theorem~3.6. Thus their imported
configuration machinery depends on the central, abelian, and fixed-$l$
consequences just distinguished, not on a final theorem asserting
full mixed-prime configuration exactness.

There is a further scope distinction. Original Mochizuki Theorem~3.10
allows arbitrary matching rational points, without a separation
hypothesis. Its proof begins with configuration cuspidally abelian
reconstruction from Theorem~1.16(iii), chooses matching splittings in
the diagonal decomposition groups, and restricts to the corresponding
cuspidally abelian fibres. Thus the reference to an outside-exceptional-set
corollary in Sa\"idi--Tamagawa is not on its own proof that the original
arbitrary-point theorem has no such argument. The needed repair is the
fibre input at the stated quotient level. Our main proof uses its own
relative fibre and transport constructions and does not assume the
published reconstruction theorem in order to prove that repair.
\citep{LA_ST2009,LA_ST2017,LA_Mochizuki2007}

\subsection{Proper support fields and the later exceptional-set theorem}

\begin{claim}{A.38}{Proved: the exact later largeness conclusion}{A10, A14}
In Sa\"idi--Tamagawa 2017, Definition~2.6 calls $\Sigma$ $J_C$-large
when the omitted-prime representation
$\rho_{J_C,\Sigma'}:G_k\to\Aut(T_{\Sigma'}J_C)$ is noninjective.
Definition~2.9 calls $C$ almost $\Sigma$-separated when its exceptional
closed points are supported over some proper algebraic subfield
$F\subsetneq\bar k$. Proposition~2.12 proves this latter conclusion
under $J_C$-largeness. It does not in general assert that the
exceptional set is finite.
\end{claim}
\begin{proof}
Put $F_0=\bar k^{\ker\rho_{J_C,\Sigma'}}$. The proof uses the difference
map
\[
 \delta:C\times C\longrightarrow J_C,\qquad (P,Q)\longmapsto[P-Q].
\]
Off the diagonal its geometric fibres have cardinality at most two,
or at most one when $C$ is not hyperelliptic. Indeed, two different
pairs in a fibre give $P+Q'\sim P'+Q$. A nonconstant degree-two linear
system is the hyperelliptic system; in that case the only possible
second pair is $(\iota Q,\iota P)$, with $\iota$ the hyperelliptic
involution. This argument is the predecessor explicitly cited from
2009 Claim~1.9(i), where the curve has first been made
nonhyperelliptic.

An exceptional pair has distinct images on $C_{\bar k}$ and its
difference lies in $J_C(\bar k)\{\Sigma'\}\subset J_C(F_0)$.
Taking the subgroup of squares of $G_{F_0}$ kills the possible
two-element permutation. Hence both coordinates are defined over a
finite extension of $F_0$, still proper in $\bar k$. This is exactly
the argument of Proposition~2.12. Proposition~2.13 then chooses a prime
$l$ for which the corresponding $\Z_l$ tower has finite intersection
with every fixed finite extension of that support field. It may choose
$l\in\Sigma\cup\{p\}$ under largeness, and $l\in\Sigma$ if the
Jacobian also has positive $p$-rank. Proposition~2.11(i) gives
$E_{C'}=f^{-1}(E_C)$ for finite \'etale $\Sigma$-covers and transports
almost separation through them.
\citep[Definitions~2.6,2.9 and Propositions~2.11--2.13]{LA_ST2017}
\end{proof}

\begin{claim}{A.39}{Proved with the relative criterion of Part A}{A10, A14}
The separation argument extends from finite exceptional support to
almost $\Sigma$-separation whenever $\Sigma\setminus\{p\}$ is nonempty,
including $p\in\Sigma$. We use the relative fixed-$r$ fibre and
transport statements of Part~A in that $\Sigma$-category, as extended
by \statementref{A.13}. Consequently the $J_C$-large
almost-separated cases of Sa\"idi--Tamagawa become separated.
\end{claim}
\begin{proof}
Fix $r\in\Sigma\setminus\{p\}$. The point-transport and difference-cocycle argument
in Part~A gives
$\kappa_r(f(R)/a)=0$ outside the exceptional support and the two
punctures, where $a\in k^\times$ and
$\operatorname{div}(f)=n(P-Q)$. Here one may take $n$ to be the
\emph{least} positive order of $[P-Q]$ in $J_C(\bar k)$; it need not
be invertible at $r$. The expression with zero residues is
$\kappa_r(f)-n\chi$, so no division by $n$ occurs. Hilbert~90 descends
the generator of this rational principal divisor to $k(C)$.

The function $f$ is not an $r$th power over $\bar k(C)$. If $r\nmid n$,
its valuations prove this; if $r\mid n$, an $r$th root would make
$(n/r)(P-Q)$ principal, contrary to minimality. Choose a finite field
$L\supset k$ containing $\mu_r$ and let $D/L$ be the smooth proper
model of $z^r=f/a$. This curve is geometrically connected. Let $h$ be
its genus and let $L_m/L$ run through the powers of the auxiliary
prime $l$ supplied by Proposition~2.13. The exceptional points rational
over these fields lie in a fixed finite set.

Write $Q_m=|L_m|$ and let $M_m$ count points of
$C(L_m)\setminus\{P,Q\}$ at which $f/a$ is an $r$th power.
Every such point has $r$ rational inverse images on $D$, so
$rM_m\leq|D(L_m)|$. The Weil bounds imply
\[
 |C(L_m)|-2-M_m\ \geq\
 (1-r^{-1})(Q_m+1)-(2g+2h/r)\sqrt{Q_m}-2,
\]
which tends to infinity. Some point outside the bounded exceptional
set therefore has nonzero mod-$r$ Kummer class, contradicting the
displayed evaluation. This proves the extension. The counting argument
does not require $l\ne r$.
\end{proof}

\paragraph{The elementary full prime-to-$p$ specialization.}
For the main theorem one can avoid the earlier Boxall-based finiteness
input entirely. Let $f_p$ be the $p$-rank of $J_C$. Choose an odd prime
$l\ne p$ not dividing $|\mathrm{GL}_{f_p}(\F_p)|$, with the evident
convention when $f_p=0$. The subgroup $\Z_l\subset G_k$ acts trivially
on $T_pJ_C$, since the congruence kernel in its automorphism group is
pro-$p$, and therefore acts trivially on all geometric $p$-primary
torsion. The preceding two-point difference-fibre bound then puts
every distinct pair with $p$-primary difference over
$\bar k^{\Z_l}$. Along the $l$-power constant-field tower these bad
labels are uniformly bounded. This is a direct special case of the
known support-field method, combined with the new criterion as proved
in Part~A; the difference-fibre lemma itself is not new.

\subsection{What is and is not supplied by the opened sources}

\begin{claim}{A.40}{Proved: the recorded predecessor comparison}{A14}
The following comparison records exact predecessors and the additional
work required in this audit. It is not a priority claim.
\end{claim}
\begin{center}\small
\begin{longtable}{p{0.17\linewidth}p{0.37\linewidth}p{0.38\linewidth}}
\toprule
Step&Opened predecessor&Additional issue resolved or retained\\\midrule
AG.40, point recognition&
Mochizuki original Remark~1.18.2 and its 2014 correction; the conclusion
is ``Any Frobenius-preserving isomorphism $\alpha$ is quasi-point-theoretic.''&
One must prove that the relative quotient still recovers the compactified
$r$-adic cohomology of every visible cover; the finite-intersection
argument must allow all relevant descended forms.\\
AG.40, Kummer evaluation&
Mochizuki Proposition~2.1(ii),(iii): residue exactness and evaluation of
a function's Kummer class at a point. Sa\"idi--Tamagawa repeats it.&
The over-identity cusp isomorphism yields a difference cocycle with
residues $+1,-1$; combining it with transported point labels gives the
contradiction. No identical full criterion was located.\\
AG.41, exactness&
Stix's relative completion; M--T Proposition~2.2 in PT formations;
Hoshi--Sawada--Tsujimura Definition~1.6(ii) and Corollary~1.14;&
Neither states the mixed proper-base, fixed-$r$ cuspidal quotient used
here. The direct proof and its finite-quotient checks remain essential.\\
AG.42, fibre transport&
Mochizuki original Remark~1.7.1 at the central level; Wakabayashi
Corollary~3.10 for a fixed prime; Sun Definition~3.1 and Corollary~3.8.&
The general pullback trick is known. Its use in the relative mixed-base
quotient also requires the local diagonal-boundary lift and compatibility.\\
Specialization and inertia&
SGA1 XIII Corollary~2.9, \S2.10, Lemma~2.11 and Proposition~5.2.&
The hypotheses require a smooth proper pair with relative normal crossings
and invertible coefficient primes; the local Kummer extension exists on
the whole regular chart, by purity.\\
Orbifold completion&
Boggi--Zalesskii Definition~1.1 and Lemma~3.1, with a torsion-free normal
subgroup and quotient in the extension-closed formation.&
The prime-to-$p$ formation meets these hypotheses when the chosen
orbifold quotient has prime-to-$p$ order; no $p$-acyclicity is required.\\
\bottomrule
\end{longtable}
\end{center}
The 2024 Hoshi--Sawada--Tsujimura result requires good log-smooth
pairs and a proper geometrically connected morphism, smooth away
from the base boundary, with smooth relative boundary there, and
log smooth as a morphism of pairs. Definition~1.6(ii) takes the
prime-to-characteristic quotient of the \emph{kernel} of the map
between full tame fundamental groups. Corollary~1.14 identifies
the geometric prime-to-characteristic fibre when it is centre-free.
The base remains the full tame group. This is a closer recent
predecessor for relative completion, and supplies no simultaneous
prime-to-$p$ completion of that base. Its proof in \S1 was checked
at these precise statements.
\citep[Definitions~1.1,1.3,1.6 and Corollary~1.14]{LA_HST2024}

The October 2026 Hoshi--Tsujimura preprint also has a point-section
injectivity theorem: Theorem~A(i), proved as Theorem~1.8, concerns
affine varieties over a characteristic-$p$ field admitting a Galois
extension of degree $p$. It uses the full \'etale fundamental group
and Artin--Schreier evaluations. It does not assert the proper-curve
prime-to-$p$ injectivity at issue here.
\citep[Theorem~A(i) and Theorem~1.8]{LA_HT2026}

The sources were opened in full or at the identified proofs, not merely
through search snippets. Searches included each named paper, both Sun
preprints, and combinations of prime-to-$p$ point separation, exceptional
sets, cuspidalization, and configuration exactness. No independently
published universal prime-to-$p$ point-separation proof was found in
that bounded search. This does not prove novelty or exclude an unknown
predecessor.\citep{LA_Mochizuki2007,LA_Mochizuki2014,LA_ST2009,LA_ST2017,LA_MT2008,LA_Stix2005,LA_Wakabayashi2016,LA_Sun2025,LA_SGA1,LA_BoggiZalesskii2019,LA_HST2024,LA_HT2026}

\subsection{A finite dependency census for the disputed exactness}

\begin{claim}{A.41}{Proved}{Which imported fibre arguments require replacement}
For the five papers specified in A10, the dependencies on the
configuration fibre assertion are as listed below. At every
listed central, abelian, or fixed-prime use, A.35 supplies the
same quotient statement with its maps and boundary data. The
literal mixed-prime fibre insertion and any argument requiring
that full insertion are lost. This is a census of the disputed
dependency, not a reproof of unrelated theorems in those papers.
\end{claim}
\begin{center}\small
\begin{longtable}{p{0.23\linewidth}p{0.35\linewidth}p{0.34\linewidth}}
\toprule
Source locus&Use of the fibre input&Verified replacement or exact loss\\\midrule
Mochizuki 2007, Proposition~1.8(ii)&
Identifies the full punctured group with a configuration inverse
image; concludes slimness of the full configuration group.&
The full mixed-prime identification is false in the specialization
of A.34. Only its central, abelian and fixed-prime versions are
proved here. Full mixed-group slimness is not inferred from those
versions or from the invalid full-fibre proof.\\
Proposition~1.8(iii),(iv)&
Specializes the fundamental section, proves surjectivity from
boundary inertia, and identifies central quotients.&
The diagonal cycle class and its point restrictions give the
same maps. Primitive boundary inertia gives surjectivity and
the universal central quotient, by A.35. The separate slimness
claim for curve groups in (iii) is proved there from curve-group
results, not from the disputed configuration insertion.\\
Paragraph before Proposition~1.12; Propositions~1.12,1.14--1.16,
1.19&
Surjectivity of symmetrized fundamental sections; inverse limits
of central quotients; reconstruction and switching of the
cuspidally abelian configuration group.&
Apply the repaired central maps on each finite proper pullback
and take their compatible inverse limits. The needed comparison
and rigidity are at the abelian level.\\
Paragraph before Proposition~3.8; Propositions~3.8--3.9;
Theorem~3.10&
Writes the full punctured fibre as a subgroup to define diagonal
actions, then uses its fixed-$l$ Lie quotients and finite proper
pullbacks.&
Replace the ambient group and the fibre from the outset by their
relative cuspidally pro-$l$ quotients. A.35 and the retained
boundary give those actions and inclusions. Every Lie object
used here factors through this quotient. The displayed full
mixed-group inclusion on original p.~63 remains invalid.\\
Mochizuki--Tamagawa 2008, Proposition~2.2 and its subsequent
configuration sequences&
Uses its own PT-formation homotopy theorem. The citation to
Mochizuki 2007 in Remark~4.8.2 concerns a proposed later
application.&
No proof in this dependency chain invokes the false mixed
formation. Its singleton and total hypotheses are retained.\\
Sa\"idi--Tamagawa 2009, Theorem~2.2, Corollary~2.3,
Proposition~2.4, Theorem~2.6&
Imports Mochizuki's abelian reconstruction, marked restrictions,
central quotient, and fixed-$l$ reconstruction, respectively.&
These are exactly the repaired quotient inputs. The conditional
full-group proposal in journal Remark~3.14 is lost; it is not a
premise of their Isom proof.\\
Sa\"idi--Tamagawa 2017, Theorem~3.2, Corollary~3.3,
Proposition~3.4, Theorem~3.6&
Reuses the preceding four constructions, with its own exceptional
sets and largeness hypotheses.&
The same four quotient repairs apply. Its support-field and
function-reconstruction arguments have independent hypotheses,
which remain as in C.5--7.\\
Wakabayashi 2016, Proposition~1.4; Theorems~3.9,4.4&
The fibre theorem comes from M--T Proposition~2.2 or Stix in
the single-prime category. The proper two-point case of 4.4
cites Mochizuki's fixed-$l$ theorem.&
No full mixed-prime fibre assertion is needed in the category
used by these proofs. Its imported fixed-$l$ theorem is in the
repaired scope, and also in the singleton case where the
homotopy theorem is valid.\\
\bottomrule
\end{longtable}
\end{center}
\begin{proof}
The full-text occurrence census for original Proposition~1.8(ii)
has three direct loci: its use in the proof of (iii), the
surjectivity paragraph before Proposition~1.12, and the fibre
insertion immediately before Proposition~3.8. The other explicit
citations to Proposition~1.8 use its proper-group slimness in
(i), its curve or central assertions in (iii),(iv), or the
definitions of the punctured groups. We checked the mathematical
content of these loci, not just the citation strings.

For the first two loci, the source's required map sends each
boundary generator to its corresponding orientation-module
generator in the extension defined by the diagonal class.
A.35 constructs this map without a full-fibre injection. Since
the central cuspidal kernel is generated by these images, the
map is surjective and has the universal central quotient kernel.
On a finite proper cover, conjugate boundary components give the
separate factors of the symmetrized extension. The labelled-cover
argument in A.35 proves compatibility with their deck actions.
Taking inverse limits identifies the full cuspidal abelian
module, hence gives the identifications in Proposition~1.14.
The proof of Theorem~1.16(iii) reconstructs these extensions and
their transition maps, and its uniqueness argument uses
Proposition~1.15 at this abelian level. The induced geometric
switching map exists before taking any quotient; its needed
abstract characterization in Proposition~1.19 likewise uses
Proposition~1.15, not a full mixed-group fibre theorem.

At the third locus the paper has $S=S_* =\{x_*\}$, hence
$S_0=\varnothing$. Its ensuing objects are constructed from the
pro-$l$ groups of the punctured finite proper covers, their
weight quotients and their Lie realizations. Consequently each
action in Propositions~3.8--3.9 factors through the relative
cuspidally pro-$l$ fibre with its proper pro-$\Sigma$ quotient.
Replace the displayed full configuration group by that relative
quotient. A.35 identifies its fibre with the relative
punctured group and identifies the image of the diagonal
decomposition group with its retained tame boundary. Conjugation
therefore supplies the same continuous diagonal actions on every
pro-$l$ object. Naturality on finite proper pullbacks supplies
their transition maps. No weight filtration, trace map or
Frobenius splitting in the subsequent proof changes. The
slimness used at the end of Proposition~3.9(v) is explicitly
slimness of the \emph{singleton-$l$} curve and configuration
groups, where M--T Proposition~2.2 applies. Thus the fixed-$l$
reconstruction of Theorem~3.10 has the repaired input, without
restoring the original full mixed-group insertion.

The two Sa\"idi--Tamagawa papers cite precisely these central,
abelian and fixed-$l$ results at the four listed theorem loci.
Their subsequent use in Green compatibility and multiplicative
function reconstruction only uses these quotients; C.5--7
give the separate consequences of removing exceptional sets.
For M--T, the full text restricts configuration exactness to
PT formations in Proposition~2.2. Its sole invocation of the
cuspidalization paper in Remark~4.8.2 is motivational, not a
proof input. Finally Wakabayashi fixes one invertible prime
throughout and gives the exact source of its fibre comparison
in Proposition~1.4; its references to Mochizuki's weight
filtration and fixed-$l$ reconstruction stay in that category.
These checks prove the stated dependency classification.
\citep[Propositions~1.8,1.12,1.14--1.16,1.19,3.8--3.9;
Theorem~3.10]{LA_Mochizuki2007}
\citep[Proposition~2.2, Remark~4.8.2]{LA_MT2008}
\citep[\S2, Remark~3.14]{LA_ST2009}
\citep[\S3]{LA_ST2017}
\citep[Proposition~1.4, Theorems~3.9,4.4]{LA_Wakabayashi2016}
\end{proof}


\section{Independent existence and reconstruction questions}
\subsection{The exact relative criterion and the missing existence assertion}

We retain the whole arithmetic group throughout. A cover is
geometrically connected, and ``constant'' refers to its deck group
over the displayed finite field. A $P$-split cover has a chosen rational
point over $P$; for a constant Galois group this makes its entire fibre
above $P$ rational.

\begin{claim}{B.1}{Proved}{The relative class-number test, with its group hypothesis}
Let $C/L$ be proper of genus at least two over a finite field of
characteristic $p$, and let $D\to C$ be a nontrivial constant Galois
\'etale $P$-split cover of degree $d$ prime to $p$. Put
$B=(\ker(J_D\to J_C))^0$. Then
\[
 \frac{\#J_D(L)}{\#J_C(L)}=\#B(L)\in\mathbf Z_{>0}.
\]
If the group has a nonidentity element of odd order, the condition
$B(L)[p]=0$ directly excludes a point-section collision at $P$.
This includes every odd prime cyclic cover used in B1.
The conclusion also holds for the remaining prime-to-$p$ groups
by the independently proved A.1, but the odd-order stabilizer
argument by itself does not prove that broader group assertion.
\end{claim}
\begin{proof}
Pullback and norm satisfy $\pi_*\pi^*=[d]$. Hence
$J_C\times B\to J_D$ is an isogeny over $L$. The equality of
finite-field point counts under isogeny gives the asserted integer.
On $J_D(L)[p^\infty]$, multiplication by $d$ is invertible and
\[
 e=d^{-1}\pi^*\pi_*=d^{-1}\sum_{\delta\in G}\delta
\]
is an idempotent. Its complementary image lies in $B[p^\infty]$:
it is the image of the homomorphism $[d]-\pi^*\pi_*$, followed
by division by $d$ on this finite $p$-primary group. Thus
$B(L)[p]=0$ implies that all rational $p$-primary classes are
fixed by $G$.

Under a collision, A.31 produces the unique lift $Q'$ above the
distinct point $Q$ with $T=[Q'-P']\ne0$ rational and $p$-primary.
If a nonidentity deck transformation fixes $T$, the effective
divisors $\delta Q'+P'$ and $Q'+\delta P'$ are distinct and
disjoint, and are linearly equivalent. They define a degree-two
pencil, so $D$ is hyperelliptic. Its hyperelliptic involution is
central. An odd-prime-order deck transformation descends to an
automorphism of the quotient projective line of that same order.
It has a geometric fixed point there; its permutation of the fibre,
which has at most two points, is trivial. It therefore has a fixed
point on $D$, contrary to the free deck action. This proves the
direct criterion. A.1, proved from the configuration, boundary,
and cuspidal criteria in Part A, supplies the final broader conclusion.
The explicit even-order control in the norm-module discussion
explains why we have kept these two justifications distinct.
\end{proof}

\begin{claim}{B.2}{Proved}{Simultaneous character criterion in every genus}
Let $a:J_C(L)\twoheadrightarrow A$ be a nontrivial finite abelian
quotient of order prime to $p$, and use the Abel map based at $P$
to pull back its constant Lang quotient. Put
\[
 \Theta_{C,P,A}=\sum_{d=0}^{g-1}
       \sum_{E\in\operatorname{Div}^{+}_d(C)(L)}
                    [a([E-dP])]\quad\text{in }\mathbf F_p[A].
\]
Its relative class number is prime to $p$ if and only if
$\overline\psi(\Theta_{C,P,A})\ne0$ for every nontrivial
character $\overline\psi:A\to\overline{\mathbf F}_p^\times$.
For $A=\mathbf Z/\ell$ this is precisely B1's simultaneous
test; it is not enough to check one reduced character unless
all primitive characters form one Frobenius orbit.
\end{claim}
\begin{proof}
For a nontrivial complex character $\psi$ put
\[
 c_d(\psi)=\sum_{E\in\operatorname{Div}^{+}_d(C)(L)}
                 \psi(a([E-dP])),\qquad
 \kappa=[K_C-(2g-2)P].
\]
The Euler product, expanded by effective divisors, gives
$L_P(\psi,T)=\sum c_d(\psi)T^d$. A line bundle $M$ of
degree $d$ has $(q^{h^0(M)}-1)/(q-1)$ effective divisors in
its class. For $d>2g-2$, Riemann--Roch makes this number
independent of the class, so the nontrivial character sum
vanishes. For $g-1\le d\le2g-2$, apply Riemann--Roch to
$M$ and $K_C-M$; the class-independent term again vanishes.
This gives
\[
 c_d(\psi)=q^{d+1-g}\psi(a(\kappa))
                       c_{2g-2-d}(\psi^{-1}).
\]
In particular the coefficients with $d\ge g$ vanish modulo
every prime over $p$. This is the stated truncation at $g-1$.

On each closed fibre of the cover, its Frobenius translation
has the character eigenvalues $\psi(a_x)$ in the regular
representation. Factoring this permutation determinant and
multiplying the Euler factors proves
\[
 \frac{Z(D,T)}{Z(C,T)}=\prod_{\psi\ne1}L_P(\psi,T),
 \qquad
 \frac{\#J_D(L)}{\#J_C(L)}=
                       \prod_{\psi\ne1}L_P(\psi,1).
\]
Since $p\nmid|A|$, all characters remain distinct after
reduction in an algebraic closure of $\mathbf F_p$ and give
all reduced characters. The product is a unit exactly when
each factor is a unit. The preceding coefficient calculation
identifies the factors with the evaluations of $\Theta$.
If $A=\mathbf Z/\ell$, the rational factors are Galois
conjugates and their product is the cyclotomic norm. Over
$\mathbf F_p$, their orbits have length
$\operatorname{ord}_\ell(p)$, which proves the last assertion.
\end{proof}

\begin{claim}{B.3}{Proved}{The corrected pair estimate and its precise limitation}
For any quotient $A$ of $J_C(L)$ let $N=\#C(L)$ and let
$N_b$ be the occupancies of the Abel-image points in its cosets.
Then
\[
 \sum_bN_b(N_b-1)
 \le \frac{N^2}{|A|}
       +\frac{|A|-1}{|A|}(2g-2)^2q-N.                    \tag{B.CL.1}
\]
If all nonzero occupancies are multiples of $p$, the left side
is at least $(p-1)(N-N_0)$. In genus two, if $p$ generates
$(\mathbf Z/\ell)^\times$, a surjection with kernel
$H\subset J_C(L)$ is a class-number witness whenever
\[
                       N\ge2|H|,\qquad \ell>N.         \tag{B.CL.2}
\]
Neither $p>(2g-2)^2$ nor these estimates alone proves universal
existence of such a quotient. For $g>2$, occupancies of points
alone are not the character test in B.2.
\end{claim}
\begin{proof}
Character orthogonality gives
\[
 \sum_bN_b^2=\frac{1}{|A|}
       \left(N^2+\sum_{\psi\ne1}
          \left|\sum_{R\in C(L)}\psi(a([R-P]))\right|^2\right).
\]
Each nontrivial Lang character pulls back to a geometrically
nontrivial unramified rank-one local system: the Abel map
induces the isomorphism on geometric abelian fundamental groups.
Its $H^0$ and $H^2$ vanish, and its $H^1$ has dimension
$2g-2$ by the unramified Euler characteristic. It is a summand
of the cohomology of the finite \emph{\'etale} Lang pullback,
so its Frobenius eigenvalues have complex absolute value
$\sqrt q$, by the opened Weil theorem
\citep[Theorem~1.6 and formula~1.5.1]{WeilCurves}.
The trace formula bounds the inner sum by $(2g-2)\sqrt q$.
Subtracting $N$ proves (B.CL.1). For a positive occupancy
divisible by $p$, $N_b(N_b-1)\ge(p-1)N_b$; sum over $b\ne0$.

For $g=2$, B.2 reads $[0]+\sum_R[a([R-P])]$.
The primitive-root assumption makes $\Phi_\ell$ irreducible
over $\mathbf F_p$. Failure of the test therefore makes all
its coefficients equal. If $\ell>N$, an empty coset exists
and is nonzero, since $P$ occupies the zero coset. All nonzero
occupancies would consequently be divisible by $p$.
For a nonzero Jacobian class, at most two ordered pairs of
distinct geometric points have that difference: two different
representations give the unique degree-two pencil, and the
second pair must be $(\iota S,\iota R)$ when the first is
$(R,S)$. This proves the bound in all characteristics.
At most $2(|H|-1)$ points therefore share a coset with a
different point. Under (B.CL.2) at least two points occupy
singleton cosets, one of them nonzero. Its occupancy is not
divisible by $p$, a contradiction. Finally the unbounded
quotient-existence requirement and the effective-divisor terms
for $g>2$ have not been removed by this reasoning; the stated
limitations follow from the actual hypotheses just used.
\end{proof}

\begin{claim}{B.4}{Proved}{A finite all-constant-extension decision for each fixed prime}
Fix $(C,P)/\mathbf F_q$ and an odd prime $\ell\ne p$.
There is a finite procedure deciding whether B1 has an
order-$\ell$ witness over \emph{any} constant extension.
For each geometric nonzero character
$\chi:\pi_1(C_{\overline k})\to\mathbf F_\ell$, let $m$
be its least Frobenius period. Construct its unique $P$-split
constant cover over $k_m=\mathbf F_{q^m}$. It suffices to test
these finitely many covers at these first fields. If $A_\chi$
is Frobenius on their new geometric $p$-torsion, the exact
extension degrees preserving the witness are
\[
  n\ge1\quad\text{such that}\quad
                  \det(1-A_\chi^n)\ne0.                \tag{B.CL.3}
\]
A failed first-field test never becomes a witness at a later
field. A successful one remains successful under all
extensions of $p$-power degree.
\end{claim}
\begin{proof}
The character space is the finite vector space
$H^1_{\et}(C_{\overline k},\mathbf F_\ell)$ of dimension
$2g$, with its computable invertible Frobenius action.
Its nonzero elements, or their one-dimensional quotients,
are finite in number. An invariant character extends uniquely
to the semidirect arithmetic group by requiring it to be zero
on the chosen section $s_P$. This is exactly its geometrically
connected constant cyclic $P$-split cover. An extension permits
that same constant character exactly when its degree is a
multiple of $m$. Abelian class field theory, or the Lang
construction and the Abel isomorphism on geometric abelian
fundamental groups, identifies it with a quotient of
$J_C(k_m)$. These are the covers in B1, with the pointing fixed.

Pullback and norm split their $p$-primary Tate data, because
$p\ne\ell$. The quotient abelian variety's nonzero rational
$p$-torsion is therefore exactly the fixed space of $A_\chi$.
This proves (B.CL.3) and the first-field assertion: an existing
fixed vector remains fixed by every power. Conversely, a
constant cover over $k_{mn}$ with this geometric character and
split point is the base change of the unique split cover just
constructed, so no further twist has escaped the test.
For $n=p^a$, the polynomial identity
$X^{p^a}-1=(X-1)^{p^a}$ shows that $A_\chi^n-1$ is
invertible exactly when $A_\chi-1$ is invertible.

For an effective procedure, one may construct the finitely
many covers and compute their zeta numerators by point counts,
or compute the induced Frobenius on coherent cohomology and
its stable part by exact finite-field linear algebra. The
Artin--Schreier sequence identifies that stable part with
geometric \(p\)-torsion cohomology; nilpotent directions do
not affect $\det(1-F)$. Thus the procedure terminates for
every fixed $\ell$. It supplies no bound on the primes that
must be tested for the unrestricted existence assertion B1.
\end{proof}

\begin{claim}{B.5}{Open}{The remaining prescribed-point existence assertion}
For every prescribed $(C,P)$ as in B1, prove or refute the
existence of a prime $\ell$ for which at least one first-field
test in B.4 succeeds. This remains undecided here, including
the universal assertion for pairs $(p,g)$ with
$p>(2g-2)^2$. B.3 gives the corrected numerical sufficient
conditions; B.28--B.29 prove unpointed existence and quantify
the prescribed-point obstruction. The R8 example has the
explicit least-prime witness proved in B.26--B.27, and B.30
gives a plane-quartic instance. None of these removes the
arbitrary-curve and prescribed-point quantifiers.
\end{claim}


\subsection{The fixed norm module, full sections, and the free-orbit question}

We fix a finite field $k$ of characteristic $p$, a smooth proper
geometrically connected curve $C/k$ of genus at least two, and
$P\in C(k)$. All arithmetic groups below retain the whole $G_k$.
Put
\[
 \Omega=\pi_1(C_{\bar k},\bar P),\qquad
 \Delta=\Omega^{(p')},\qquad
 N=\ker(\Omega\longrightarrow\Delta),\qquad
 V=N^{\mathrm{ab},p}.
\]
Write $F$ both for arithmetic Frobenius through $s_P$ and for its
actions, and write $\alpha$ for its action on $\Delta$. For each
open normal $H\triangleleft\Delta$, choose the compatible pointed
geometric cover $D_H\to C_{\bar k}$. Its Jacobian norm maps define
\[
 M_P=\varprojlim_{H\triangleleft_o\Delta}
                 J_{D_H}(\bar k)[p^\infty].
\]
The action of $F$ on this entire inverse system permutes its
indices and satisfies $F(gm)=\alpha(g)F(m)$. We never assume that
it commutes with the geometric deck action.

\begin{claim}{B.6}{Proved}{The integral description of the fixed norm module}
There is a natural exact sequence of compact abelian pro-$p$ groups
\[
 0\longrightarrow V\xrightarrow{F-1}V
   \xrightarrow{\varepsilon}M_P^F\longrightarrow0.
\]
For every $F$-stable normal level $H$, projection gives a continuous
surjection
\[
 M_P^F\longrightarrow J_{D_H}(k)[p^\infty].
\]
In particular $M_P^F$ need not vanish. The first term in the exact
sequence is a Tate-module inverse limit; its Frobenius invariants
vanish, whereas its Frobenius coinvariants need not vanish.
\end{claim}
\begin{proof}
The $\alpha$-stable open normal subgroups form a cofinal directed
system. A finite pointed geometric cover descends to a finite
constant extension, so its subgroup has a finite $\alpha$-orbit;
intersecting that orbit gives a stable refinement. Finite
intersections give common refinements. For such $H$, the subgroup
generated by $H$ and $s_P(G_k)$ defines a model $D_H/k$ with
compatible point $P_H\in D_H(k)$. Put
\[
 U_H=\operatorname{pr}^{-1}(H)\subset\Omega,
 \qquad T_H=T_pJ_{D_H}.
\]
The pointed Abel map identifies $U_H^{\mathrm{ab},p}$ with $T_H$,
including the characteristic prime. The exact geometric comparison
is \cite[Proposition 9.1, Lemmas 9.4--9.5, and Corollary 9.6]{MilneJac}:
for a smooth complete curve over a separably closed field, the Abel
map identifies its full abelian fundamental group with that of its
Jacobian. The characteristic-$p$ part is explicitly included, using
Artin--Schreier cohomology and its Witt-vector extensions.

One can describe the Tate module in this comparison without treating
$[p^a]$ as an étale isogeny. Quotient the source of $[p^a]$ by
the connected component of its kernel. The resulting isogeny to the
Jacobian is étale, with geometric kernel
$J_{D_H}[p^a](\bar k)$. These covers are cofinal among the
connected étale abelian $p$-covers of the Jacobian. Their inverse
limit is $T_H$, and all constructions commute with $G_k$.
For $H'\subset H$, the map induced by
$U_{H'}\hookrightarrow U_H$ is the Jacobian norm on Tate modules:
the pointed Abel maps commute with the curve morphism and its
Jacobian norm. This is an inclusion-induced map, rather than a
transfer multiplied by the covering degree.

We next prove the $F$-equivariant identification
\[
                    V\simeq\varprojlim_H T_H.                 \tag{NM.1}
\]
Let $c:N\to\Q_p/\Z_p$ be a continuous character. Its image is
finite because $N$ is compact and the target is discrete. Choose
an open normal $O\triangleleft\Omega$ with
$N\cap O\subset\ker c$. The formula
\[
                  \widetilde c(no)=c(n)\quad(n\in N,o\in O)
\]
defines a continuous character on $NO$. It is well-defined since
two decompositions differ by $N\cap O$, and is additive since
$[O,N]\subset O\cap N$. The subgroup $NO$ is open, contains $N$,
and corresponds to an open normal subgroup of $\Delta$. A stable
refinement therefore extends $c$ to some $U_H$. Conversely, a
character of $U_H$ vanishing on $N$ factors through the
pro-prime-to-$p$ group $H$ and has $p$-primary image, so it is
zero. Restriction consequently identifies
\[
 \Hom_{\mathrm{cont}}(N,\Q_p/\Z_p)
 =\varinjlim_H\Hom_{\mathrm{cont}}(U_H,\Q_p/\Z_p).
\]
Pontryagin duality proves (NM.1), with its actual norm transition
maps and $F$-action.

At a stable level set $W_H=T_H\otimes_{\Z_p}\Q_p$. Its quotient
$W_H/T_H$ is the group of geometric $p$-primary points of the
Jacobian. The fixed subgroup $T_H^F$ is zero. Indeed an invariant
compatible $p$-division sequence would lie in the finite group
$J_{D_H}(k)$; multiplying sufficiently far down that sequence by
the exponent of its $p$-primary subgroup kills every coordinate.
Thus $F-1$ is injective on the finite free module $T_H$, and
invertible on $W_H$. There is an exact sequence
\[
 0\longrightarrow T_H\xrightarrow{F-1}T_H
 \xrightarrow{\varepsilon_H} J_{D_H}(k)[p^\infty]
 \longrightarrow0,\qquad
 \varepsilon_H(a)=(F-1)^{-1}a\bmod T_H.                       \tag{NM.2}
\]
For the image assertion, a coset $w+T_H$ is $F$-fixed exactly
when $(F-1)w\in T_H$. Its fixedness is exactly rationality over
$k$. Formula (NM.2) is also the finite-field Kummer comparison,
with cocycle convention $F(\widetilde x)-\widetilde x$.

All maps in (NM.2) commute with norms. Cofinality of the stable
levels gives
\[
 M_P^F=\varprojlim_H J_{D_H}(k)[p^\infty].                    \tag{NM.3}
\]
This supplies its compact pro-$p$ topology, without imposing a
bounded exponent on $M_P$.

For completeness, inverse-limit surjectivity requires no uniform
denominator. Given a compatible family $m_H$ in (NM.3), the sets
$\varepsilon_H^{-1}(m_H)\subset T_H$ are nonempty and compact.
Any finite collection of compatibility equations can be solved by
choosing a preimage at a common refinement and taking its norms.
The product of these compact fibres therefore has the finite
intersection property for all compatibility conditions. A compatible
lift exists. If a compatible family $a_H$ maps to zero, its unique
preimages $b_H$ under $F-1$ are compatible, since the norms commute
with $F$. This proves the kernel equality; left injectivity holds
at each level. Together with (NM.1), this proves the asserted
exact sequence.

Finally, for a transition of degree $d$ prime to $p$,
$\operatorname{Nm}\circ\pi^*=[d]$. Both maps are defined over
$k$. On the finite rational $p$-primary groups, multiplication by
$d$ is an automorphism. Thus norm has the section
$d^{-1}\pi^*$ and is surjective. An inverse system of nonempty
finite fibres with these compatible maps has surjective projection
to each coordinate, again by compactness. This proves the final
surjectivity assertion.
\end{proof}

\begin{claim}{B.7}{Proved}{All fixed classes have full arithmetic section realizations}
Every $m\in M_P^F$ is the compatible $p$-adic abelianized
difference of a continuous full étale arithmetic section from
$s_P$. Its projection to the geometrically prime-to-$p$ quotient
is exactly $s_P$. If $m\ne0$, the realizing section is not
geometrically conjugate to a point section. This conclusion uses
point separation proved by the cuspidal criterion and transport
above; it does not assert free deck orbits for $m$.
\end{claim}
\begin{proof}
Choose $v\in V$ with $\varepsilon(v)=m$ and choose $n\in N$
mapping to $v$. In the full arithmetic group write
$\mathfrak f=s_P(\operatorname{Frob}_k)$ and consider
$n\mathfrak f$. Every element $a$ of a profinite group defines a
continuous homomorphism $\widehat\Z\to\overline{\langle a\rangle}$:
the ordinary integer power map extends, because its composition
with each finite quotient has a finite period. Apply this to
$a=n\mathfrak f$. Its projection to $G_k\simeq\widehat\Z$ sends
the dense generator $1$ to $1$, so the composite is the identity.
It therefore gives a continuous section $s_n$. Since $n\in N$,
its prime-to-$p$ projection is exactly $s_P$.

At every stable pointed level, $n\in U_H$ and its image in
$U_H^{\mathrm{ab},p}=T_H$ is $v_H$. Hence the abelianized
section-difference cocycle takes the value $v_H$ at Frobenius.
Its class is $v_H\bmod(F-1)T_H$, and (NM.2) identifies this with
$m_H$. These identities commute with norm and prove realization
in the whole inverse limit.

Every section whose prime-to-$p$ projection is exactly $s_P$ has
this form: its value at Frobenius is uniquely $n\mathfrak f$ for
some $n\in N$, and that value determines the section. Conjugation
by $a\in N$ replaces $n$ by $a n F(a)^{-1}$, hence changes its
image in $V$ by $(1-F)\bar a$ and leaves $m$ unchanged. We obtain
a surjection from these $N$-conjugacy classes of full sections to
$V/(F-1)V=M_P^F$. No injectivity of this abelianization map is
being asserted.

Suppose $m\ne0$ and $s_n$ were geometrically conjugate to the
section of $Q\in C(k)$. Projection and point separation give
$Q=P$. If $a\in\Omega$ conjugates $s_n$ to $s_P$, its image in
$\Delta$ centralizes the common section $s_P(G_k)$. The geometric
centralizer is trivial by \statementref{A.17}. Thus $a\in N$,
and the preceding coboundary calculation would give $m=0$, a
contradiction. The dependencies on the point-separation proof are
the exact cuspidal transport and \statementref{A.19}; the latter
may use its independent support-field proof
\statementref{A.21}.
\end{proof}

\begin{claim}{B.8}{Proved}{An exact finite criterion for free orbits at every normal level}
For an $F$-stable open normal $H\triangleleft\Delta$, let $S_H$
be the finite set of $a\in J_{D_H}(k)[p^\infty]$ such that, for
every normal $K\triangleleft\Delta$ containing $H$, the norm of
$a$ to $D_K$ has trivial stabilizer in $\Delta/K$. The quotient
$K/H$ and the curve $D_K$ need not be $F$-stable. Then the set of
fixed norm-module vectors with free orbits at every normal level
is naturally
\[
                         \varprojlim_H S_H.
\]
It is nonempty if and only if every $S_H$ is nonempty. Any such
vector is nonzero. Consequently, nonexistence of an all-level free
vector has a finite certificate of this specific form: one stable
pointed normal cover for which $S_H$ is empty.
\end{claim}
\begin{proof}
There are only finitely many intermediate normal subgroups between
$H$ and $\Delta$, so $S_H$ is a finite, explicitly specified set.
If $H'\subset H$ are stable, norm maps $S_{H'}$ into $S_H$:
its composite with each norm to $D_K$ is the norm from $D_{H'}$
to that same quotient. This proves preservation of all defining
stabilizer conditions, whether or not $K$ is stable.

An all-level free fixed vector projects to every $S_H$. Conversely,
a compatible family in these sets gives a vector of $M_P^F$ by
(NM.3). For an arbitrary normal level $K$, choose a stable
$H\subset K$; the definition of $S_H$ then proves freeness of
that coordinate. Thus the displayed inverse limit is exact.

If every $S_H$ is nonempty, impose the norm compatibility equations
on the compact product $\prod_H S_H$. A finite list is satisfied
by choosing an element at a common refinement and projecting to
the finitely many levels involved. Compactness supplies a point
of the inverse limit. This argument does not require the restricted
maps between the sets $S_H$ to be surjective. The converse is
immediate. There exists a nontrivial prime-to-$p$ geometric normal
cover, for the abelian geometric group surjects onto $\Z/r$ for
any $r\ne p$. At such a level the zero class has nontrivial
stabilizer, so an all-level free vector cannot be zero.

The criterion can also be written directly in terms of a section
representative from \statementref{B.7}. Let $v_H\in T_H$ be
its abelianized value at Frobenius and put
$w_H=(F-1)^{-1}v_H\in W_H$. For $K\supset H$ let
$\nu_{HK}:W_H\to T_K\otimes\Q_p$ be the rationalized norm.
For $g\in\Delta/K$ its action fixes the coordinate precisely when
\[
                         (g-1)\nu_{HK}(w_H)\in T_K.           \tag{NM.4}
\]
Indeed that coordinate is $\nu_{HK}(w_H)\bmod T_K$.
Formula (NM.4), tested for every $g\ne1$ and every intermediate
normal quotient, is equivalent to membership in $S_H$. It
respects the semilinear action: we have not commuted $g$ with
$F-1$, which would generally be incorrect.
\end{proof}

\begin{claim}{B.9}{Proved}{A relative class-number witness obstructs free vectors}
Suppose there are a finite extension $L/k$ and compatible pointed
étale prime-to-$p$ covers $Z\to Y\to C_L$, both geometrically
normal over $C$, such that $Z\to Y$ is nontrivial Galois and
\[
                           p\nmid\frac{\#J_Z(L)}{\#J_Y(L)}.
\]
Then $M_P^F$ contains no vector with free deck orbit at every
normal level. The geometric deck groups need not be constant over
$L$. This is a sufficient condition, not a claimed equivalence
with the existence of one cyclic witness over $C$.
\end{claim}
\begin{proof}
The chosen rational lifts identify these covers with levels in the
pointed tower whose descent is induced by a power of $s_P(F)$.
Thus the coordinate of any $m\in M_P^F$ on $Z$ belongs to
$J_Z(L)[p^\infty]$. Write $d=\deg(Z/Y)$ and
$\pi:Z\to Y$. Since $d$ is invertible on $p$-primary groups,
the identities $\operatorname{Nm}\pi^*=[d]$ give a direct-sum
decomposition of finite abelian groups
\[
 J_Z(L)[p^\infty]
 =\pi^*J_Y(L)[p^\infty]\oplus
       \ker\!\left(\operatorname{Nm}:J_Z(L)[p^\infty]
                                  \to J_Y(L)[p^\infty]\right).
\]
The order of the second summand is the $p$-part of the displayed
class-number ratio, so the hypothesis makes that summand zero.
Every coordinate is consequently a pullback from $Y$. All elements
of the nontrivial geometric group $\Gal(Z_{\bar k}/Y_{\bar k})$
fix that pullback, so it fails the required free-orbit condition
at the normal level $Z$. Norm and pullback are defined over $L$,
which is all that this argument needs; individual deck
transformations need not be defined over $L$.
\end{proof}

\begin{claim}{B.10}{Proved}{Nonzero fixed classes on an explicit genus-two curve}
For the smooth proper model over $\F_2$ of
\[
                     C:\quad y^2+y=x+\frac1x+\frac1{x+1},
\]
one has $J_C(\F_2)\simeq(\Z/2)^2$. For every rational base point
$P$, the group $M_P^F$ therefore has nonzero classes and these
have non-point-theoretic full arithmetic section realizations.
This assertion does not provide an all-level free vector.
\end{claim}
\begin{proof}
The three poles at $0,1,\infty$ are simple. The Artin--Schreier
extension is geometrically connected, since a function of the form
$b^2+b$ cannot have an odd pole order. It is étale away from these
poles. At a simple pole, $1/y$ is a local uniformizer upstairs and
the involution $y\mapsto y+1$ changes it by a quantity of valuation
two. Thus the local different exponent is two. Riemann--Hurwitz
gives $2g-2=2(-2)+3\cdot2=2$, so $g=2$. Each pole gives one
rational point, denoted $W_0,W_1,W_\infty$.

There are no other $\F_2$-points. Over
$\F_4=\F_2[t]/(t^2+t+1)$, the only additional possible
$x$-values are $t,t+1$. Their Artin--Schreier right sides have
absolute trace one, so neither gives a point. Hence $N_1=N_2=3$.
The genus-two Frobenius identities, or
\cite[Theorem 11.1 and Corollary 11.4]{MilneJac}, then give
\[
              P_C(T)=1-T^2+4T^4,\qquad \#J_C(\F_2)=P_C(1)=4.
\]
The divisor identities
\[
 \operatorname{div}(x)=2(W_0-W_\infty),\qquad
 \operatorname{div}(x+1)=2(W_1-W_\infty)
\]
give two distinct nonzero classes of order two. Nonzero and
distinct follow because the difference of two distinct points on
a positive-genus curve cannot be principal. Thus these classes
give all of $J_C(\F_2)\simeq(\Z/2)^2$.

The retained exact computation in
\path{code/criterion_norm_module.py} checks the field arithmetic,
both point counts, and the Frobenius polynomial. Now apply
\statementref{B.6} and \statementref{B.7}. All actual point
sections are distinct by the point-separation theorem; the new
sections constructed from a nonzero fixed class are different
from all of them.
\end{proof}

\begin{claim}{B.11}{Refuted}{Freeness at one normal level need not survive norm}
The implication from a free vector in a finite coprime deck module
to a free norm at every normal quotient is false, even for an
abelian odd group acting on an $\F_2$-vector space.
\end{claim}
\begin{proof}
Take $G=(\Z/3)^2$, $\F_4=\F_2[\omega]$ with
$\omega^2+\omega+1=0$, and $V=\F_4\oplus\F_4$. Let
\[
          (a,b)(u,v)=(\omega^a u,\omega^b v).
\]
The vector $(1,1)$ has trivial stabilizer. For the diagonal
subgroup $H=\langle(1,1)\rangle$, its norm is zero, since
$1+\omega+\omega^2=0$ in both summands. The quotient $G/H$ is
nontrivial and fixes this zero norm. This is an exact finite-module
counterexample to that inference; it is not claimed to be the
Jacobian module of a particular curve. The program just cited
enumerates the nine group elements, the orbit, and the norm.
It explains why the finite sets in \statementref{B.8} must test
every normal quotient of each chosen cover.
\end{proof}

\begin{claim}{B.12}{Open}{The universal free-orbit assertion of B2}
It remains undetermined here whether $M_P^F$ has no nonzero
all-level free vector for every genus-two curve in characteristic
two over every finite field of definition, and whether an
all-level free vector exists for any proper hyperbolic curve in
any characteristic. The general fixed classes and their full
section realizations have been determined above, but these results
do not decide nonemptiness of the inverse limit in
\statementref{B.8} for all such curves.
\end{claim}
\paragraph{The exact remaining step.}
To prove the universal nonexistence statement one must produce,
for each pointed curve, a stable normal level $H$ with
$S_H=\varnothing$. A witness as in \statementref{B.9} suffices,
but point separation supplies no such witness: it excludes classes
represented compatibly by point differences, whereas the
surjection in \statementref{B.7} produces many sections which
are not point sections. To refute the statement one must instead
prove $S_H\ne\varnothing$ for every stable normal level of one
fully specified pointed curve, not merely compute regular vectors
at a finite list of levels. Compactness would then give a genuine
all-level counterexample. Either sort of all-level free vector,
if it exists, automatically has a non-point-theoretic section
realization by \statementref{B.7}; none can come from a pair of
distinct points after the main point-separation theorem.


\subsection{An exact class-number criterion at every finite abelian norm level}

\begin{claim}{B.13}{Proved}{Finite abelian norm obstructions, including nonconstant deck groups}
Let $D\to C$ be a pointed geometrically connected finite
\emph{\'etale} cover over $k=\mathbf F_q$, with abelian geometric
Galois group $A$ of order prime to $p=\operatorname{char}k$.
The chosen point of $D$ lies rationally above $P$.
Let $F$ act on $A$ by the pointed arithmetic descent, and put
\[
 S_D(k)=\{a\in J_D(k)[p^\infty]:
    \operatorname{Nm}_{D/D_H}(a)
    \text{ has trivial stabilizer in }A/H
                   \text{ for every }H\le A\},                 \tag{B.AB.1}
\]
where $D_H=D_{\overline k}/H$, with its pointed descent whenever
the subgroup is stable under a power of $F$.

For each subgroup $H$ for which $A/H$ is a nontrivial cyclic
prime-power group $\mathbf Z/\ell^a$, let $m_H$ be its orbit
length under $F$, and let $H^+/H$ be the unique subgroup of
order $\ell$ in $A/H$. Both quotient curves have their pointed
models over $k_{m_H}=\mathbf F_{q^{m_H}}$. Then
\[
 S_D(k)\ne\varnothing
 \quad\Longleftrightarrow\quad
 p\mid
 \frac{\#J_{D_H}(k_{m_H})}{\#J_{D_{H^+}}(k_{m_H})}
 \quad\text{for every such }H.                                \tag{B.AB.2}
\]
Thus an empty set at a finite abelian normal level is equivalent
to a relative class-number unit in one cyclic prime-power tower,
at that quotient's first pointed field of definition. The step
$D_H\to D_{H^+}$ has prime degree $\ell$; its lower curve need
not be $C$.
\end{claim}
\begin{proof}
We use the geometric $p$-primary torsion group
\[
                   T=J_D(\overline k)[p^\infty].
\]
For a subgroup $H\le A$, the element
\[
                   e_H=\frac1{|H|}\sum_{h\in H}h
                         \in\mathbf Z_p[A]
\]
is an idempotent. Pullback from $D_H$ is injective on
$p$-primary torsion, and pullback followed by norm is
$\sum_{h\in H}h=|H|e_H$. Its image is $e_HT$; this follows
also by applying norm and dividing by $|H|$. Consequently
the stabilizer in $A$ of $e_Ha$ is the inverse image of the
stabilizer of the norm coordinate in $A/H$. Freeness at this
quotient means exactly
\[
                         \operatorname{Stab}_A(e_Ha)=H.         \tag{B.AB.3}
\]

Here is the needed character decomposition, including its
integrality. Over a finite unramified extension of $\mathbf Z_p$
containing all roots of unity of exponent $\exp(A)$, the
idempotents of the characters $\chi$ diagonalize $A$.
For each subgroup $K$ with $A/K$ cyclic, sum these idempotents
over all characters whose kernel is exactly $K$, and denote
the result by $\varepsilon_K$. That sum is invariant under
the unramified Galois group and hence belongs to
$\mathbf Z_p[A]$. The $\varepsilon_K$ are pairwise orthogonal
and sum to one. Moreover
\[
       e_H=\sum_{K\supset H}\varepsilon_K.                     \tag{B.AB.4}
\]
Every nonzero element of $\varepsilon_KT$ has stabilizer
exactly $K$. To see this after the faithful unramified scalar
extension, if $g\notin K$, every corresponding value
$\chi(g)-1$ is a unit: a nontrivial root of unity of order
prime to $p$ remains nontrivial in the residue field. Thus
$g-1$ is invertible on that entire summand. This reasoning
holds for arbitrary finite $p$-power torsion, not only for
torsion killed by $p$.

If $A/H$ is cyclic of order $\ell^a$, its primitive packet is
particularly simple:
\[
                         \varepsilon_H=e_H-e_{H^+}.            \tag{B.AB.5}
\]
Every strict overgroup of $H$ contains $H^+$. Hence (B.AB.3)
forces $\varepsilon_Ha\ne0$: otherwise the whole norm
coordinate would be fixed by $H^+/H$. Conversely, suppose
$a$ has nonzero primitive component at \emph{every} subgroup
with cyclic prime-power quotient. For an arbitrary $K\le A$
and $g\notin K$, a prime-primary character of the finite
abelian group $A/K$ detects $g$. Its kernel lifts to a subgroup
$H\supset K$ with $A/H$ cyclic of prime-power order and
$g\notin H$. The nonzero component $\varepsilon_Ha$ survives
in $e_Ka$, and $g$ cannot fix it. Thus (B.AB.3) holds for
every $K$. We have proved the exact equivalence
\[
 a\text{ satisfies all norm-freeness conditions}
 \quad\Longleftrightarrow\quad
 \varepsilon_Ha\ne0\text{ for every cyclic prime-power }A/H.
                                                                  \tag{B.AB.6}
\]

We next impose Frobenius, without assuming that the deck group
is constant. Conjugation carries $\varepsilon_H$ to
$\varepsilon_{F(H)}$. If $a\in T^F=J_D(k)[p^\infty]$,
then $\varepsilon_Ha$ is fixed by $F^{m_H}$. Necessity in
(B.AB.6) therefore requires
\[
                         (\varepsilon_HT)^{F^{m_H}}\ne0.        \tag{B.AB.7}
\]
Conversely assume (B.AB.7) for every indicated $H$.
Choose one representative in each Frobenius orbit of these
subgroups and a nonzero $b_H$ in the displayed fixed group.
Then
\[
                  a_H=\sum_{i=0}^{m_H-1}F^i b_H
\]
is fixed by $F$. Its components in the distinct packets
$\varepsilon_{F^i(H)}T$ are all nonzero. Summing $a_H$ over
the finitely many subgroup orbits gives a rational torsion
class satisfying (B.AB.6). No division by $m_H$ occurs;
the proof is valid even when $p\mid m_H$.

Finally both $H$ and $H^+$ are stable under $F^{m_H}$, and
their quotient maps and rational lifts of $P$ descend to
$k_{m_H}$. Put
\[
 B_H=\bigl(\ker(J_{D_H}\longrightarrow J_{D_{H^+}})\bigr)^0.
\]
Pullback identifies $B_H[p^\infty]$ with the primitive packet
$\varepsilon_HT$, compatibly with $F^{m_H}$. Norm and
pullback split the $p$-primary groups because their degree
$\ell$ is prime to $p$. The isogeny
$J_{D_{H^+}}\times B_H\to J_{D_H}$ shows
\[
 \#B_H(k_{m_H})
       =\frac{\#J_{D_H}(k_{m_H})}{\#J_{D_{H^+}}(k_{m_H})}.
\]
Thus (B.AB.7) is exactly divisibility by $p$ of the indicated
integer. This proves (B.AB.2).
\end{proof}

\begin{claim}{B.14}{Proved}{The elementary-prime case and the stronger higher-power test}
If the group $A$ in \statementref{B.13} is constant, then all
$m_H=1$. If in addition $A=(\mathbf Z/\ell)^d$, then
$S_D(k)$ is nonempty exactly when every degree-$\ell$ quotient
$D_H\to C$ has relative class number divisible by $p$.
Thus at a constant elementary-abelian level, an obstruction
is exactly a prime-degree class-number witness among those
quotients.

For higher-exponent abelian groups, the exact criterion also
tests the primitive parts at cyclic prime-power quotients.
These tests are strictly stronger as finite module conditions
than testing only prime-degree quotients. This strictness does
not assert that a particular curve has no prime-degree witness
over all constant extensions.
\end{claim}
\begin{proof}
For elementary $A$, every nontrivial cyclic quotient has degree
$\ell$, and $H^+=A$. The relative factor in (B.AB.2) is
therefore exactly that of $D_H/C$. This proves the first
assertion and identifies the relationship with B1 at the
displayed finite field and level.

Here is a complete finite-module control for the higher-power
assertion. Take characteristic two and
\[
        A=\mathbf Z/9\times\mathbf Z/3,
                     \qquad M=\mathbf F_4^4.
\]
Let $\omega\in\mathbf F_4$ have order three. An element
$(a,b)$ acts diagonally on the four summands with scalars
\[
          \omega^a,\qquad\omega^b,\qquad
          \omega^{a+b},\qquad\omega^{a+2b}.                    \tag{B.AB.8}
\]
These are all four character kernels of prime-degree
quotients. A vector with all four coordinates nonzero has
free norm orbit at every such quotient. There are $3^4=81$
such vectors. Nevertheless the subgroup
\[
                   \{(0,0),(3,0),(6,0)\}
\]
fixes every element of $M$, so no vector has free orbit at
the full level. The three missing primitive packets have
cyclic quotient of order nine.

For comparison, choose a primitive ninth root
$\zeta\in\mathbf F_{64}$ and add three summands
$\mathbf F_{64}$ with characters
\[
                \zeta^{a+3jb},\qquad j=0,1,2.
\]
Now all seven cyclic prime-power kernels are present.
The vector with every coordinate one has free norm orbit
at all ten subgroups of $A$. In fact the exact number of
all-quotient free vectors in this module is
\[
                       3^4\cdot63^3=20\,253\,807.
\]
The packet criterion proves the formula; the accompanying
program checks every subgroup and every norm for the displayed
vector. It exhausts the $256$ vectors of the first module and
independently verifies its count $81$ and its absence of an
all-quotient free vector.

There is also a check of the nonconstant Frobenius step when
the orbit length is divisible by the characteristic. Set
$A=(\mathbf Z/3)^2$, retain (B.AB.8), and define
\[
 F(a,b)=(b,a),\qquad
             F(u,v,w,z)=(v,u,w,z^2).
\]
These actions satisfy $F(gm)=F(g)F(m)$. The fixed module
has $32$ vectors, of which exactly nine have free norm orbit
at every subgroup. The first two character kernels form an
orbit of length two, and the other two are stable. The sum
\[
                  (1,0,0,0)+F(1,0,0,0)=(1,1,0,0)
\]
is fixed and is nonzero in both moving packets, despite
working in characteristic two. Thus the Frobenius-orbit
argument in \statementref{B.13} does not hide an averaging
division by an inadmissible integer.

All finite-field operations, subgroup enumeration, norms,
stabilizers and counts are independently implemented in
\path{code/criterion_abelian_norm.py}; its complete results
are in the adjacent JSON certificate. These are finite
representation examples, with no claimed Jacobian realization.
The universal genus-two characteristic-two assertion in B2
still requires an empty level for every pointed curve. The
present theorem exactly reduces every \emph{abelian} candidate
level to the displayed relative class-number tests; it does
not settle the existence of such a level, or rule out a
stronger obstruction at a nonabelian level.
\end{proof}


\subsection{Solvable deck groups: a single common-stabilizer obstruction}

\begin{claim}{B.15}{Proved}{Simultaneous norm freeness for solvable coprime actions}
Let $G$ be a finite solvable group of order prime to $p$, acting
on a finite abelian $p$-group $M$. Let $F$ be an automorphism
of $M$, and let $\sigma\in\operatorname{Aut}(G)$ satisfy
$F(gm)=\sigma(g)F(m)$. For $K\triangleleft G$ put
\[
                       e_K=|K|^{-1}\sum_{g\in K}g.
\]
There exists $a\in M^F$ with
\[
                 \operatorname{Stab}_G(e_Ka)=K
                      \quad(K\triangleleft G)                 \tag{B.SV.1}
\]
if and only if
\[
        \bigcap_{b\in M^F}\operatorname{Stab}_G(e_Kb)=K
                      \quad(K\triangleleft G).                \tag{B.SV.2}
\]
Neither $M^F$ nor an individual norm image is assumed to be
$G$-stable. The order of $F$ or of $\sigma$ may be divisible
by $p$.
\end{claim}
\begin{proof}
The implication from (B.SV.1) to (B.SV.2) is immediate.
We prove the converse by induction on the derived length
of $G$. If $G=1$, choose $a=0$. Otherwise let $N$ be the
last nontrivial derived subgroup of $G$. It is nontrivial,
abelian, normal, and characteristic, so it is $\sigma$-stable.
There is a $G$- and $F$-stable direct sum
\[
                    M=e_NM\oplus W,\qquad W=(1-e_N)M.          \tag{B.SV.3}
\]
The first summand is a module for $G/N$. Its fixed group is
$e_N(M^F)$, since the idempotent commutes with $F$.
For every normal subgroup $K\supset N$, its common
stabilizer after $e_K$ is $K/N$ by (B.SV.2). The induction
hypothesis supplies $v\in(e_NM)^F$ such that
\[
                \operatorname{Stab}_G(e_Kv)=K
                        \quad(K\triangleleft G,\ N\subset K).
                                                                  \tag{B.SV.4}
\]

We need one simultaneous support lemma for the abelian group
$N$. For each subgroup $L\le N$ with $N/L$ cyclic, let
$\varepsilon_L^N\in\mathbf Z_p[N]$ be the sum of the
character idempotents with exact kernel $L$, including the
trivial character when $L=N$. As in \statementref{B.13},
these are integral orthogonal idempotents summing to one;
every nonzero element in their $L$-summand has stabilizer
exactly $L$ in $N$. The reason is that $\chi(n)-1$ is a
unit for every nontrivial prime-to-$p$ root $\chi(n)$.

All the $e_K$, for $K\triangleleft G$, are central in
$\mathbf Z_p[G]$, hence commute with each other and with
these $N$-idempotents. Form their simultaneous orthogonal
idempotents by taking all products
\[
 \varepsilon_L^N
       \prod_{K\in\mathcal S}e_K
       \prod_{K\notin\mathcal S}(1-e_K),                       \tag{B.SV.5}
\]
where $\mathcal S$ ranges through subsets of the finite set
of normal subgroups of $G$, and discard zero summands.
They decompose $W$. On each resulting summand, every $e_K$
is either zero or the identity, and every nonzero vector has
one fixed stabilizer $L$ in $N$.

The automorphism $F$ permutes these summands: it carries
$e_K$ to $e_{\sigma(K)}$ and $\varepsilon_L^N$ to
$\varepsilon_{\sigma(L)}^N$. Decompose $W$ further into
the direct sums belonging to these $F$-orbits. In each such
orbit whose fixed subgroup is nonzero choose one nonzero
fixed vector. Its projection to every summand in that orbit
is nonzero, since $F$ cyclically identifies all projections.
Sum the chosen vectors over the finitely many available
orbits. The resulting $w\in W^F$ has a nonzero component
in exactly every joint summand that occurs in any fixed
vector. Consequently, simultaneously for every normal $K$,
\[
 \operatorname{Stab}_N(e_Kw)
       =\bigcap_{b\in W^F}\operatorname{Stab}_N(e_Kb).           \tag{B.SV.6}
\]
Indeed $e_K$ retains a specified collection of the joint
summands; on each retained nonzero component the stabilizer
is its exact $N$-character kernel. Thus both sides are the
intersection of the same kernels. This argument uses no
averaging over $F$ and works for arbitrary $p$-power torsion.

The summand $e_NM$ is fixed by $N$, so (B.SV.2) and
(B.SV.3) identify the right side of (B.SV.6) with $N\cap K$.
Set $a=v+w$. If $g\in G$ fixes $e_Ka$, the $G$-stable
decomposition (B.SV.3) shows that it fixes $e_Kv$ and
$e_Kw$ separately. Since
\[
                         e_Kv=e_{KN}v,
\]
equation (B.SV.4) implies $g\in KN$. Write $g=nk$, with
$n\in N$ and $k\in K$. The group $K$ fixes every $e_K$
image. Hence $g$ fixes $e_Kw$ exactly when $n$ does.
Equation (B.SV.6) gives $n\in N\cap K$, so $g\in K$.
The reverse inclusion is automatic. This proves (B.SV.1)
for every normal $K$ and completes the induction.
\end{proof}

\begin{claim}{B.16}{Proved}{Constant solvable covers have exactly the relative-unit obstruction}
Let $D\to C$ be a constant pointed normal prime-to-$p$
cover over a finite field $k$, with solvable deck group $G$.
The finite all-normal-quotient norm-free set $S_D(k)$ is empty
if and only if there exist normal subgroups
$K\subsetneq L\triangleleft G$ such that
\[
                         p\nmid
                   \frac{\#J_{D_K}(k)}{\#J_{D_L}(k)}.          \tag{B.SV.7}
\]
In particular, for constant solvable deck groups, the relative
class-number-unit sufficient criterion is also necessary at
that finite level. The lower curve $D_L$ can differ from $C$,
and the degree of the step is not required to be prime.
\end{claim}
\begin{proof}
Apply \statementref{B.15} with
$M=J_D(k)[p^\infty]$ and $F=1$; the actual arithmetic
Frobenius has already been imposed by taking rational points.
This is a $G$-module because the deck group is constant.
The common stabilizer $L$ of $e_KM$ is a normal subgroup of
$G$, since $e_KM$ is $G$-stable, and contains $K$.

Pullback identifies $J_{D_K}(k)[p^\infty]$ with $e_KM$:
the averaging idempotent is defined over $k$, and pullback
and norm split because $|K|$ is prime to $p$.
If $L\supsetneq K$ is its common stabilizer, all of this
group is invariant under $L/K$. Pullback from $D_L$ is
therefore all of $J_{D_K}(k)[p^\infty]$, again by averaging.
The $p$-parts of the two Jacobian orders agree, giving
(B.SV.7). Conversely that relative unit makes every rational
$p$-primary class on $D_K$ a pullback from $D_L$, so its
stabilizer contains the nontrivial group $L/K$. No class can
satisfy all normal quotient conditions. The equivalence in
\statementref{B.15} proves that these are the only ways the
finite set can be empty.
\end{proof}

\begin{claim}{B.17}{Proved}{A single-element finite certificate in characteristic two}
For any proper hyperbolic pointed curve $(C,P)/k$ in
characteristic two, $M_P^F$ has no all-level free vector if
and only if there are a pointed $F$-stable normal cover
$D\to C$, a normal subgroup $K$ of its geometric deck group
$G$, and an element $g\in G\setminus K$ such that
\[
 g\,\operatorname{Nm}_{D/D_K}(a)
       =\operatorname{Nm}_{D/D_K}(a)
        \qquad\text{for every }a\in J_D(k)[2^\infty].          \tag{B.SV.8}
\]
The quotient $D_K$ need only be defined over its pointed
finite constant extension. Thus in characteristic two a
finite obstruction can always be witnessed by one quotient
and one element fixing every candidate norm coordinate.
It cannot require only a union of different stabilizers
covering the candidates.

This reformulation does not prove that such a finite
certificate exists for every genus-two curve. That universal
existence assertion in B2 remains open here.
\end{claim}
\begin{proof}
The geometric deck group of a prime-to-two cover has odd
order. The exact theorem of Feit--Thompson is that every
finite group of odd order is solvable; these are its only
hypotheses \citep[Chapter~I, \S1, p.~775, unnumbered main
theorem]{FeitThompson1963}. Hence \statementref{B.15} applies
to every finite level under consideration.

Choose $2^e$ killing the finite group $J_D(k)[2^\infty]$
and apply that theorem to
\[
                      M=J_D[2^e](\overline k).
\]
The geometric deck group acts on $M$; arithmetic Frobenius
normalizes that action and has fixed group exactly
$J_D(k)[2^\infty]$. Pullback identifies the stabilizer of
each norm coordinate with the stabilizer of its averaging
idempotent, so \statementref{B.15} says that $S_D(k)$
is empty precisely when (B.SV.8) holds for some $K,g$.
By the compact finite criterion \statementref{B.8}, absence
of an all-level free vector is equivalent to emptiness at
one stable finite level. This proves the stated equivalence.

The distinction between this statement and
\statementref{B.16} matters. With nonconstant deck group,
the Frobenius-fixed group need not be $G$-stable, so the
common stabilizer in (B.SV.8) need not be normal and cannot
automatically be interpreted as a relative class-number unit
over $k$. The semilinear common-stabilizer criterion is the
assertion proved in this generality.
\end{proof}

\paragraph{Exhaustive finite checks of the induction.}
The program \path{code/criterion_solvable_norm.py} verifies the
theorem on a nonabelian solvable group with a Frobenius of order
divisible by the characteristic. Over $\mathbf F_3$, put
\[
 i=\begin{pmatrix}0&1\\2&0\end{pmatrix},\qquad
 j=\begin{pmatrix}1&1\\1&2\end{pmatrix},\qquad
 T=\begin{pmatrix}1&1\\0&1\end{pmatrix}.
\]
The first two generate $Q_8$, since $i^2=j^2=-1$ and
$ij=-ji$. The matrix $T$ has order three and normalizes this
group. Let $V=\mathbf F_3^2$ be the natural module and add
the three nontrivial one-dimensional characters of
$Q_8/\{\pm1\}$, giving $M=V\oplus\mathbf F_3^3$.
Define Frobenius by $T$ on $V$ and by the induced cyclic
permutation on the three character summands. It satisfies
the semilinear compatibility in \statementref{B.15}.

The program enumerates all six normal subgroups, all $243$
vectors, and every norm and stabilizer. With $F=1$, exactly
$64$ vectors satisfy all norm-free conditions. With the
displayed order-three Frobenius, exactly nine vectors are
fixed and four satisfy all conditions; every common-stabilizer
test is exact. On the natural module alone and on the three
character summands alone, the all-level set is empty, and
the program identifies a failed common-stabilizer condition
in each case, both with constant and nonconstant Frobenius.
These finite representation checks are recorded in
\path{code/criterion_solvable_norm.json}; no geometric
realization is asserted by the checks.


\subsection{A genuinely nonabelian finite obstruction mechanism}

\begin{claim}{B.18}{Proved}{A finite norm obstruction with no vanishing normal-step module}
The equivalence of \statementref{B.13} cannot be extended to
arbitrary finite groups using only the vanishing of relative
$p$-primary modules at normal quotient steps. There is a
faithful action of a finite nonabelian simple $7'$-group $G$
on a nonzero finite $7$-group $M$ such that the set of vectors
with free norm orbit at every normal quotient is empty,
although the relative module at every nontrivial normal step
is nonzero. This is an exact finite representation example,
not a claimed realization by a Jacobian or a counterexample
to universal B1.
\end{claim}
\begin{proof}
Let $M=\mathbf F_7^5$, and define $G\le\operatorname{GL}(M)$
by the three matrices
\[
 U=\operatorname{diag}(1,4,1,4,4),\qquad
 P=\begin{pmatrix}
 1&0&0&0&0\\
 0&1&0&0&0\\
 0&0&0&0&1\\
 0&0&1&0&0\\
 0&0&0&1&0
 \end{pmatrix},\qquad
 Q=\begin{pmatrix}
 2&4&4&4&4\\
 2&5&4&5&5\\
 2&4&5&5&5\\
 2&5&5&5&4\\
 2&5&5&4&5
 \end{pmatrix}.                                                \tag{B.NA.1}
\]
Here $U^3=P^3=Q^2=1$, and $UP\ne PU$.
The following finite certificate proves both the order and
simplicity needed in the argument, without assuming an
identification with a named finite classical group.

Starting with the identity, close the set of matrices under
right multiplication by $U,P,Q$. Exact arithmetic modulo
seven gives
\[
                        |G|=25920=2^6\cdot3^4\cdot5.
\]
Conjugation by the same generators partitions this complete
matrix list into twenty classes, with sizes
\[
\begin{gathered}
 1,40,40,45,240,270,360,360,480,540,720,720,1440,\\
                   2160,2160,2160,2880,2880,3240,5184.
\end{gathered}                                                   \tag{B.NA.2}
\]
For each nonidentity class representative $a$, impose the
equivalence relation $1\sim a$ on the matrix list and close
it under left and right multiplication by $U,P,Q$. This is
the group congruence whose identity class is the normal
closure of $a$. In all nineteen cases its identity class
has all $25920$ elements. Every nontrivial normal subgroup
contains one of these representatives up to conjugacy, so
$G$ is simple.

The data are fully reproducible. The generating program saves
every matrix in
\path{code/criterion_nonabelian_norm_group.json}. The independent
checker \path{code/criterion_nonabelian_norm_verify.py} reads
the three generators and regenerates the complete group using
a separate general matrix multiplication routine. It checks
the conjugacy partition and the normal closures just described.
Congruence closure uses a finite union-find calculation: after
a new pair is identified, its six images under left and right
generator multiplication are queued. This computes the least
invariant equivalence relation and therefore the asserted
normal closure. The checker also verifies that the common
fixed space of the three generators has dimension zero and
that $\sum_{g\in G}g$ is the zero operator on $M$.

The action is faithful by its definition as a matrix group,
and it is coprime because $7\nmid25920$. Nevertheless
\[
                         |M|=7^5=16807<|G|.
\]
No orbit in $M$ can therefore have size $|G|$. In particular
the norm-freeness condition already fails at the identity
subgroup, where the norm coordinate is the vector itself.

On the other hand, simplicity leaves only the normal
subgroups $1$ and $G$. The old module at the latter level is
$M^G=0$, and the relative module for the sole nontrivial
normal step is the entire nonzero module $M$. Its order is
$7^5$, not a unit. Thus there is no vanishing relative module
at any nontrivial normal step, yet the full finite norm-free
set is empty. The obstruction is the union of nontrivial
stabilizers, rather than one absent abelian character packet.

For an actual cover in B2, the separate task is to determine
whether its Frobenius-fixed Jacobian module has such a
configuration. This example proves that a search for
nonabelian finite obstructions has a mechanism not captured
by the relative class-number-unit sufficient criterion. It
does not supply that geometric realization or settle the
universal assertion.
\end{proof}


\subsection{Point-isolating covers and exact constructions}

\begin{claim}{B.19}{Proved}{A pointed cover isolating one rational point}
For every $C/k$ as in \statementref{A.1} and every $P\in C(k)$,
there is a connected finite \emph{\'etale} cover $D\to C$ with
geometrically Galois group of order prime to $p$, and a specified
$P'\in D(k)$ over $P$, such that the image of $D(k)$ in $C(k)$ is
exactly $\{P\}$. The geometric deck group need not be constant over $k$.
\end{claim}
\begin{proof}
Split $\Pi_C=\Delta_C\rtimes s_P(G_k)$. For $Q\ne P$, the equations
expressing geometric conjugacy of $s_Q$ and $s_P$ have no solution
in the compact group $\Delta_C$, by \statementref{A.1}. If they had
a solution modulo every open characteristic subgroup $N$, the
closed solution sets in $\Delta_C$ would have the finite-intersection
property and yield a solution in $\Delta_C$. Thus some such $N_Q$
detects the failure. We can take characteristic subgroups cofinally:
a finitely generated profinite group has only finitely many open
subgroups of any prescribed index, and their finite intersection
is characteristic.

Intersect the finitely many $N_Q$, for $Q\in C(k)-\{P\}$, to obtain
$N$. The subgroup $N\rtimes s_P(G_k)$ defines the desired pointed
cover. A rational lift of $Q$ would give a fixed coset in its finite
fibre, equivalently a conjugacy of the two section subgroups modulo
$N$, contradicting the choice of $N_Q$. Its geometric deck group
is $\Delta_C/N$, of prime-to-$p$ order. This proof uses precisely
the relative exactness, boundary transport, and cuspidal criterion
audited in Part A; it gives no class-number unit assertion.
\end{proof}

\begin{claim}{B.20}{Proved}{A constructive two-stage odd Lang criterion}
Let $C/L$ be proper of genus at least two over a finite field of
characteristic two, $P\in C(L)$, and let $D\to C$ be a connected
constant cyclic odd-degree \emph{\'etale} cover with a point
$P_D\in D(L)$ above $P$. Suppose $h(D)/h(C)$ is odd.
Let $E\to D$ be the Lang cover associated with the maximal odd
quotient of $J_D(L)$ and based at $P_D$. Then no point of $C(L)$
other than $P$ lifts rationally to $E$. Taking the geometric
Galois closure and its $P$-pointed twist preserves this property
and gives an odd geometric Galois cover of $C$.
\end{claim}
\begin{proof}
For $G=\operatorname{Gal}(D/C)$ of odd order, the idempotent
$e=|G|^{-1}\sum_{g\in G}g$ acts on the two-primary rational torsion
of $J_D$. The invariant part is the image of the pullback from
$J_C$: pullback and norm are inverse after division by $|G|$.
The complementary abelian factor has rational point count
$h(D)/h(C)$, since the rational Frobenius polynomial factors as
the product of the base and complementary polynomials. Therefore
the oddness assumption says that all rational two-primary points
of $J_D$ are fixed by $G$.

If $Q_D\ne P_D$ lies above $Q\ne P$, its difference class cannot be
fixed by a nonidentity odd deck transformation $a$. Otherwise
$aQ_D+P_D\sim Q_D+aP_D$. These are different effective degree-two
divisors, so the curve is hyperelliptic and uniqueness of its
degree-two pencil gives $aQ_D=\iota P_D$ and $aP_D=\iota Q_D$.
Every automorphism commutes with the hyperelliptic involution
$\iota$, whence $a^2$ fixes $Q_D$. A free deck action of odd order
makes this impossible. Consequently $[Q_D-P_D]$ is not
two-primary. In a Lang cover the rational lift criterion is
exactly vanishing of this class in the chosen quotient of
$J_D(L)$: the Lang endomorphism is $F-1$, with kernel $J_D(L)$.
Thus $Q_D$ does not lift to $E$.

The two stages have constant odd Galois groups. The geometric
Galois closure of their composite embeds in
$\operatorname{Gal}(E/D)^{|G|}\rtimes G$, hence still has odd order.
The geometric core of the pointed subgroup is stable under the
action of $s_P(G_L)$, so its pointed twist has an $L$-point over
$P$ and refines the original cover. A refinement cannot acquire
a rational lift absent from the original cover.
\end{proof}

\begin{claim}{B.21}{Proved by exact computation}{The two requested genus-two classes over $\mathbf F_4$}
Write $\alpha^2+\alpha+1=0$. The two conjugate curves
\[
 C_\alpha:\quad Y^2+(x^2+x+\alpha)Y
   =x^5+\alpha x^3+x+1
\]
are the curves denoted by the pair of class labels $203/302$ in
the supplied return. Each has seven $\mathbf F_4$-points.
For every one of these fourteen pointed cases, the file
\path{code/isolating_lang_covers.json} specifies a connected
pointed cover over $\mathbf F_4$, with odd geometric Galois group,
whose rational-point image is just the chosen point.
The representation is an exact Lang-cover and geometric-core
construction; it does not expand enormous affine equations.
\end{claim}
\begin{proof}
The program \path{code/isolating_lang_covers.py} works in
$L=\mathbf F_{64}=\mathbf F_2[t]/(t^6+t+1)$, with elements encoded
by polynomial bits; the two roots $\alpha$ have codes $58,59$.
It implements polynomial Euclidean division and the full Mumford
addition and reduction formulas for this equation, checking the
defining divisibility after every addition. Direct counts are
\[
 \#C(\mathbf F_4)=7,\qquad \#C(\mathbf F_{16})=15,\qquad
 \#C(\mathbf F_{64})=91.
\]
They give the Frobenius polynomial
$X^4+2X^3+X^2+8X+16$ and
$\#J_C(L)=6076=4\cdot7^2\cdot31$. The two independent classes
specified in the JSON generate the full $7$-primary subgroup.
Its 49-element addition table is checked against
$\mathbf F_7^2$ and all $49^3$ associativity triples.
Multiplication by $124$ projects rational point classes onto it.

For every original point $P$, all eight quotient kernels
$\chi:\mathbf F_7^2\to\mathbf F_7$ are tested. If $N_b$ is the
computed occupancy of the coset $b$ for the $P$-based Abel map,
and $c=\chi([P-\infty])$, Riemann--Roch gives exactly
\[
 L_P(\chi,T)=1+\left(\sum_{b=0}^6N_b\zeta_7^b\right)T
                    +64\zeta_7^{-2c}T^2 .
\]
The integer resultant of $L_P(\chi,1)$ and
$\Phi_7$ is the relative class number. The selected relative
values are
\[
 46713473689,\qquad120974743723,\qquad103441382829,
\]
all odd. The full occupancy vectors, character coordinates,
base-point twists, and every attempted value are retained, not
just their residues modulo two.

For $\alpha=58$ the selected characters in the recorded basis are
$(1,0)$ at $\infty,(58,1),(58,58)$ and $(1,1)$ at the other four
original points. For $\alpha=59$ they are $(1,0)$ at
$\infty,(59,1),(59,59)$ and $(1,1)$ at the other four points.
The odd degrees $m$ of the second Lang stage are respectively
\[
 70957766533591,\qquad183760635715237,\qquad157127460517251.
\]
The preceding theorem isolates $P$ over $L$.

To specify the descent, use the splitting $s_P$ of the geometric
fundamental extension over $\mathbf F_4$. Let $N$ be the geometric
core for the two-stage cover over $L$ and let $F$ be conjugation
by $s_P(\operatorname{Frob}_4)$. It satisfies $F^3N=N$. Put
$N_0=N\cap F(N)\cap F^2(N)$ and take the cover corresponding to
$N_0\rtimes s_P(G_{\mathbf F_4})$. It is an exact construction
from the computed Lang data, stable under $F$, and still of odd
geometric Galois order. After extension to $L$ it refines the
isolating cover, so it isolates $P$ already over $\mathbf F_4$.
Its degree is at most $(7m^7)^3$; we do not claim equality.
\end{proof}

\begin{claim}{B.22}{Proved}{A direct isolating cover for the curve over $\mathbf F_{2^{38}}$}
For the curve and point $P=\infty$ specified in
\statementref{B.23} below, put $m=2^{19}+1=524289$.
The pullback of $[m]:J_C\to J_C$ by
$i_P:C\to J_C$, $Q\mapsto[Q-P]$, is a connected
\emph{\'etale} Galois cover of degree $m^4$ over
$\mathbf F_{2^{38}}$. Its constant deck group is
$(\mathbf Z/m)^4$, and its rational-point image in $C$ is
exactly $\{P\}$.
\end{claim}
\begin{proof}
The explicit elliptic maps in \statementref{B.23} show that
Frobenius on $J_C$ is $[-2^{19}]$. Thus
$J_C(k)=\ker(F-1)=J_C[m](\bar k)$; all $m$-torsion is rational,
and multiplication by $m$ kills every rational point.
The Abel map induces a surjection onto the geometric abelian
fundamental group modulo $m$, so its pullback of the multiplication
cover is connected. It is \emph{\'etale}, since $m$ is odd.
A rational point $Q$ lifts precisely when $[Q-P]\in[m]J_C(k)=0$.
The Abel map is injective on geometric points, so this means $Q=P$.
The point $(P,0)$ is the specified rational lift. The exact integer
degree and the field and curve data are recorded in
\path{code/supersingular_isolating_cover.json}.
\end{proof}



\subsection{The nontrivial primitive-cover counterexample}

\begin{claim}{B.23}{Refuted}{Universal nontrivial primitive-zero cyclic covers}
The assertion that every genus-two curve in characteristic two has
a connected cyclic \emph{\'etale} cover of odd degree $n>1$ whose
primitive Jacobian part has two-rank zero is false. A completely
specified counterexample is the smooth proper model over
$k=\mathbf F_2[t]/(M(t))$ of
\[
\begin{split}
 M(Z)&=Z^{38}+Z^{14}+Z^{10}+Z^9+Z^8+Z^5+Z^2+Z+1,\\
 d&=t^3+1,\qquad
 C:\ y^2+y=A(u)^3,\\
 A(u)&=(u^2+t^2u+t)/d,\qquad
 B(u)=(tu^2+u+t^2)/d .
\end{split}
\]
The polynomial $M$ is irreducible, and $C$ is supersingular of
genus two with a unique rational point $P$ above infinity.
For every such cover over $\bar k$, the primitive part has positive
two-rank. This independently confirms the particular counterexample
AS-A.86. It does not refute the version allowing degree one.
\end{claim}
\begin{proof}
We supply the geometry, the finite cutoff, and the exact certificates.

\emph{The curve and its Jacobian.}
Successive squaring modulo $M$ gives
$Z^{2^{38}}\equiv Z$ and
$\gcd(M,Z^{2^2}-Z)=\gcd(M,Z^{2^{19}}-Z)=1$.
These are precisely the irreducibility tests for degree 38; every
residue is retained in \path{code/primitive_certificates/manifest.json}.
In particular $t\notin\mathbf F_4$ and $d\ne0$.
After Artin--Schreier changes removing even powers, $A(u)^3$ has
highest odd term $(t^2/d^3)u^5$, with nonzero coefficient.
The degree-two map to $\mathbf P^1_u$ is ramified only at infinity,
with different exponent six. Riemann--Hurwitz gives genus two,
one rational point $P$ at infinity, and canonical divisor $2P$.
The Artin--Schreier $p$-rank formula, or its elementary one-pole
Cartier calculation, gives two-rank zero. In this case a basis of
regular differentials is $du,u\,du$ after the reduced quintic
equation; Cartier sends $du$ to zero and $u\,du$ to $du$, so its
stable rank is zero.

On $E:y^2+y=x^3$ there are two morphisms
\[
 (x_1,y_1)=(A(u),y),\qquad
 (x_2,y_2)=(B(u),y+(u^3+1)/d).
\]
The polynomial identity
\[
 (u^2+t^2u+t)^3+(tu^2+u+t^2)^3
       =d(u^3+1)^2+d^2(u^3+1)
\]
checks the second equation. Both maps send $P$ to the origin.
The joint map is birational onto its image, since
$u=x_2+t x_1$ and $y=y_1$. Its image therefore has normalization
of genus two and cannot lie in an elliptic subvariety or its
translate. The induced map $J_C\to E^2$ is consequently an isogeny
defined over $k$. On $E$, the chord-and-tangent formula gives
$[2](x,y)=(x^4,y^4+1)$; hence its $\mathbf F_4$-Frobenius is $[-2]$.
The $\mathbf F_{2^{38}}$-Frobenius of $E^2$, and therefore that of
$J_C$, is $[-2^{19}]$. For the last implication compose
$F_{J_C}+[2^{19}]$ with the isogeny: its connected image would lie
in a finite kernel, so it is zero. This proves supersingularity
and the exact scalar Frobenius, without assuming a torsion-group
isomorphism from an unspecified isogeny.

\emph{The theta-orbit criterion.}
A cyclic cover of odd degree $n$ corresponds to a line bundle $L$
of exact order $n$ with a trivialization of $L^n$. Its coherent
cohomology decomposes as
$\bigoplus_{i\in\mathbf Z/n}H^1(C,L^i)$, and absolute Frobenius
sends the $i$-summand semilinearly to the $2i$-summand. The
nontrivial summands have dimension one by Riemann--Roch.
For a multiplication-by-two orbit $O$ of primitive indices, its
contribution to the two-rank is $|O|$ times the eventual rank of
the orbit composite. This follows from the Artin--Schreier
sequence: the Frobenius-stable subspace has dimension equal to
$\dim_{\mathbf F_2}H^1_{\et}(D,\mathbf F_2)$. The idempotent
primitive character projector, or the image of
$\prod_{\ell\mid n}(1-\sigma^{n/\ell})$, selects exactly these
primitive summands.

Let $\mathcal B$ be the cokernel of
$\mathcal O_{C^{(1)}}\to F_*\mathcal O_C$. In characteristic two it
is a line bundle with square the canonical bundle: locally write
a change of uniformizer as $v=a^2+u b^2$; its class in the
cokernel changes by $b$, whose square is $dv/du$. The cohomology
sequence tensored with $L^i$ identifies the kernel of the
Frobenius arrow with $H^0(\mathcal B\otimes L^i)$.
There is just one geometric theta characteristic because
$J_C[2](\bar k)=0$, so $\mathcal B=\mathcal O(P)$.
The corresponding theta locus is precisely the symmetric Abel
curve $\Theta=C-P$. Thus zero primitive two-rank requires at
least one theta point in every multiplication-by-two orbit of
primitive points of the cyclic subgroup $\langle L\rangle$.

\emph{A complete finite bound.}
Write $f=\operatorname{ord}_n(2)$,
$e=\gcd(f,19)$, and $M_n=\varphi(n)/e$.
Scalar Frobenius and symmetry of $\Theta$ show that one point
in a required orbit brings its entire multiplication-by-$2^{19}$
orbit into $\Theta$. Consequently there are at least $M_n$
primitive theta points, and also zero. For a nonzero Jacobian
difference, at most two ordered pairs of distinct points of
$\Theta$ realize it: two representations yield two linearly
equivalent degree-two divisors, and the unique hyperelliptic
pencil gives the sole alternative representation. Hence
\[
 M_n(M_n+1)\le2(n-1).
\]
For odd $n$, prime-power multiplication proves
$\varphi(n)\ge(2/3)n^{5/6}$. Indeed
$(\ell-1)^6\ge\ell^5$ for every prime $\ell\ge5$, and the possible
factor at three is at least $64/243>(2/3)^6$.
It follows that $n<10\cdot19^3=68590$.

The standard-library program \path{code/primitive_conductors.py}
exhaustively enumerates all odd integers in that interval and gives
\[
191,229,361,419,457,571,573,647,687,955,1083,1257,1371.
\]
It then imposes the full difference constraint on the mandatory
packets $H=\langle 2^{19},-1\rangle\subset(\mathbf Z/n)^\times$.
Already $\{0\}\cup H$ excludes all but $191,229,457,1257$.
For $457$, every coset of $U=\langle2,-1\rangle$ must contain an
$H$-packet. There are six such cosets. Normalize the first packet
to $H$ by changing a cyclic generator, and enumerate every
remaining packet. The surviving search depths are
$1,17,226,879,262,0$, so there is no admissible selection.
For $1257$, there are two $U$-cosets and each mandatory packet
has 22 elements, half in each nonzero residue class modulo three.
With zero, the residue counts are $(1,22,22)$, forcing
$924$ ordered distinct differences divisible by three. Only
$2(1257/3-1)=836$ are permitted. Thus only $191,229$ remain.
The script retains every input, packet, candidate list and count.

\emph{The two full polynomial certificates.}
For $E$ put $r_n(X)=x([n]R)$ when $X=x(R)$. The addition formula gives
\[
 r_1=X,\qquad r_{2m}=r_m^4,\qquad
 r_{2m+1}=X+(r_m+r_{m+1})^{-2}.
\]
Compute these as reduced rational functions over $\mathbf F_2$.
Independently form
\[
 \psi_1=\psi_2=\psi_4=1,\quad\psi_3=X^4+X,\quad
 \psi_{2m}=\psi_m^4,\quad
 \psi_{2m+1}=\psi_{m+2}\psi_m^3+\psi_{m-1}\psi_{m+1}^3.
\]
For each required $n$ the program verifies that the reduced
denominator of $r_n$ is exactly $\psi_n^2$, with coprime numerator.
Thus its zeros are exactly the finite coordinates of nonzero
$n$-torsion, without importing a division-polynomial formula
that excludes characteristic two.
\[
\begin{array}{c|r|r|r}
 n&\deg\psi_n&\deg\operatorname{num}r_n&
                  \deg\operatorname{den}r_n\\ \hline
191&18240&36481&36480\\
229&26220&52441&52440
\end{array}
\]
Exact arithmetic in the stated field gives
\[
 U_n(u)\psi_n(A(u))+V_n(u)\psi_n(B(u))=1
                    \quad(n=191,229).
\]
The composed polynomial degrees are respectively $36480,52440$;
both B\'ezout coefficients have degree four less in each case.
All eight polynomials are supplied as complete binary coefficient
arrays, with ascending unsigned little-endian 64-bit coefficients;
bit $i$ represents $t^i$. The manifest gives SHA-256 hashes.
The producer checks both whole identities. The separately written
\path{code/primitive_verify.py} loads the arrays, verifies their
hashes and degrees, independently composes the polynomials over
$\mathbf F_2$, and verifies the two whole identities again.

If a finite point $Q$ of $C$ gave a nonzero Abel class killed by
one of these integers, both of its finite elliptic images would
be killed by the integer. Its coordinate $u$ would therefore be a
common zero of the two composed polynomials, contradicting the
B\'ezout identity. The point at infinity is the origin and cannot
be primitive. Every possible conductor has now been excluded,
over $\bar k$ itself. This proves the counterexample.
\end{proof}

\begin{claim}{B.24}{Open}{The literal degree-one-permitted primitive-cover assertion}
For every proper curve of genus at least two over
$\overline{\mathbf F}_p$, does there exist a connected cyclic
\emph{\'etale} cover of odd degree prime to $p$, allowing degree one,
whose primitive part has $p$-rank zero, with the convention
$P_1=J_C$?
\end{claim}
The preceding counterexample does not settle this statement.
The identity cover proves it for every curve of $p$-rank zero,
including that example. For positive $p$-rank, degree one fails.
The exact Frobenius-orbit formulation above extends to arbitrary
genus with character spaces of dimension $g-1$: one needs
nilpotence of each orbit composite, not just a noninvertible
arrow. The conductor bound proved here uses scalar Frobenius
and the one-dimensional genus-two character spaces, so it does
not bound the positive-$p$-rank cases. We tried transferring the
scalar obstruction, using a single nonordinary character, and
using the identity convention; each fails for this stated reason.
The smallest unresolved scope in this work is positive-two-rank
genus-two curves in characteristic two. No such curve is claimed
to be a counterexample.



\subsection{Rational witnesses on the explicit primitive-cover counterexample}

The geometric counterexample in \statementref{B.23} is not a counterexample
to the rational class-number criterion. We give two exact witnesses,
including the least possible prime over all constant extensions.
The computations certify the class-number parity; they do not claim
to compute the full characteristic-zero zeta numerator.

\begin{claim}{B.25}{Proved}{A normalized Cartier test in the case used here}
Let $K=\mathbf F_{2^d}$, and let
\[
 C:\quad Y^2+hY=f(x),\qquad h\in K^\times,
 \quad f\text{ monic of degree five},
\]
be smooth, with marked point $P=\infty$. Let $\ell$ be an odd prime
such that $\mu_\ell\subset K$, and let $a\in J_C(K)$ have exact
order $\ell$. Suppose every nonzero $ia$ has reduced Mumford
representative $(u_i,v_i)$ of degree two. Write $j=2i\bmod\ell$ and
\[
 \begin{split}
 w_i&=(v_i^2+hv_i+f)/u_i,\\
 V_i&=v_i+u_i\bigl((w_i/h)\bmod u_i\bigr),\\
 c_i&=(V_i)_3\,(u_j)_1.
 \end{split}
\]
Here subscripts on polynomials indicate coefficients. Then the
unique constant cyclic, $P$-split cover with character subgroup
$\langle a\rangle$ has odd relative class number if and only if
\[
 A_i:=\prod_{s=0}^{d-1}c_{2^s i}^{\,2^{d-1-s}}\ne1
       \qquad(1\le i<\ell).
\]
The index on each $c$ is reduced modulo $\ell$. The assertion includes
the specified pointing; there is no unspecified constant twist.
\end{claim}
\begin{proof}
Write $ia=[D_i-2P]$ and $L_i=\mathcal O(D_i-2P)$.
The Artin--Schreier model has a local parameter $t=x^2/Y$ at $P$;
monicity gives leading residues one for $xt^2$ and $Yt^5$.
Since $dx/h$ has divisor $2P$, Riemann--Roch and Serre duality
give dual bases
\[
 \omega_i=u_i\,dx/h\quad\text{of }H^0(C,\Omega_C\otimes L_i^{-1}),
 \qquad e_i\quad\text{of }H^1(C,L_i).
\]
For example, the divisor of $u_i\,dx/h$, after subtracting
$D_i-2P$, is $\iota D_i$, which is effective. Each space has
dimension one because $L_i$ is nontrivial of degree zero.

The displayed interpolation formula is the unique polynomial of
degree at most three congruent to $v_i$ modulo $u_i$ for which
$u_i^2\mid V_i^2+hV_i+f$. Indeed, modulo $u_i^2$ the correction
$u_i b$ changes the norm by $hu_i b$. Since the reduced double
has degree two, $V_i$ is cubic and Cantor reduction gives
\[
 V_i^2+hV_i+f=(V_i)_3^2u_i^2u_j,
 \qquad v_j\equiv V_i+h\pmod {u_j}.
\]
Consequently
\[
 g_i=\frac{Y+V_i}{(V_i)_3u_j},\qquad
 \operatorname{div}(g_i)=2D_i-D_j-2P.
\]
Its leading unit at $P$ is one. Multiplication by $g_i$ therefore
identifies $L_i^2$ with $L_j$ respecting their rigidifications at $P$.

The residue definition of Cartier gives
$\langle g_iF(z),\eta\rangle
 =\langle z,\mathcal C(g_i\eta)\rangle^2$.
Using $Y/h=(Y/h)^2+f/h^2$, the even-power terms have Cartier zero,
whereas the odd-power part of $f+hV_i$ equals that of
$(V_i)_3^2u_i^2u_j$. Thus
\[
 \mathcal C\bigl((Y+V_i)dx/h\bigr)
      =(V_i)_3\sqrt{(u_j)_1}\,u_i\,dx/h,
 \qquad
 \mathcal C(g_i\omega_j)=\sqrt{c_i}\,\omega_i.
\]
It follows that absolute Frobenius sends $e_i$ to $c_i e_{2i}$
square-semilinearly. Iterating $d$ times gives $A_i e_i$, because
$2^d\equiv1\pmod\ell$ and every representative is defined over $K$.

For completeness, the cover and the normalization can be described
without a choice of a cohomological scalar. Since $a$ has exact
order $\ell$, a rational function with divisor $\ell(D_1-2P)$
exists over $K$. Scale it to leading unit one at $P$. Adjoining
its $\ell$th root gives a connected cyclic \emph{\'etale} cover;
all its valuations are divisible by $\ell$, its character has
exact order $\ell$, and its leading unit makes the fibre over
$P$ split completely. Any two descents with the same geometric
character differ by a constant character, whose evaluation at
$P$ proves uniqueness of the split descent. Equivalently,
rigidifying all the $L_i$ at $P$ makes all their multiplication
isomorphisms canonical: any associativity discrepancy is a
constant with value one at $P$ and therefore is one. The $g_i$
used above are exactly these canonical doubling isomorphisms.

Finally, the Artin--Schreier sequence identifies the geometric
\'etale $\mathbf F_2$ cohomology with the Frobenius-stable part
of coherent cohomology; see also the characteristic-$p$
comparison in \cite[Section 9]{MilneJac}. A finite-dimensional
$K$-space with square-semilinear operator $T$ has a nonzero
$T$-fixed vector over $K$ exactly when $T^d$ has eigenvalue one:
on $\ker(T^d-1)$ the operator defines descent from $K$ to
$\mathbf F_2$, and ordinary finite-field semilinear descent
has the same dimension. Applying this on the nontrivial
character summands shows that rational new two-torsion exists
exactly when at least one $A_i$ is one. Dualizing the finite
\'etale two-torsion representation does not change the dimension
of its Frobenius-fixed space. Since the cover degree is odd,
the invariant and new parts split by the deck averaging
idempotent, and $\operatorname{Nm}\circ\pi^*=[\ell]$ is invertible
on two-primary groups. The relative class number is odd
exactly when that new rational two-primary group is zero.
\end{proof}

\begin{claim}{B.26}{Proved}{Exact degree-191 and degree-229 rational witnesses}
For the curve and point $(C,P)$ of \statementref{B.23}, a constant
cyclic $P$-split cover of degree $191$ with odd relative class
number exists over $\mathbf F_{2^{380}}$. A second such cover,
of degree $229$, exists over $\mathbf F_{2^{152}}$.
Their genera are respectively $192$ and $230$, and their new
Jacobian parts are ordinary of dimensions $190$ and $228$.
\end{claim}
\begin{proof}
We specify the fields, the normalization, and the exact finite
certificates. No search heuristic is used as evidence of their parity.
Put $d_0=t^3+1$ and use $A(u),B(u)$ from \statementref{B.23}.
Write $A(u)^3=\sum_{j=0}^6 a_j u^j$. Successively eliminate the
terms of degree six, four and two by adding $S(u)^2+S(u)$, with
$S$ cubic. The resulting polynomial is
\[
 f_0(u)=a_5u^5+b_3u^3+b_1u+b_0,
 \qquad a_5=t^2/d_0^3\ne0.
\]
No field extension is required for these square roots in a finite
field. Set
\[
 u=a_5x,\qquad y=a_5^3Y+S(a_5x),\qquad
 h=a_5^{-3},\qquad f(x)=a_5^{-6}f_0(a_5x).
\]
Then $Y^2+hY=f(x)$ is monic of degree five and has the same marked
point at infinity. In particular the normalization is an isomorphism
over the stated field, not a scalar extension omitted from the result.

The two certificate fields are
\[
 \begin{array}{c|c|c}
 \ell&K&[K:k]\\ \hline
 191&\mathbf F_2[z]/(z^{380}+z^{47}+1)&10\\
 229&\mathbf F_2[z]/(z^{152}+z^6+z^3+z^2+1)&4.
 \end{array}
\]
The following hexadecimal strings encode elements of $K$: binary
bit $j$ is the coefficient of $z^j$. Consecutive lines belonging
to one entry are concatenated. They give the embedding of $k$ and
the reduced representative $u=x^2+u_1x+u_0$, $v=v_1x+v_0$ of the
chosen class $a$ on the normalized model.
\begin{quote}\small
\begin{verbatim}
191:
 t  = cebd4b3dac18387ef98c70a8bc69661f941f1ab70330e1fb
      19cea93138d8ab1686e3f52a7f31d78c0d349454dd0a7a3
 u0 = 35800adb325893dcca9e4a4670173048e3e683b2ccd8383d
      42236023341647d76e22bcb0bf5df97910948458ab1fd37
 u1 = e9edfad1fe30d89fd0afb8a9cacda03aca4c455c01b4c42
      a0f679640b5f557a33700a882ea77325babe828b006e95a4
 v0 = fea3805d1c99d0cdcad5dd87afc416eb89ff67d3976e543
      0580f7152b05537b2ad904bc24804bfa1cc3fdd706890fec
 v1 = e8156100e879ff743228cc826dbe666030ae4636505cf28
      83274f15cd61a2e7e3cdd52bacc8324e95038057b70470e5
229:
 t  = 80fafae83a3d67cbb2033e0b744dcc9de8aa86
 u0 = 8f276acfd0d6f09be86235d08a7ea29b9771cb
 u1 = 47138bd65385caa096e9e047779586fa5eda15
 v0 = 8f83e3c8bccd0c1f58d6f6168617be21fab882
 v1 = 568c198f174a8035b7a6d31f2c28bce57c2db0
\end{verbatim}
\end{quote}

The files \path{code/r8_rational_witness_191.json} and
\path{code/r8_rational_witness_229.json} retain these data and
every reduced nonzero multiple of $a$, every $c_i$, and every $A_i$.
Their SHA-256 hashes are respectively
\begin{quote}\small\ttfamily
2ccc98fdd5ae31557e6f1977ceb2dd1f51face99f9a08b48f11e6d7808bee256
\par
73cee7d15efe7b6620c01660605dc4d49ab128856163b5375ea6a0fb1e93851c.
\end{quote}
The independent verifier \path{code/criterion_r8_verify.py} reads
these files and imports no producer routines. A separate binary
polynomial implementation proves both field polynomials irreducible
by Rabin's criterion, and proves the degree-38 base polynomial
irreducible. It verifies $M(t)=0$, the displayed normalization, and
the elliptic-map identity from \statementref{B.23}.

For every supplied pair $(u_i,v_i)$, the verifier checks reducedness,
$u_i\mid v_i^2+hv_i+f$, the tangent interpolation, and both the
$x$- and $Y$-coordinates of its reduced double. Independently of
the producer's general Cantor addition formula, it reconstructs
$ia+a$ by coprime Chinese-remainder interpolation for
$2\le i\le\ell-2$. All these supports are coprime in the supplied
data. It also checks $2a$ by tangent interpolation and
$(\ell-1)a=-a$. Thus $a$ is nonzero of exact prime order $\ell$,
and all the supplied indices are certified. The verifier additionally
checks the displayed class as the sum of the two saved elliptic
pullbacks, although the order and divisor-sequence verification
already suffice for the existence assertion.

It recomputes every $c_i$ and recomputes each product $A_i$ directly
from separately accumulated squares, rather than reusing the
producer's Horner recurrence. The exact results are
\[
 \begin{array}{c|r|r|r|r}
 \ell&\text{nonzero }(u_i,v_i)&\text{CRT additions}&
                 \text{checked }c_i&\#\{i:A_i=1\}\\ \hline
 191&190&188&190&0\\
 229&228&226&228&0.
 \end{array}
\]
Every $c_i$ and $A_i$ is nonzero. Extra identities, not used by
the producer, give useful checks on the indices and rigidifications:
\[
 A_{2i}=A_i^2,\qquad A_{-i}=A_i,\qquad
 A_i^{2^{19}}=A_i,\qquad
 A_{-2^{19}i}=A_i^{2^{38}}.
\]
As a compact exact record, let $m_b(X)=\sum_j b_jX^j$ for the
binary digits $b_j$ of the positive integer $b$. The characteristic
polynomial of the computed linear coherent Frobenius on the new
part is
\[
 \begin{cases}
 m_{892407}(X)^{10},&\ell=191,\\
 m_{713865}(X)^4m_{785567}(X)^4m_{913207}(X)^4,&\ell=229.
 \end{cases}
\]
Each displayed $m_b$ has degree nineteen, is the minimal polynomial
of a checked multiplier, and has $m_b(1)=1$. The verifier obtains
it by multiplying all nineteen distinct Frobenius conjugates and
checking that every coefficient is in $\mathbf F_2$.
The full results and input hashes are in
\path{code/criterion_r8_verify.json}; the concise run transcript is
\path{code/criterion_r8_verify.out}.

The scalar Frobenius of \statementref{B.23} is $s=-2^{19}$.
Direct modular multiplication gives
\[
 \operatorname{ord}_{191}(2)=95,
 \quad\operatorname{ord}_{191}(s)=10,
 \qquad
 \operatorname{ord}_{229}(2)=76,
 \quad\operatorname{ord}_{229}(s)=4.
\]
Thus the stated fields contain all the corresponding Jacobian
torsion and all the roots of unity; the constant-deck descent
condition $\sigma_{\#K}(a)=(\#K)a$ holds. Statement
\statementref{B.25} now proves both odd relative class numbers.
Nonvanishing of every $c_i$ makes the absolute Frobenius maps
invertible on all nontrivial coherent character lines, so the
new parts are ordinary. Riemann--Hurwitz gives genus $\ell+1$,
and their dimension is $\ell-1$ as asserted.
\end{proof}

\begin{claim}{B.27}{Proved}{The exact least prime and the two field minima}
For the fixed marked curve $(C,P)$ of \statementref{B.23}, the
least prime degree of a constant cyclic $P$-split cover satisfying
the rational class-number criterion over any finite extension
of $k=\mathbf F_{2^{38}}$ is exactly $191$.
At degree $191$ the least possible constant extension degree is ten;
at degree $229$ it is four, and each minimum is attained by
\statementref{B.26}. The second assertion is a minimum for that
specified prime, not a claimed minimum over all primes.
\end{claim}
\begin{proof}
We include the pointed-involution argument, since a geometric
positive-rank obstruction alone would not prove a rational
lower bound. Fix any geometric connected cyclic cover of odd
prime degree $\ell$. Scalar Frobenius preserves its character
line, so its geometric isomorphism class is already defined
over $k$. Choosing a point above $P$ gives its unique pointed
descent. On its deck group Frobenius acts by $s=-2^{19}$:
Cartier duality changes the scalar on $J_C[\ell]$ into
$(\#k)/s=s$. Therefore its first constant-deck field has degree
$m=\operatorname{ord}_\ell(s)$ over $k$; every later
constant-deck field has degree divisible by $m$.

The hyperelliptic involution fixes $P$ and acts as minus one
on $J_C$. Its unique lift fixing the selected point above $P$
inverts the deck group. Uniqueness makes the lift commute with
the arithmetic Frobenius of the pointed descent. Let
\[
 R=\mathbf F_2[X]/\Phi_\ell(X),\qquad
 W=H^1_{\et}(D_{\bar k},\mathbf F_2)_{\rm new}.
\]
Every nontrivial coherent character space has dimension one.
The injection of geometric \emph{\'etale} cohomology into its
Frobenius-stable coherent part consequently makes $W$ free
of rank one over a direct factor $R'$ of the finite \emph{\'etale}
algebra $R$. A zero-rank character factor is simply omitted.

Frobenius $F$ is semilinear for the automorphism $X\mapsto X^s$
of $R'$. Since $s^m=1$, its $m$th power is $R'$-linear, hence
multiplication by a unit $b\in R'$. Commuting $F^m$ with $F$
and with the lifted involution forces
\[
 b\in\bigl(R'^{\langle s,-1\rangle}\bigr)^\times.
\]
Put $f_\ell=\operatorname{ord}_\ell(2)$. If $19\nmid f_\ell$,
then $\langle2\rangle\subset\langle2^{19},-1\rangle
=\langle s,-1\rangle$. On the finite \emph{\'etale} algebra $R$,
the automorphism $X\mapsto X^2$ is the absolute Frobenius.
Its fixed ring is a product of copies of $\mathbf F_2$, whose
unit group is trivial. Therefore $b=1$. Every geometric new
two-torsion point is then rational at the first constant-deck
field and remains rational over all subsequent ones.

By \statementref{B.23}, the new part of every prime-degree
geometric cover has positive two-rank. Thus none of these
covers can have odd relative class number when
$19\nmid\operatorname{ord}_\ell(2)$. If $\ell<191$ and
$19\mid\operatorname{ord}_\ell(2)$, then $19\mid\ell-1$;
the only possibilities for an odd prime are
$\ell\in\{39,77,115,153\}$, all composite. This excludes every
odd prime below $191$ over every finite extension. Degree two
is excluded by the prime-to-characteristic requirement.

The degree-$191$ certificate proves attainment, and the necessary
constant-deck degrees computed above prove the two field minima.
For clarity, allowing a different unpointed constant twist does
not evade the lower bound: splitting above the specified point
selects the unique pointed twist, to which the commuting
involution argument applies.
\end{proof}

\paragraph{Consequences and scope.}
The same explicitly specified geometric curve has positive primitive
two-rank for every nontrivial odd cyclic cover, yet possesses rational
class-number witnesses. It therefore separates the geometric
primitive-zero condition from the rational criterion on a fixed
curve, with exact field and prime data. This is a positive instance
of B1, not a proof of its universal assertion, including the proposed
range $p>(2g-2)^2$.


\subsection{What varying the constant twist does and does not prove}

\begin{claim}{B.28}{Proved}{A uniform count of bad constant twists}
Let $C/K$ be smooth proper of genus $g\ge2$ over a finite field
of characteristic $p$. Let $D\to C$ be a geometrically connected
\emph{\'etale} cover with constant cyclic deck group of odd prime
order $\ell\ne p$. Put $f=\operatorname{ord}_\ell(p)$.
Among the $\ell$ constant twists of this cover, at most
\[
                \frac{(g-1)(\ell-1)}{f}
\]
have relative class number divisible by $p$. In particular, if
$f\ge g-1$, at least one twist has $p$-unit relative class number.
For genus two this last condition is automatic.
\end{claim}
\begin{proof}
Fix a generator $\delta$ of the deck group. The nontrivial factor
of its group algebra is
\[
 R=\mathbf F_p[X]/\Phi_\ell(X)
       \simeq\prod_{j=1}^{(\ell-1)/f}\mathbf F_{p^f}.
\]
Let $W$ be the new part of
$H^1_{\et}(D_{\bar K},\mathbf F_p)$. Over an algebraic closure,
the nontrivial coherent character spaces are $H^1(C,L)$ for
nontrivial degree-zero line bundles $L$ of order $\ell$; each
has dimension $g-1$. The Artin--Schreier comparison embeds
\'etale cohomology tensored with $\bar K$ into the Frobenius-stable
coherent part, equivariantly for the deck action. Consequently the rank $d_j$ of $W$ over each field
factor of $R$ satisfies $0\le d_j\le g-1$.

The deck group is constant, so arithmetic Frobenius $F$ commutes
with its action and is linear over every field factor. Twisting
by the constant cocycle with Frobenius value $\delta^t$, for
$t\in\mathbf Z/\ell$, replaces its action on $W$ by
$\delta^tF$, up to using $-t$ consistently in the contravariant
cohomology convention. On the $j$th field factor, write the image
of $\delta$ as the primitive root $\zeta_j$. A new fixed vector
exists on this factor exactly when
\[
               \det(\zeta_j^{-t}-F\mid W_j)=0.
\]
This polynomial has degree $d_j$, and its $\ell$ arguments are
distinct. At most $d_j$ twists are therefore bad on this factor.
Taking the union over the $(\ell-1)/f$ factors gives the claimed
bound. Nonsemisimplicity of $F$ causes no problem: a fixed vector
is detected exactly by this determinant. The rational two-primary
argument in \statementref{B.25} applies with $p$ in place of two:
norm and pullback split the $p$-primary groups because
$p\nmid\ell$, and rational new $p$-torsion is equivalent to
divisibility of the relative class number. Finally,
$f\ge g-1$ makes the bound at most $\ell-1$, strictly below the
number of twists.
\end{proof}

\begin{claim}{B.29}{Proved}{Existence with an unspecified base point}
For every smooth proper curve $C/k$ of genus $g\ge2$ over a
finite field, there are a finite extension $K/k$, a point
$P_K\in C(K)$, and a constant cyclic $P_K$-split cover of odd
prime degree, different from the characteristic, with $p$-unit
relative class number. The point $P_K$ is not prescribed in
advance. More quantitatively, for a constant degree-$\ell$ cover
as in \statementref{B.28}, put
$B=(g-1)(\ell-1)/f$ and $g_D=1+\ell(g-1)$. If $f\ge g-1$, the
number of $P\in C(K)$ for which its unique split twist is a
class-number witness is at least
\[
       \frac{\ell-B}{\ell}
             \bigl(\#K+1-2g_D\sqrt{\#K}\bigr).
\]
\end{claim}
\begin{proof}
For each twist $D_t$, let $N_t$ be the number of rational points
of $C$ that split in it. Every geometric fibre has $\ell$ points
and Frobenius is a deck translation. Thus each $P\in C(K)$
splits in precisely one twist, and
$\#D_t(K)=\ell N_t$. At least $\ell-B$ twists have $p$-unit
relative class number. Applying the Weil lower bound to each
of these geometrically connected curves gives the displayed
inequality \cite[Theorem 1.6 and formula 1.5.1]{WeilCurves}.
If its right-hand side is negative, the inequality is automatic;
otherwise summing over at least $\ell-B$ good twists proves it.

There exist odd primes $\ell\ne p$ with
$\operatorname{ord}_\ell(p)\ge g-1$: the primes with order less
than $g-1$ divide the fixed nonzero integer
$\prod_{j=1}^{g-2}(p^j-1)$ and are finite in number.
Choose such an $\ell$ and a nonzero geometric cyclic
$\ell$-cover, which exists because the prime-to-$p$ abelian
fundamental group has rank $2g$ at $\ell$.
It, its deck group, and a model of its base curve descend to a
finite extension. Enlarge that field until
$\#K+1>2g_D\sqrt{\#K}$. The quantitative assertion then produces
a good twist with a rational point above some $P_K$.
All descent and field enlargements precede the choice of a good
twist; no assertion that a previously good twist remains good
under extension is needed.
\end{proof}

\paragraph{The remaining pointing obstruction in B1.}
For a fixed prescribed $P$, its split twist is one specified
element of the $\ell$ twists. The counting theorem can leave
that element among the bad twists. There is no way to repair
this by simply extending a good unpointed twist until $P$
splits: this extension may make its nontrivial $p$-torsion
eigenvalues equal to one.

This failure occurs on the explicit R8 curve, not just on an
abstract representation. Consider the cyclic degree-three
character obtained by pulling back the point $(0,0)$ of exact
order three on $E:y^2+y=x^3$ along its first degree-two map
$C\to E$. Norm followed by pullback is multiplication by two,
so this character is nonzero. Its pointed cover is constant over
$k=\mathbf F_{2^{38}}$, since $s=-2^{19}\equiv1\pmod3$.
Its new geometric two-rank is positive by \statementref{B.23};
the nontrivial group algebra is the single field $\mathbf F_4$,
and the coherent character dimension is one, so its new
\'etale cohomology has rank one over that field. The commuting
pointed hyperelliptic involution used in \statementref{B.27}
forces Frobenius to be the identity. Both nontrivial constant
twists instead have primitive third-root Frobenius multiplier,
and hence odd relative class number. They do not split over $P$.
The cubic constant extension first makes that fibre split and
simultaneously cubes the multiplier to one, making the relative
class number even again. Thus the unpointed existence theorem
does not settle B1, including the range $p>(2g-2)^2$.

These three covers can be given explicitly in the original R8
coordinates by
\[
       w^3=y,\qquad w^3=ty,\qquad w^3=t^2y.
\]
On the elliptic curve, $\operatorname{div}(y)=3(0,0)-3O$;
its pullback gives the exact degree-three character above.
If $\tau=x/y$ is the elliptic local parameter at $O$, then
$y\tau^3$ has residue one. The degree-two map has
$\tau=b\eta^2+\cdots$ at $P$ for a rational local parameter
$\eta$ and $b\in k^\times$. The leading unit of $y$ at $P$
is therefore $b^{-3}$, a cube, proving that the first cover
is precisely the split twist.
The pure binary-polynomial program
\path{code/criterion_twist_obstruction.py} verifies irreducibility
of $M$ and computes
\[
        t^{(2^{38}-1)/3}=588964846
\]
in the polynomial-bit encoding of $k$. This value is a nontrivial
cube root of unity, so $t$ is a noncube and the other two covers
first split above $P$ over the cubic extension. The same program
checks the two-dimensional new cohomology matrices
\[
 I,\qquad
 \begin{pmatrix}0&1\\1&1\end{pmatrix},\qquad
 \begin{pmatrix}1&1\\1&0\end{pmatrix},
\]
up to interchanging the last two by the chosen cohomology
convention. Their determinants $\det(1-F)$ are $0,1,1$ in
$\mathbf F_2$, and all three matrices cube to $I$.
The full finite data and the geometric scope of this check are
recorded in its JSON and run transcript.


\subsection{Exact extension-field tests for plane quartics}

\begin{claim}{B.30}{Proved}{A least-prime-five quartic witness over \(\mathbf F_{3^8}\)}
Let $C/\mathbf F_3$ be the smooth plane quartic
\[
                W^4=X^3Z-XZ^3,\qquad P=(1:0:0).
\]
There is a $P$-split constant cyclic degree-five cover over
$\mathbf F_{3^8}$ whose relative class number is prime to three.
Thus five is the least odd prime different from the characteristic
which gives a witness for this pointed curve over some constant
extension. Every constant degree-five cover of this curve requires
a constant degree divisible by eight; the displayed field is
therefore also the least one for a degree-five witness.
The certificates give six such covers. For each of them the
relative class number stays prime to three exactly under further
constant extensions of odd degree.
\end{claim}
\begin{proof}
The three partial derivatives of the displayed homogeneous
polynomial have no common projective zero: in characteristic
three they are $-Z^3,W^3,X^3$, up to nonzero constants.
The unique point at infinity is $P$. On the affine chart put
$w=W/Z$ and $x=X/Z$. The degree-two morphism
\[
       C\longrightarrow E:\ y^2=x^3-x,
                    \qquad (x,w)\longmapsto(x,w^2)
\]
is branched, in particular at infinity. The elliptic curve has
four points over $\mathbf F_3$, so its Frobenius satisfies
$F^2=[-3]$. Consequently $E[5]$ is constant over
$K=\mathbf F_{3^8}$.

The program \path{code/quartic_cyclic5.py} works in
\[
 K=\mathbf F_3[z]/
       (z^8+2z^5+z^4+2z^2+2z+2).
\]
It specifies a basis of $E[5]$ and all six cyclic subgroups.
For a generator $T=(a,b)$ and $2T=(c,d)$ put
$h=(x-a)(x-c)$. The normalized isogeny has
\[
 r(x)=N(x)/h(x)^2,
 \qquad r'(x)=g(x)/h(x)^3,\qquad
       g=N'h-2Nh',
\]
where
\[
 N=xh^2+
       \sum_{(u,v)\in\{T,2T\}}
        \left(\frac{h^2}{x-u}+\frac{v^2h^2}{(x-u)^2}\right).
\]
These are polynomials. The addition formula on $E$ proves
this isogeny formula by summing $x(R+S)-x(S)$ over the four
nonzero kernel points. The program verifies its defining
identity
\[
 (x^3-x)g^2=N^3+A'Nh^4+B'h^6.
\]
It gives $s,r_0\in K$ with $s^4=-A'$ and
$r_0^3+A'r_0+B'=0$, hence the isogeny to $E$ is
\[
       x_E=(N/h^2-r_0)/s^2,
       y_E=yg/(s^3h^3).
\]
All its kernel points are $K$-rational. Pulling back
$C\to E$ gives the cover whose normalization has equation
\[
       v^4=F(x),\qquad
       F=(x^3-x)g^2h^2/s^6,\qquad
       (X_C,W_C)=((N/h^2-r_0)/s^2,v/h^2).              \tag{B.Q.1}
\]
It is connected: the branched quadratic extension defining
$C/E$ is disjoint from this degree-five unramified isogeny.
Base change makes its map to $C$ finite \emph{\'etale}, cyclic
of degree five, and constant over $K$. Infinity is a rational
point above $P$, and its deck translates split that fibre.
Its genus is $1+5(3-1)=11$.

We give the entire finite cohomological test, not a point-count
guess. Here $\deg F=19$, and $f=x^3-x$, $g$, and $h$ are
pairwise coprime and squarefree, of degrees three, six, and two.
A basis of regular differentials on the normalization is
\[
 \left\{\frac{x^jdx}{v}:0\le j\le3\right\},\quad
 \left\{\frac{gh\,dx}{v^2}\right\},\quad
 \left\{\frac{ghx^jdx}{v^3}:0\le j\le5\right\}.       \tag{B.Q.2}
\]
At a root of $f$ the ramification index is four; at a root of
$gh$ it is two. Checking the orders there proves regularity.
At infinity, $v_\infty(x)=-4$, $v_\infty(v)=-19$ and
$v_\infty(dx)=-5$, giving exactly the displayed ranges.
The eleven forms are independent by their distinct cyclic
characters and powers of $x$, so the genus calculation makes
them a basis.

For a differential $H(x)dx/v^a$, choose $b\in\{1,2,3\}$
with $3b\equiv a\pmod4$. Cartier is computed exactly by
\[
 \frac{1}{v^b}\,\mathcal C_{\mathbf P^1}
       \bigl(H(x)F(x)^{(3b-a)/4}dx\bigr),\qquad
 \mathcal C_{\mathbf P^1}\left(\sum c_ix^i dx\right)
       =\sum c_{3j+2}^{1/3}x^jdx.
\]
Division by $gh$ in the relevant character spaces is checked
with zero remainder. If $M$ is its inverse-Frobenius-semilinear
matrix, eight iterations produce its $K$-linear matrix. Every
entry of $M$, that eighth power, and the cover polynomials is
included in \path{code/quartic_cyclic5.json}. For each of the
six covers the result is
\[
     \det(T-\mathcal C^8)=T^3(T+1)^8.                 \tag{B.Q.3}
\]
The same calculation on $C$, with basis
$dx/w^2,dx/w^3,x\,dx/w^3$, gives Cartier zero.
Serre duality identifies the linear eighth powers of Cartier
and coherent Frobenius as transposes. The Artin--Schreier
sequence identifies the Frobenius-stable coherent subspace
with geometric three-torsion cohomology. Thus the cover's
new three-primary part has rank eight and Frobenius eigenvalue
$-1$ only; its relative class number is a unit. Equation
(B.Q.3) also gives the exact odd-extension condition. This
conclusion is insensitive to possible Jordan blocks.

For completeness, exact enumeration over $\mathbf F_3$,
$\mathbf F_9$, and $\mathbf F_{27}$ gives $4,4,28$ points
on $C$. Newton identities and the reciprocal functional
equation give its Frobenius polynomial
\[
                     (T^2+3)(T^2-3)^2.
\]
For an abelian variety over a finite field,
\citep[Chapter~II, Lemma~1.3 and Remark~1.4]{MilneAV2008}
prove that the algebra generated by Frobenius is a product
of fields. Its minimal polynomial is therefore squarefree,
so $F^4=[9]$ on $J_C$;
equality on a Tate module implies equality of these
endomorphisms. On every order-five character this says
$F^4=-1$ and $F^8=1$. Hence every nonzero character has
period exactly eight, proving the constant-field minimum.
Finally three is excluded in characteristic three and five
is the first admissible odd prime. The independently loaded
certificates verify the finite arithmetic in this argument.
\end{proof}

\begin{claim}{B.31}{Proved}{A four-cover all-extension negative test over \(\mathbf F_5\)}
Let $C/\mathbf F_5$ have the same equation and point as in B.30.
Pull its quadratic map to $E:y^2=x^3-x$ back along the duals
of the four cyclic degree-three isogenies out of $E$, over
$K=\mathbf F_{5^8}$. These four constant $P$-split covers
have genus seven. Every one has relative class number divisible
by five over $K$ and over every further constant extension.
This assertion concerns these four covers; it does not exclude
the other order-three characters of $J_C$ or other primes.
\end{claim}
\begin{proof}
The program \path{code/quartic_cyclic3_f5.py} computes the four
kernel pairs on $E$ from the roots of its three-division
polynomial $3x^4-6x^2-1$. For a kernel pair generated by
$(a,b)$, the general odd-characteristic isogeny formula is
\[
 h=x-a,\quad N=xh^2+2(3a^2+A)h+4b^2,
 \quad g=N'h-2Nh',
\]
where $E$ has equation $y^2=x^3+Ax+B$. The identity
$fg^2=N^3+A'Nh^4+B'h^6$ is checked coefficient by
coefficient. The image of a different three-torsion point
generates the dual kernel on $E'$. Applying the formula there
gives an isogeny with target coefficients $3^4A,3^6B$;
scaling by three identifies that target with the original $E$.
Thus we have used the dual isogeny, without identifying an
ordinary isogenous elliptic curve with $E$ without proof.

With $f'=x^3+A'x+B'$ on $E'$, the pullback cover is
\[
        v^4=f'g^2h^2/3^6,
        \qquad (X_C,W_C)=(N/(9h^2),v/h^2).
\]
It is a connected constant \emph{\'etale} degree-three
$P$-split cover by the same disjointness and base-change
argument as B.30. Now $\deg F=11$, and $\deg(gh)=4$.
The regular basis consists of $x^jdx/v$ for $j=0,1$,
$gh\,dx/v^2$, and $ghx^jdx/v^3$ for $0\le j\le3$.
The same valuation calculation proves it is a full basis.
Cartier extracts the coefficients numbered $5j+4$ from
$H(x)F(x)^a$ in character $a$, and takes fifth roots.

The exact eighth powers have characteristic polynomials
$(T-1)^3$ on the old genus-three cohomology and
$(T-1)^7$ on the genus-seven cover. Polynomial division
therefore gives $(T-1)^4$ for the new factor, with value
zero at one. The old factor has been divided as a polynomial;
we have not divided the two zero values at one.
The new geometric five-torsion has a Frobenius fixed vector.
It remains fixed by every power, proving the all-extension
failure. All four input kernels, both isogenies, the cover
polynomials, and the matrices are in the accompanying JSON.
\end{proof}

\begin{claim}{B.32}{Open}{The unbounded quartic classification in B6}
The least witness over arbitrary constant extensions for
\emph{every} smooth plane quartic over $\mathbf F_3$ and
$\mathbf F_5$ remains undecided here. B.30 is an exact
least-prime and first-field result for one pointed quartic;
B.31 excludes four actual covers at all later fields.
The independent program \path{code/pointed_quartic_tests.py}
also computes the point counts and Weil polynomials of
$X^4+Y^4+Z^4=0$ over $\mathbf F_3$, at $(1:1:1)$,
and $X^4+Y^4-Z^4=0$ over $\mathbf F_5$, at $(1:0:1)$.
Their base-field Jacobian orders are $64$ and $256$,
respectively, so they have no base-field odd-prime Lang
quotient. This is not an all-extension exclusion.

The previous return's large base-field census programs were
not attached and have not been rerun here. We do not treat
their numerical tables as newly verified certificates.
B.4 reduces each fixed-prime all-extension question to a
finite calculation, but gives no terminating bound on the
prime search. It also remains undecided whether some curve
has a nonabelian witness of the stated type and has no
class-number witness over any extension. Failure of finitely
many characters, primes, or fields cannot prove that assertion.
\end{claim}


\subsection{A first-extension witness for the pointed Fermat quartic in characteristic five}

\begin{claim}{B.33}{Proved}{A constant cyclic thirteen-cover over \(\mathbf F_{125}\)}
Let
\[
 C/\mathbf F_5:\quad W^4=X^3Z-XZ^3,\qquad P=(1:0:0).
\]
Over $K=\mathbf F_{125}$ there is a constant cyclic, $P$-split,
\emph{\'etale} cover $D\to C$ of degree thirteen with
\[
  g_D=27,\qquad
  \frac{\#J_D(K)}{\#J_C(K)}\equiv1\pmod5.
\]
Its new Jacobian part is ordinary of dimension twenty-four.
Under an extension of $K$ of degree $n$, its relative class
number stays prime to five if and only if $4\nmid n$.
\end{claim}
\begin{proof}
All the following defining polynomials lie in $\mathbf F_5[x]$:
\[
\begin{split}
 h={}&x^6+x^2+4,\\
 N={}&x^{13}+3x^{11}+x^9+3x^7+2x^5+x,\\
 g={}&x^{18}+2x^{16}+x^{14}+3x^{12}+3x^{10}+2x^2+4.
\end{split}                                                    \tag{B.QM.1}
\]
Direct polynomial multiplication gives
\[
 g=N'h-2Nh',\qquad
       (x^3-x)g^2=N^3-Nh^4.                                    \tag{B.QM.2}
\]
The polynomials $x^3-x,g,h$ are squarefree and pairwise coprime,
and $N$ is coprime to $h$. Consequently
\[
 E:\ y^2=x^3-x\longrightarrow E,\qquad
       (x,y)\longmapsto
       \left(\frac{N(x)}{4h(x)^2},\frac{y g(x)}{2h(x)^3}\right)
                                                                  \tag{B.QM.3}
\]
is a separable degree-thirteen isogeny. Indeed (B.QM.2) verifies
the target equation, the map takes infinity to infinity, and
the rational function defining its first coordinate has degree
thirteen. A nonconstant morphism between elliptic curves taking
the origins to each other is an isogeny; its degree is thirteen
because the $x$-coordinate maps on source and target both have
degree two.

The kernel is constant over the displayed cubic field. A fully
explicit presentation and generator are
\[
 K=\mathbf F_5[z]/(z^3+4z^2+3z+1),\qquad
                    T=(z,z+3)\in E(K).                           \tag{B.QM.4}
\]
The cubic has no root in $\mathbf F_5$, so this is a field.
The elliptic addition law gives $13T=0$, with $T\ne0$.
The twelve nonzero multiples of $T$ have precisely the six
distinct roots of $h$ as their $x$-coordinates. These assertions
are also a short exact arithmetic certificate: its thirteen
points, including the identity, are listed in
\path{code/criterion_quartic13_verify.json}, in the basis
$1,z,z^2$. Thus they constitute the entire kernel of (B.QM.3).

Pull back (B.QM.3) along $C\to E$, $(X,W)\mapsto(X,W^2)$.
The normalization of the resulting connected curve is
\[
 D:\quad w^4=F(x),\qquad
 F(x)=4(x^3-x)g(x)^2h(x)^2,\qquad
 (X_C,W_C)=\left(\frac{N(x)}{4h(x)^2},\frac{w}{h(x)^2}\right).
                                                                  \tag{B.QM.5}
\]
One can verify connectedness directly: $F$ has simple zeros
at the three roots of $x^3-x$, so the fourth-root extension
has degree four even over an algebraic closure. Alternatively,
the degree-two extension $C/E$ and the odd-degree isogeny are
linearly disjoint. The map to $C$ is finite \emph{\'etale} by
base change of the isogeny; it is a constant cyclic cover over
$K$ by (B.QM.4). The fibre of $P$ is the rational kernel of
(B.QM.3), so it splits completely. Riemann--Hurwitz gives
$g_D=1+13(3-1)=27$.

We give the cohomological calculation explicitly enough to
reproduce it without an isogeny-search program. Here $\deg F=51$
and $\deg(gh)=24$. A basis of regular differentials is
\[
 \left\{\frac{x^jdx}{w}:0\le j\le11\right\},\qquad
 \left\{\frac{gh\,dx}{w^2}\right\},\qquad
 \left\{\frac{ghx^jdx}{w^3}:0\le j\le13\right\}.                \tag{B.QM.6}
\]
At a simple zero of $F$ the ramification index is four, while
at a double zero it is two. The factors $gh$ exactly remove
the poles in characters two and three at the latter zeros.
At infinity the orders of $x,w,dx$ are $-4,-51,-5$;
these give exactly the ranges in (B.QM.6). The forms are
independent by their characters and their powers of $x$, and
their number is the genus, proving that this is a full basis.

Put $h_1=1$ and $h_2=h_3=gh$. Since $5\equiv1\pmod4$,
Cartier in character $s$ is computed by
\[
 \mathcal C\!\left(\frac{h_s x^jdx}{w^s}\right)
   =\frac1{w^s}\,
      \mathcal C_{\mathbf P^1}(h_sx^jF^s dx),\qquad
 \mathcal C_{\mathbf P^1}\!\left(\sum a_i x^i dx\right)
        =\sum a_{5b+4}x^b dx.                                  \tag{B.QM.7}
\]
No coefficient roots need to be represented here because
all coefficients lie in $\mathbf F_5$. Divide the extracted
polynomial by $h_s$ and expand in the relevant part of
(B.QM.6); every division has zero remainder. This gives a
$27\times27$ matrix $M$ over $\mathbf F_5$. Its independently
computed characteristic polynomials are
\[
\begin{split}
 \det(U-M)
   &=(U-2)^5(U+2)^6
       (U^2-2U-1)^4(U^2+2U-1)^4,\\
 \det(U-M^3)&=(U-2)^{14}(U+2)^{13}.
\end{split}                                                    \tag{B.QM.8}
\]
For the old curve, the basis $dX/W^2,dX/W^3,X\,dX/W^3$ gives
\[
 M_C=\begin{pmatrix}3&0&0\\0&0&4\\0&1&0\end{pmatrix},\qquad
                \det(U-M_C^3)=(U-2)^2(U+2).                    \tag{B.QM.9}
\]
The entire matrix, its cube, and all intermediate polynomials
are in the certificate. The independent checker
\path{code/criterion_quartic13_verify.py} implements its own
prime-field polynomial arithmetic, the cubic field, the
elliptic group law, and (B.QM.7); it then computes the
characteristic polynomials by integer linear algebra reduced
modulo five. It imports no search-producer code.

Pullback and norm split old and new five-primary cohomology,
since thirteen is invertible there. Cartier and coherent
Frobenius are dual after a full finite-field cycle, and the
Artin--Schreier sequence identifies the stable part with
geometric five-torsion cohomology. Polynomial division of
(B.QM.8) by (B.QM.9) gives the new factor
\[
                  (U-2)^{12}(U+2)^{12}.                        \tag{B.QM.10}
\]
It has degree twenty-four and no zero root, proving the
claimed ordinarity. Its value at one is
$(-1)^{12}3^{12}=1$ in $\mathbf F_5$, so the relative
class number is a unit, with the asserted residue.
Both $2$ and $-2$ have multiplicative order four in
$\mathbf F_5$. The eigenvalues of a further $n$th Frobenius
power contain one exactly when $4\mid n$; possible Jordan
blocks do not affect this criterion. This proves the final
assertion.
\end{proof}

\begin{claim}{B.34}{Proved}{The exact first-field minimum and the pointed Fermat identification}
For the pointed quartic of \statementref{B.33}, the least
constant extension degree over which an odd-prime constant
cyclic $P$-split class-number witness exists is exactly three.
At that first field, thirteen is the least possible prime,
and \statementref{B.33} attains it. The same assertions hold
for the pointed Fermat quartic
\[
       U^4+V^4=R^4\quad\text{over }\mathbf F_5,
                   \qquad (U:V:R)=(1:0:1).
\]
This is a field-first minimum. The least prime when arbitrary
constant extension degrees are allowed is not determined here.
\end{claim}
\begin{proof}
Direct enumeration gives
\[
 \#C(\mathbf F_5)=8,\qquad
 \#C(\mathbf F_{25})=44,\qquad
 \#C(\mathbf F_{125})=104.
\]
Newton identities and the reciprocal functional equation
give the full Jacobian Frobenius polynomial
\[
          H_C(U)=(U^2+2U+5)^2(U^2-2U+5).                       \tag{B.QM.11}
\]
The independent certificate enumerates these points, verifies
the first three Newton identities, and checks the resulting
polynomial. Its consequences for the first three fields are
\[
\begin{split}
 \#J_C(\mathbf F_5)&=8^2\cdot4=256,\\
 \#J_C(\mathbf F_{25})&=32^3=32768,\\
 \#J_C(\mathbf F_{125})&=104^2\cdot148
                         =2^8\cdot13^2\cdot37.
\end{split}                                                    \tag{B.QM.12}
\]
A constant cyclic $P$-split cover of odd prime degree
$\ell\ne5$ gives a quotient $J_C(\mathbf F_{5^m})\twoheadrightarrow
\mathbf Z/\ell$, by the pointed Abel and Lang construction.
The first two orders in (B.QM.12) exclude every such odd prime
over degrees one and two. The third order permits only
thirteen and thirty-seven. The thirteen-cover already proved
therefore establishes both claimed minima.

Finally the characteristic-five identity
\[
       (X+2Z)^4-(X-2Z)^4=X^3Z-XZ^3
\]
shows that the invertible projective linear map
\[
       (X:W:Z)\longmapsto(X-2Z:W:X+2Z)
\]
carries $C$ to the displayed Fermat curve and $P$ to
$(1:0:1)$. Thus the result concerns that exact pointed
Fermat example, not merely another curve with the same
Jacobian point counts.
\end{proof}


\subsection{Supersingular correspondences: a larger proved class}

The general assertion in B5 concerns a finite \'etale cover whose
geometric Galois closure has order prime to $p$. This condition is
stronger than requiring the degree of a particular non-Galois cover
to be prime to $p$. It is retained throughout the constructions below.

\begin{claim}{B.35}{Proved: placing finitely many values over prime-to-$p$ torsion}{B5}
Let $E/\overline{\F}_p$ be an elliptic curve, let
$\pi:E\to E/\{\pm1\}\simeq\mathbf P^1$ be its quotient map, and
let $S\subset\mathbf P^1(\overline{\F}_p)$ be finite. There exists
$\phi\in\mathrm{PGL}_2(\overline{\F}_p)$ such that every point of
$\phi(S)$ is the image under $\pi$ of a prime-to-$p$ torsion point
of $E$. The points of $\phi(S)$ may also be required to avoid the
branch locus of $\pi$.
\end{claim}
\begin{proof}
Choose coordinates putting both $S=\{t_1,\ldots,t_N\}$ and the
branch locus $B$ of $\pi$ in $\mathbf A^1$. This is permissible even
in characteristic two; the quotient map is then separable of degree
two, although ramification may be wild. Choose $c\ne0$ outside the
finite set determined by
\[
 c(t_i-t_j)=b-b'\qquad(i\ne j,\ b,b'\in B).
\]
Thus the finite sets $B-ct_i$ are pairwise disjoint. Descend all this
data to a finite field $k=\F_q$.

If $E$ is supersingular, put $a=0$ and $v=\operatorname{id}_E$.
If $E$ is ordinary, its $p$-adic Tate module has rank one. After
enlarging $k$, write the Frobenius action as multiplication by
$\alpha\in1+p\Z_p$, and put $a=v_p(\alpha-1)<\infty$.
For every $n$ prime to $p$,
\[
 v_p(\alpha^n-1)=a,
\]
because $(\alpha^n-1)/(\alpha-1)\equiv n\pmod p$.
Thus the $p$-primary subgroup of $E(\F_{q^n})$ is killed by $p^a$.
Let $v:E'=E^{(p^a)}\to E$ be the iterated Verschiebung. It is an
\'etale isogeny of degree $p^a$ in the ordinary case, and
\[
 v(E'(\F_{q^n}))=[p^a]E(\F_{q^n})
                 =E(\F_{q^n})\{p'\},\qquad p\nmid n.
\]
Here relative Frobenius is a bijection on points over a perfect
finite field and $v\circ F^a=[p^a]$. In the supersingular case all
geometric torsion already has order prime to $p$.

The composite $h=\pi\circ v:E'\to\mathbf P^1$ is Galois with
group $G=\ker(v)\rtimes\{\pm1\}$ and branch locus $B$.
Its branch inertia groups normally generate $G$: the quotient by
their normal closure would be a connected finite \'etale cover of
$\mathbf P^1$, which must have degree one by Riemann--Hurwitz.
Consider the $N$ covers of the $s$-line given by
\[
 h(R_i)=s+ct_i.
\]
Each is connected and has branch locus $B-ct_i$. Their compositum
has geometric Galois group $H\subseteq G^N$ projecting onto each
factor. At a branch point belonging only to the $i$th factor, its
inertia subgroup is supported in that factor and projects onto the
corresponding inertia of $G$. The other covers are unramified there.
Conjugating by $H$ and using the surjection to the $i$th factor shows
that $H$ contains the normal closure of every such inertia in that
factor. These normal closures generate $G$. Consequently $H$ contains
every factor and equals $G^N$. This also proves geometric
irreducibility of the normalized fibre product; it does not assume
linear disjointness without checking it.

Let $Y/k$ be its smooth proper model. Fix a prime $l\ne p$ and take
$n=l^m$. The Weil bound makes $|Y(\F_{q^n})|$ tend to infinity.
Only finitely many geometric points lie above the branch points or
infinity, so for sufficiently large $m$ there is a point outside that
set. It supplies $s,R_1,\ldots,R_N$ over $\F_{q^n}$ with
$s+ct_i=\pi(v(R_i))$. Each $v(R_i)$ has prime-to-$p$ order by the
previous paragraph. The transformation $t\mapsto s+ct$, composed
with the initial coordinate change of $S$, is the desired $\phi$.
\end{proof}

\paragraph{A finite point-count bound.}
The curve in the proof has degree $d^N$ over the $s$-line, where
$d=2p^a$. Since each individual cover has genus one and their branch
loci are disjoint, the different formula gives
$g(Y)=1+(N-1)d^N$, also when $p=2$. Thus the final existence step
can be made effective by choosing an $l$-power extension for which
the lower Weil bound exceeds the finite number of excluded geometric
points. The $p$-power isogeny in this auxiliary proof is not one of
the \'etale covers of the original curve in the next construction.

\begin{claim}{B.36}{Proved: every odd-characteristic hyperelliptic source}{B5}
Let $p$ be odd and let $C/\overline{\F}_p$ be hyperelliptic of genus
at least two. For every hyperelliptic curve $H/\overline{\F}_p$ there
is a connected finite \'etale cover $D\to C$ with prime-to-$p$
geometric Galois closure and a nonconstant map $D\to H$.
In particular B5 holds for every such $C$, including characteristic
three, with a supersingular target of genus at least two.
\end{claim}
\begin{proof}
Let $C\to\mathbf P^1$ be the degree-two map, with branch set $T$
of cardinality $2g(C)+2$. Choose four points of $T$ and let
$E\to\mathbf P^1$ be the elliptic double cover branched there.
The normalization $C_0$ of $C\times_{\mathbf P^1}E$ is connected:
the two quadratic function fields are different because their
branch sets are different. By tame ramification cancellation,
$C_0\to C$ is \'etale of degree two. The map $C_0\to E$ is a
double cover ramified above the points lying over the remaining
members of $T$. Translate the group law of $E$ so that one such
branch point is its origin. Then the entire fibre of $C_0\to E$
over zero has ramification index two.

Choose the quotient map $\pi:E\to E/\{\pm1\}$ and write $S$
for the branch set of a hyperelliptic map $H\to\mathbf P^1$.
By B.35, after changing the coordinate on that target, its branch
points are $\pi(P_i)$ for finitely many prime-to-$p$ torsion points
$P_i\in E$. Choose an integer $m$ prime to $p$ killing all $P_i$.
Pull $C_0\to E$ back along $[m]:E\to E$ and normalize, obtaining
$C_1$. The map $C_1\to C_0$ is finite \'etale with group $E[m]$.
The double cover $C_1\to E$ is still connected, since its nonempty
branch divisor is the inverse image of the original branch divisor.
Above every point of $E[m]$ its ramification index is two.

Consequently every point of $C_1$ above a branch value $\pi(P_i)$
has even ramification index in the composite $C_1\to E\to\mathbf P^1$:
the two inverse images under $\pi$ are $P_i,-P_i\in E[m]$.
The normalization of $C_1\times_{\mathbf P^1}H$ is therefore
\'etale over $C_1$, by Abhyankar's lemma. Any connected component
$D$ dominates both $C_1$ and $H$, and has degree one or two over
$C_1$. The tower $D\to C_1\to C_0\to C$ consists of finite
\'etale Galois steps of prime-to-$p$ orders. Its Galois closure has
prime-to-$p$ order: inductively its group is a subgroup of an
extension of the preceding Galois group by a product of conjugate
step groups, equivalently of the corresponding finite wreath
product. This verifies the stronger Galois-closure condition.

For an explicit supersingular target choose an odd prime power
$Q=p^b\ge5$ and let
\[
 H_Q:\quad y^2=x^Q-x.
\]
Its genus is $(Q-1)/2$. For $x\in\F_{Q^2}\setminus\F_Q$, the
nonzero value $x^Q-x$ has trace zero. All such values have the same
quadratic character
$\epsilon=(-1)^{(Q+1)/2}$ in $\F_{Q^2}$: writing a trace-zero
element as $a$, one has $a^{Q-1}=-1$ and hence
$a^{(Q^2-1)/2}=(-1)^{(Q+1)/2}$; elements of $\F_Q^\times$ are
squares in $\F_{Q^2}$. Therefore
\[
 |H_Q(\F_{Q^2})|=Q^2+1+\epsilon(Q^2-Q).
\]
This attains one of the Weil bounds. All Frobenius eigenvalues
over $\F_{Q^2}$ are consequently $-\epsilon Q$, so all Newton
slopes are $1/2$. Thus $H_Q$ is supersingular. For $p=3$, take
$Q=9$, yielding genus four; for $p\ge5$, take $Q=p$.
\end{proof}

The first double-cover construction is the one used in
Bogomolov--Tschinkel's hyperelliptic argument. The extra assertion
proved above is that the finite torsion positioning can be arranged
with prime-to-$p$ torsion, which keeps the final covering tower in
the required formation. We make no priority claim for this
strengthening. The ramification cancellation is the precise
one-dimensional form of SGA1 XIII Proposition~5.2 and Poonen
Lemma~2.1.\citep{LA_BogomolovTschinkel2002,LA_Poonen2005,LA_SGA1}

\begin{claim}{B.37}{Proved: a characteristic-two ramification criterion}{B5}
Let $p=2$. Suppose a curve $C/\overline{\F}_2$ admits a finite
separable map to $\mathbf P^1$ with three distinct fibres in which
every ramification index is divisible by three. Then a connected
finite \'etale cover of $C$ of degree one or three maps nonconstantly
to a supersingular elliptic curve. In particular this applies to
every cyclic degree-three cover of $\mathbf P^1$ of genus at least
two.
\end{claim}
\begin{proof}
Move the three values to $0,1,\infty$ and use the cyclic cover
$E:y^3=x(x-1)$. Its smooth projective model is elliptic, branched
with index three at those values. Normalized base change along
$C\to\mathbf P^1$ is \'etale by Abhyankar's lemma, and any
connected component has cyclic degree one or three over $C$.
After interchanging the coordinates, the elliptic curve is
$z^2+z=t^3$. It has nine points over $\F_4$: for each of the four
values of $t$, the right side has trace zero to $\F_2$, so there
are two values of $z$, in addition to the point at infinity.
Thus its two $\F_4$-Frobenius eigenvalues are $-2$ and it is
supersingular. A cyclic cubic cover has uniform ramification
index three over each branch point, and has at least three such
points when its genus is at least two.
\end{proof}

\begin{claim}{B.38}{Proved: the first three low-genus Galois cases}{B5}
Suppose $p$ is odd, $C/\overline{\F}_p$ has genus at least two,
and $G\subset\Aut(C)$ has order prime to $p$. Put $Y=C/G$ and
suppose $g(Y)\le2$. If at least one of the following holds, then
$C$ has a connected finite \'etale cover with prime-to-$p$ Galois
closure mapping to a supersingular curve of genus at least two:
\begin{enumerate}
\item $g(Y)=1$ or $2$;
\item $G$ is solvable;
\item two distinct branch values of $C\to Y$ have ramification
indices with a common prime factor.
\end{enumerate}
In particular this proves B5 for every tamely superelliptic curve
in every odd characteristic, including characteristic three.
\end{claim}
\begin{proof}
We verify the extra prime-to-$p$ condition in Poonen's
Theorem~1.7, cases~(1)--(3), rather than infer it from the word
``unramified.'' His theorem assumes precisely a prime-to-$p$
Galois action, genus at least two for the source, and quotient
genus at most two, with one of the listed extra hypotheses.
The proof below follows his \S\S3.1--3.3 and records the orders
of the covering groups.\citep[Theorem~1.7 and \S3]{LA_Poonen2005}

First suppose $Y$ has genus two. It is already hyperelliptic,
so no \'etale covering is needed to dominate a hyperelliptic curve.
Suppose instead that $Y=E$ is elliptic. The map $C\to E$ is
ramified, since $g(C)>1$. Translate a branch value to the origin,
and choose a prime $e\ne p$ dividing its ramification index.
The Galois hypothesis ensures that this divisibility holds at
every point above the origin. Pull back along $[l]:E\to E$ for
an odd prime $l\ge5$ distinct from $p$, and choose a connected
component $C'$. Then $C'\to C$ is finite \'etale Galois with
group a subgroup of $E[l]$. Every point of $C'$ above $E[l]$
has ramification index divisible by $e$ in $C'\to E$.

Under $E\to E/\{\pm1\}=\mathbf P^1$, the nonzero $l$-torsion
points give $(l^2-1)/2$ distinct values $a_i$. If $e\ge3$, take
the cyclic cover
\[
 H:\quad z^e=\frac{(x-a_1)(x-a_2)}{(x-a_3)(x-a_4)}.
\]
It has genus $e-1\ge2$, and its $z$-map has degree two, so it
is hyperelliptic. If $e=2$, use three factors in the numerator
and three in the denominator, giving a genus-two hyperelliptic
curve instead. At every branch value of $H\to\mathbf P^1$,
all ramification indices of $C'\to\mathbf P^1$ are divisible
by $e$. A connected component of the normalized pullback of
$H$ is therefore \'etale over $C'$, with cyclic group of order
one or $e$, and dominates $H$.

Now let $Y=\mathbf P^1$ and $G$ be solvable. Induct on $|G|$.
For a nontrivial proper normal subgroup $N$, if $g(C/N)\ge2$
apply induction to $C/N\to\mathbf P^1$ and pull its \'etale
cover back to $C$. If $g(C/N)=1$, apply the elliptic case to
$C\to C/N$. If $g(C/N)=0$, apply induction to $C\to C/N$.
It remains to treat $G$ cyclic of prime order $e\ne p$.
For $e=2$, $C$ itself is hyperelliptic. For $e\ge3$, pull
back along a cyclic degree-$e$ map $\mathbf P^1\to\mathbf P^1$
branched at two branch values of $C$. The two degree-$e$
extensions are different (the original curve has positive genus),
so the pullback is connected and gives an \'etale cyclic
degree-$e$ cover $C'\to C$. Its new cyclic map to $\mathbf P^1$
has at least four branch values, since its genus has increased.
Using four of these values in the preceding displayed equation
again gives an \'etale cyclic pullback dominating a hyperelliptic
curve.

Finally suppose two branch indices have a common prime divisor
$e$. Pull back along a cyclic degree-$e$ map of projective lines
branched at those values. A third branch value has at least two
inverse images; a prime factor $e'$ of its ramification index
therefore occurs at two new branch values. Repeat this construction.
Every step on the source is cyclic \'etale of order dividing the
chosen prime, hence prime to $p$. Every new map to the projective
line is Galois of degree at most $|G|$, with ramification indices
dividing $|G|$. The degrees of the composite maps of projective
lines tend to infinity, forcing the degrees of the covers of $C$
and hence their genera to tend to infinity. Riemann--Hurwitz then
forces the number of branch values to tend to infinity. There are
only finitely many possible indices, so at some stage at least six
branch values have a common index $a>1$. Choose a prime $e\mid a$
and a cyclic degree-$e$ cover of $\mathbf P^1$ branched at these
six points, using three exponents $1$ and three exponents $-1$.
It has genus $2(e-1)\ge2$. Its normalized pullback is \'etale
over the stage curve, and its solvable cyclic group puts the target
in the preceding case.

Thus in every case a hyperelliptic curve is dominated after a
finite tower of \'etale Galois prime-to-$p$ covers and their base
changes. A connected base change of such a Galois cover has a
subgroup of the original group; finite composition has prime-to-$p$
Galois closure by the wreath-product argument in B.36. Composing
with B.36 (and pulling its cover back) gives the claimed
supersingular target.
\end{proof}

\paragraph{The modular fourth case, with its group condition.}
Poonen's fourth condition asks for three branch indices divisible
by $2,3,l$, respectively, where $l\le89$ is prime or
$l\in\{101,103,107,131,167,191\}$. For $l=2,3$ the preceding
common-factor case applies. For $l\ge5$ his proof uses
$X(l)\to\mathbf P^1$ with geometric group
$\mathrm{PSL}_2(\F_l)$, then solvable covers of $X_0(l)$ and the
Atkin--Lehner quotient $X_0^*(l)$ of genus at most two.
Consequently the same prime-to-$p$ conclusion follows if additionally
$p\nmid l(l^2-1)$: the modular pullback, its Borel-subgroup steps,
and the order-two Atkin--Lehner step all have prime-to-$p$ Galois
closure. In the $l=5$ case, the new Galois map to $X(5)\simeq
\mathbf P^1$ has repeated branch indices and the third case applies.
Without the additional group-order check, tameness of the local
indices $2,3,l$ alone is insufficient to certify the required
Galois closure.\citep[Lemma~2.3 and \S3.4]{LA_Poonen2005}

\paragraph{The general reduction still required.}
B.36 also applies whenever a prime-to-$p$ \'etale cover of the
given odd-characteristic curve maps to a hyperelliptic curve.
B.38 establishes this under the listed low-genus Galois
hypotheses, but a general curve is not assumed to have such a map.

For a general curve, a tame map to $\mathbf P^1$ does not suffice.
Its ramification over a branch value may include unramified points,
where a branched target cover would remain ramified after base
change. Nor need the Galois closure of a tame map have
prime-to-$p$ order. Thus the tame Belyi theorem does not provide
the missing divisible fibres or the missing prime-to-$p$ covering
tower. The general universal-correspondence assertion is explicitly
Conjecture~2.6 in Bogomolov--Tschinkel; the theorem there is for
hyperelliptic sources and $p\ge5$. Neither that conjecture nor the
present special cases settle B5 for every nonhyperelliptic curve
or for all curves in characteristic two.
\citep[Theorem~1.7 and \S3]{LA_Poonen2005}
\citep[Theorem~1.7 and Conjecture~2.6]{LA_BogomolovTschinkel2002}
\citep[Theorem~1.5]{LA_KedlayaLittWitaszek2023}

The ordinarity literature does not supply the missing conclusion.
Raynaud's Theorem~2 constructs, for a proper smooth curve of genus
at least two in characteristic $p$, an \'etale cover with a solvable
prime-to-$p$ group which is nonordinary. Nonordinarity is weaker than
having a supersingular isogeny factor, and still weaker than the
required map to a supersingular curve. Thus this result cannot be
substituted for B5. The Prym results of Ozman--Pries likewise give
$p$-rank information, with no universal supersingular-target assertion.
\citep[Theorem~2]{LA_Raynaud2000}
\citep[Theorem~1.1 and \S8.2]{LA_OzmanPries2019}


\subsection{A positive-rank primitive-zero cover without a supersingular factor}

\begin{claim}{B.39}{Proved}{An exact cyclic-thirteen distinction between B4 and B5}
Let $C/\F_2$ be the smooth proper model
\[
 C:\quad Y^2+X^2Y=X^5+X,
 \qquad P=\infty.
\]
It has genus two, $2$-rank one, and no nonzero supersingular
isogeny factor in its geometric Jacobian. Over $k=\F_{64}$
there is a connected $P$-pointed cyclic \'etale cover
$D\to C_k$ of degree $13$, with zeta numerator
\[
 P_D(T)=
 (1-T-80T^2-64T^3+4096T^4)
 (1+12T+64T^2)^{12}.
 \tag{B.LIT.7}
\]
Its primitive degree-$13$ factor has Newton slopes
$1/3$ and $2/3$, each with multiplicity $12$, and hence
has $2$-rank zero. Its relative class number is
\[
 \frac{|J_D(k)|}{|J_C(k)|}=77^{12}
       =43439888521963583647921,
\]
which is odd. Nevertheless $D$ has no nonconstant map,
even over $\overline{\F}_2$, to any supersingular curve
of positive genus. Thus this is a B1 witness and a
nontrivial positive-$2$-rank case of B4, but this particular
cover does not supply the correspondence asked for in B5.
\end{claim}

\begin{proof}
The substitution $x=X^{-1}$, $y=Y/X^2$ identifies the
function field with
\[
 y^2+y=x^3+x^{-1}.
\]
The Artin--Schreier poles have odd orders three and one.
Their different exponents are four and two, so
Riemann--Hurwitz gives genus two. In the polynomial model
the regular differential $dX/X^2$ has divisor $2P$.
Direct counts give $|C(\F_2)|=4$ and $|C(\F_4)|=4$;
the genus-two functional equation therefore gives the
Frobenius characteristic polynomial
\[
 H_2(U)=U^4+U^3+2U+4.
\]
It is irreducible modulo three and hence over $\Q$.
Its Newton slopes are $0,1/2,1/2,1$, which also verifies
the asserted $2$-rank. Taking sixth powers of its roots
gives
\[
 H_{64}(U)=U^4-U^3-80U^2-64U+4096,
 \qquad |J_C(k)|=H_{64}(1)=3952=16\cdot13\cdot19.
\]

Here is a complete definition of the cover used in the
computation. Write $k=\F_2[z]/(z^6+z+1)$ and encode
elements by the binary coefficients in the basis
$1,z,\ldots,z^5$. In reduced Mumford coordinates for the
displayed polynomial model, put
\[
 G=\bigl(X^2+25X+14,\;17X+26\bigr)\in J_C(k).
\]
The numbers in this formula are binary field encodings,
not integers embedded in a characteristic-two field.
The certificate verifies $G\ne0$ and $13G=0$.
Since the $13$-primary subgroup has order $13$, there is
a unique surjection
\[
 \chi:J_C(k)\longrightarrow\Z/13,
 \qquad [304]A=\chi(A)G.
\]
Pull back the Lang isogeny $F_{64}-1:J_C\to J_C$
along $R\mapsto[R-P]$ and quotient its deck group
$J_C(k)$ by $\ker\chi$. This defines $D$. The point
above $P$ represented by zero makes it pointed. The
Abel map induces the isomorphism from the abelianized
geometric fundamental group of $C$ to the fundamental
group of $J_C$; thus the pullback of this connected
cyclic isogeny cover is geometrically connected.
This is the precise Jacobian covering statement of
Milne, \S9, Proposition~9.1, applied after base change
to $\overline{k}$: that proposition assumes a separably
closed field, a smooth proper curve, and a chosen point.
\citep[\S9, Proposition~9.1]{MilneJac}

For every $R\in C(k)$ the program computes
$\chi([R-P])$ by exact Cantor arithmetic. Its complete
table, including every point and every element of
$\langle G\rangle$, gives
\[
 \#\{R:\chi([R-P])=0\}=16,
 \qquad
 \#\{R:\chi([R-P])=j\}=4
 \quad(1\le j\le12).
\]
Every polynomial division in the group law is checked
to be exact, and every output Mumford pair is checked
against the curve equation. The producer
\path{code/b5_char2_13_certificate.py} writes these finite
tables and the entire expanded degree-$28$ numerator to
\path{code/b5_char2_13_certificate.json}; its textual
output is also retained.

The independent verifier
\path{code/c0_cyclic13_verify.py} imports no Cantor
implementation. It uses fixed-point addition by the
Chinese remainder theorem and tangent Taylor expansion,
and independently reproduces all $64$ point projections,
all $13$ subgroup elements, and $H_{64}$ by a Newton-sum
recurrence. Its exact output and certificate are retained
beside the producer.

Let $\zeta$ be a primitive thirteenth root of unity.
The complete twisted polynomial has degree two. Its
linear coefficient, for every $i\ne0$, is
\[
 \sum_{R\in C(k)}\zeta^{i\chi([R-P])}
       =16+4\sum_{j=1}^{12}\zeta^{ij}=12.
\]
Because the canonical divisor is $2P$, its quadratic
coefficient is $64$. These are exactly the effective-divisor
and functional-equation formulas proved in B.2, so no
single evaluation at $T=1$ is substituted for a Frobenius
polynomial. Multiplying the twelve complete twisted
polynomials and the numerator of $C_k$ proves (B.LIT.7).
It also gives the relative class number and
$g(D)=13(g(C)-1)+1=14$.

The two roots of $U^2+12U+64$ have $2$-adic valuations
two and four: the point $(1,2)$ lies strictly below the
segment joining $(0,6)$ and $(2,0)$. Dividing valuations
by six gives slopes $1/3$ and $2/3$. Thus the primitive
factor has $2$-rank zero and no supersingular factor.
For completeness, $J_C$ has no geometric supersingular
factor either. If a Frobenius root $\alpha$ belonged to
one after a finite constant extension, then
$\alpha/\sqrt2$ would be a root of unity. One elementary
way to see this criterion is that a supersingular Weil
number divided by the appropriate square root of the
field size is an algebraic unit at every finite place
and has absolute value one at every complex place;
Kronecker's argument then makes it a root of unity.
The property descends from $\alpha^m$ to $\alpha$.
Irreducibility of $H_2$ would make all four of its roots
have this property, contradicting its slope-zero root.
This also proves geometric simplicity here: a geometric
elliptic decomposition of a $2$-rank-one surface would
include one supersingular elliptic factor.

Finally $J_D$ is isogenous to $J_C\times\Pr(D/C)$
by pullback and norm. A nonconstant map $D\to S$ of
smooth proper curves induces pullback and norm with
composition multiplication by its degree on $J_S$.
Consequently a positive-dimensional supersingular
$J_S$ would give a supersingular factor of $J_D$,
which the preceding two paragraphs exclude. This holds
for inseparable as well as separable maps, since a
nonzero multiplication map is invertible in the
isogeny category.
\end{proof}

\paragraph{The scope of the search.}
This does not refute the existence of a different odd
\'etale cover of $C$ with a supersingular factor. The
independent theta-Taylor elimination in
\path{code/b5_theta_division.py} excludes torsion orders
$3,5,7,9,11$ on the Raynaud theta curve of this example;
the order-$13$ candidate locus is $r^2+r+1=0$.
The companion probe records complete character sums
for exponent-one prime quotients over all
$\F_{2^m}$ with $1\le m\le10$. These bounded searches
are retained as evidence about the attempted route,
not as a classification of all odd covers or all
constant extensions.

\subsection{Two exact limits of the proposed correspondence constructions}

\begin{claim}{B.40}{Proved}{Ordinary cubic Pryms in two small genera}
Let $C$ be a smooth hyperelliptic curve over an algebraically
closed field of characteristic two, of genus $g\in\{2,3\}$
and $2$-rank zero. For every connected cyclic \'etale
cover $D\to C$ of degree three, the Prym variety is
ordinary of dimension $2(g-1)$. In particular,
$f_D=2(g-1)$. In genus three its genus-two reflection
quotient is ordinary, so this cubic-reflection construction
does not yield a supersingular target.
\end{claim}

\begin{proof}
The hyperelliptic Artin--Schreier map has exactly one
branch point, since its number of branch points is
$f_C+1=1$. Put this point at infinity, denote its
unique preimage by $P$, and use an equation
$y^2+y=f(x)$ with $\deg f=2g+1$. Then
\[
 \operatorname{div}(dx)=(2g-2)P,
 \qquad B^1\simeq\mathcal O_C((g-1)P),
\]
where $B^1$ is Raynaud's line bundle of locally exact
differentials on the Frobenius twist. We identify that
twist with its corresponding equation over the perfect
ground field; all arguments below are unchanged on it.

Suppose a nonzero three-torsion class $L$ satisfies
$H^0(B^1\otimes L)\ne0$. Removing copies of $P$ and
hyperelliptic conjugate pairs from an effective divisor
of degree $g-1$ gives
\[
 L\simeq\mathcal O_C(E-dP),\qquad 1\le d\le g-1,
\]
with $E$ reduced in the hyperelliptic divisor sense:
it contains no $P$ and no conjugate pair. There is a
rational function $h$ with
$\operatorname{div}(h)=3E-3dP$. For $g=2$ or $3$,
\[
 3d\le3(g-1)\le2g<2g+1.
\]
The pole orders of $x,y$ at $P$ are $2,2g+1$.
Consequently every function whose only pole is at
$P$ of order at most $2g$ belongs to $k[x]$.
Every finite zero of such a function occurs together
with its hyperelliptic conjugate and with the same
multiplicity; there are no finite branch points.
This contradicts the form of $E$. Thus neither $L$
nor $L^2$ lies on the theta divisor of $B^1$.

For clarity, the ordinarity implication uses the exact
sequence
\[
 0\longrightarrow L\longrightarrow F_*F^*L
   \longrightarrow B^1\otimes L\longrightarrow0.
\]
Since $H^0(L)=H^0(L^2)=0$, the just-proved vanishing
makes Frobenius an isomorphism between the two
$g-1$ dimensional twisted $H^1(\mathcal O)$ spaces.
These are precisely the nontrivial character summands
of $H^1(D,\mathcal O_D)$, as
$\pi_*\mathcal O_D=\mathcal O_C\oplus L\oplus L^2$.
Their stable Frobenius rank is $2(g-1)$, which proves
the assertion. This is also the rank-one instance of
Raynaud's ordinarity criterion in \S1.
\citep[\S1, before Proposition~1]{LA_Raynaud2000}

Here are the details of the reflection assertion in
characteristic two. Let $\delta$ generate the cyclic
deck group, and let $\iota$ be the hyperelliptic
involution on $C$. Since $\iota^*L=L^{-1}$,
$\iota$ lifts to an automorphism $\tau$ of $D$ with
$\tau\delta\tau^{-1}=\delta^{-1}$.
Write $\tau^2=\delta^a$. The element $\tau$ commutes
with its square and also sends $\delta^a$ to
$\delta^{-a}$, so $a=-a$ in $\Z/3$ and $a=0$.
Thus $\tau$ is an involution. It fixes exactly one
of the three points over $P$, since it acts there
as a reflection of the three-element deck torsor.
It has no other fixed points, as $\iota$ has none.
The \'etale map $D\to C$ identifies the completed
local rings at that fixed point and at $P$.
The local different for $D\to H:=D/\langle\tau\rangle$
therefore has exponent $2g+2$, the same as that of
$C\to\mathbb P^1$. Consequently
\[
 6g-6=2g_D-2=2(2g_H-2)+(2g+2),
 \qquad g_H=g-1.
\]

Let $A,A'$ be the pullback images in $J_D$ of the
Jacobians of the quotients by $\tau$ and by a distinct
conjugate reflection. Both lie in the Prym: their
norm to $J_C$ is fixed by the induced hyperelliptic
action $[-1]$, hence has finite image and is zero
on these connected abelian varieties. Their
intersection is finite. Indeed an element fixed by
the two reflections is fixed by $\delta$ as well.
Its norm $z$ to $J_C$ satisfies $2z=0$, while the
pullback of $z$ equals $3$ times that element; hence
the intersection is killed by six. Pullback from
each quotient has finite kernel by its pullback--norm
identity. Both images thus have dimension $g-1$,
and their addition gives an isogeny
\[
 J_H\times J_H\longrightarrow\Pr(D/C).
\]
For $g=3$ the quotient therefore has genus two and
is ordinary.
The inequality used above fails for $g\ge4$; no
assertion in those genera follows from this proof.
\end{proof}

\begin{claim}{B.41}{Proved}{A fixed simply branched function cannot supply the missing fibres}
Let $k$ be algebraically closed, let $f:C\to\mathbb P^1_k$
be finite separable, and assume that for every
$t\in\mathbb P^1(k)$,
\[
 \gcd\{e_P(f):P\in f^{-1}(t)\}=1.
\]
Let $C'\to C$ be any finite \'etale cover and
$h:\mathbb P^1\to\mathbb P^1$ any finite separable map.
The composite $h\circ f\circ\pi$ cannot have three
distinct fibres in each of which all ramification
indices are divisible by three. It cannot have four
distinct fibres in each of which they are all even.
More generally, uniform fibre divisibilities $m_i$
at distinct values require
\[
 \sum_i(1-1/m_i)<2.
\]
\end{claim}

\begin{proof}
The \'etale pullback preserves every ramification index
over the algebraically closed ground field. For a
fixed $t$ over $u$ the gcd of the indices of the
composite at points over $t$ is therefore $e_t(h)$.
Taking the gcd over all $t\in h^{-1}(u)$ shows that
any uniform divisor $m_i$ of the composite fibre
already divides every $e_t(h)$ in that fibre of $h$.
Write $n=\deg h$. There are at most $n/m_i$ geometric
points in such a fibre. Their contribution to the
different is at least
$\sum_{t\mid u_i}(e_t(h)-1)\ge n(1-1/m_i)$,
also in the presence of wild ramification.
Riemann--Hurwitz for $h:\mathbb P^1\to\mathbb P^1$
now gives
\[
 n\sum_i(1-1/m_i)\le2n-2<2n.
\]
The claimed patterns each give a sum of two and are
impossible. In particular, the hypothesis holds for
a simply branched map of degree at least three:
every branch fibre has an unramified point. Thus
rational postcomposition and \'etale pullback of that
fixed function cannot complete the general B5 route;
a new function on the covering curve would be needed.
\end{proof}

\begin{claim}{B.42}{Proved}{Three uniformly divisible fibres with a supersingular Fermat quotient}
Let $C/\overline{\F}_p$ be a smooth proper curve admitting
a finite separable map to $\mathbb P^1$ with three
distinct fibres in which every ramification index is
divisible by the same odd integer $r\ge3$.
Suppose $r\mid p^a+1$ for some $a\ge1$.
Then $C$ has a connected cyclic \'etale cover of degree
dividing $r$ which dominates a supersingular curve of
genus $(r-1)/2$. Its geometric Galois closure therefore
has order prime to $p$.
In characteristic two this applies to every cyclic
prime-degree cover of $\mathbb P^1$ of genus at least
two whose degree $\ell$ satisfies that
$\operatorname{ord}_\ell(2)$ is even.
\end{claim}

\begin{proof}
Send the three specified values to $0,1,\infty$.
The smooth proper model
\[
 H_r:\quad y^r=x(x-1)
\]
is a cyclic cover of degree $r$, totally ramified
at those three values and nowhere else. Since $r$
is odd, its genus is $(r-1)/2$. Normalizing
$C\times_{\mathbb P^1}H_r$ and taking any connected
component gives a cyclic cover $D\to C$ of degree
dividing $r$. At each of the three relevant fibres,
the tame local ramification index $r$ divides that
of $C\to\mathbb P^1$, so the normalized pullback is
unramified. It is unramified elsewhere as well.
Thus $D\to C$ is \'etale, and $D\to H_r$ is dominant.

Put $q=p^a$. The Fermat curve of degree $r$ maps
to $H_r$ by
\[
 U^r+V^r=1,\qquad x=U^r,\quad y=-UV.
\]
It is in turn dominated by the Hermitian Fermat
curve $U^{q+1}+V^{q+1}=1$, by taking coordinate
powers $(q+1)/r$. This latter curve has genus
$q(q-1)/2$ and exactly $q^3+1$ points over
$\F_{q^2}$. Indeed the norm map to $\F_q$ has
one preimage of zero and $q+1$ preimages of each
nonzero value. Counting affine solutions gives
$2(q+1)+(q-2)(q+1)^2=q^3-q$, and there are
$q+1$ points at infinity. This attains the upper
Weil bound, so every Frobenius root equals $-q$.
The curve is therefore supersingular. Pullback
and norm inject the Jacobian of every dominated
curve in the isogeny category; hence $H_r$ is
supersingular too.

For the last assertion, a cyclic cover of odd prime
degree $\ell$ has all ramification indices equal
to $\ell$ at its branch values. Genus at least two
forces at least three such values. An even order
of $2$ modulo $\ell$ gives
$2^{\operatorname{ord}_\ell(2)/2}=-1$ modulo
$\ell$, so the proved criterion applies. Degrees
$5,11,13,17,19$ give examples beyond the cubic
case. Degree $7$ has order three and is outside
this criterion; no inference for arbitrary
characteristic-two curves follows.
\end{proof}


% Configuration audit contribution to Part B. AI-produced by GPT-6 Astra.
% Provisional statement identifiers B.CF*; parent editor renumbers.
\subsection{Ordinary pro-reconstruction in finite transcendence degree}

We use an ordinary category here, not a category in which coherence is an
additional hypothesis. Let \(\mathsf T\) be the category whose objects are
Grothendieck topoi and whose morphisms are isomorphism classes of geometric
morphisms. Its ordinary pro-category has
\[
 \operatorname{Hom}_{\operatorname{Pro}(\mathsf T)}
       (\{\mathcal X_i\},\{\mathcal Y_j\})
   =\varprojlim_j\varinjlim_i
       \operatorname{Hom}_{\mathsf T}(\mathcal X_i,\mathcal Y_j).
                                                        \tag{B.CF.1}
\]
Fix a field \(k\) of characteristic zero and models of finite-type
\(k\)-schemes \(X,Y\) over a finitely generated subfield \(F_0\).
For finitely generated \(F\subset k\) containing \(F_0\), put
\[
 \mathcal X=\{X_{F,\mathrm{et}}\}_F,
 \qquad \mathcal Y=\{Y_{F,\mathrm{et}}\}_F,
 \qquad \mathcal B=\{\operatorname{Spec}(F)_{\mathrm{et}}\}_F.
\]
These are pro-objects of \(\mathsf T\), with structural arrows to
\(\mathcal B\). A morphism in the slice
\(\operatorname{Pro}(\mathsf T)_{/\mathcal B}\) induces a map
\(|X|\to |Y|\): apply the underlying-space functor and take inverse
limits. The identification with the scheme spaces follows from
\citep[Lemma~32.4.2, Tag~0CUF]{StacksProject}; its hypotheses are a directed
inverse system of schemes with affine transition maps, as here.
We call this pro-morphism \emph{pinned} when that induced map sends
closed points to closed points. This definition does not impose any
same-field representatives or coherence data.

\begin{claim}{B.43}{Proved}{Ordinary pro-reconstruction for every field of finite transcendence degree}
Suppose \(\operatorname{trdeg}_{\Q}k<\infty\). For finite-type
\(k\)-schemes, the functor to the preceding pro-systems is fully faithful
after universal homeomorphisms are inverted, with pinned morphisms in
\(\operatorname{Pro}(\mathsf T)_{/\mathcal B}\) as target morphisms.
Equivalently, for absolutely weakly normal finite-type \(X,Y\), every
such pinned ordinary pro-morphism is induced by a unique \(k\)-morphism
\(X\to Y\). This includes all finite-transcendence-degree fields of B8,
regardless of whether reconstruction from the single topos over \(k\)
holds.
\end{claim}
\begin{proof}
Enlarge \(F_0\) to contain a transcendence basis of \(k/\Q\).
Then \(k/F_0\) is algebraic. Every finitely generated subfield of \(k\)
is contained in one containing this \(F_0\), and every such field is
finite over \(F_0\). Thus we have made a cofinal change of the indexing
system, and all its field transitions are finite separable. In
particular \(X_E\) and \(Y_F\), for \(E\supset F\supset F_0\),
are finite-type schemes over the same finitely generated field \(F\).

We first verify that forgetting structural 2-cells causes no ambiguity
at these levels. For every finitely generated characteristic-zero
field \(F\), its absolute Galois group \(G_F\) is slim, meaning that
all its open subgroups have trivial centre. Indeed, if
\(\operatorname{trdeg}_{\Q}F>0\), apply
\citep[Corollary~4.5(1)]{HoshiTsujimura2023}; for a number field apply
Corollary~4.5(3) using any of its discrete mixed-characteristic valuation
rings. The same source, Definition~4.3, states the equivalent condition
that the centralizer in \(G_F\) of every normal open subgroup is trivial.
It follows also for every open subgroup, by taking its normal core.
For each connected component \(Z\) of a nonempty finite-type
\(F\)-scheme, the image of \(\pi_1(Z)\to G_F\) is open: it is the
stabilizer of a geometric connected component, of which there are
finitely many. An automorphism of the structural geometric morphism
\(Z_{\mathrm{et}}\to BG_F\) is an element of the centralizer of this
image. This follows by evaluating its action on every finite
\(G_F\)-set at a geometric base point; naturality and transport require
and suffice for commutation with the fundamental-group image.
Consequently that automorphism group is trivial. The same conclusion
holds component by component. Therefore, whenever the structural
triangle commutes as isomorphism classes, its commutation 2-cell is
unique. Moreover, any absolute isomorphism between two geometric
morphisms over \(BG_F\) automatically respects their structural 2-cells.

Let \(\alpha:\mathcal X\to\mathcal Y\) be a pinned pro-morphism over
\(\mathcal B\). At a target level \(F\), formula~(B.CF.1) represents
it by a geometric morphism
\[
             a:X_{E,\mathrm{et}}\longrightarrow Y_{F,\mathrm{et}}
                                                        \tag{B.CF.2}
\]
for some \(E\supset F\). The equality of structural pro-morphisms,
applied at the target level \(BG_F\), means that after enlarging \(E\)
once more, (B.CF.2) is a morphism over \(BG_F\). This is precisely the
equality criterion in the filtered colimit in (B.CF.1); it uses no
uniform choice over all target levels.

The representative (B.CF.2) is pinned. To prove this, let \(x\) be a
closed point of \(X_E\). The nonzero finite \(k\)-algebra
\(\kappa(x)\otimes_E k\) has a closed point, so \(x\) lifts to a
closed point \(\widetilde x\) of \(X\). The point
\(\widetilde y=|\alpha|(\widetilde x)\) is closed by assumption.
Its image in \(Y_F\) is closed: its residue field is generated by
finitely many coordinates over \(F\), is contained in a finite
extension of \(k\), and is algebraic over \(F\), since \(k/F\) is
algebraic. It is therefore finite over \(F\). Compatibility of the
underlying-space pro-morphism identifies that image with \(|a|(x)\).
This proves pinnedness at the actual finite stage.

We now apply the opened theorem
\citep[Theorem~0.1 and the equivalent formulation~(0.2)]{CarlsonHaineWolf2024}:
for a field finitely generated over its prime field, with the additional
hypothesis of positive transcendence degree only in positive
characteristic, the functor from finite-type schemes localized at
universal homeomorphisms to pinned relative geometric morphisms is
fully faithful. The field here is \(F\), of characteristic zero;
both schemes in (B.CF.2) are finite type over it. It follows that
(B.CF.2) comes from a unique localized \(F\)-scheme morphism
\(g_F:X_E\to Y_F\). Pairing this morphism with the structural map
to \(\operatorname{Spec}E\) gives a localized \(E\)-scheme morphism
\(X_E\to Y_E\), and then a localized \(k\)-scheme morphism
\(f_F:X\to Y\).

Choose this construction first at \(F_0\), and write \(f=f_{F_0}\).
For another target field \(F\), compatibility of the pro-morphism
means that \(g_F\) followed by \(Y_F\to Y_{F_0}\), and \(g_{F_0}\),
induce isomorphic geometric morphisms after a common finite source
extension \(H\subset k\). They are morphisms of finite-type
\(F_0\)-schemes. The uniqueness of structural 2-cells just proved
makes this a relative isomorphism; the faithfulness in Theorem~0.1
therefore makes the two localized scheme morphisms equal.
Their lifts to \(Y_F=Y_{F_0}\times_{F_0}F\) agree as well, since they
also have the same structural morphism to \(\operatorname{Spec}F\).
It follows that \(f_F=f\), and the pro-morphism induced by \(f\)
is exactly \(\alpha\), at every target level.

Conversely a localized \(k\)-scheme morphism descends to some finite
extension of \(F_0\): a representative finite diagram, its equations,
and its inverse universal homeomorphisms involve only finitely many
coefficients. Being a universal homeomorphism descends under the
faithfully flat field extension, and any necessary enlargement is
still finite over \(F_0\). Thus it defines the required ordinary
pro-morphism, and its map on limits is pinned. If two such morphisms
induce the same pro-morphism, compare them at target \(F_0\), refine
once, and apply the same finite-stage faithfulness and base change.
They are equal. This proves full faithfulness in the localized
category. The final formulation uses the equivalence of that
localization with absolutely weakly normal schemes, as in
\citep[Theorem~1.13, Lemma~1.32, Proposition~2.28]{CarlsonHaineWolf2024}.
\end{proof}

\begin{claim}{B.44}{Proved}{The remaining ordinary pro-category issue is an infinite-transcendence issue}
The mechanism in B.43 does not extend by the same argument to an
arbitrary characteristic-zero field: transitions corresponding to
transcendental field extensions need not preserve closed points.
Specifically, for \(F\) of characteristic zero and \(u\) transcendental,
the point \(T=u\) of \(\mathbb A^1_{F(u)}\) is closed, whereas its
image in \(\mathbb A^1_F\) is generic. Consequently the full model
system is not in general an object of the ordinary pro-category of
pinned geometric morphisms. This is an obstruction to that categorical
formulation, and is not a counterexample to reconstruction for pinned
morphisms in \(\operatorname{Pro}(\mathsf T)_{/\mathcal B}\).
\end{claim}
\begin{proof}
The contraction of the ideal \((T-u)\) from \(F(u)[T]\) to \(F[T]\)
is zero, because \(u\) is transcendental. This proves the assertion
on points. The scheme projection induces the transition geometric
morphism of the model system and has that same underlying-space map.
An inverse system in a category must have all its transition arrows
in that category. Finally, no assertion in this example constructs
a pinned ordinary pro-morphism failing to come from a scheme map;
thus it supplies no such counterexample.
\end{proof}

\begin{claim}{B.45}{Proved}{The infinite-transcendence obstacle and its resolution}
For \(\operatorname{trdeg}_{\Q}k=\infty\), ordinary relative
pro-reconstruction is proved by B.62--5 below, culminating in
B.66. The shorter argument of B.43 alone does not prove it:
a representative may have the form
\(X_{E,\mathrm{et}}\to Y_{F,\mathrm{et}}\) with \(E/F\)
transcendental. Pinnedness of the map on limits does not, by the
argument of B.43, give pinnedness of a finite-stage representative.
The coherent same-field result AS-FF.13 supplies neither that
finite-stage assertion nor an equivalence with this ordinary
pro-category. The new proof resolves the obstacle by deriving
coherence from compactness, comparing whole Kummer systems after
one finitely generated refinement, and reconstructing every étale
slice over a fixed source field. It does not assert that those
finite-stage transition morphisms preserve closed points.
\end{claim}

\subsection{A henselian field with no centralizing coefficient automorphisms}

\begin{claim}{B.46}{Proved}{A formal arc with no nontrivial algebraic self-correspondence}
There exists \(u\in t\Q[[t]]\), transcendental over \(\Q(t)\), with
the following property. Let \(\Omega\) be any complete rank-one valued extension
of \(\Q((t))\), with the valuation trivial on \(\Q\), and let
\(h\in\Omega\) satisfy \(v(h)=v(t)=1\). If both \(h\) and
\(u(h)\) are algebraic over \(\Q(t,u(t))\), and \(h,u(h)\) are
algebraically independent over \(\Q\), then
\[
                         h=t.                            \tag{B.CF.3}
\]
The series \(u\) can be specified by a recursive sequence of rational
finite jets; the proof gives a finite elimination certificate at every
stage. The property holds simultaneously in all finite ramified
extensions and their completions.
\end{claim}
\begin{proof}
Let \(Z\subset\mathbb A^2_{\Q}\times\mathbb A^2_{\Q}\) be an
irreducible closed subvariety of dimension two, dominant under both
projections, and different from the diagonal. We first prove that
in every cylinder of prescribed rational coefficients in \(t\Q[[t]]\)
there is a smaller cylinder none of whose series has a pair
\[
                 (t,u(t),h,u(h))\in Z(\Omega),\qquad v(h)=1.
                                                        \tag{B.CF.4}
\]
Write the prescribed jet as \(J(t)\) modulo \(t^N\), \(N\geq1\).
Consider the two-parameter family
\[
                       P(t)=J(t)+a t^N+b t^{N+1}.
\]
A generic member has finite intersection
\[
                         Z\cap(\Gamma_P\times\Gamma_P), 
                                                        \tag{B.CF.5}
\]
where \(\Gamma_P=\{(x,P(x))\}\). Here is a dimension proof, including
the loci where evaluation parameters coincide. On
\(xx'(x-x')\ne0\), evaluation at \(x,x'\) is a surjective linear
map from the two parameters \((a,b)\) to two values; the determinant
is \(x^N(x')^N(x'-x)\). The incidence variety of (B.CF.5) consequently
has dimension at most two over this locus. When \(x=x'\ne0\), the
point must lie in \(Z\cap\mathrm{Diag}\), of dimension at most one;
one independent evaluation condition again leaves total dimension at
most two. When \(x=0\), \(x'\ne0\), the first point is the fixed
point \((0,J(0))\). A fibre of the first dominant projection of the
irreducible surface \(Z\) has dimension at most one, and evaluation
at \(x'\) imposes one independent parameter condition. The same
argument applies with the factors exchanged. When both \(x=x'=0\),
only the single pair of fixed points is possible. Thus the full
incidence variety has dimension at most the two-dimensional parameter
space. Generic fibres are zero-dimensional or empty. We may choose
\(a,b\in\Q\), because rational points are Zariski dense in that
parameter space, including after passage to an algebraic closure.
This proves the assertion for the entire affine intersection; no
localization or denominator is needed.

Choose polynomial generators \(Q_j(x,y,x',y')\) of the ideal of
\(Z\). Since (B.CF.5) is zero-dimensional, its coordinate algebra is
finite-dimensional over \(\Q\). Elimination therefore supplies a
nonzero polynomial \(R(x)\in\Q[x]\) and polynomials
\(A_j(x,x')\in\Q[x,x']\) with
\[
 R(x)=\sum_j A_j(x,x')Q_j\bigl(x,P(x),x',P(x')\bigr).
                                                        \tag{B.CF.6}
\]
If \(u\equiv P\pmod{t^M}\), \(v(h)=1\), and (B.CF.4) holds,
then each summand on the right after \((x,x')=(t,h)\) has valuation
at least \(M\). Indeed \(u(t)-P(t)\) and \(u(h)-P(h)\) have
valuation at least \(M\), all rational coefficients have valuation
zero, and all remaining arguments have nonnegative valuation.
Taking \(M>\operatorname{ord}_{t=0}R(t)\), as well as
\(M>\deg P\), contradicts (B.CF.6). The same finite coefficient
cylinder is excluded for every valued extension \(\Omega\), not
just for one choice of Puiseux series field.

There are countably many such rational varieties \(Z\). Each
forbidden set is nowhere dense by the cylinder argument. The series
algebraic over \(\Q(t)\) also form a countable union of nowhere dense
sets: for every nonzero \(A(t,U)\in\Q[t,U]\), a polynomial \(P\)
with any prescribed jet can be chosen so that \(A(t,P(t))\ne0\),
and a sufficiently long further jet then preserves its nonzero
leading coefficient. The complete metric space \(t\Q[[t]]\), with
the coefficient topology, is a Baire space. Choose \(u\) outside
all these forbidden sets.

For the claimed pair \(h,u(h)\), take the rational Zariski closure of
\((t,u(t),h,u(h))\). Algebraicity and the two algebraic-independence
conditions make this an irreducible surface dominant on both factors.
If it were different from the diagonal it would be among the excluded
varieties. It is therefore the diagonal, which gives (B.CF.3).

For a recursive specification, enumerate rational polynomial ideals
which define the indicated irreducible dominant surfaces, and
interleave all nonzero polynomials \(A(t,U)\). At each step choose
the first rational pair \((a,b)\) in a fixed enumeration for which
(B.CF.5) is zero-dimensional, compute an elimination polynomial in
(B.CF.6), and extend the fixed jet beyond its order. Increase the
length strictly at every step. Polynomial ideal dimension,
irreducibility and elimination over \(\Q\) are finite algorithms;
the preceding proof guarantees that every search terminates. The
nested jets define a single series and prove its property. No finite
initial segment alone is asserted to establish all the infinitely
many exclusions.
\end{proof}

\begin{claim}{B.47}{Proved}{Trivial centralizers at every finite extension of a transcendence-degree-two henselian field}
Choose a series \(u\) as in B.46, and let
\[
             K=\Q(t,u(t))^h\subset\Q((t))               \tag{B.CF.7}
\]
be the henselization for the induced \(t\)-adic valuation. Then
\(K\) is countable, has transcendence degree two over \(\Q\), value
group \(\Z\), residue field \(\Q\), and completion \(\Q((t))\).
For every finite extension \(E/K\), in a fixed algebraic closure
\(\overline K\), one has
\[
              C_{\operatorname{Aut}(\overline K)}(G_E)=1.
                                                        \tag{B.CF.8}
\]
Thus the centralizer hypothesis of B7 is attained among the
residue-characteristic-zero fields with positive transcendence defect
specified by B8.
\end{claim}
\begin{proof}
The series is transcendental, so the initial field has transcendence
degree two. It lies in \(\Q((t))\), contains \(\Q(t)\), and has
residue field \(\Q\) and value group \(\Z\). Henselization is
algebraic and immediate. These properties, including its universal
property as a valued-field extension, are given in the opened
\citep[Proposition~3.3.25 and Corollary~3.3.26]{ADH2017}.
The polynomial truncations show that \(\Q(t)\) is dense in
\(\Q((t))\); hence the completion of \(K\) is that same field.
Countability follows from algebraicity over a countable field.

We record explicitly the group calculation used below. Inside a fixed
valued algebraic closure, put
\[
                  K^\dagger=\bigcup_{n\geq1}
                         K\overline{\Q}(t^{1/n}),
                                                        \tag{B.CF.9}
\]
with compatible roots. It is henselian, has algebraically closed
residue field \(\overline\Q\), and has divisible value group
\(\Q\). By the opened
\citep[Corollary~3.3.22]{ADH2017}, a valued field of equicharacteristic
zero is algebraically closed exactly when it is henselian with
algebraically closed residue field and divisible value group.
Thus (B.CF.9) is \(\overline K\). The constants \(\Q\) are
algebraically closed in \(K\), since \(K\subset\Q((t))\).
Adjoining a number field changes only the residue field, and adjoining
\(t^{1/n}\) then has degree \(n\), by its distinct valuation classes.
These facts show that restriction to constants and action on the
compatible roots give
\[
        G_K\simeq\widehat\Z(1)\rtimes G_\Q,             \tag{B.CF.10}
\]
where a splitting acts on the constants and fixes the chosen roots
of \(t\). The action on inertia is cyclotomic.

The centralizer in \(G_K\) of every open subgroup \(H\) is trivial.
An element of that centralizer projects to the centralizer of an
open subgroup of \(G_\Q\), which is trivial by
\citep[Corollary~4.5(3) and Definition~4.3]{HoshiTsujimura2023}.
It consequently lies in \(\widehat\Z(1)\). In each \(\ell\)-adic
coordinate, an element fixed by an open subgroup of \(G_\Q\) is
zero: the \(\ell\)-adic cyclotomic character is nontrivial on the
absolute Galois group of every number field, because a number field
contains only finitely many roots of unity, and multiplication by a
nonzero element of \(\Z_\ell\) is injective on \(\Z_\ell\).
This proves the asserted centralizer calculation without using the
conclusion (B.CF.8).

Now let \(E/K\) be finite and let
\(\sigma\in\operatorname{Aut}(\overline K)\) centralize \(G_E\).
Its fixed-field action preserves \(E\). The unique prolongation of
our valuation to \(E\) is henselian and discrete, with residue field
a number field, so \(E\) is not separably closed. The opened
\citep[Lemma~2.1]{KhassaKhanduja2010} states that every automorphism
of a non-separably-closed henselian valued field is continuous for
its valuation topology; its proof applies the theorem of F.\ K.
Schmidt to the original and transported valuations. For a discrete
rank-one valuation, equivalent valuation topologies give the same
valuation ring. An automorphism of its ordered value group is the
identity. Hence \(v(\sigma a)=v(a)\) on \(E\), and uniqueness of
prolongation over a henselian field gives the same equality on
\(\overline K\).

The field \(\overline\Q\subset\overline K\) is invariant under
\(\sigma\). The restriction of \(G_E\) to it is the open group
\(G_{E\cap\overline\Q}\); this follows from the Galois
correspondence, and the intersection is a number field since
\(E/K\) is finite and \(K\cap\overline\Q=\Q\). Therefore
\(\sigma|_{\overline\Q}\in G_\Q\) centralizes an open subgroup
of \(G_\Q\), and is the identity. In particular all coefficients
of \(u\) are fixed.

Embed \(\overline K\) as a valued field in an algebraic closure of
\(\Q((t))\), and work in its completion. Continuity and convergence
of the rational polynomial truncations give
\[
          \sigma u(t)=u(\sigma t),\qquad v(\sigma t)=1.
                                                        \tag{B.CF.11}
\]
Both \(\sigma t\) and \(\sigma u\) are algebraic over
\(\Q(t,u)\), because \(\overline K\) is algebraic over that
field; they are algebraically independent over \(\Q\), because
\(\sigma\) is a field automorphism. B.46 applied to
\(h=\sigma t\) yields \(\sigma t=t\), and (B.CF.11) gives
\(\sigma u=u\). The uniqueness of a valued embedding of a
henselization now makes \(\sigma\) the identity on \(K\).
It lies in \(G_K\) and centralizes its open subgroup \(G_E\).
The calculation following (B.CF.10) forces \(\sigma=1\), proving
(B.CF.8) for every finite \(E\).
\end{proof}

\subsection{A residue-characteristic-zero obstruction beyond coefficient automorphisms}

\begin{claim}{B.48}{Proved}{Compatible logarithmic reduction under a supported infinitesimal scaling}
Let \(R\) be a henselian discrete valuation ring of residue
characteristic zero, with fraction field \(K\), uniformizer \(t\),
and the unique extension \(v\) of its valuation to \(\overline K\).
Put \(\lambda=1+t\), and define
\[
 F(a)=
 \begin{cases}
   \lambda a,&v(a)>0,\\
   a,&v(a)\leq0,
 \end{cases}
 \qquad F(0)=0,
 \qquad F(\infty)=\infty.                                \tag{B.CF.12}
\]
For every finite \(G_K\)-stable set \(S\subset\mathbb P^1(\overline K)\)
containing \(0,1,\infty\), the stable marked genus-zero models for
\(S\) and \(F(S)\), over a finite Galois extension splitting their
marks, have isomorphic logarithmic special fibres over the same
standard log point. These isomorphisms preserve the labels and are
compatible with forgetting marks, finite extensions of the field of
definition, and the Galois descent data. At a mark \(a\), the induced
comparison of the arithmetic cusp group preserves its inertia and
induces the identity on its quotient \(G_{K(a)}\), with
\(K(F(a))=K(a)\) as actual subfields of \(\overline K\).
\end{claim}
\begin{proof}
The map (B.CF.12) is a \(G_K\)-equivariant bijection, whose inverse
scales the same open valuation ball by \(\lambda^{-1}\).
For two distinct finite points,
\[
          \frac{F(a)-F(b)}{a-b}\in1+\mathfrak m_{\overline K}.
                                                        \tag{B.CF.13}
\]
When both points belong to the ball this ratio is \(\lambda\), and
when neither does it is one. In the remaining case, say
\(v(a)>0\geq v(b)\), the perturbation is \(ta\), which has
valuation strictly larger than \(v(a-b)=v(b)\). This proves
(B.CF.13). In particular every difference has the same valuation and
the same initial form in the associated graded valued field before
and after applying \(F\). The assertion also holds for differences
in the reciprocal coordinate near infinity: \(F\) fixes all points
in that branch.

Choose a finite Galois extension \(L/K\) containing all the marks,
and write \(\pi\) for a uniformizer of its valuation ring. Distinct
rational sections of \(\mathbb P^1_{R_L}\) can be separated by
successive blowups at their common reductions. Each blowup is at a
smooth point of the special fibre. Its exceptional component has an
affine coordinate
\[
                          z=(x-a)/\pi^m,                 \tag{B.CF.14}
\]
for a centre \(a\) in the corresponding cluster of marks, with
\(m\) the depth of that cluster. Equivalently one may use the
reciprocal coordinate on the infinity branch. The finite tree of
components, their depths, and the locations of the marked and nodal
points on each component are obtained from the valuations and
normalized initial forms of the differences between the marks.
Contracting components with fewer than three special points gives
the stable marked model. This description is also a direct
construction of it over \(R_L\).

Equation (B.CF.13) identifies the two trees, their depths, and all
component configurations. In the coordinates (B.CF.14), the
component isomorphisms have the same reduced coordinates. They are
unique on stable components because three labelled special points
fix an automorphism of \(\mathbb P^1\). Equip the models with the
fs divisorial log structures defined by the special fibre and the
marked sections. A node on the stable model has a toric chart
\(uv=\epsilon\pi^d\), with \(d\geq1\) and \(\epsilon\) a
unit. The logarithmic special fibre retains \(d\) and the residue
of the smoothing unit, with its two tangent directions. In cluster
coordinates before contraction one may write
\(uv=\pi^{n-m}\); after changing normalized coordinates the
smoothing unit is a product and quotient of initial forms of the
same differences. Thus (B.CF.13) preserves the smoothing data and
the component maps lift to log isomorphisms over the fixed log point.
These toric charts show at the same time that the total marked model
is log regular; contracting a chain merely replaces thickness one
by a positive integral thickness.

We give the compatibility calculation, since uniqueness of the
isomorphisms by itself would not suffice for contractions. If
\(z=(x-a)/\pi^m\) is a coordinate for a target component and
\(w=(x-b)/\pi^n\) is one for a source component, the transition is
\[
       z=\frac{b-a}{\pi^m}+\pi^{n-m}w.                   \tag{B.CF.15}
\]
Extract the smaller valuation of the two coefficients. Its first
nonzero normalized reduction is a nonzero constant, \(w\), or
\(c+w\); it depends only on the initial form of \(b-a\) and the
two depths. The transformed expression has exactly the same
exponent and the same reduced factor, by (B.CF.13). This checks
both the characteristic monoid and the unit part of the pullback
of each logarithmic generator, including a generator which
vanishes identically on a contracted tail. All node, mark, and
infinity charts are obtained by products, quotients, and
reciprocals of these linear expressions. Hence the logarithmic
isomorphisms commute with every mark-forgetting contraction.
Equality of the unit factors may be checked on the irreducible
components, because the special fibre is reduced.

For completeness, the lift of a fixed component isomorphism is
unique once the characteristic monoids and log base are fixed.
At a node its possible change multiplies the two generators by
units \(u_1,u_2\) with \(u_1u_2=1\), \(u_1x=x\), and
\(u_2y=y\) in the reduced nodal ring. Thus
\(u_1-1\in(x)\cap(y)=0\). On the smooth locus the vertical
generator comes from the fixed log base, and a horizontal generator
maps to a non-zero-divisor. There is no additional freedom there.
The same coordinate check as in (B.CF.15) works under finite base
extension: if \(\pi=\epsilon'(\pi')^e\), the same unit
\(\epsilon'\) occurs on both sides, and initial forms commute
with the map between the associated graded valued fields.
This proves base-extension compatibility. Since (B.CF.12) is
\(G_K\)-equivariant, the uniquely labelled component maps and their
log lifts commute with \(\operatorname{Gal}(L/K)\).

At a marked section (a), the local horizontal parameter is
\(x-a\). Its image is \(x-F(a)\); inside the ball their
comparison is the scaling \(\lambda\), and outside it the
identity. The vertical log structure and the field \(K(a)\)
are unchanged, since multiplication by the nonzero element
\(\lambda\in K\) preserves \(K(a)\). The local Kummer chart
therefore identifies the horizontal inertia \(\widehat\Z(1)\)
and its arithmetic quotient \(G_{K(a)}\), with the identity on
the latter. This is the usual cusp description obtained from the
punctured strict henselian normal trait, or directly from the
Kummer chart in SGA~1, XIII, Proposition~5.2. Choices of geometric
fibre functors change the resulting group maps by the corresponding
cuspidal conjugacy, not by an unspecified automorphism of the
residue field.
\end{proof}

\begin{claim}{B.49}{Refuted}{Reconstruction over an equicharacteristic-zero henselian discretely valued field}
For every field \(K\) admitting a nontrivial henselian discrete
valuation with residue characteristic zero, closed-point-preserving
reconstruction of finite-type \(K\)-schemes from their relative
small étale topoi is false, even with universal homeomorphisms
inverted. The supported scaling (B.CF.12) is the geometric-point
function of a pinned autoequivalence
\[
          (\mathbb A^1_K)_{\mathrm{et}}
             \longrightarrow(\mathbb A^1_K)_{\mathrm{et}}
                  \quad\hbox{over }BG_K                 \tag{B.CF.16}
\]
which is not induced by a morphism of \(K\)-schemes.
\end{claim}
\begin{proof}
We first state precisely the opened comparison theorem.
\citep[Theorem~1.6]{Lepage2013} assumes that \(\mathscr X\) is a
log-regular fs log scheme and every prime in \(L\) is invertible
on it. Restriction identifies its finite Kummer-étale covers of
pro-\(L\) type with the corresponding finite étale covers of its
log-trivial locus. Theorem~1.8 in the same paper assumes a
noetherian henselian base \(B\) with closed point \(s\), a
connected fs log scheme \(\mathscr X\), and properness of its
underlying scheme over \(B\); restriction identifies finite
Kummer-étale covers of \(\mathscr X\) with those of its log
special fibre. These are equivalences of covering categories,
functorial in the log schemes. For our models all primes are
invertible: their valuation rings contain \(\Q\). Taking
\(L\) to be the set of all primes gives, for each marked model,
\[
 \operatorname{FEt}(\mathbb P^1_L\setminus S)
       \simeq\operatorname{KCov}(\mathscr X_S)
       \simeq\operatorname{KCov}((\mathscr X_S)_s).
                                                        \tag{B.CF.17}
\]
Here \(L/K\) now denotes a finite splitting field for the marks;
the prime-set letter in the quoted theorem has served its purpose.
The model is proper over a henselian DVR and log regular by the
charts in B.48. Thus every hypothesis has been checked. This
application retains all finite covers, with no prime omitted.

Apply (B.CF.17) to the logarithmic special-fibre isomorphisms of
B.48. We obtain compatible inverse-image equivalences
\[
 \operatorname{FEt}(\mathbb P^1_L\setminus F(S))
           \longrightarrow
 \operatorname{FEt}(\mathbb P^1_L\setminus S).            \tag{B.CF.18}
\]
Their compatibility with mark-forgetting, finite constant
extensions, and Galois descent was checked in (B.CF.15), so
finite étale descent gives the same system over \(K\).
These equivalences are over \(\operatorname{FEt}(K)\), rather
than merely geometric: apply the two quoted theorems also to
\(\operatorname{Spec}R\) with its vertical log structure.
They identify covers of the fixed log point with finite étale
\(K\)-algebras, and every special-fibre isomorphism used above
lies over that same log point. At each puncture (B.CF.18)
identifies the full decomposition and inertia pair, and acts as
the identity on the actual quotient \(G_{K(a)}\), as proved
at the end of B.48.

We now pass from this compatible system of finite covering
categories to the entire small étale site. This step is necessary;
an isomorphism at one punctured fundamental group alone would
not suffice. A quasi-compact separated étale object
\(V\to\mathbb P^1_K\) is a disjoint union of normal curves.
Zariski's Main Theorem identifies it with an open subscheme of
the finite normal compactification \(\overline V\to\mathbb P^1_K\)
in its generic finite separable algebras. Its complement is a
finite set of points. Choose a finite set \(S\) containing the
branch locus, the images of every point of
\(\overline V\setminus V\), and \(0,1,\infty\).
Over its complement the compactification is finite étale.
Use (B.CF.18) to transport this
finite étale cover, and normalize \(\mathbb P^1_K\) in the
resulting generic algebra. The full cusp pairs give the
bijections between the points over the omitted base points:
ramification is measured by inertia, and the unramified residue
extension by the quotient decomposition group over the fixed
\(G_{K(a)}\). Therefore they transport the removed points of
\(\overline V\setminus V\), preserve the unramified locus, and
produce a quasi-compact separated étale object again.
The construction is independent, up to its compatible natural
isomorphisms, of enlarging \(S\), by the restriction
compatibility in (B.CF.18).

A morphism between such étale objects is determined by its map
of generic finite algebras, extended uniquely to their normal
compactifications, together with the condition that its image
lies in the chosen open subsets. The covering-category
comparisons preserve these generic maps and the conditions on
points. Thus they give an equivalence of these categories of
étale objects. They preserve covering families: joint
surjectivity at a geometric generic point is detected by the
finite covering category; at a geometric closed point it is
detected by the same local decomposition/inertia comparison.
These objects form a generating small étale site. We have
therefore constructed an autoequivalence of
\((\mathbb P^1_K)_{\mathrm{et}}\), over \(BG_K\), whose
function on \(\overline K\)-valued points is exactly (B.CF.12). It is pinned,
since that function is a \(G_K\)-equivariant bijection and
preserves the finite field \(K(a)\) at every algebraic point.
It fixes infinity and hence restricts to (B.CF.16).

Every nonzero rational number is a valuation unit and is fixed
by (B.CF.12), whereas \(t\) is sent to \((1+t)t\ne t\).
If a \(K\)-scheme morphism \(\mathbb A^1_K\to\mathbb A^1_K\)
induced (B.CF.16), its polynomial would agree with \(T\) on the
infinite subset \(\Q^\times\), and would consequently be \(T\)
identically. Its value at \(t\) gives the contradiction.
The scheme \(\mathbb A^1_K\) is already absolutely weakly
normal, so inverting universal homeomorphisms introduces no
additional morphism which could realize this point function.
\end{proof}

\begin{claim}{B.50}{Refuted}{The centralizer condition in B7 is insufficient}
There exists a characteristic-zero field \(K\), countable and of
transcendence degree two over \(\Q\), such that
\[
 C_{\operatorname{Aut}(\overline K)}(G_E)=1
           \quad\hbox{for every finite }E/K,
\]
but closed-point-preserving reconstruction of finite-type
\(K\)-schemes from their relative étale topoi fails. One may take
exactly the recursively specified field (B.CF.7). The obstruction
is a compatible nonalgebraic permutation of points and their
whole covering categories, not a centralizing field automorphism.
\end{claim}
\begin{proof}
B.47 proves the displayed condition for this field, with all
finite extensions included. Its valuation is henselian and
discrete, with residue field \(\Q\), so B.49 supplies the
pinned nonalgebraic autoequivalence of its affine-line topos.
These are the same field and the same entire arithmetic base
\(BG_K\). This is a counterexample to the asserted implication.
We also verify the surviving necessary direction, rather than
accepting AS-FF.6 on authority. If a field satisfies reconstruction,
then every finite extension does: finite étale slicing identifies
objects and relative morphisms over that extension, and faithfulness
forces the reconstructed map to preserve the extension's structural
map. Let \(E/K\) be finite and let \(\sigma\) centralize
\(G_E\). It preserves \(E\). The coefficient automorphism on
\(\mathbb G_{m,E}\), together with the identification of its
base Galois action supplied by \(\sigma G_E\sigma^{-1}=G_E\)
pointwise, gives a pinned relative topos automorphism whose geometric
point function is \(a\mapsto\sigma(a)\). This includes the case
where \(\sigma\) fixes \(E\): then it is the identity absolute
morphism with its relative structural 2-cell twisted by the central
element. The point-function description follows directly by pulling
back to finite separable extension points, as in the opened
\citep[Proposition~3.6]{CarlsonHaineWolf2024}.
Reconstruction over \(E\) would make this an \(E\)-automorphism
of \(\mathbb G_m\), hence \(T\mapsto cT^{\pm1}\).
Evaluation at one gives \(c=1\); evaluation at two excludes
inversion. Consequently \(\sigma(a)=a\) for every
\(a\in\overline K^\times\), so \(\sigma=1\).
This proves necessity for all finite extensions. B.50 shows that
excluding this particular obstruction is not sufficient.
A precise valid sufficient-and-necessary replacement is the
point-function condition of the opened
\citep[Theorem~4.18(5)]{CarlsonHaineWolf2024}: for every affine
regular connected finite-type source with a rational point,
every pinned relative map to \(\mathbb G_m\) must have the
geometric-point function of a regular unit. The supported
scaling demonstrates why a condition solely on coefficient
field automorphisms cannot replace that requirement.
\end{proof}

\begin{claim}{B.51}{Refuted}{The specified transcendence-defect field in AS-FF.18 does not reconstruct}
Let
\[
     u_*(t)=\sum_{n\geq1}t^{n!},\qquad
     K_* =\Q(t,u_*(t))^h\subset\Q((t)).                  \tag{B.CF.19}
\]
This field has transcendence degree two, residue field \(\Q\),
value group \(\Z\), and positive transcendence defect one.
It does not satisfy closed-point-preserving étale reconstruction.
This settles the particular residue-characteristic-zero case
left undecided in AS-FF.18, without claiming that its centralizers
have the property proved for the different series in B.46.
\end{claim}
\begin{proof}
For a nonzero \(P(t,U)\in\Q[t,U]\), write \(e=\deg_U P\)
and \(m=\deg_t P\), and put
\(u_N=\sum_{n=1}^N t^{n!}\). If \(P(t,u_*)=0\), then a
nonzero \(P(t,u_N)\) has order at least \((N+1)!\) but degree
at most \(m+eN!\), impossible for sufficiently large \(N\).
Thus infinitely many distinct \(u_N\) would be roots of the
nonzero polynomial \(P(t,U)\) over \(\Q(t)\), again impossible.
The series is transcendental. The field and its henselization
therefore have the asserted transcendence degree, and
henselization preserves residue field and value group.
The defect is \(2-0-1=1\). B.49 applies directly to this
henselian discrete valuation and refutes reconstruction.
\end{proof}

\begin{claim}{B.52}{Open}{Exact remaining scope of the henselian-field question B8}
The universal quantifier in B8 is not fully settled. B.49--9,
together with B.54--6, refute reconstruction for every field of
finite transcendence degree admitting a nontrivial henselian
valuation of residue characteristic zero. This includes all the
residue-characteristic-zero transcendence-defect cases in AS-FF.18,
the specified field \(K_*\), and fields with a mixed-characteristic
valuation admitting such a nontrivial residue-characteristic-zero
coarsening. The following mixed-characteristic cases remain from
the original \texttt{agent\_F/F\_section.tex}.
\begin{enumerate}
\item General finite-transcendence-degree henselian fields of mixed
residue characteristic outside the positive and negative classes
already separated in AS-FF.16, AS-FF.22, and AS-FF.26--39.
In particular, an unbounded wild ramification index is not by itself
a negative criterion: the cyclotomic, maximal abelian, and each fixed
bounded-nilpotency class in AS-FF.35, 37, 39 were already positive
claims in the return.
\item For every prime \(p\), the specified maximal metabelian field
\[
 K_p=\operatorname{Frac}(\Z_{(p)}^{sh}),\qquad
 I=G_{K_p},\qquad
 B_p=\overline\Q^{\,\overline{[I',I']}},\quad
 I'=\overline{[I,I]},
\]
and general finite extensions of it. AS-FF.40 singles out this
field; AS-FF.18 names \(B_2\) as the remaining degree-zero test.
Neither its asserted centralizer condition nor the affine and
selected-cover rigidity assertions AS-FF.44--50 prove reconstruction
for arbitrary pinned morphisms or arbitrary covers. The negative
residue-characteristic-zero argument here has no application to its
mixed-characteristic valuation.
\item AS-FF.39 explicitly makes no assertion about the union of
the fixed nilpotency-class fields
\(\bigcup_c\overline\Q^{\,\gamma_{c+1}(I)}\), or about a general
solvable closure. The former union remains undecided here. If the
latter phrase means the maximal solvable extension of this precise
strict-local field \(K_p\), B.53 below identifies it with
\(\overline\Q\) and refutes reconstruction. It does not identify
the lower-central union with the derived-series union, and makes no
assertion about solvable closures of arbitrary other base fields.
B.57 determines the stable rational Kummer kernel at the
lower-central union for \(p\geq5\), with explicit exceptions
at \(p=2,3\), but its full reconstruction problem remains open.
\end{enumerate}
\end{claim}
\begin{proof}[Scope verification]
The original pending list is stated at AS-FF.18 and its subsequent
updates: AS-FF.27 handles algebraic fields with a discrete
mixed-characteristic valuation; AS-FF.29--30 handle the specified
bounded-wild fields; AS-FF.35, 37, 39 handle the cyclotomic, maximal
abelian, and fixed bounded-nilpotency extensions. AS-FF.40--50 retain
the full reconstruction problem over \(B_p\), while proving only the
more restricted assertions indicated above. None of these scope
statements converts a failed proof mechanism into a counterexample.
Theorem B.49 uses a discrete residue-characteristic-zero valuation
for its noetherian stable logarithmic models. B.54 removes
discreteness by the verified valuative-logarithmic approximation
and covering-category continuity; B.55 removes the rank-one
restriction in finite transcendence degree by henselian coarsening.
These additions settle the entire residue-characteristic-zero
portion. They do not treat positive residue characteristic and
its wild covers. The full classification demanded by B8 therefore
remains open in the mixed-characteristic respects listed above.
\end{proof}

\begin{claim}{B.53}{Refuted}{The unbounded solvable closure of the strict-local field}
Fix an embedding of the strict henselization
\(A=\Z_{(p)}^{sh}\) in \(\overline\Q\), put
\(K_p=\operatorname{Frac}(A)\), and write \(I=G_{K_p}\).
Let \(I^{(0)}=I\) and
\(I^{(d+1)}=\overline{[I^{(d)},I^{(d)}]}\) be the closed
derived series. Then
\[
 K_p^{\mathrm{sol}}
   :=\bigcup_{d\geq 0}\overline\Q^{\,I^{(d)}}
   =\overline\Q.                                      \tag{B.CF.20}
\]
Consequently this maximal solvable closure does not satisfy
closed-point-preserving \(\acute etale\) reconstruction. The statement
concerns the actual field in AS-FF.39, not a substitution of a
completed local field. It leaves open the bounded derived-length
field \(B_p=\overline\Q^{I^{(2)}}\), and the lower-central union
\(\bigcup_c\overline\Q^{\gamma_{c+1}(I)}\).
\end{claim}
\begin{proof}
The ring \(A\) is a strictly henselian discrete valuation ring
with residue field \(\overline{\F}_p\). For a finite Galois
extension \(M/K_p\), henselianity gives a unique extension of
the valuation. Thus its decomposition group is the whole
\(\operatorname{Gal}(M/K_p)\). Its residue extension is trivial,
since it is finite over \(\overline{\F}_p\), so the whole
Galois group is also its inertia group. The opened
\citep[Lemma~15.114.5, Tag~09EE]{StacksProject} gives an exact
sequence with a finite \(p\)-group as kernel and a cyclic
prime-to-\(p\) group as quotient. Both are solvable, hence every
finite quotient of \(I\) is solvable. Equivalently one may apply
the profinite version, Lemma~58.13.5, Tag~0BUA, to this same DVR.

Every nonidentity element of a profinite group survives in some
finite quotient. If that quotient is solvable of derived length
\(d\), the image of \(I^{(d)}\) in it is trivial. It follows
that \(\bigcap_d I^{(d)}=1\). Alternatively, every element of
\(\overline\Q\) lies in a finite Galois extension of \(K_p\),
whose finite solvable Galois group has some finite derived length;
that extension lies in \(\overline\Q^{I^{(d)}}\). This proves
(B.CF.20), without asserting any uniform bound on derived length.

Choose a nonidentity automorphism \(\sigma\) of
\(\overline\Q/\Q\), for example complex conjugation after
choosing a complex embedding. Acting on coefficients induces an
autoequivalence of \(\mathbb A^1_{\overline\Q,\mathrm{et}}\)
over \(B G_{\overline\Q}=\mathsf{Sets}\); it preserves closed
points and acts on their values by \(a\mapsto\sigma(a)\).
If induced by a morphism of \(\overline\Q\)-schemes, its
polynomial would fix every rational number, and hence would be
\(T\). It moves an algebraic number, a contradiction. Inverting
universal homeomorphisms does not change this conclusion for the
normal characteristic-zero affine line. This is a genuine
negative reconstruction result for the exact unbounded derived
tower, rather than a failure of a proposed proof of reconstruction.
\end{proof}


\subsection{Nondiscrete valuations of residue characteristic zero}

The logarithmic argument can be extended beyond DVRs. The full
valuative log structure is essential here: replacing it by a fixed
finitely generated submonoid would, in general, change its Kummer
coverings and its arithmetic fundamental group.

\begin{claim}{B.54}{Proved}{All-prime logarithmic comparison over a henselian rank-one valuation ring}
Let $R$ be a henselian rank-one valuation ring with fraction field $K$
and residue field $\kappa$ of characteristic zero. Let
$S\subset\mathbb P^1(K)$ be a finite set containing $0,1,\infty$,
and let $\mathscr X/R$ be its stable marked genus-zero model.
Endow $\operatorname{Spec}R$ with the canonical valuative log
structure $R\setminus\{0\}\to R$, and $\mathscr X$ with the
log structure of its vertical boundary and its marked sections.
Write $\mathscr X_s$ for its log closed fibre and
$U=\mathbb P^1_K\setminus S$. Then restriction gives equivalences
\[
 \operatorname{FEt}(U)
 \ \longleftarrow\
 \operatorname{KCov}(\mathscr X)
 \ \longrightarrow\
 \operatorname{KCov}(\mathscr X_s).                 \tag{B.LIT.1}
\]
Here $\operatorname{KCov}$ denotes all finite Kummer \'etale
coverings in the saturated, not necessarily fs, sense of
Achinger--H\"ubner--Lara--Stix. In that definition the underlying
scheme map is integral and need not be finite. The equivalences
are over the entire $BG_K$, and are compatible with finite valued
extensions and with morphisms between stable models obtained by
forgetting marks. Neither discreteness nor completeness is assumed.
\end{claim}
\begin{proof}
We state the exact source inputs. AHLS Lemma~5.3.2 says that a
microbial valuation ring of residue characteristic zero is a
filtered colimit of local excellent regular subalgebras
$\Lambda_i$ with local maps to $R$, containing a fixed
pseudouniformizer $\pi$, for which $V(\pi)_{\mathrm{red}}$ is
a strict normal crossings divisor. If $R$ is henselian the
$\Lambda_i$ may be chosen henselian. A nontrivial rank-one
valuation is microbial, because $R[1/\pi]=K$ for every
$0\ne\pi\in\mathfrak m_R$. Put
\[
 M=R\setminus\{0\},\qquad
 M_i=\Lambda_i\cap\Lambda_i[1/\pi]^\times.
\]
The paragraph following that lemma gives
$(M\to R)=\varinjlim_i(M_i\to\Lambda_i)$ as saturated prelog
rings; the associated log schemes $B_i$ are fs and log regular.
AHLS Proposition~4.2.6 identifies finite Kummer \'etale covering
categories with their filtered colimit for an affine charted
inverse system of qcqs saturated log schemes, in the sense of its
Definition~3.4.4. These are the precise continuity and approximation
theorems used below.\citep[Definition~3.4.4,
Proposition~4.2.6, Lemma~5.3.2]{LA_AHLS2024}

For completeness, the required stable family exists over $R$,
without a noetherian assumption. Hassett Theorem~2.1, with genus
zero and all weights one, gives a proper moduli stack and its
projective coarse scheme over $\mathbb Z$. Every stable labelled
genus-zero curve has trivial automorphism group, so in this case
the stack is the fine projective moduli scheme
$\overline M_{0,|S|}$. Its universal family and the valuative
criterion give $\mathscr X$. The finite-presentation moduli
problem over arbitrary bases is obtained from the initially
noetherian formulation by descent of its finite presentation.
The model and its generic identification are unique, and its base
change to an extension of valuation rings is the corresponding
stable model.\citep[Theorem~2.1]{LA_Hassett2003}

One can describe its finite chart data directly. For a finite set
of split marks, the finite tree of branching clusters is determined
by the valuations of pairwise differences. At a branching radius,
choose a difference $d\ne0$ of that valuation and blow up the ideal
$(x-a,d)$ in the appropriate smooth coordinate. Its charts are
$x-a=dz$ and $(x-a)w=d$. Repeat at the finitely many branching
radii, omitting or contracting components with fewer than three
special points. This is the stable model just described. Its node
charts have the form $xy=a$, for $0\ne a\in\mathfrak m_R$;
its marked charts have a smooth horizontal parameter. All nodes
and their branches are split, since the marks are split and their
dual graph is a labelled tree.

Over the base monoid $M$, the log charts are $M$, $M\oplus\mathbb N$
at a mark, and
\[
       M[e_1,e_2]/(e_1+e_2=a)
                                                        \tag{B.LIT.2}
\]
at a node. AHLS Example~2.4.5 proves that these node monoids are
saturated and that the maps from $M$ are smooth, vertical and
saturated for every saturated $M$ and every $a\in M$.
The standard vertical log structure agrees with these charts by
AHLS Proposition~5.3.1; the disjoint smooth horizontal sections
add the displayed $\mathbb N$ factors. The log-trivial open is $U$.

Descend the marked stable family, its finite chart cover, and the
chart gluings to a sufficiently large $\Lambda_i$. This is ordinary
finite-presentation descent; equivalently, descend the $R$-point
of the fine moduli scheme and pull back its universal family.
After increasing $i$, all the finitely many node parameters and
chart identities descend. For each such nonzero parameter $a$,
rank one gives $N$ with $\pi^N/a\in R$. Require $a$ and
$\pi^N/a$ to belong to $\Lambda_i$. Then $a\in M_i$; on the
regular snc base it is a unit times a monomial in the local
equations of $V(\pi)_{\mathrm{red}}$. Include also the inverses
of the finitely many unit coefficient differences which separate
sections and chart overlaps. All these requirements persist on a
cofinal subsystem.

Let $\mathscr X_i/B_i$ denote the resulting marked log curve.
The proper stable underlying scheme is an ordinary base change
of one fixed descended family. Its local log charts are the
versions of (B.LIT.2) and $M_i\oplus\mathbb N$ over $M_i$.
It is thus fs and log smooth over the log-regular $B_i$, hence
log regular by the log-regularity theorem recalled immediately
after AHLS Lemma~5.3.2. All integers are invertible: the
$\Lambda_i$ can be taken as $\mathbb Q$-algebras in its
characteristic-zero construction. Let $U_i$ be the log-trivial
open; its limit is $U$.

The node and marked monoid maps are saturated. Their saturated
log base changes consequently have the same underlying schemes
as ordinary base changes. The compatible finite strict \'etale
chart covers make $\{\mathscr X_i\}$ an affine charted system
in exactly AHLS Definition~3.4.4, with limit $\mathscr X$.
The same applies to the log closed fibres
$\mathscr X_{i,s_i}$. Indeed, all the maps $\Lambda_i\to R$
and the transition maps are local inclusions. Their induced
residue-field maps are injective: the inverse image of
$\mathfrak m_R$ in each ring is its maximal ideal. Those
residue fields therefore have colimit $\kappa$. Taking the same charts modulo their maximal
ideals therefore gives the log fibre $\mathscr X_s$ as their
saturated inverse limit.

Lepage Theorem~1.6 identifies finite Kummer \'etale covers of an
fs log-regular scheme with finite \'etale covers of its log-trivial
open, for primes invertible on that scheme. His Theorem~1.8
identifies finite Kummer \'etale covers of a connected fs log
scheme proper over a noetherian henselian base with those of its
log closed fibre. Apply both theorems to every $\mathscr X_i$,
with all primes. They give
\[
 \operatorname{FEt}(U_i)\ \longleftarrow\
 \operatorname{KCov}(\mathscr X_i)\ \longrightarrow\
 \operatorname{KCov}(\mathscr X_{i,s_i}).
\]
Ordinary finite \'etale continuity on the left and AHLS
Proposition~4.2.6 on the other two terms prove (B.LIT.1).
This is a limit argument using the full valuative monoid, not an
application of the fs theorems directly to $R$.
\citep[Theorems~1.6 and 1.8]{Lepage2013}

Apply the same argument to the log base, proper over itself.
It identifies the covering categories of that base and its log
closed point with $\operatorname{FEt}(K)$. Independently, AHLS
Corollary~4.4.9 says that the Kummer \'etale fundamental group of
the canonical log valuation ring is the maximal tame quotient of
$G_K$. Every finite extension is tame in residue characteristic
zero, so this is all of $G_K$. Since all the restriction functors
commute with pullback of base covers, (B.LIT.1) is over $BG_K$.

For a finite valued extension $L/K$, AHLS Proposition~4.4.3
identifies the canonical valuative log map as a finite Kummer
\'etale map in the stated integral sense. Stable models base
change as schemes, and their saturated node charts base change
as just checked. This yields compatibility with the actual field
extension. For a finite diagram of mark-forgetting morphisms,
descend the diagram and all its charts to a common stage; the
restriction squares commute at that stage and hence at the limit.
Equivalently, a transported cover with its specified restriction
has a contractible groupoid of extensions by full faithfulness.
The resulting comparison isomorphisms therefore satisfy their
cocycle identities. This proves the asserted functoriality.
\citep[Proposition~4.4.3, Corollary~4.4.9]{LA_AHLS2024}
\end{proof}

\begin{claim}{B.55}{Refuted}{Reconstruction over all finite-transcendence-degree henselian fields of residue characteristic zero}
Every field $K$ admitting a nontrivial henselian rank-one valuation
of residue characteristic zero fails closed-point-preserving
reconstruction of finite-type $K$-schemes from their \'etale topoi
over $BG_K$. In particular this holds for every field of finite
transcendence degree over $\mathbb Q$ with a nontrivial henselian
valuation of residue characteristic zero, regardless of its rank,
discreteness, or transcendence defect.
\end{claim}
\begin{proof}
We first work in rank one. Choose $0\ne t\in\mathfrak m_K$ and
put $\lambda=1+t$. On $\mathbb P^1(\overline K)$ set
\[
 F(a)=\begin{cases}\lambda a,&v(a)>0,\\a,&v(a)\leq0,\end{cases}
 \qquad F(0)=0,\qquad F(\infty)=\infty.
                                                        \tag{B.LIT.3}
\]
Henselianity makes the extended valuation $G_K$-invariant.
Thus $F$ is an equivariant bijection, with inverse scaling the
same ball by $\lambda^{-1}$. As in \statementref{B.48},
\[
       \frac{F(a)-F(b)}{a-b}\in1+\mathfrak m_{\overline K}
       \quad(a\ne b).                                  \tag{B.LIT.4}
\]
When both points lie in the ball the ratio is $\lambda$; when
neither does it is one. In the remaining case
$v(a)>0\geq v(b)$, the change $ta$ has strictly greater
valuation than $a-b$. This proves (B.LIT.4) for an arbitrary
ordered rank-one value group. Hence every difference has the same
valuation and initial form before and after $F$.

Take a finite $G_K$-stable marked set containing $0,1,\infty$
and a finite Galois extension $L/K$ splitting it. The stable
models for $S$ and $F(S)$ have the same labelled cluster tree.
For a cluster with centre $a$ choose a scalar $d\in L^\times$
realizing its radius; use that \emph{same} scalar for the
transformed cluster with centre $F(a)$. Instead of the integral
power coordinates of the DVR proof, write $z=(x-a)/d$. If
$w=(x-b)/e$ is another component coordinate, the transition is
\[
                  z=\frac{b-a}{d}+\frac ed w.            \tag{B.LIT.5}
\]
Extract the coefficient of smaller valuation. Its first nonzero
normalized reduction is a nonzero constant, $w$, or $c+w$.
Equation (B.LIT.4) preserves this reduced factor and its value.
It therefore preserves both the ghost-monoid element and the
unit part of every logarithmic pullback. The smoothing parameter
$e/d$ can be chosen literally the same on the two sides.
Products, quotients and reciprocal coordinates give all node,
mark and infinity charts, including generators whose pullback
vanishes identically on a contracted tail. This proves
compatibility with every mark-forgetting contraction.

The labelled component maps lift uniquely once these ghost maps
and the log base are fixed. At a node a possible difference of
two lifts multiplies the branch generators by units $u_1,u_2$
with $u_1u_2=1$, $u_1x=x$, and $u_2y=y$. Thus
$u_1-1\in(x)\cap(y)=0$ in the reduced nodal ring. On the
smooth locus the vertical log structure comes from the fixed
base, while the horizontal generator is a non-zero-divisor.
Initial forms commute with finite extensions of valued fields,
so the comparisons commute with those extensions; their unique
labelled description and equivariance of $F$ give Galois descent.
At a mark $a$, the horizontal parameter changes by $\lambda$
or one. In particular $K(F(a))=K(a)$ as actual subfields,
the cusp inertia is preserved, and the quotient $G_{K(a)}$ is
identified by the identity. This is the same local Kummer-chart
calculation as in \statementref{B.48}; principal units have all
their roots by Hensel's lemma, since every integer is invertible.

Apply \statementref{B.54} to these split models and descend.
It gives compatible equivalences between the finite \'etale
cover categories of $\mathbb P^1_K\setminus S$ and
$\mathbb P^1_K\setminus F(S)$ over the entire $BG_K$.
They respect all inclusions of marked sets and the actual
arithmetic cusp residue extensions. We recall why this transports
the whole small \'etale site, as required in
\statementref{B.49}. A quasi-compact separated \'etale object
over $\mathbb A^1_K$ is an open subscheme of the normalization
in its generic finite algebra. Its data consist of that finite
algebra, its branching over finitely many closed points, and the
retained unramified points with their residue extensions.
The cover equivalences and cusp comparisons carry all these
data. Maps extend uniquely from their generic algebra maps,
with the condition on retained points preserved. Geometric
point surjectivity, which detects covering families, is preserved
as well. Such objects form a basis of the small \'etale site.
The compatibility already proved therefore yields a pinned
topos self-equivalence over $BG_K$, whose function on
$\overline K$-valued points is (B.LIT.3).

Residue characteristic zero implies $v(q)=0$ for every
$q\in\mathbb Q^\times$. Thus $F$ fixes all rational numbers
and moves $t$. No $K$-polynomial has that point function: a
polynomial fixing infinitely many rational numbers is the
identity. The constructed topos morphism is consequently not
induced by a $K$-scheme morphism. This proves failure of
reconstruction in rank one.

Finally let $\operatorname{trdeg}_{\mathbb Q}K=d<\infty$
and let the original valuation have residue characteristic zero.
It is trivial on $\mathbb Q$. If
$v(a_1),\ldots,v(a_m)$ are rationally independent, distinct
monomials in the $a_i$ have distinct valuations, so no nonzero
$\mathbb Q$-polynomial in them can vanish. Hence $m\leq d$:
the value group has finite rational rank. Any strict chain of
convex subgroups gives rationally independent elements chosen
in successive differences; the highest nonzero coefficient in
a rational relation cannot vanish modulo the preceding convex
subgroup. The ordered group therefore has finite rank and a
maximal proper convex subgroup. Its quotient is a nonzero
rank-one ordered group. Coarsen by this subgroup.
ADH Lemma~3.4.2 states that a valued field is henselian if and
only if both a coarsening and its induced residue valuation are
henselian. The coarsening is therefore henselian; its residue
characteristic is still zero because $\mathbb Q^\times$
remains in its unit group. The rank-one result applies to the
same field $K$.\citep[Lemma~3.4.2]{ADH2017}
\end{proof}

This extension removes the nondiscrete residue-characteristic-zero
cases from the unresolved portion of B8. It does not replace the
all-prime comparison by a tame quotient in mixed characteristic.
Fields admitting no nontrivial henselian residue-characteristic-zero
coarsening require a separate argument.


% Configuration audit continuation: exact mixed-characteristic frontiers.
% AI-produced by GPT-6 Astra; provisional statement identifiers B.56--13.
\subsection{Wild Kummer obstructions and the unbounded nilpotent tower}

Throughout this subsection,
\[
 K_p=\operatorname{Frac}(\Z_{(p)}^{sh})\subset\overline\Q,
 \qquad I=G_{K_p}=I_{\Q_p},\qquad
 F_p=\overline\Q^{I'},\qquad B_p=\overline\Q^{I''}.
                                                        \tag{B.MIX.1}
\]
The identification with local inertia is the henselization--completion
comparison, with the chosen prolongation of the valuation. It concerns
absolute Galois groups; the global field \(K_p\) has not been replaced
by \(\Q_p^{\mathrm{ur}}\). For clarity, a finite collection of global
algebraic numbers and its local Galois conjugates lie in a number field
\(D\) whose completion \(M\) is finite Galois over \(\Q_p\).
Then \(I\cap G_M=I_M\) fixes the compositum \(K_pD\).
These are the identifications used below.

\begin{claim}{B.56}{Proved}{Wild covers obstruct supported scaling at the metabelian field}
For a finite extension \(E/B_p\), let \(\kappa_E\) be the rational
\(p\)-primary Kummer map. Put
\[
 W_p=\{a\in B_p^\times:a^m\in F_p^\times
                         \text{ for some }m\geq1\}.
\]
Then
\[
 \ker\kappa_E=\bigcap_{n\geq1}(E^\times)^{p^n}=W_p,
 \qquad W_p\cap(1-W_p)=F_p\setminus\{0,1\}.             \tag{B.MIX.2}
\]
In particular the supported scaling used in B.48--7 cannot be the
point function of a pinned relative autoequivalence over \(B_p\).
Precisely, for \(\lambda=1+p\), the bijection
\[
 f(a)=\begin{cases}\lambda a,&v(a)>0,\\a,&v(a)\leq0,
                         \end{cases}\qquad f(0)=0
                                                        \tag{B.MIX.3}
\]
of \(\overline\Q\) is excluded by two one-puncture
\(p\)-Kummer comparisons. This is an obstruction to this proposed
counterexample, and not a proof of reconstruction over \(B_p\).
\end{claim}
\begin{proof}
We independently reprove the kernel input of AS-FF.42. Suppose
\(\kappa_E(a)=0\). Enlarge to \(E'=B_pD\), with \(D\) a number
field stable under the local Galois group, containing \(a\) and the
finitely many generators for \(E/B_p\). Its local completion
\(M/\Q_p\) is finite Galois. Write
\[
 J=I\cap G_M=I_M,\quad H=I''\cap J,
 \quad\Gamma=J/H,\quad C=(I'\cap J)/H.
\]
The Kummer class of \(a\) on \(J\) inflates from
\(H^1(\Gamma,\Q_p(1))\). The group \(C\) and the quotient
\(\Gamma/C\) are abelian; the cyclotomic character is trivial
on \(C\) and has infinite image on \(\Gamma/C\).
For an abelian group \(A\) with nontrivial scalar character \(\chi\),
the cocycle equation gives
\((\chi(g)-1)c(h)=(\chi(h)-1)c(g)\). Choosing one
\(g\) with \(\chi(g)\ne1\) proves \(H^1(A,\Q_p(\chi))=0\).
Inflation--restriction therefore injects
\(H^1(\Gamma,\Q_p(1))\) into \(H^1(C,\Q_p(1))\).
Every \(g\in I'\) centralizes \(C\), because the commutator
lies in \(I''\cap J\), and fixes its coefficient module.
Thus it fixes the class, so
\(\kappa_J(g(a)/a)=0\).

Both \(a\) and \(g(a)\) belong to \(M\). Restriction from
\(H^1(G_M,\Q_p(1))\) to \(H^1(I_M,\Q_p(1))\) is injective,
since \(\Q_p(1)^{I_M}=0\). Over the finite local field \(M\),
the rational Kummer kernel is exactly \(\mu(M)\): valuation
detects the uniformizer component, and the logarithm identifies a
sufficiently small principal-unit group with an additive lattice
in \(M\). Hence \(g(a)/a\in\mu(M)\) for every \(g\in I'\).
An integer \(m\) killing the finite group \(\mu(M)\) makes
\(a^m\) fixed by \(I'\), so \(a^m\in F_p\).
Since \(F_p\) contains all roots of unity,
\(F_p(a)/F_p\) is cyclic and belongs to its maximal abelian
extension \(B_p\). Thus \(a\in W_p\).

Conversely, if \(a^m\in F_p\), every \(p^n\)-th root of
\(a\) is a radical of an element of \(F_p\). Its extension
over \(F_p\) is cyclic, so it belongs to \(B_p\).
Compatible roots exist. Since \(G_E\) fixes \(\mu_{p^\infty}\),
integral Kummer classes are homomorphisms to the torsion-free
module \(\Z_p(1)\); rational vanishing is the same as integral
vanishing. This proves the first equality in (B.MIX.2).
If \(a,1-a\in W_p\), then for each \(g\in I'\) one has
\(g(a)=\zeta a\), \(g(1-a)=\eta(1-a)\), with
\(\zeta,\eta\in\mu_\infty\subset F_p\).
If \(g(a)\ne a\), their sum gives
\(a=(1-\eta)/(\zeta-\eta)\in F_p\), a contradiction;
the case \(\zeta=\eta\) forces both to be one.
This proves the second equality.

Choose \(x\in B_p\setminus F_p\) with \(v(x)>0\), and
\(c\in K_p\) with both \(c\) and \(1-c\) units. Such an
\(x\) exists because \(I'\) is a nontrivial closed subgroup
of the free wild pro-\(p\) group, hence has nontrivial
abelianization; multiplying an element of \(B_p\setminus F_p\)
by a sufficiently large power of \(p\) ensures positive value.
This is also the nontriviality argument in AS-FF.40. The residue
field \(\overline\F_p\) supplies \(c\), including when \(p=2\).
Put \(y=1+c(x-1)\). Then \(y\) is a unit,
\(f(1)=1\), \(f(y)=y\), and \(f(x)=\lambda x\).

Suppose the stated relative autoequivalence existed. On the line
punctured at a point \(a\in B_p\), it sends the rank-one geometric
\(p\)-Kummer line to itself by a scalar
\(\alpha_a\in\Z_p^\times\). The remaining summand is an
arithmetic Kummer class, which cancels on subtracting two point
evaluations. The puncture at \(1\), evaluated at \(y,x\), gives
\[
 \kappa_{B_p}\!\left(\frac{y-1}{\lambda x-1}\right)
     =\alpha_1\kappa_{B_p}\!\left(\frac{y-1}{x-1}\right)
     =\alpha_1\kappa_{B_p}(c)=0.
\]
The puncture at \(x\), evaluated at \(y,1\), similarly gives
\[
 \kappa_{B_p}\!\left(\frac{y-\lambda x}{1-\lambda x}\right)
     =\alpha_x\kappa_{B_p}(1-c)=0.
\]
The two left arguments are \(A\) and \(1-A\), for
\(A=c(x-1)/(\lambda x-1)\). Equation (B.MIX.2) forces
\(A\in F_p\). But this fractional linear transformation has
coefficients in \(F_p\) and nonzero determinant
\(c(\lambda-1)\), so it cannot carry
\(x\notin F_p\) into \(F_p\). This contradiction pinpoints the
wild arithmetic obstruction absent from the residue-zero proof.
\end{proof}


\begin{claim}{B.57}{Proved}{A prime-dependent Kummer theorem for the unbounded nilpotent tower}
Let
\[
 N=\bigcap_{c\geq1}\gamma_{c+1}(I),\qquad
 K_{p,\mathrm{nil}}=\overline\Q^N
      =\bigcup_{c\geq1}\overline\Q^{\gamma_{c+1}(I)}.
                                                        \tag{B.MIX.4}
\]
For every \(p\geq5\) and every finite extension
\(E/K_{p,\mathrm{nil}}\),
\[
 \ker\!\left(E^\times\longrightarrow H^1(G_E,\Q_p(1))\right)
       =\bigcap_{n\geq1}(E^\times)^{p^n}
       =\mu_\infty.                                    \tag{B.MIX.5}
\]
For \(p=2\) and \(p=3\), the corresponding kernel is already
strictly larger than \(\mu_\infty\) over
\(K_{p,\mathrm{nil}}\) itself. Explicit nontorsion elements in it are
\[
       2\quad(p=2),\qquad
       \frac{2+\sqrt{-3}}{2-\sqrt{-3}}\quad(p=3).         \tag{B.MIX.6}
\]
These conclusions decide this exact Kummer input to AS-FF.39 at its
unbounded union. They do not decide reconstruction at that union.
\end{claim}
\begin{proof}
Write \(K=K_{p,\mathrm{nil}}\) and \(Q=I/N\).
The quotient \(Q\) is the maximal pronilpotent quotient of
\(I\): a finite quotient is nilpotent exactly when it kills
some lower-central term, and any finite compositum of such
quotients is again nilpotent. Its pro-\(p\) Sylow subgroup
\(Q_{(p)}\) is normal, and
\[
                    Q=Q_{(p)}\times T,
                                                        \tag{B.MIX.7}
\]
where \(T\) is the tame prime-to-\(p\) procyclic group.
Indeed every finite quotient of \(I\) has a normal wild
\(p\)-subgroup and cyclic prime-to-\(p\) quotient, by the
inertia sequence verified in B.53. In a nilpotent quotient
the different Sylow subgroups commute. Taking inverse limits
gives (B.MIX.7); the full tame quotient survives, since it
is already abelian. All roots of unity belong to \(K\).
For odd \(p\), the cyclotomic character on \(T\) surjects
onto \(\mu_{p-1}\subset\Z_p^\times\), and its pro-\(p\)
part comes from \(Q_{(p)}\).

Let \(a\in E^\times\) have zero rational Kummer class.
As above, \(G_E\) fixes the coefficient module, so this is
equivalent to having all compatible \(p\)-power roots of
\(a\) in \(E\). Enlarge to a finite extension \(E'/K\)
normal under the local absolute Galois group, containing the
finite local orbit \(a_1=a,\ldots,a_m\). This is possible
by adjoining the finitely many local conjugates of generators
for \(E/K\); the field \(K\) is stable under \(G_{\Q_p}\),
since \(N\) is characteristic in its normal inertia subgroup.
Choose a number field \(D\subset E'\), stable under
\(G_{\Q_p}\), containing these generators and conjugates,
and let \(M=D_v\), a finite Galois extension of \(\Q_p\).
Put
\[
 F=K_pD(\mu_{p^\infty}),\qquad
 R=F(a_j^{1/p^n}:1\leq j\leq m,\ n\geq1).
\]
Then \(R/K_p\) is Galois and \(R\subset E'\).
The group \(D_R=\operatorname{Gal}(R/F)\) embeds by
Kummer characters into \(\Z_p(1)^m\); in particular it is
torsion-free. The image \(C_R\) of \(N\) in
\(\operatorname{Gal}(R/K_p)\) is finite, because
\(E'/K\) is finite.

Choose \(t\in T\) whose cyclotomic character is a primitive
\((p-1)\)-st root \(\zeta\), and choose a lift
\(\tau\in I\) of \((1,t)\in Q_{(p)}\times T\).
The cyclotomic character factors through \(Q\), so
\(\chi(\tau)=\zeta\) exactly. Its image is central modulo
\(C_R\). Consequently
\[
             [\tau,D_R]\subset C_R\cap D_R=1.           \tag{B.MIX.8}
\]
Here the last equality uses finite versus torsion-free, not
an assertion that the entire radical extension is abelian
over \(K_p\).

Set \(J=I_M=I\cap G_M\) and \(J_0=J\cap\ker\chi\).
The group \(J_0\) fixes \(K_p,D,\mu_{p^\infty}\) and
maps onto \(D_R\). Kummer conjugation and (B.MIX.8) give
\[
       \kappa_{J_0}(\tau a)=\zeta\,\kappa_{J_0}(a).
                                                        \tag{B.MIX.9}
\]
There is no lost subgroup here: \(J\) already fixes the
strict-local field \(K_p\). Restriction
\(H^1(J,\Q_p(1))\to H^1(J_0,\Q_p(1))\) is injective,
because \(J/J_0\) is an abelian group with nontrivial
cyclotomic character and its first cohomology vanishes by
the displayed cocycle calculation in B.56. Restriction
from \(G_M\) to \(I_M\) is injective as well, since the
inertia invariants of \(\Q_p(1)\) vanish. Both \(a\) and
\(\tau a\) belong to \(M\), so (B.MIX.9) implies
\[
                \kappa_M(\tau a)=\zeta\kappa_M(a).
                                                        \tag{B.MIX.10}
\]

Valuation in (B.MIX.10) gives \(v(a)=\zeta v(a)\);
as \(\zeta\ne1\), the element \(a\) is a unit.
The principal-unit logarithm then gives
\(\log(\tau a)=\zeta\log(a)\).
The opened Baker--Brumer theorem, with proof in
\citep[Théorème, pp.~1--2]{Rousseau1969}, states that
\(p\)-adic logarithms of algebraic units which are linearly
independent over \(\Q\) remain linearly independent over
\(\overline\Q\). If \(\log(a)\ne0\), the preceding
algebraic relation therefore makes the two logarithms
\(\Q\)-dependent, and forces \(\zeta\in\Q\).
For \(p\geq5\), a primitive \((p-1)\)-st root is not
rational. Thus \(\log(a)=0\); the unit \(a\) in the
finite local field \(M\) is a root of unity. Conversely
all roots of unity have all their divisions in \(K\).
This proves (B.MIX.5), with every finite extension included.

For \(p=2\), each splitting field over \(K_2\) of
\(X^{2^n}-2\) has a finite \(2\)-group as Galois group:
adjoin the \(2^n\)-th roots of unity and then a Kummer
root. These are nilpotent groups, so all those roots lie
in \(K_{2,\mathrm{nil}}\), proving the first exception.

For \(p=3\), put \(z=\sqrt{-3}\),
\(F_0=K_3(z)=K_3(\mu_3)\), and
\(a=(2+z)/(2-z)\). The nonidentity element of
\(\operatorname{Gal}(F_0/K_3)\) sends \(a\) to \(a^{-1}\).
Let \(R_n=F_0(\mu_{3^n},b_n)\), where \(b_n^{3^n}=a\).
This is Galois over \(K_3\), and its kernel over \(F_0\)
is a \(3\)-group \(P_n\). Cauchy's theorem supplies an
involution \(s\) in the Galois group. It inverts
\(\mu_{3^n}\). Writing \(s(b_n)=\xi b_n^{-1}\), its
order two forces \(\xi^2=1\), hence \(\xi=1\).
For \(g\in P_n\), write
\(g(b_n)=\eta b_n\) and \(g(\zeta_{3^n})=\zeta_{3^n}^{u}\).
Direct substitution on these two generators gives
\(sgs^{-1}=g\). Thus
\(\operatorname{Gal}(R_n/K_3)=P_n\times\langle s\rangle\)
is nilpotent. All \(3\)-power roots of \(a\) lie in
\(K_{3,\mathrm{nil}}\). Finally
\(a+a^{-1}=2/7\), which is not an algebraic integer,
so \(a\) is not a root of unity. This proves the second
exception and completes the stated Kummer classification.
\end{proof}

\paragraph{Why the new Kummer theorem does not invoke an existing reconstruction theorem.}
The amenability hypothesis in
\citep[Definition~5.15(5.15.1)]{CarlsonHaineWolf2024}
requires each torus to have no nonzero divisible rational
points, with all primes included in characteristic zero.
It fails already for \(\mathbb G_m\), since
\(\mu_\infty\subset K_{p,\mathrm{nil}}\).
The independent condition~(5.15.2), requiring the Tate module
of \(\operatorname{Pic}(X)\) to vanish for every regular
finite-type \(X\), fails as well. For an explicit example take
\(A:y^2+xy=x^3-p\), with \(c_4=1\) and
\(\Delta=p(1-432p)\). Its node has the two rational tangent
lines \(y=0\), \(y+x=0\), hence split multiplicative reduction.
The opened Tate uniformization and torsion sequence
\citep[Theorem~1, author p.~6; pp.~18--19]{Tate1995}
give a compatible torus Tate line \(\Z_p(1)\subset T_pA\).
The subgroup \(I'\) fixes this line pointwise. Torsion points
of the globally defined elliptic curve are algebraic over
\(\Q\), so their coordinates lie in \(F_p\), and consequently
in \(K_{p,\mathrm{nil}}\). Their divisor classes give
\(\Z_p\subset T_p\operatorname{Pic}(A_{K_{p,\mathrm{nil}}})\).
Thus even the Picard condition has an actual counterexample;
it is not merely absent from the proof of (B.MIX.5).

The explicit positive argument for a fixed bounded nilpotency
class in AS-FF.39 also uses an inertia-invariant lift of every
visible generalized-Albanese functional. Passing to a
pronilpotent quotient does not retain its bounded-nilpotency
Lie-algebra argument: arbitrary pro-\(p\) groups are
pronilpotent. Formula (B.MIX.5) alone supplies neither that
lift nor compatible morphisms on all finite étale covers.
The unbounded tower therefore still requires a separate
reconstruction argument. The larger kernels at \(p=2,3\)
likewise refute a particular input, not reconstruction.
No single-topos reconstruction conclusion is deduced from
ordinary pro-reconstruction.

\subsection{Affine point functions over the unbounded tower when \texorpdfstring{\(p\geq5\)}{p >= 5}}

\begin{claim}{B.58}{Proved}{An actual Kummer value detects a rational scalar}
Let \(p\geq5\), \(K=K_{p,\mathrm{nil}}\), and let \(E/K\)
be finite. If \(b\in E^\times\) and \(\alpha\in\Q_p\)
satisfy
\[
                 \kappa_E(b)=\alpha\kappa_E(p),           \tag{B.MIX.11}
\]
then \(\alpha\in\Q\).
\end{claim}
\begin{proof}
Pass to a finite extension \(E'/K\) normal under
\(G_{\Q_p}\), as in B.57, and put \(N'=G_{E'}\).
It is normal in \(I\), contained in \(N=G_K\), and of
finite index in \(N\). For \(g\in I\), conjugate (B.MIX.11)
on \(N'\). The restricted Kummer class of the rational
number \(p\) is \(I\)-invariant, since it comes from
\(H^1(I,\Q_p(1))\). Hence
\(\kappa_{E'}(g(b)/b)=0\). B.57 makes every
\(g(b)/b\) a root of unity. The finite local orbit of
the globally algebraic number \(b\) lies in one number
field, whose roots of unity form a finite group. A single
integer \(m>0\) therefore gives
\(b^m\in\overline\Q^I=K_p\).

The class
\(\kappa_I(b^m)-m\alpha\kappa_I(p)\) restricts to zero
on \(N'\), so it inflates from \(H^1(I/N',\Q_p(1))\).
This group is zero. Indeed the kernel \(N/N'\) of
\(I/N'\to I/N=Q\) is finite and fixes the coefficient
module, so its rational first cohomology is zero.
Inflation--restriction identifies the group with
\(H^1(Q,\Q_p(1))\). The group \(Q\) has a central
tame element \(t\) with \(\chi(t)\ne1\).
For every cocycle \(c\), centrality gives
\((\chi(t)-1)c(g)=(\chi(g)-1)c(t)\), so \(c\) is a
coboundary. This proves the claimed vanishing.

The element \(b^m\in K_p\) lies in a finite unramified
local extension \(M/\Q_p\), with \(I_M=I\).
The local Kummer class comparison is injective upon
restriction to inertia, because \(\Q_p(1)^I=0\).
Thus the preceding zero class already vanishes on
\(G_M\). Normalizing its valuation by \(v_M(p)=1\)
now yields
\(m\alpha=v_M(b^m)\in\Z\). Hence \(\alpha\) is rational.
\end{proof}

\begin{claim}{B.59}{Proved}{Affine actual point functions for all pinned affine-line automorphisms}
Let \(p\geq5\), and let \(E\) be any finite extension
of \(K_{p,\mathrm{nil}}\). Every pinned autoequivalence
of \(\mathbb A^1_{E,\mathrm{et}}\) over \(BG_E\), with
its relative structural isomorphism retained, has actual
geometric point function
\[
                         a\longmapsto ua+v,
                 \qquad u\in E^\times,\ v\in E.          \tag{B.MIX.12}
\]
The coefficients are unique. This conclusion concerns the
actual point function; it does not identify the remaining
pointwise-trivial automorphism kernel on the full étale site,
or assert existence for arbitrary noninvertible pinned maps.
\end{claim}
\begin{proof}
Let \(f:\overline\Q\to\overline\Q\) be the actual point
function. Relative point evaluation
\citep[Proposition~3.6]{CarlsonHaineWolf2024}, together with
the inverse equivalence, makes it a \(G_E\)-equivariant
bijection. On every finite extension \(L/E\), it gives a
bijection of \(L\). Compose with the unique affine scheme
automorphism taking \(f(0),f(1)\) to \(0,1\); henceforth
\(f(0)=0\), \(f(1)=1\).

For \(a\in\overline\Q\), the equivalence after a finite
scalar extension restricts to the lines punctured at
\(a\) and \(f(a)\). On their rank-one geometric
\(p\)-Kummer lines it acts by
\(\alpha_a\in\Z_p^\times\). The arithmetic summand
cancels when two point evaluations are subtracted. Thus,
over any common finite extension containing \(a,b,x\),
\[
 \kappa_L\!\left(\frac{f(x)-f(a)}{f(b)-f(a)}\right)
       =\alpha_a\kappa_L\!\left(\frac{x-a}{b-a}\right)
       \qquad (b,x\ne a).                              \tag{B.MIX.13}
\]
This only uses the one-puncture cohomology calculation;
no preservation theorem for the place at infinity and
no equality of different inertia exponents is assumed.
Set \(b=a+1\), \(x=a+p\). B.58 applies to the left
argument in (B.MIX.13), and proves
\(\alpha_a\in\Q^\times\) for every \(a\).

Fix a complex embedding of \(\overline\Q\), and write
\(|\cdot|\) for its ordinary complex absolute value.
Clear the rational denominator of \(\alpha_a\) in
(B.MIX.13). B.57 says that the resulting multiplicative
error is a root of unity. Its complex absolute value is one.
Consequently there is a positive real constant \(C_a\)
such that
\[
                 |f(x)-f(a)|=C_a|x-a|^{\alpha_a}
                       \qquad(x\ne a).                  \tag{B.MIX.14}
\]
One can take \(C_a=|f(a+1)-f(a)|\).

We spell out how (B.MIX.14) forces the exponents and constants
to coincide. If \(\alpha_a<0\), the sequence \(f(n)\), for
positive rational integers \(n\to\infty\), tends to \(f(a)\).
For another centre \(b\ne a\), its distance from \(f(b)\)
would therefore tend to the finite positive number
\(|f(a)-f(b)|\). The formula at \(b\) makes that distance
tend to either zero or infinity, since \(\alpha_b\ne0\).
This contradiction proves \(\alpha_a>0\) for every \(a\).
The reverse triangle inequality gives
\[
 \left|C_a|n-a|^{\alpha_a}-C_b|n-b|^{\alpha_b}\right|
                      \leq |f(a)-f(b)|.
\]
Divide by the larger power of \(n\) and let \(n\to\infty\).
Unequal positive exponents give a contradiction, so all
\(\alpha_a\) equal one positive number \(\alpha\).
Dividing again by \(n^\alpha\) gives \(C_a=C_b\).
Normalization at zero and one makes the common constant one.
Thus
\[
                       |f(x)-f(y)|=|x-y|^\alpha
                       \qquad(x,y\in\overline\Q).       \tag{B.MIX.15}
\]
The triangle with vertices \(0,1,2\) gives
\(2^\alpha\leq2\), hence \(\alpha\leq1\).
Since \(f\) is bijective, its inverse satisfies the same
identity with exponent \(1/\alpha\), and gives
\(1/\alpha\leq1\). Hence \(\alpha=1\).

An isometry of \(\overline\Q\subset\mathbb C\) fixing
zero and one is the identity or complex conjugation.
Indeed distances to zero and one give \(f(i)=i\) or
\(f(i)=-i\); distances to the three noncollinear points
\(0,1,i\) then determine the image of every point, with
the second case giving the reflected point.

The conjugation case cannot be \(G_E\)-equivariant.
First \(E\) is not separably closed. For \(p\geq5\),
the splitting field of \(X^p-p\) over \(K_p\) has
non-nilpotent Galois group
\(\mathbb F_p\rtimes\mathbb F_p^\times\): the radical
degree is \(p\), the cyclotomic tame degree is \(p-1\),
and the natural scalar action is faithful. Thus \(N\ne1\).
It is a closed subgroup of the free wild pro-\(p\) group,
so is torsion-free and infinite; every finite-index subgroup
\(G_E\) is nontrivial. The field \(E\) carries its
nontrivial rank-one henselian mixed-characteristic valuation.
If complex conjugation centralized \(G_E\), it would
preserve \(E\). The opened
\citep[Lemma~2.1]{KhassaKhanduja2010} makes this automorphism
continuous for the valuation; in rank one it preserves the
valuation ring. Uniqueness of prolongation then makes it
preserve the chosen valuation on \(\overline\Q\), so it
lies in the local decomposition group \(G_{\Q_p}\).
Its element of order two maps trivially to the unramified
quotient \(\widehat\Z\), then trivially to the torsion-free
tame inertia quotient. It would lie in wild pro-\(p\)
inertia, which has no element of order two for \(p\geq5\).
This is impossible. The normalized \(f\) is therefore the
identity. Undoing normalization gives (B.MIX.12).
\end{proof}

\begin{claim}{B.60}{Proved}{Actual Kummer values force rational boundary residues}
Let \(p\geq5\), let \(E/K_{p,\mathrm{nil}}\) be finite,
and let \(U\) be a smooth geometrically connected curve over
\(E\), with smooth proper compactification \(\overline U\).
Suppose \(\xi\in H^1(U,\Z_p(1))\) evaluates at every point
\(x\in U(L)\), for every finite extension \(L/E\), to the
Kummer class of an actual element of \(L^\times\).
Every geometric boundary residue of \(\xi\) belongs to
\(\Q\cap\Z_p\).
\end{claim}
\begin{proof}
Split the finitely many boundary points over a finite extension
and fix one of them, \(q\). Riemann--Roch supplies a nonconstant
rational function \(g\) with its only pole at \(q\) and with a
zero at every other boundary point: prescribe those finitely
many zeros and allow a sufficiently large pole at \(q\).
The map \(g:\overline U\to\mathbb P^1\) is finite.
Remove zero, infinity, and all branch values from the target,
obtaining \(V\); its inverse image \(U'\) lies in \(U\)
and is finite étale over \(V\). Corestrict \(\xi|_{U'}\)
to a class \(\eta\in H^1(V,\Z_p(1))\).

Its geometric residues vanish away from zero and infinity.
At a removed branch value the original class has zero
residues at all its preimages. The corestriction residue
is the sum of source residues multiplied by residue-field
degrees, not by ramification indices. This formula follows
at finite Kummer levels from valuation of the field norm,
and then by passage to the compatible inverse limit.
The residue at infinity is precisely the residue at the
sole point \(q\), since \(q\) is rational over the chosen
base. Thus the class has the form
\(\eta=\eta_0+\alpha\kappa(T)\), with \(\eta_0\) arithmetic
and \(\alpha=-\operatorname{res}_q\xi\). This also follows
directly from the Kummer sequence on the punctured rational
line, whose Picard group is zero; zero residues at the
additional punctures remove their basis terms.

At every \(t\in V(E)\), the evaluation of \(\eta\) is
an actual Kummer value. Indeed the fibre is a product of
finite separable extensions, and corestriction evaluates
as the product of the field norms of the actual values
of \(\xi\) on these points. Choose \(t_0\in V(E)\) and
a positive integer \(n\) for which \(p^nt_0\in V(E)\);
only finitely many choices are excluded. Subtracting the
two evaluations cancels \(\eta_0\) and gives an actual
value \(b\in E^\times\) with
\(\kappa_E(b)=n\alpha\kappa_E(p)\).
B.58 gives \(n\alpha\in\Q\), hence the desired
rationality of the residue. The argument works for each
boundary point and is unchanged by the finite scalar
extension used to split them.
\end{proof}

\begin{claim}{B.61}{Proved}{Polynomial actual point functions for arbitrary affine-line morphisms}
Let \(p\geq5\), and let \(E\) be any finite extension
of \(K_{p,\mathrm{nil}}\). Every pinned geometric morphism
\[
       \mathbb A^1_{E,\mathrm{et}}
              \longrightarrow\mathbb A^1_{E,\mathrm{et}}
\]
over \(BG_E\) has actual geometric point function
\(x\mapsto P(x)\) for a unique polynomial \(P\in E[T]\).
Invertibility of the geometric morphism is not assumed.
This does not identify the whole inverse-image functor
with the functor induced by \(P\), and does not prove
existence for an arbitrary source scheme.
\end{claim}
\begin{proof}
Write \(f:\overline\Q\to\overline\Q\) for the actual
point function supplied by relative pinning. If the
underlying generic point maps to a closed point, continuity
makes the entire image that closed point. Evaluation at
zero makes it an \(E\)-rational point, and the assertion
holds with a constant polynomial. Otherwise the generic
point maps to the generic point. Each geometric fibre
\(S_a=f^{-1}(a)\) is then finite: after a finite constant
extension containing \(a\), it is a proper closed subset
of the source affine line. In particular the image of
\(f\) is infinite.

For \(a\in\overline\Q\), restrict to
\(U_a=\mathbb A^1\setminus S_a\) and pull back the
Kummer class of \(T-a\). Split \(S_a\) after a finite
extension. The resulting class is an arithmetic class
plus
\(\sum_{r\in S_a}q_{a,r}\kappa(T-r)\), because a
localization of the polynomial ring has trivial Picard
group. B.60 gives \(q_{a,r}\in\Q\cap\Z_p\).
Subtract the evaluation at any fixed \(x_0\notin S_a\),
and clear these finitely many rational denominators.
B.57 identifies the resulting multiplicative error
with a root of unity, over every finite field of
point definition. In a fixed complex embedding, this gives
\[
 |f(x)-a|^2=P_a(x),\qquad
 P_a(z)=C_a^2\prod_{r\in S_a}|z-r|^{2q_{a,r}},
 \quad C_a>0,
 \qquad x\in\overline\Q\setminus S_a.                   \tag{B.MIX.16}
\]
No sign or integrality of the exponents has been assumed.
The function \(P_a\) is positive and real analytic off
\(S_a\), and \(\log P_a\) is harmonic there.

Put \(S=S_0\cup S_1\cup S_i\) and
\(\Omega=\mathbb C\setminus S\). Define a real-analytic
function on this connected open set by
\[
 F(z)=\frac{P_0(z)-P_1(z)+1}{2}
          +i\frac{P_0(z)-P_i(z)+1}{2}.                  \tag{B.MIX.17}
\]
The elementary squared-distance identities show that
\(F(x)=f(x)\) at every algebraic \(x\in\Omega\).
Density and continuity then give
\(|F(z)-a|^2=P_a(z)\) on \(\Omega\setminus S_a\)
for every algebraic \(a\). The function \(F\) is
nonconstant, since otherwise \(f\) would have finite image.

Here is the analytic implication in full. At a point
\(z\in\Omega\), write
\(w=F(z)\), \(A=F_z(z)\), \(B=F_{\bar z}(z)\), and
\(C=F_{z\bar z}(z)\). For all algebraic \(a\ne w\)
for which \(z\notin S_a\), harmonicity yields
\[
 \operatorname{Re}\left(\frac{C}{w-a}
                  -\frac{AB}{(w-a)^2}\right)=0.         \tag{B.MIX.18}
\]
There is at most one excluded algebraic centre: all
\(S_a\) consist of algebraic points and are disjoint.
Thus density extends (B.MIX.18) to every complex
\(a\ne w\). Substituting \(b=1/(w-a)\) gives
\(\operatorname{Re}(Cb-ABb^2)=0\) for every nonzero
\(b\in\mathbb C\). The two polynomial coefficients
must vanish, so \(C=0\) and \(AB=0\).
Consequently \(F\) is harmonic and conformal.
More explicitly, if \(F_z\) is nonzero anywhere,
\(F_{\bar z}\) vanishes on an open neighbourhood, and
real-analytic continuation makes it vanish on all of
\(\Omega\). Otherwise \(F_z\) vanishes identically.
Thus \(F\) is holomorphic or antiholomorphic on \(\Omega\).

For \(r\in S\), choose an algebraic centre
\(a\ne f(r)\). Its fibre omits \(r\), so \(P_a\)
is finite and positive near \(r\). Equation (B.MIX.16)
makes \(F\) bounded on a punctured neighbourhood.
The singularity is removable. Choose three noncollinear
such centres; their distance identities at \(r\),
including the original identities for \(f(r)\), show
that the extended value \(F(r)\) equals \(f(r)\).
This gives an entire holomorphic or antiholomorphic
function agreeing with \(f\) on every algebraic point.
For any fixed centre \(a\), (B.MIX.16) bounds its growth
at infinity by a finite power of \(|z|\). Cauchy's
estimates make it a polynomial or an antipolynomial:
\(F(z)=P(z)\) or \(F(z)=P(\bar z)\).
Interpolation at distinct rational integers, whose
images lie in \(E\), gives \(P\in E[T]\) in either case.

The nonconstant antipolynomial case is impossible.
Let \(c\) denote complex conjugation on \(\overline\Q\).
For \(g\in G_E\), equivariance would give
\(P(cg x)=P(gc x)\) for every \(x\). Thus the field
commutator \(h=cgcg^{-1}\) satisfies \(P(hz)=P(z)\)
for every \(z\). Apply this to \(z+n\), with
\(n\in\Z\). If \(d>0\) is the degree of \(P\),
comparison of the coefficients of \(n^{d-1}\) in
\(P(hz+n)=P(z+n)\) gives \(hz=z\).
So \(c\) centralizes \(G_E\), which is excluded by
the valuation and ramification argument in B.59.
Hence \(F=P\), proving the asserted actual point
function. Its polynomial is unique because the infinite
set of rational inputs already determines it.
\end{proof}


% Pro-category completion, developed during the independent configuration audit.
% Provisional B.PRO identifiers; parent editor may renumber.
\subsection{Completion of ordinary pro-reconstruction in arbitrary transcendence degree}

We retain the ordinary category \(\mathsf T\), the model pro-objects,
the slice over \(\mathcal B\), and the limit-point definition of
pinnedness in B.43. In particular no compatible structural 2-cells
are assumed in the definition of a pro-morphism. The following
argument supplies the coherence that is used, and then reconstructs
at a fixed source field. This avoids trying to make the transition
maps pinned.

\begin{claim}{B.62}{Proved}{Compact coherence for ordinary pro-morphisms}
Let \(k\) have characteristic zero, let \(X,Y\) be finite-type
\(k\)-schemes, and let
\(\alpha:\mathcal X\to\mathcal Y\) be an ordinary pro-morphism
over \(\mathcal B\). There is a geometric morphism
\[
                 \alpha_k:X_{\mathrm{et}}\longrightarrow
                           Y_{\mathrm{et}}                 \tag{B.PRO.1}
\]
over \(BG_k\), compatible with the finite-level representatives of
\(\alpha\). Its underlying-space map is the map on inverse limits
defined by \(\alpha\), so it is pinned if \(\alpha\) is pinned.
Moreover finite constant base change, precomposition by a scheme
morphism, restriction to open subobjects, and slicing by a
quasi-compact separated étale object preserve the class of maps
obtained from ordinary pro-morphisms. For a slice defined at a
particular target level, the source model can be taken to be its
pullback by that level's chosen representative.
\end{claim}
\begin{proof}
Nilpotents do not affect the étale sites, so reduce the source.
At each target field \(F\), choose a representative
\(a_F:X_{E_F,\mathrm{et}}\to Y_{F,\mathrm{et}}\), where
\(E_F\supset F\) is finitely generated. Equality over the pro-base
permits one further source refinement making this a relative
morphism over \(BG_F\). Its structural 2-cell is unique:
\(G_F\) is slim by the opened Hoshi--Tsujimura Corollary~4.5,
and the image of each connected component's fundamental group in
\(G_F\) is open. For the latter assertion first map to \(G_{E_F}\);
the relative algebraic closure of \(F\) in the finitely generated
field \(E_F\) is finite. The same argument applies after every
finitely generated source refinement. Thus every absolute
comparison isomorphism at such a refinement is automatically
relative to \(BG_F\).

The opened \citep[Theorem~A.8 and Corollary~A.9]{CarlsonHaineWolf2024}
apply: the source is topologically noetherian and the target qcqs.
Every representative is coherent, and its inverse image takes
quasi-compact separated étale schemes to quasi-compact separated
étale schemes. Compose with \(X\to X_{E_F}\), obtaining functors
from the qcs étale sites of \(Y_F\) to that of \(X\).

Here is the finite-choice point that permits coherence. For two
qcs separated étale schemes \(U,V\) over a reduced noetherian
scheme \(X\) with finitely many irreducible components,
\(\operatorname{Isom}_X(U,V)\) is finite. An isomorphism is
determined on the finite fibres over the generic points: equality
there implies equality on \(U\), since \(U\) is reduced, those
fibres are dense, and \(V\to X\) is separated. The possible maps
of the finite generic étale algebras form a finite set.
Consequently the set of natural isomorphisms between two of the
preceding inverse-image functors is a closed subset of a product
of finite sets. It is compact, and the relative condition over
\(BG_F\) is closed as well.

For \(F\subset H\), ordinary pro-compatibility supplies an
isomorphism between the functor at \(F\) and the restriction of
the functor at \(H\), after a single finitely generated source
refinement and hence after base change to \(k\). The compact
space of relative comparison isomorphisms is therefore nonempty.
In each such compact space use the \emph{closure of the subset
of isomorphisms which descend after a finitely generated source
refinement}. This refinement is important: an arbitrary isomorphism
over \(k\) need not descend as an isomorphism of the entire
inverse-image functors. Membership in this closure says that on
every finite list of étale objects its components agree with those
of one whole comparison isomorphism descending at a finite stage.
Take the product of these closed spaces over all pairs of fields, and
impose naturality, identity, and the triple cocycle equations.
These are closed conditions. Every finite collection of them is
satisfiable: choose one common upper target field, compare its
representative with all the finitely many representatives after
one common finitely generated source refinement, and obtain all
comparisons by passing through this one representative. They
then satisfy every required finite cocycle equation. Compactness
gives a simultaneous coherent choice. This proves existence of
coherence; it has not been inserted into the hypothesis.

Every qcs étale object, every morphism between such objects, and
every finite diagram over \(Y\) descends to a finitely generated
field, by finite presentation. Finite covering families descend
too; the exact site-continuity statement is the opened
\citep[Lemma~59.51.2, Tag~0EYL, together with
Lemma~32.10.1, Tag~01ZM]{StacksProject}. We use the canonical
associative base-change identifications, or equivalently first
strictify this fixed diagram of sites. The coherently identified
inverse-image functors consequently
define a left exact covering-preserving functor on this site of
\(Y\), hence (B.PRO.1). The identical construction on finite étale
constant objects gives its structural isomorphism over \(BG_k\).
The comparison on open subobjects proves the assertion on spaces.

For clarity, slicing entails no additional coherence assumption.
Suppose \(V_F\to Y_F\) is qcs separated étale and put
\(U_{E_F}=a_F^*V_F\). For any higher target level, the preceding
comparisons identify the pullback of \(V\) with \(U\) over \(X\).
This is an isomorphism between finitely presented étale schemes,
so it descends after one finitely generated source refinement.
The resulting slice map represents a map of the model systems
\(\mathcal U\to\mathcal V\). For ordinary pro-compatibility,
approximate a selected comparison on the finite list consisting
of \(V\) and any specified arrows between slice objects by a
\emph{whole} comparison isomorphism descending at a finite stage.
This is exactly the closure property imposed above. Its component
on \(V\) agrees with the selected identification, so it induces
the required comparison of the entire slice geometric morphisms
after one refinement. Thus descending an individual isomorphism
of étale objects is not being substituted for descending a
comparison of functors. The same construction
over a finite étale arithmetic base gives finite constant base
change; cofinal fields have the form \(F'L_0\) for a fixed finite
field extension \(L_0/F_0\) defining that base change. Open
restriction and scheme precomposition are immediate. If the
original map is pinned, so are the slice maps: an étale
finite-type scheme has only closed points above a closed point.
\end{proof}

\begin{claim}{B.63}{Proved}{Pro-points and a single-stage Kummer comparison}
For every characteristic-zero \(k\), pinned ordinary pro-morphisms
\(\mathcal B\to\mathcal Y\) over \(\mathcal B\) are precisely
the \(k\)-rational points of \(Y\). This statement is compatible
with finite extensions of \(k\).

In particular, suppose a fixed representative of a pinned
pro-morphism pulls the universal Kummer class back to
\(\xi\in H^1(X_E,\widehat\Z(1))\). If \(x\in X_E(L)\), with
\(L/E\) finite, and \(b\) is the value of its induced pro-point
in \(\mathbb G_m\), then there is \emph{one} finitely generated
field \(M\) containing \(L\) and \(b\) for which
\[
        \operatorname{res}_{M/L}(x^*\xi)=\kappa_M(b)
                \quad\hbox{in }\widehat{M^\times}.       \tag{B.PRO.2}
\]
The field \(M\) works simultaneously for all Kummer levels.
\end{claim}
\begin{proof}
Let the limit point of a pinned pro-point be \(y\in Y\).
It is closed, and a finite closed subscheme containing it descends
to some \(Y_{F_0}\). At every sufficiently large target field
\(F\), its image \(y_F\) is therefore closed, with
\(\kappa(y_F)/F\) finite. A representative
\(BG_E\to Y_{F,\mathrm{et}}\) has underlying image \(y_F\).
It factors through the closed subtopos
\(B G_{\kappa(y_F)}\): the inverse image of the complementary
open subobject is empty. After one source refinement it is
relative over \(BG_F\), with the unique structural 2-cell
established in B.62. Since \(\kappa(y_F)/F\) is finite
separable, the finite étale slice description says that this
relative morphism is exactly an \(F\)-embedding
\(\kappa(y_F)\hookrightarrow E\). Thus it is an actual
\(E\)-valued point of \(Y_F\). Compatibility at a common
source refinement makes these points the same \(k\)-point;
the finite residue extension of the limit point consequently
has degree one. This also proves uniqueness of the pro-morphism.

For the second assertion base change the pro-morphism to the
finite extension \(kL/k\), using B.62, and compose with the
point \(x\). The first assertion gives an actual value
\(b\in(kL)^\times\). Choose a finitely generated target field
\(H\subset kL\) containing \(b\) and all fields needed for the
displayed point. At target \(\mathbb G_{m,H}\), the pro-point's
representative has the closed underlying image \(T=b\).
After one refinement making it relative to \(BG_H\), it is
the constant-point morphism, as just proved. Its comparison
with the original representative at the original target is an
isomorphism of \emph{whole geometric morphisms} after one common
finitely generated source refinement \(M\). Pulling back all
the compatible universal Kummer torsors gives (B.PRO.2).
This proves simultaneous equality in the completed multiplicative
group, rather than a sequence of equalities over fields depending
on \(n\), or an equality only after restriction to \(G_{kL}\).
\end{proof}

\begin{claim}{B.64}{Proved}{Intersection with a finitely generated field's Kummer completion}
Let \(A\subset B\) be finitely generated fields of characteristic
zero. If \(\xi\in\widehat{A^\times}\), \(b\in B^\times\), and
\[
                    \operatorname{res}_{B/A}\xi=\kappa_B(b),
                                                               \tag{B.PRO.3}
\]
then \(b\in A^\times\) and \(\xi=\kappa_A(b)\).
\end{claim}
\begin{proof}
Let \(A'\) be the relative algebraic closure of \(A\) in
\(B\); it is finite over \(A\). Take a normal projective
\(A'\)-model of the finitely generated regular extension
\(B/A'\). Every codimension-one valuation of this model is
trivial on \(A\). Applying its completed valuation to
(B.PRO.3) shows that the ordinary integer valuation of \(b\)
is zero. Normality then makes \(b\) a global invertible function
on this proper model. Such functions belong to its constant
field \(A'\), so \(b\in A'^\times\).

The map \(\widehat{A'^\times}\to\widehat{B^\times}\) is
injective: for every \(n\), an \(n\)-th root in \(B\) of an
element of \(A'\) is algebraic over \(A'\), hence already in
\(A'\). Thus (B.PRO.3) already holds over \(A'\). Taking the
norm to \(A\), with \(d=[A':A]\), gives
\[
                       d\xi=\kappa_A(N_{A'/A}b).          \tag{B.PRO.4}
\]

We record the elementary saturation argument, so no unsupported
finite-generation claim about \(A^\times\) is needed. In the
completion \(\widehat{A^\times}=\varprojlim_n
A^\times/(A^\times)^n\), if \(d\eta=0\), then its component
modulo \(n\) is represented by a \(d\)-th root of unity.
Indeed lift to level \(nd\): if \(a^d=c^{nd}\), then
\(a=c^n\zeta\) with \(\zeta^d=1\). The finite set
\(\mu_d(A)\) and compatibility over common multiples imply
\(\eta=\kappa_A(\zeta)\) for one \(\zeta\in\mu_d(A)\).
If \(d\xi=\kappa_A(a)\), reduction modulo \(d\) first gives
\(a=c^d\) with \(c\in A^\times\). Apply the preceding
torsion assertion to \(\xi-\kappa_A(c)\); hence \(\xi\)
is the Kummer class of an actual element \(c'\in A^\times\).
Apply this to (B.PRO.4). Finally Kummer maps of finitely generated
characteristic-zero fields are injective, by
\citep[Proposition~5.24 and Observation~6.9]{CarlsonHaineWolf2024}.
Equation (B.PRO.3) therefore makes \(c'=b\) in \(B\), proving
both assertions.
\end{proof}

\begin{claim}{B.65}{Proved}{Reconstruction at one fixed source field}
Let \(\alpha\) be a pinned ordinary pro-morphism and choose a
relative representative
\[
                   a:X_{E,\mathrm{et}}\longrightarrow
                     Y_{F,\mathrm{et}}                 \tag{B.PRO.5}
\]
over \(BG_F\), with \(F\subset E\subset k\) finitely
generated. If \(X_E\) is normal, then (B.PRO.5) is induced by a
unique \(F\)-scheme morphism \(f_E:X_E\to Y_F\).
Its lift to \(Y_E\), after extension to \(k\), realizes the
point function of \(\alpha_k\). The same conclusion holds with
the absolutely weak normalization of \(X_E\) as source.
\end{claim}
\begin{proof}
\emph{1. Units over the fixed field.}
First take \(Y_F=\mathbb G_{m,F}\), and suppose \(X_E\) is
affine, geometrically integral, normal, and has an \(E\)-point.
The pullback \(\xi\) of the compatible Kummer classes lies in
\(\widehat{\mathcal O(X_E)^\times}\). For normal schemes this
follows from the Kummer sequence and finite generation of
\(\operatorname{Pic}(X_E)\), as stated in the opened
\citep[Proposition~5.22]{CarlsonHaineWolf2024}.
For every finite \(L/E\) and \(x\in X_E(L)\), B.63 gives
one finitely generated \(M\supset L\) and a point value \(b\)
satisfying (B.PRO.2). B.64 gives
\[
                         b\in L^\times,\qquad
                       x^*\xi=\kappa_L(b).              \tag{B.PRO.6}
\]
In particular the value at the chosen \(E\)-point is an actual
element of \(E^\times\); divide it out.

The normalized unit group is finitely generated free, by
\citep[Corollary~6.18]{CarlsonHaineWolf2024}. The evaluation of
the normalized completed unit at every algebraic point is an
actual algebraic unit by (B.PRO.6). The arithmetic argument of
Proposition~6.22 of that source now makes the completed unit
an actual unit of \(X_E\). Its hypotheses are precisely a
finitely generated characteristic-zero field (hence amenable
and admitting a nontrivial discrete valuation), the stated
finitely generated free unit group, and these actual point
evaluations. The proof works for normal sources as well:
Picard finite generation is Proposition~5.22, the unit assertion
is Corollary~6.18, and the nonconstant unit maps used in Claim~6.23
have image containing a nonempty open of \(\mathbb G_m\).
The source explicitly records this normal-source extension in
Remark~4.20. Equivalently, its full-rank integer valuation matrix
makes an integer multiple of each profinite exponent integral;
\(\widehat\Z\cap\Q=\Z\) then makes each exponent integral.
Thus \(\xi=\kappa(f)\) for \(f\in\mathcal O(X_E)^\times\).

This \(f\) realizes the value at \emph{every} point of
\(X(\overline k)\), not only the points algebraic over \(E\).
For such a point choose a finitely generated field of definition
\(L\supset E\) inside \(\overline k\). This extension need
not be finite. First base change the pro-system to the finite
extension \(kL/k\), then refine its source model from \(E\)
to \(L\). The point is now \(L\)-rational, so B.63 applies
with the identity finite extension of the new source field
\(L\); B.64 applies to the resulting finitely generated
comparison field. Its value \(b\) satisfies
\(\kappa_L(b)=\kappa_L(f(x))\); injectivity gives
\(b=f(x)\). The finite extension of \(k\) needed to contain
the point is permitted by B.62.

Remove the rational-point and geometric-connectedness assumptions
by a finite Galois extension of \(E\) splitting the finitely
many geometric components and supplying a point on each. The
preceding proof gives units on these components. Their Kummer
classes are the pullback of the one fixed \(\xi\), so Galois
conjugates agree by Kummer injectivity. They descend to a unit
over \(E\). Passing to this finite extension is justified by
finite constant base change to its compositum with \(k\).
Finally a finite affine open cover handles a nonaffine normal
source; equality of the point functions glues the units uniquely.

\emph{2. Coordinates and the point function.}
For general \(Y_F\), inverse images of Zariski opens under
\(a\) are Zariski opens of \(X_E\), and the corresponding pro
restrictions are pinned. For a regular function \(h\) on an
affine target open, cover that open by \(D(h)\) and \(D(h-1)\).
Step~1 reconstructs the pulled-back units \(h\) and \(h-1\)
on their respective source opens over the same field \(E\).
On the overlap their difference is one, since this holds on
all \(\overline k\)-points and the source is reduced. They
therefore give a regular function on the entire inverse-image
open. The ring operations and all defining relations hold on
\(\overline k\)-points, hence as functions. The resulting maps
on affine target opens glue uniquely to
\(f_E:X_E\to Y_F\). Their point function is that of
\(\alpha_k\). The open inverse-image conditions hold as well,
because finite-type schemes over a field are Jacobson and
membership in an open is determined by the closed points of
its complement. This argument does not require the target to
be separated.

\emph{3. The entire functor at the same source level.}
For each qcs separated étale \(V_F\to Y_F\), the sheaf
\(a^*V_F\) is represented by a qcs separated étale scheme
\(U_E\to X_E\), by CHW A.8--A.9. It is normal. Slice the
pro-morphism, as in B.62, using the fixed representative
\(a_V:U_{E,\mathrm{et}}\to V_{F,\mathrm{et}}\).
Step~2 gives a scheme map \(g_V:U_E\to V_F\) over this
same field \(E\), with the slice point function. This step
does not require the whole topos map of the slice to have been
reconstructed already. Composing to \(Y_F\) and comparing
point functions gives
\[
                 g_V\text{ over }Y_F
                    = f_E\circ(U_E\to X_E).
\]
It induces \(h_V:U_E\to X_E\times_{Y_F}V_F\).

The coherent realization (B.PRO.1) supplies the usual étale
slice fibre bijection: above every \(x\in X(\overline k)\),
the points of \(U\) are precisely the points of \(V\) above
\(\alpha_k(x)\). The map \(h_V\) realizes that bijection.
An étale morphism between étale schemes over \(X\) which is
bijective on every geometric closed fibre is an isomorphism:
the failure loci are constructible, so any nonempty such locus
over a field would contain a closed point. Thus \((h_V)_k\)
is an isomorphism, and faithful flatness makes \(h_V\) an
isomorphism over \(E\). This argument requires no additional
extension depending on \(V\).

For a morphism \(V_F\to W_F\), the two maps between the
reconstructed pullbacks agree on \(\overline k\)-points by
the natural slice functor of \(\alpha_k\); reducedness makes
them equal as scheme maps. The \(h_V\) consequently form a
natural isomorphism \(a^*\simeq f_E^*\) on the entire qcs
étale site. That site is a basis, so (B.PRO.5) is the induced
geometric morphism. This is equality at the fixed source field,
not just an equality of limit point functions or an objectwise
collection of equalities after unrelated extensions.

\emph{4. Faithfulness, automorphisms, and a seminormal source.}
We spell out the descent needed for the last assertion. For two
\(F\)-scheme maps from a reduced finite-type \(E\)-scheme to
\(Y_F\), an isomorphism of their relative geometric morphisms
implies equality of their pullbacks of all open subobjects and
all Kummer classes of units on target opens. Unit Kummer maps
on the source opens are injective: evaluate at closed points,
whose residue fields are finitely generated of characteristic
zero, and use their injectivity, followed by reducedness.
The coordinate argument of Step~2 therefore makes the two
scheme maps equal. This proves the required faithfulness even
when \(E/F\) is transcendental.

An actual such scheme map has no automorphisms as a geometric
morphism. At any source point \(x\), its residue field and
that of its image are finitely generated over \(\Q\); the
image of the source residue Galois group in the target residue
Galois group is open. The latter group is slim. Thus the
centralizer describing an automorphism on that stalk is trivial.
Geometric stalks are conservative, proving the assertion for
the whole inverse-image functor.

Replace \(X_E\) by its absolute weak normalization; in
characteristic zero this is finite-type seminormal. Its finite
normalization \(Z_E\to X_E\) is an h-cover. Apply the normal
case to \(Z_E\) and to a normalization of the reduced fibre
product \(Z_E\times_{X_E}Z_E\). The preceding faithfulness
makes the two maps agree after this normalization. Since that
normalization is an h-cover, h-separatedness in CHW Theorem~4.1
makes them agree on \((Z_E\times_{X_E}Z_E)^{\mathrm{awn}}\)
before applying the h-sheaf equalizer. The opened
\citep[Theorem~4.1]{CarlsonHaineWolf2024}, which states that
\(T\mapsto\operatorname{Hom}(T^{\mathrm{awn}},Y_E)\)
is an h-sheaf, descends them to a scheme map on \(X_E\).
Their topos isomorphisms are unique by the just-proved absence
of automorphisms and hence satisfy the cocycle condition.
Étale-topos h-descent, the opened Proposition~4.8 of the same
source, descends the isomorphism to (B.PRO.5) itself. This proves
the last assertion and uniqueness.
\end{proof}

\begin{claim}{B.66}{Proved}{The ordinary pro-form for every characteristic-zero field}
For every field \(k\) of characteristic zero, the functor from
finite-type \(k\)-schemes localized at universal homeomorphisms
to their model pro-systems is fully faithful, with pinned
morphisms in the ordinary slice
\(\operatorname{Pro}(\mathsf T)_{/\mathcal B}\) as defined
above. In particular B9 is proved in arbitrary transcendence
degree; no coherent-relative morphism condition has been added
to its input.
\end{claim}
\begin{proof}
For a pinned ordinary pro-morphism \(\alpha\), choose a target
field \(F_0\) and a relative representative at a finitely
generated source field \(E\). B.65 reconstructs it as a
localized scheme map \(X_E\to Y_{F_0}\). Pairing with the
structural map to \(E\) and extending to \(k\) gives
\(f:X^{\mathrm{awn}}\to Y\). At every other target field
\(F\), B.65 reconstructs that level as well. Pro-compatibility
compares the two maps after one common finitely generated
source extension. The faithfulness proved in Step~4 of B.65
makes their projections to \(Y_{F_0}\) equal; their structural
maps agree, so their lifts to \(Y_F\) and to \(Y\) agree.
At each target field the ordinary pro-class is therefore that
of \(f\). This is surjectivity on the Hom set in the ordinary
pro-category, rather than only on the corresponding limit
point functions.

Conversely, every localized finite-type \(k\)-scheme morphism
and its finite diagram of equations descend to a finitely
generated subfield, and hence define the model pro-morphism.
Its limit map preserves closed points. If two scheme morphisms
give the same pro-morphism, compare at a common target field,
refine once, and use the same fixed-field faithfulness; their
base changes to \(k\) are equal. Absolute weak normalizations
identify the localization, as in CHW Theorem~1.13 and
Proposition~2.28. This proves full faithfulness.
\end{proof}


\subsection{The remaining universal supersingular-target assertion}

\begin{claim}{B.67}{Open}{The exact unsettled quantifier in B5}
For every prime $p$ and every proper smooth connected curve
$C/\overline{\mathbf F}_p$ of genus at least two, there exists
a connected finite \emph{\'etale} cover $D\to C$ whose geometric
Galois closure has order prime to $p$, and a nonconstant morphism
$D\to H$ to a positive-genus supersingular curve $H$.
This assertion remains undecided here for arbitrary $C$.
\end{claim}

We proved the assertion for every odd-characteristic hyperelliptic
source, including characteristic three, and for the further tame
Galois classes in B.38. In characteristic two we proved the
specified cyclic-three and uniformly triple-ramified cases in
B.37. These are statements about actual morphisms of curves and
prime-to-$p$ Galois closures, not only about nonordinary Jacobians.

We tried the corrected torsion-position construction on an elliptic
quotient, then the low-genus quotient and repeated-branch-index
constructions. The first requires such a quotient or an already
constructed correspondence; the latter require the explicit
Galois and index hypotheses in B.38. For a general tame map,
an unramified point in a branch fibre prevents ramification
cancellation, and tameness alone does not control the prime
divisors of the Galois closure. The general tame Belyi theorem
therefore does not finish this step. Nonordinary covering theorems
do not furnish a supersingular target. These failed reductions
are not counterexamples to the displayed assertion.

The first unhandled genus is two in characteristic two for sources
outside B.37; in odd characteristic the first unhandled genus is
three for nonhyperelliptic sources outside B.38. A next needed
step is a prime-to-$p$ covering construction that creates the
uniformly divisible branch fibres used in B.36--4 for one of
these remaining classes, or a genuine obstruction to every such
supersingular correspondence. No finite search is asserted to
exclude all covers of a remaining source.


\section{Consequences and limits}
\subsection{The manuscript's statements and relative class numbers}

\begin{claim}{C.1}{Proved}{The two curve assertions in cote 149, scan 36}
Over every field finitely generated over its prime field, the
point-section map for a smooth hyperbolic curve is injective:
use the full geometric group in characteristic zero, and its
prime-to-$p$ quotient in characteristic $p$. For a proper curve
of genus at least two over a finite field, the covering condition
of scan 36 forces $P=Q$.
\end{claim}
\begin{proof}
In positive characteristic this is exactly
\statementref{A.1} and \statementref{A.33}. Their dependencies
are the same-prime compactification lemma, the cuspidal criterion
with its orientation and Kummer evaluation, the relative
configuration exactness, and the diagonal boundary transport;
each was independently proved in Part A. The covering condition
is equivalent to section conjugacy by \statementref{A.31}, and
the main theorem then gives equality. Characteristic zero is the
all-prime Kummer argument in \statementref{A.30} and
\statementref{A.33}.

Thus the sentence in scan 36, \emph{«en déduire $P=Q$»}, holds in
the specified curve scope. The phrase
\emph{«sans hyp. de caractéristique»} requires retaining the
finite-generation and hyperbolicity hypotheses. It does not make
the corresponding abelianized assertion true in positive
characteristic, and it does not include proper elliptic curves.
\end{proof}

\begin{claim}{C.2}{Refuted}{The characteristic-$p$ prime-to-$p$ extension of Théorème 3 in all dimensions}
The asserted characteristic-$p$ extension of cote 149, scan 32,
Théorème 3, for every geometrically connected finite-type scheme
embedded in a semiabelian variety, is false, already in dimension
one. Over $\mathbf F_2$, take
\[
 E:\quad y^2+xy=x^3+1,\qquad O=\infty,\quad T=(0,1).
\]
The two points have identical prime-to-two section classes in
the whole arithmetic extension. For an arbitrary semiabelian
$G/K$ over a finitely generated characteristic-$p$ field, the
exact replacement is
\[
 \ker\bigl(G(K)\to\hbox{prime-to-}p\hbox{ section classes}\bigr)
       =G(K)[p^\infty]
\]
after using the identity point to identify this section set with
abelian Galois cohomology. A subscheme embedded in $G$ has
injective point sections whenever no difference of distinct
rational points lies in $G(K)[p^\infty]$.
\end{claim}
\begin{proof}
The explicit checks in \statementref{A.24} verify that $E$ is
smooth and $T$ has exact order two. Its geometric fundamental
group is abelian, its prime-to-two part is the product of its
prime-to-two Tate modules, and the difference of the two point
sections is the Kummer class of $T$, which is zero. This is a
conjugacy assertion in the whole group, not a finite-level test.

For a semiabelian variety the same abelian fundamental-group
description and \statementref{A.30} give the exact kernel.
Functoriality of fundamental groups for an embedding $X\to G$
shows that a collision on $X$ maps to a point difference in this
kernel. Its avoidance is therefore a sufficient condition on
$X$. We do not assert the converse for $X$: hyperbolic curves
may still separate points nonabelianly when their generalized
Jacobian Kummer map fails to do so, as the negative control
\statementref{A.25} and the main theorem show. Products
$E\times\mathbf G_m^{d-1}$ give the same counterexample in every
positive dimension.

This matches the warning on scan 34 about translation by
$p$-torsion and the failure of the abelianized map (**) on scan
35. Removing a point from an elliptic curve can restore the
nonabelian hyperbolic assertion, as proved in Part A; it does
not invalidate the warning about (**).
\end{proof}

\paragraph{Readings and the limits of the transcript check.}
We checked the supplied second-pass transcriptions, not absent
images of the original scans. The alternatives
\emph{à l'aise / l'air} on scan 35 do not affect a hypothesis.
The margin's uncertain reading involving an \emph{anneau} is
mathematically consistent with spreading to a finite-type model
and specializing at a closed point, as in \statementref{A.32};
the proof does not depend on selecting that reading. The
alternative \emph{ou / en} in the preceding sentence also changes
no mathematical quantifier used here. In scan 36, a cover $Y'$
in the displayed covering condition is a cover of $Y$, as its
chosen point and generalized Jacobian require; the transcription's
$X$ is the noted variable slip. We interpret the torsion there
as $p$-primary, rather than as necessarily killed by $p$.
No statement in this paper depends on resolving a doubtful word.

\begin{claim}{C.3}{Proved}{The precise class-number consequences and a sharp limitation}
Point separation gives the isolating pointed cover of
\statementref{B.19}; it gives no numerical value for the
$p$-part of its relative class number. There is no universal
assertion that every nontrivial pointed prime-to-$p$ cover has
$p$-unit relative class number, or that every such cover has
$p$-divisible relative class number. Both behaviours occur for
the same curve, point, field, and prime degree in our computation.

For $C_\alpha/\mathbf F_{64}$ of \statementref{B.21},
$\alpha$ with bit code $58$, and $P=\infty$, the degree-seven
characters $(1,0)$ and $(1,1)$ in the recorded basis give
respectively
\[
 h(D_{(1,0)})/h(C)=46713473689,\qquad
 h(D_{(1,1)})/h(C)=55586549824.
\]
The first is odd and the second is even. Nevertheless the
point-section map of $C$ is injective over this field and every
constant extension.
\end{claim}
\begin{proof}
The exact occupancy vectors for these two characters are
$(9,12,17,12,12,17,12)$ and $(11,18,14,8,8,14,18)$.
The degree-two twisted polynomials, including leading coefficient
64 and the infinity pointing, are supplied in
\statementref{B.21}. Their full integer resultants give exactly
the two displayed integers; both are retained in the computed
JSON, and neither is a conjectural value. The main theorem and
its checked inputs apply to this same curve.

A hypothetical collision would force the relative divisibility
in the applicable class-number criterion. Its nonexistence
removes that premise; contraposition does not turn it into
existence of a cover with nondivisible relative class number.
Our explicit constructions prove such existence for the
fourteen requested $\mathbf F_4$ pointed cases and for the
supersingular test curve at primes 191 and 229, as shown in
Part B. Those are additional calculations, not formal
consequences of the point-separation theorem. The universal
quantifier in B1 therefore remains a separate question.
\end{proof}

\begin{claim}{C.4}{Proved}{What transfers to the other research line}
The characteristic-zero construction of a virtual puncture
does not prove the surjectivity of the section map for the
tripod over $\mathbf Q$, and does not construct a birational
lift of an arbitrary section. The present finite-field
injectivity theorem supplies no such assertion. Likewise a
non-point-theoretic full section over a finite field obtained
from a norm-module class in Part B is not a counterexample
over a field finitely generated over $\mathbf Q$.
\end{claim}
\begin{proof}
The pullback construction in the next subsection supplies an
extension for every section, without giving a rational point or
a splitting of its virtual cusp. Recognizing it as the punctured
group of an actual point is an additional condition. The
characteristic-zero separation criterion compares two already
geometric point sections; it has no arbitrary section as a
second input. The finite-field norm-module realization changes
both the characteristic and the arithmetic Galois group, so no
base-change map transfers it to $\mathbf Q$. These are exact
limitations of the maps constructed, not assertions that the
characteristic-zero surjectivity or birational-lifting questions
are false. They address the separate LM6 line referred to in
the supplied intake.
\end{proof}



\subsection{The precise separation consequences for the published theorems}

In this subsection $\Pi_Z^\Sigma$ always retains the entire arithmetic
quotient $G_k$. For a proper curve it is the geometrically pro-$\Sigma$
\'etale group; for an open curve it is the geometrically pro-$\Sigma$
tame group. The set $\Sigma$ may contain $p$, but must contain some
$r\ne p$. The relative comparison used in this scope is
\statementref{A.13}, not a claim that the full pro-$\Sigma$
configuration fibre sequence is exact.

\paragraph{The exceptional set in the source.}
Sa\"idi--Tamagawa 2017, Definitions~1.1, 1.2 and 1.4, first takes
the exceptional locus of the map from points of the universal
pro-cover to \emph{commensurability classes} of decomposition
subgroups, then takes its image $E_Z$ in the closed points of the
compactification. Its ``almost injective in the strong sense'' means
finite exceptional locus for a map of sets. In contrast,
Definition~2.9 calls a proper curve almost $\Sigma$-separated when
the residue fields of its exceptional points lie in a proper
algebraic subfield of $\bar k$. Proposition~2.12 assumes
$J_Z$-largeness and proves this latter support statement; it does
not assert that $E_Z$ is finite. The earlier cofinite result is
ST2009 Proposition~1.8(vi), with almost injectivity of the
closed-point decomposition map in Corollary~1.11(iii).
These are distinct assertions.\citep[Definitions~1.1, 1.2, 1.4,
2.6, 2.9; Proposition~2.12]{LA_ST2017}
\citep[Proposition~1.8(vi), Corollary~1.11]{LA_ST2009}

\begin{claim}{C.5}{Proved: the refined proper theorem with every point labelled}{C1, C5}
Let $X/k_X$ and $Y/k_Y$ be proper smooth geometrically connected
curves of genus at least two over finite fields. Let $\Sigma_X,
\Sigma_Y$ each contain a prime different from the respective
characteristic, and suppose
\[
 \ker\!\left(G_{k_X}\longrightarrow
       \Aut\!\left(\prod_{s\notin\Sigma_X}T_sJ_X\right)\right)
 \ne1,
 \qquad
 \ker\!\left(G_{k_Y}\longrightarrow
       \Aut\!\left(\prod_{s\notin\Sigma_Y}T_sJ_Y\right)\right)
 \ne1.
\]
At the characteristic prime $T_p$ denotes the \'etale Tate module.
For every isomorphism $\alpha:\Pi_X^{\Sigma_X}\simeq
\Pi_Y^{\Sigma_Y}$, the two prime sets, characteristics and field
cardinalities coincide under the group-theoretic identifications,
and $\alpha$ preserves Frobenius. Moreover:
\begin{enumerate}
\item $E_X=E_Y=\varnothing$, and the same is true on every finite
connected \'etale covering arising from an open subgroup;
\item the decomposition-group construction gives a unique bijection
$\phi:X^{\cl}\simeq Y^{\cl}$, preserving residue-field degrees and
carrying $D_x$ to $D_{\phi(x)}$ up to conjugacy at every point;
\item every corresponding pair of finite marked sets $S,T$ may be
used in the cuspidally abelian reconstruction of ST2017
Corollary~3.3; the Green-compatibility assertion in Proposition~3.12
holds at all points and at every corresponding finite covering;
\item $\alpha$ is induced by the unique isomorphism of curves with
their corresponding pro-coverings in ST2017 Theorem~4.22.
\end{enumerate}
The fourth assertion was already the global conclusion of that
theorem. The improvement is the removal of the exceptional-point
restrictions in the preceding reconstruction statements.
\end{claim}
\begin{proof}
The displayed kernel conditions are exactly $J_X$-largeness and
$J_Y$-largeness in Definition~2.6. Proposition~1.7 recovers the
geometric subgroup and the prime sets. Lemma~2.7 gives largeness
relative to the constant fields, and Proposition~2.5 gives equality
of the characteristics and cardinalities and preservation of
Frobenius. Proposition~2.12 gives almost $\Sigma$-separation.
\statementref{A.23}, using the arbitrary-prime-set comparison
\statementref{A.13} and the retained boundary, upgrades it to
$E=\varnothing$. By ST2017 Proposition~2.11(i), an \'etale
$\Sigma$-cover has exceptional set equal to the inverse image of
the original exceptional set, so emptiness holds at every such
level, including constant extensions. No largeness hypothesis on
the new Jacobian is silently added.

Theorem~1.9 reconstructs the set of decomposition subgroups from
the Frobenius group. Emptiness makes the association with points
injective, hence gives $\phi$ on all closed points. The arithmetic
image of $D_x$ is $G_{k(x)}$, which proves the degree assertion.
Corollary~3.3 and Proposition~3.12 now have all their original
hypotheses with $X^{\cl}\setminus E_X=X^{\cl}$ and similarly
for $Y$ and for every covering. Substitution proves precisely the
marked-point and Green-compatibility assertions. Finally the
hypotheses of Theorem~4.22 are exactly those displayed here,
so its curve and pro-covering isomorphism applies.
\citep[Proposition~1.7, Theorem~1.9, Proposition~2.5,
Lemma~2.7, Propositions~2.11--2.12, Corollary~3.3,
Proposition~3.12, Theorem~4.22]{LA_ST2017}
\end{proof}

The $J$-largeness assumption includes every cofinite prime set and
some infinite-complement sets described in ST2017 Remark~2.8.
The present argument does not replace this condition by an arbitrary
nonempty $\Sigma$. In particular, the singleton-prime collisions
of \statementref{A.26} prevent such a replacement. Nor does
$E=\varnothing$ remove the omitted-prime torsion
$J(k)\{\Sigma'\}$ or the quotient by
$k^\times\{\Sigma'\}$ in the pseudo-function reconstruction
(Definition/Remark~4.1 and Theorem~4.2). The finite-index step in
the proof of Theorem~4.22 uses Remark~4.7's separate control of
complementary-prime torsion along a constant-field tower. Point
separation alone does not provide that control.

\begin{claim}{C.6}{Proved: the exact affine hypothesis is retained}{C1, C5}
Let $U,V$ be hyperbolic curves over finite fields, with an
isomorphism of geometrically pro-$\Sigma$ tame groups. Assume,
as in ST2017 Theorem~4.23, that there are corresponding open
subgroups defining covers $U',V'$ whose smooth compactifications
$X',Y'$ have genus at least two, with $\Sigma$ respectively
$J_{X'}$-large and $J_{Y'}$-large. Then $E_U=E_V=\varnothing$,
all closed-point decomposition labels are unique, and the unique
scheme isomorphism with its pro-coverings in Theorem~4.23 has
no exceptional-point restriction. The matched-open-subgroup
hypothesis is part of this assertion.
\end{claim}
\begin{proof}
By C.5, $E_{X'}=E_{Y'}=\varnothing$. If two interior points
of the universal pro-cover of $U$ had commensurate decomposition
groups, restrict to their common open arithmetic subgroup and to
the given open subgroup for $U'$. After a finite constant
extension they give distinct rational points of $U'$ with equal
section subgroups. They remain distinct on $X'$, and projection
to $\Pi_{X'}^\Sigma$ contradicts its all-extension separation.
Two distinct upstairs points above the same geometric point are
already excluded by ST2017 Proposition~1.3. Cusp decomposition
groups are recognized by their nontrivial inertia, and their
commensurability separation is the cusp assertion of ST2009
Proposition~1.8(i),(ii), recalled in ST2017 Theorem~1.8. Thus
no interior or cuspidal exception remains. The same argument
applies to $V$. Theorem~4.23 then applies with its original
matched-cover hypotheses. This proves both the new separation
clause and the already published geometric conclusion.
\citep[Proposition~1.3, Theorems~1.8 and 4.23]{LA_ST2017}
\end{proof}

\begin{claim}{C.7}{Proved: the prime-to-characteristic theorem and its simplified restrictions}{C2, C5}
Let $U,V$ be hyperbolic curves over finite fields, and let
$\Sigma_U^\dagger$ and $\Sigma_V^\dagger$ consist of every prime
different from their respective characteristics. Every isomorphism
of their geometrically pro-$\Sigma$ tame groups arises from a unique
isomorphism of curves and corresponding pro-coverings. All their
closed-point decomposition maps are injective, including after
every constant extension and every open-subgroup covering. For
proper $X,Y$, the restrictions to complements of exceptional sets
in ST2009 Corollary~2.3, Theorem~3.6 and Proposition~3.8 can be
replaced by all points.
\end{claim}
\begin{proof}
For $\Sigma=p'$, \statementref{A.1} gives proper point
separation after every finite extension and
\statementref{A.22} gives $E_X=\varnothing$. If $p$ is also
retained, a collision would survive in the prime-to-$p$ quotient,
so it is excluded as well. At every proper \'etale level the
curve still has genus at least two; the same proof applies there.
For a nonproper hyperbolic curve choose a corresponding finite
tame cover whose compactification has genus at least two, as in
the proof of ST2009 Corollary~3.10, and use the argument of
C.6. The cusp separation is unchanged.

The exact published Isom statements are ST2009 Theorem~3.9
for proper curves and Corollary~3.10 for arbitrary hyperbolic
curves. They permit either all primes or all primes except $p$.
Their hypotheses are met here. In the proper reconstruction,
Corollary~2.3 now applies to arbitrary finite $S,T$;
$\mathcal O_{E_X}=K_X$ and $\mathcal O_{E_Y}=K_Y$ in
Theorem~3.6; and Proposition~3.8 has no excluded points at any
finite covering. The proof of Theorem~3.9 is the formal
combination of Theorem~3.7 with Proposition~3.8, and the proof
of Corollary~3.10 descends compatible proper isomorphisms from
the cofinal family of such covers. These are the precise places
where the separation result removes restrictions.
\citep[Corollary~2.3, Theorems~3.6--3.7, Proposition~3.8,
Theorem~3.9, Corollary~3.10]{LA_ST2009}

This simplification does not set $J_X(k_X)[p^\infty]$ equal
to zero. The $p$-power correction to functions in Theorem~3.6(ii)
therefore remains necessary. Similarly, the reconstruction of
additive structure is still a separate part of the proof; an
injective point-label map alone is not an Isom theorem.
\end{proof}

\paragraph{What becomes of the proposed configuration route.}
ST2009 arXiv Remark~3.13, journal Remark~3.14, would first obtain
an isomorphism of the \emph{full} prime-to-$p$ punctured groups
from the proposed full fibre exactness, then apply the published
Isom theorem to infer equality of the punctures. The full
exactness premise is false by Part~A. Replacing the group by its
cuspidally pro-$r$ quotient does not permit that last application
of the full-group Isom theorem. What does work is the distinct
chain
\[
 \text{relative exactness and diagonal transport}
 \Longrightarrow
 \text{the cuspidal criterion of Part A}
 \Longrightarrow E_X=\varnothing.
\]
This chain is independent of the published Isom conclusion.
Combining it with the explicitly retained cuspidally abelian
and fixed-$r$ reconstruction inputs gives C.7 by the source's
function-reconstruction argument. It is a corrected route to
point separation followed by the known Isom proof, rather than
a restoration of the false full fibre premise.

\begin{claim}{C.8}{Proved: specific Mochizuki separation hypotheses become automatic}{C5}
For proper hyperbolic curves over finite fields and
$\Sigma^\dagger=p'$, the $\Sigma$-separation assumption in
Mochizuki's original Definition~1.18, Proposition~2.6(ii),(iii)
and Corollary~3.11 is automatic. Thus a Frobenius-preserving
isomorphism is point-theoretic, arbitrary finite corresponding
marked sets are allowed in Theorem~2.5(i), the indicated proper
and cuspidally abelian cusp decomposition groups are
commensurably terminal, and the isomorphism is totally globally
Green-compatible. The quotient repairs of A.35
are used at the fibre-identification inputs.
\end{claim}
\begin{proof}
Definition~1.18(i) asks for closed-point decomposition
injectivity on every connected finite \'etale covering arising
from an open subgroup. C.7 proves exactly this requirement.
Decomposition-group recognition and this injectivity change
quasi-point-theoreticity into point-theoreticity under
Definition~1.18(ii), as in Remark~1.18.2(ii). Theorem~2.5(i)
assumes point-theoreticity and gives the marked-set assertion.
Proposition~2.6(ii),(iii) now have their separation assumption
satisfied and give the two commensurable terminality assertions.
Finally Corollary~3.11 assumes finite fields,
$\Sigma^\dagger=p'$ and separation; all three conditions hold.
Its fixed-prime fibre uses are in the repaired scope of A.35.
These are applications of the stated source results with their
actual assumptions, not an assertion that every conclusion of
the entire paper has been reverified.
\citep[Definition~1.18, Remark~1.18.2, Theorem~2.5(i),
Proposition~2.6(ii),(iii), Corollary~3.11]{LA_Mochizuki2007}
\end{proof}

For reference, the original/LaTeX numbering pairs in the last
claim are: Definition~1.18/Definition~1.5;
Theorem~2.5/Theorem~2.1; Proposition~2.6/Proposition~2.3;
Corollary~3.11/Corollary~3.1. The source's
Theorem~2.5(ii) assumes $\Sigma$ is \emph{all} primes in order
to recover principal divisors exactly. Separation does not remove
that hypothesis. Its original Theorem~3.10 (LaTeX Theorem~3.1)
already permits arbitrary matching rational points; the new
fact makes the matching point unique, rather than supplying a
previously absent arbitrary-point existence assertion.

\paragraph{Bounds on the consequence claim.}
Wakabayashi's reconstruction theorems use a single invertible
prime. Their pro-$l$ point labels cannot be declared injective
from a theorem about the larger prime-to-$p$ group; Part~A gives
actual singleton-prime collisions. The virtual-curve construction
in Sun 2025 already applies to arbitrary sections, and Sun 2026's
specified open-injection criterion has its own geometricity
hypotheses. Point separation neither constructs a nongeometric
section nor realizes its virtual extension as a punctured
algebraic curve. Finally no arbitrary-homomorphism anabelian
theorem follows merely by replacing exceptional sets with the
empty set in the checked Isom statements. These limits are part
of the scope of C1, C2 and C5.

\subsection{The separation-hypothesis census in the remaining named papers}

\begin{claim}{C.9}{Proved}{Exact removal of the named separation conditions}
In the original Mochizuki 2007 text, the explicit separation or
point-theoreticity conditions relevant to C5 are the loci in the
following table. For finite fields with $\Sigma^\dagger=p'$,
point-theoreticity and separation at these loci follow from
\statementref{A.1} and its all-cover corollaries, as applied in
C.7--8. Other hypotheses in the same statements remain.
Wakabayashi 2016 and the two Sun preprints have no additional
unproved finite-field prime-to-$p$ separation condition whose
removal strengthens their stated theorems. The distinctions in
the table are part of the conclusion.
\end{claim}
\begin{center}\small
\begin{longtable}{p{0.25\linewidth}p{0.33\linewidth}p{0.34\linewidth}}
\toprule
Locus&Condition or use&Consequence and retained restriction\\\midrule
Mochizuki Definition~1.18; Remark~1.18.2(i)--(iii)&
Separation on every finite open-subgroup cover; passage from
quasi-point-theoretic to point-theoretic; possible failure of
maximality or terminality of point decomposition groups.&
Separation and unique point labels hold at every such level.
For a Frobenius-preserving isomorphism the matching point is
unique. The terminality conclusion is Proposition~2.6 below.\\
Proposition~2.3(ii)&
Proves separation when $\Sigma$ is all primes.&
Its finite-field point-separation conclusion now also holds for
$\Sigma=p'$. Its assertion about detecting principal divisors
in (i) is a different statement and is not extended by this
substitution.\\
Definition~2.4; Theorem~2.5(i)&
Point-theoreticity supplies matching point sets used in Green
compatibility and marked cuspidally abelian reconstruction.&
All finite corresponding sets are allowed, functorially, with
the same cusp and proper-quotient data. The auxiliary definition
of Green compatibility is unchanged.\\
Theorem~2.5(ii)&
Adds $\Sigma=\mathrm{Primes}$ to reconstruct principal
divisors and the stated push-forward of multiplicative groups.&
Point-theoreticity is automatic in the finite-field setting,
but the all-prime hypothesis is retained. Separation alone does
not remove complementary-primary torsion.\\
Proposition~2.6(ii),(iii)&
Assumes $\Sigma$-separation for commensurable terminality in
the proper group and in the cuspidally abelian punctured group.&
Both terminality conclusions hold without an additional
separation assumption in the finite-field prime-to-$p$ cases.\\
Remark~2.6.1; Corollary~2.7&
The remark compares an already given punctured-group
isomorphism with the reconstructed abelian one. The corollary
assumes all primes, point-theoreticity, and total local principal
Green compatibility.&
The point-label condition is automatic as above. Existence of
the punctured-group isomorphism in the remark is not deduced
from separation. In the finite-field all-prime corollary,
Corollary~3.11 supplies Green compatibility as in C.8;
the Isom conclusion was already part of the source's program.\\
Corollary~2.8; Proposition~2.10(iii)&
The former is a point-theoretic affine Isom theorem with all
primes; the latter supplies point-theoreticity for quasi-Belyi
curves over nonarchimedean local fields.&
The finite-field all-prime premise of the former is automatic;
the prime-to-$p$ version is obtained by the separately checked
ST2009 argument, C.7. The local-field result has no new
finite-field consequence from R1.\\
Theorem~3.10; Corollary~3.11&
The theorem assumes matching rational point decomposition
groups, not separation. The corollary assumes finite fields,
$\Sigma^\dagger=p'$, and separation.&
The matching point in the theorem becomes unique; arbitrary
matching-point existence was already present. The corollary's
separation condition disappears and total global Green
compatibility follows.\\
Wakabayashi Theorems~3.9,4.4; Corollary~3.10 and
Remark~3.10.1&
One invertible prime $l$; matching decomposition groups in the
punctured theorem; the remark explicitly allows several
matching closed points.&
The source does not assume a mixed-prime separation condition.
R1 cannot make its single-prime point correspondence unique;
the collisions in Part~A prohibit that inference.\\
Sun 2025, Definition~3.1, Proposition~3.7, Corollary~3.8,
Theorems~A--C&
Virtual pullbacks are defined for arbitrary sections in
characteristic zero. Geometric sections identify their actual
fibres; the main reconstruction statements use number fields
or mixed-characteristic local fields.&
There is no conditional finite-field separation clause to
remove. Point separation does not prove that a virtual curve
has an algebraic realization or that its section is geometric.\\
Sun 2026, Lemma~5.6; Theorem~5.7 and Corollary~5.8&
Lemma~5.6(i) uses the already known injectivity of points into
full sections over an arithmetic field; (ii) uses an additional
open-immersion geometricity theorem. The later results impose
specific virtual covering and cusp conditions.&
No such condition is supplied merely by R1. The source's
arithmetic fields are number fields or mixed-characteristic
local fields, and its groups are full \'etale groups.
None of these assertions becomes a finite-field prime-to-$p$
section-surjectivity theorem.\\
\bottomrule
\end{longtable}
\end{center}
\begin{proof}
We searched the complete opened texts for the source's
separation, point-theoreticity, point-injectivity and geometricity
terms, then read the surrounding numbered statements and their
proof uses. In Mochizuki's original text, the explicit
$\Sigma$-separation occurrences are Definition~1.18,
Remark~1.18.2, Proposition~2.3(ii), Proposition~2.6(ii),(iii)
and Corollary~3.11. The point-theoreticity occurrences add
Definition~2.4, Theorem~2.5, Remark~2.6.1, Corollaries~2.7--2.8
and Proposition~2.10(iii). Theorem~3.10 requires corresponding
point groups, but not separation. This produces the displayed
finite list; no unexamined whole-paper assertion is substituted
for it.

C.7 proves separation on every required cover, C.8 proves
the resulting point-theoreticity and terminality, and A.35
and A.41 supply the repaired quotient fibre inputs. Substituting
these established assertions into the displayed hypotheses gives
the consequences in the first eight rows. In particular, the
proof of Theorem~2.5(i) restricts the reconstructed cuspidally
abelian configuration extension to the matched point groups;
the repaired abelian fibre is exactly what this restriction
requires. The all-prime requirement in (ii) is separate and
has not been removed. The proof of Corollary~3.11 uses the
fixed-$l$ reconstruction for each $l\in\Sigma^\dagger$
and recognizes noncuspidal decomposition groups there. A.41
repairs that input, while all-cover separation makes each of
the matching labels unique. This is precisely total global
Green compatibility.

Wakabayashi's Remark~3.10.1 expressly distinguishes existence
of a matching point from uniqueness. Its surrounding theorems
remain in the single-prime category described in A.37, so
the larger prime-to-$p$ result does not alter their hypotheses.
Sun 2025 constructs the virtual pullback before making any
assumption that its section is a point. Sun 2026 defines an
arithmetic field to mean a number field or a mixed-characteristic
local field; its Lemma~5.6(i) invokes the injectivity portion
of the section conjecture in that already known full-group
setting. Its part (ii) and later geometricity criteria have
additional inputs, not a missing point-separation premise.
These readings prove the last three rows and complete the
stated separation-condition census. They also explain why no
new arbitrary-homomorphism theorem or section-surjectivity
theorem has been inferred here.
\citep[Definition~1.18, \S2, Theorem~3.10,
Corollary~3.11]{LA_Mochizuki2007}
\citep[Theorems~3.9,4.4, Remark~3.10.1]{LA_Wakabayashi2016}
\citep[Definition~3.1, Proposition~3.7, Corollary~3.8,
Theorems~A--C]{LA_Sun2025}
\citep[\S2, Lemma~5.6, Theorem~5.7, Corollary~5.8]{LA_Sun2026}
\end{proof}


\subsection{Virtual punctures and the characteristic-zero fibre construction}

\begin{claim}{C.10}{Proved: the exact pullback construction}{C4}
Let $k$ have characteristic zero, let $C/k$ be hyperbolic, and consider
the projection $C^{(2)}=(C\times_k C)\setminus\Delta_C\to C$.
Work in a specified quotient category in which there is an exact
sequence
\[
 1\longrightarrow F\longrightarrow E\xrightarrow{q}\Pi_C
 \longrightarrow1.
\]
For every section $s:G_k\to\Pi_C$ the profinite pullback
\[
 \Pi_s=E\times_{\Pi_C,s}G_k
\]
is an extension $1\to F\to\Pi_s\to G_k\to1$. If $s=s_P$ is the
section of $P\in C(k)$, the natural fibre map identifies $\Pi_s$ with
the corresponding quotient of $\Pi_{C\setminus P}$, whenever the
assumed exact fibration statement includes that fibre identification.
\end{claim}
\begin{proof}
The kernel of the pullback projection is $F\times\{1\}$. Surjectivity
follows from that of $q$, proving exactness. For a geometric section
the map from the arithmetic group of the actual fibre has the specified
geometric fibre map and is the identity on $G_k$. If those geometric
groups agree, the extension isomorphism criterion proves that the
arithmetic groups agree. These arguments are valid for profinite groups
and continuous maps, without requiring the section to come from a point.
\end{proof}

\paragraph{The quotient qualification is necessary.}
An exact sequence in a quotient cannot recover information already
discarded by that quotient. For example, killing the entire diagonal
inertia gives the product group with its exact projection sequence;
its fibre is the unpunctured geometric group. Consequently the phrase
``the punctured group in every exact quotient'' means the corresponding
punctured quotient and its retained boundary data. It cannot literally
mean the full punctured fundamental group in a quotient that kills
the new cusp. This is a limitation of that wording, not a failure of
the pullback construction.

\begin{claim}{C.11}{Proved: retaining the distinguished cusp}{C4}
Assume the chosen configuration quotient also retains a diagonal
boundary decomposition group $D_\Delta$ mapping surjectively to
$\Pi_C$, with kernel the surviving inertia $I_\Delta$. Then
\[
 D_s=D_\Delta\times_{\Pi_C,s}G_k\subset\Pi_s
\]
is the distinguished virtual cusp extension. The pair $(\Pi_s,D_s)$
depends only on the geometric conjugacy class of $s$, up to the
corresponding inner ambiguity. For $s_P$ it is the actual punctured
fibre group with its new cusp in the chosen quotient.
\end{claim}
\begin{proof}
Suppose $s'=\Inn(\gamma)s$ for $\gamma\in\Delta_C$. The diagonal
decomposition map has a geometric lift $d$ of $\gamma$; conjugation
by $d$ identifies the inverse images of the two section subgroups and
their intersections with $D_\Delta$. Thus it transports both the
extension and the cusp. If the second proper projection is retained,
it changes that projection by $\Inn(\gamma)$, which may be removed
by a geometric inner automorphism in the target fibre because its
map to the other proper factor is surjective. This leaves the cusp
conjugacy class unchanged. For $s_P$, a local parameter at $P$ gives
the ordinary boundary decomposition description, and the exact fibre
comparison gives the asserted identification.
\end{proof}

The diagonal lift is extra geometric information. Exactness of
$E\to\Pi_C$ alone does not specify which lift preserves a cusp.
Also $\Pi_s$ is an extension, not automatically a split extension;
the construction has not supplied a section of $\Pi_s\to G_k$.
In particular the pullback of the diagonal cusp extension along a
nongeometric section is a virtual boundary datum, whose realizability
as a rational puncture has not been proved by writing the fibre product.

\begin{claim}{C.12}{Proved: precise relation to Sun's virtual curves}{A14, C4}
The construction in C.10 is already Sun's virtual fundamental group.
The two opened preprints do not assert that every virtual puncture is
the puncture at a rational point.
\end{claim}
\begin{proof}
Sun 2025 Definition~3.1 starts with a smooth surjective morphism over
a characteristic-zero field and a Galois section of the base, and defines
the virtual group by this pullback. Corollary~3.2 records its extension
and the independence of the geometric kernel from the section.
For a hyperbolic-curve family, Proposition~3.7 gives the fibre comparison,
and Corollary~3.8 identifies the virtual group with the actual fibre
when the section is geometric. This is a direct predecessor, rather
than a new construction of this paper.

Sun 2026 develops additional reconstruction and geometricity criteria.
Its introduction's Theorem~D corresponds to Corollary~5.8 in the body
(the introduction calls the reference a lemma). It concerns a number
field or a mixed-characteristic local field, virtual curves arising
from curves of types $(0,3)$ and $(1,3)$, and an open injection of their
virtual groups respecting their geometric kernels. Under those
hypotheses it forces the geometric index to be two and both sections
to be geometric. These restrictive hypotheses are not available for
an arbitrary section over a finitely generated extension of $\Q$.
\citep[Definition~3.1, Corollary~3.2, Proposition~3.7 and
Corollary~3.8]{LA_Sun2025}
\citep[Theorem~D and Corollary~5.8]{LA_Sun2026}
\end{proof}

\begin{claim}{C.13}{Proved: the characteristic-zero recognition criterion}{C4}
Suppose $k$ is finitely generated over $\Q$, $C/k$ is hyperbolic,
and the base quotient retains the full arithmetic fundamental group,
or at least the point-separating Kummer maps used below. A cusp-preserving
isomorphism over the identity of this base quotient between the
punctures at $P,Q\in C(k)$ forces $P=Q$.

More generally, in the full base quotient, if the virtual puncture
$(\Pi_s,D_s)$ is isomorphic over $\Pi_C$, preserving the distinguished
cusp, to the actual puncture at $P$, then $s$ and $s_P$ are geometrically
conjugate. Conversely a geometric section has that puncture by C.10.
\end{claim}
\begin{proof}
Projecting the cusp decomposition groups to $\Pi_C$ identifies the
point sections up to geometric conjugacy. If the smooth proper
compactification $\bar C$ has positive genus, abelianization gives
zero Kummer class for $[P-Q]\in J_{\bar C}(k)$. The
Mordell--Weil--Lang--N\'eron finite-generation theorem applies to
$J_{\bar C}(k)$. For a finitely generated
abelian group $A$, the map to all its finite-quotient completions is
injective: $\bigcap_{n\ge1}nA=0$. The Kummer exact sequences therefore
make the all-prime Kummer map injective on $J_{\bar C}(k)$. Hence $[P-Q]=0$.
A nonconstant function with divisor $P-Q$ would have degree one and
identify $\bar C$ with $\mathbf P^1$, impossible for positive genus; thus
$P=Q$.

If $\bar C$ has genus zero, extend $k$ finitely so that two boundary
points $a,b$ are rational. The existence of $P$ identifies $\bar C$
with $\mathbf P^1$. A degree-one coordinate with zero at $a$ and pole
at $b$ gives an injective map $C\to\mathbf G_m$ on rational points.
The full Kummer map on $k^\times$ is injective. To see this, let $K$
be the algebraic closure of $\Q$ in $k$, a number field. If an element
of $k^\times$ has an $n$th root in $k$ for every $n$, all its divisorial
valuations on a normal projective $K$-model are zero, so it lies in
$K^\times$. Its roots also lie in $K$, because they are algebraic over
$K$. Their valuations at every finite prime are zero; the finitely
generated unit group of $K$ has trivial intersection of all its
$n$th-power subgroups. Thus the element is one. Applying this to the
ratio of the coordinate values at $P,Q$ again proves $P=Q$.

For the virtual statement, the image of its distinguished cusp in
$\Pi_C$ is $s(G_k)$ by definition. An over-$\Pi_C$, cusp-preserving
comparison with the actual puncture makes this subgroup conjugate
to $s_P(G_k)$. Their projections to $G_k$ identify the corresponding
sections. Conversely C.10 and C.11 give the comparison when $s$
is geometric.
\end{proof}

This gives a recognition equivalence for the open characteristic-zero
question; it neither constructs a nongeometric section nor proves
that none exists. The finite-field proof of AG.40 used finiteness of
$J_C(k)$ to obtain a principal multiple of $P-Q$, constant-field point
counts, and the weight calculation eliminating nonconstant proper
Kummer classes. Those steps cannot simply be reused over $k$.
In a singleton-$l$ quotient, even the abelian Kummer map on $J_C(k)$
has the prime-to-$l$ torsion subgroup in its kernel. Thus exactness of
a chosen configuration quotient is not by itself a proof of point
separation for its base quotient. The full-group case, including
affine curves, is covered by the genus-zero and positive-genus
arguments above.\citep[\S1.2]{LA_CiperianiStix2015}

\bibliographystyle{plainnat}
\bibliography{ANABELIAN_A1_references}
\end{document}
